% Appendix F body (mass-gap module) -- integrated into the main manuscript.
% Mass-gap module (single-path presentation; earlier internal route labels removed).

% --- Local macros for Appendix F ---
\providecommand{\Nb}{N_b}
\providecommand{\Ccount}{C_{\mathrm{count}}}
\providecommand{\Cint}{C_{\mathrm{int}}}
\providecommand{\Cact}{C_{\mathrm{act}}}
\providecommand{\Kzero}{K_0}
\providecommand{\Ktwo}{K_2}
\providecommand{\Krep}{K_{\mathrm{rep}}}
\providecommand{\Kproj}{K_{\mathrm{proj}}}
\providecommand{\Kact}{\overline{K}_{\mathrm{act}}}
\providecommand{\gb}{g_b}
\providecommand{\qkp}{q}
\providecommand{\dist}{\mathrm{dist}}
\providecommand{\E}{\mathbb{E}}

\begin{remark}[Appendix~F deliverable]
The unique paper-level target theorem is Theorem~\ref{thm:main}.
This appendix proves three Route~1 outputs on the tuned dyadic spine.
First, it proves the one-shot entrance package at the explicit block scale $B(\beta)$.
Second, it closes QE2$_{\mathrm{lat}}$ by combining exact terminal-block conditional independence with entrance-scale Dobrushin decay and feeding the resulting uniform local-observable decay into the fixed-lattice OS gap theorem.
Third, it closes QE3 by transporting that full fixed-lattice gap through the dyadic bounded-state limits and the continuum-time OS reconstruction furnished by Lane~A.
A dense-core Euclidean-clustering route is retained later only as a downstream reformulation feeding Lane~B.
\end{remark}


\begin{quote}\small
\noindent\textbf{Notation and constants (quick dictionary).}\par\smallskip
\noindent$u_{\mathrm{ren}}$: fixed weak-window target coupling used to define the tuned dyadic trajectory.\par
\noindent$\ell_0$: fixed physical length unit; $a_n=\ell_0 2^{-n}$ and $\beta_n:=\beta(a_n;u_{\mathrm{ren}})$.\par
\noindent$q_\ast$: entrance-scale Dobrushin contraction constant (Theorem~\ref{thm:h60_one_shot_entrance}).\par
\noindent$m_{\mathrm{blk}}$: block-mixing rate, $m_{\mathrm{blk}}:=-\log q_\ast$ (see~\eqref{eq:mstar_explicit}).\par
\noindent$C_{\mathrm{ent}}$: entrance bound for the one-shot scale: $a_nB_n\le C_{\mathrm{ent}}\ell_0$ along the tuned dyadic sequence (Corollary~\ref{cor:route1_bounded_physical_entrance}).\par
\noindent$\mathfrak m_\ast$: target physical gap constant produced directly from $m_{\mathrm{blk}}$ and the entrance bound $C_{\mathrm{ent}}$; see~\eqref{eq:mstar_explicit}.\par
\noindent$\Delta_{\mathrm{lat}}(\beta_n)$: full fixed-lattice OS/transfer-matrix gap in lattice units along the tuned dyadic path; physical rate is $\Delta_{\mathrm{lat}}(\beta_n)/a_n$.\par
\noindent$m_{\mathrm{lat}}(\beta_n)$: auxiliary probe-family/channel clustering rate, retained only as a diagnostic and not used on the primary QE2/QE3 spine.
\end{quote}

\subsection*{Proof roadmap (load-bearing vs.\ exploratory)}\label{sec:massgap_roadmap}
\addcontentsline{toc}{subsection}{Proof roadmap}

\noindent\textbf{What this appendix proves.}
The unique paper-level target theorem is Theorem~\ref{thm:main}.
This appendix proves the Route~1 entrance package (QE1), the tuned-sequence full fixed-lattice OS gap lower bound (QE2$_{\mathrm{lat}}$), and the continuum vacuum-gap transport step (QE3) on the tuned dyadic spine.
The dense-core Euclidean-clustering / Lane~B chain is recorded only as a downstream reformulation of the already established Route~1 gap.

\medskip
\noindent\textbf{Primary proof path.}
The front-facing Route~1 package now has the form:
\begin{enumerate}[leftmargin=*]
\item one-shot entrance into the KP/Dobrushin basin together with the tuned-sequence physical entrance scale (proved here);
\item exact terminal-block covariance transfer plus entrance-scale block decay, fed into the fixed-lattice OS gap theorem, yielding QE2$_{\mathrm{lat}}$ unconditionally;
\item dyadic bounded-state convergence plus Lane~A continuum-time OS transport, yielding the continuum spectral gap by non-collapse.
\end{enumerate}

\medskip
\noindent\textbf{Secondary material.}
A dense-core Euclidean-clustering route feeding Lane~B is retained later in the appendix as an alternative reformulation.
Measure-level domination and variance/CLT reductions remain exploratory side-branches and are not part of the front-facing chain unless explicitly cited.


\subsection{Front-facing closure chain (QE1/QE2-lat/QE3)}
\label{sec:main_theorem_package}

This subsection isolates the frozen Route~1 proof spine.
The front-facing package is
\begin{center}
\fbox{\parbox{0.94\linewidth}{\small\centering
\textbf{QE1:} one-shot entrance + tuned-sequence physical entrance bound (proved here)\\[2pt]
$\Longrightarrow$\\[2pt]
\textbf{QE2$_{\mathrm{lat}}$:} full fixed-lattice OS gap from exact entrance-scale block-to-microscopic transfer (proved here)\\[2pt]
$\Longrightarrow$\\[2pt]
\textbf{QE3:} continuum OS spectral gap by non-collapse, closed by dyadic-to-continuum OS transport through Lane~A.
}}
\end{center}
Thus QE1, QE2$_{\mathrm{lat}}$, and QE3 are all proved on the primary Route~1 spine.
A dense-core Route~1 / Lane~B bridge is recorded later, but it is downstream of the already established gap and is not part of this primary package.
Continuum sharp-local existence, universality, and scheme-independence are proved in the Lane~A / Lane~C modules and are cited directly in Theorem~\ref{thm:main}; they are not repackaged here.
Throughout we work in lattice dimension $d=4$ and with a fixed compact connected simple gauge group $G$ together with the Wilson representation $\rho$ of Theorem~\ref{thm:main}.

\subsubsection{Standing choices and notational conventions}

\paragraph{Wilson lattice action.}
On each finite lattice region $\Lambda\subset\mathbb Z^4$ (torus or box with admissible boundary conditions) we consider the standard Wilson gauge action
\begin{equation}
 S_{\beta}(U):=\beta\sum_{p\in P(\Lambda)} A(U_p),
 \qquad A(g):=1-\frac{1}{d_{\pi}}\Re\,\mathrm{Tr}_{\pi}(g),
 \label{eq:wilson_action_main}
\end{equation}
with $\pi=\rho$ the fixed Wilson representation (equivalently, after descent to the effective quotient $G_{\mathrm{eff}}:=G/\ker(\rho)$; cf.\ \Cref{rem:effective_group_convention,def:effective_group,lem:center_blindness_descent}).
The corresponding finite-volume Gibbs measure is
\(d\mu_{\beta,\Lambda}(U)\propto e^{-S_{\beta}(U)}\,dU\), with $dU$ Haar on edges and tree gauge used when derivatives are taken.

\begin{remark}[Global form, effective quotient, and line operators]
For a fixed pair $(G,\pi)$, the Wilson plaquette weight depends on $G$ only through the effective quotient
$G_{\mathrm{eff}}:=G/\ker(\pi)$ (\Cref{def:effective_group,lem:center_blindness_descent}).
In particular, for the $\Spin$ families with $\pi$ equal to the defining (vector) representation, $G_{\mathrm{eff}}$ is the corresponding $\SO$ group.
The sharp local algebra treated in this manuscript is generated by strictly local gauge-invariant observables (flowed local composites of curvature),
and therefore cannot detect the difference between global forms with the same Lie algebra once $\pi$ is fixed (\Cref{lem:adjoint_local_generators}).
Differences between global forms are encoded in \emph{nonlocal} line/surface operators (e.g. Wilson lines in spinor representations and the dual \,''t~Hooft defects),
which correspond to superselection data for representations of the same local net; see \S\ref{subsec:spin_vs_so}.
\end{remark}

\paragraph{Admissible action class and schemes.}
We allow, when stated, local gauge actions in the universality class of Definition~\ref{def:action_universality_class} and renormalization schemes in the sense of Definition~\ref{def:admissible_scheme}.
The Wilson action \eqref{eq:wilson_action_main} is a distinguished member of that class.

\paragraph{Tuned dyadic scaling trajectory.}
Fix once and for all a weak-window target coupling $u_{\mathrm{ren}}\in(0,u_\ast]$ and the physical length unit $\ell_0>0$.
The load-bearing Route~1 path uses the dyadic sequence
\begin{equation}
 a_n:=\ell_0 2^{-n},
 \qquad
 \beta_n:=\beta(a_n;u_{\mathrm{ren}}),
 \qquad
 g_{\mathrm{GF}}^2(\ell_0;a_n,\beta_n)=u(a_n),
 \label{eq:a_beta_main}
\end{equation}
where $u(a_n)=u_{-n}$ is the backward step-scaling trajectory from Proposition~\ref{prop:h59_uv_tuning_exists} and
$\beta(a_n;u_{\mathrm{ren}})$ is selected by Theorem~\ref{thm:h49_bare_flow_window}.
This tuned dyadic sequence is the only renormalization prescription used on the paper-level Route~1 mass-gap spine.
Older $t_0$-based scale-setting formulas are retained only in auxiliary/archival subsections and are not cited in the final theorem chain.

\paragraph{One-shot entrance scale.}
For the load-bearing Route~1 chain, the entrance mechanism is the explicit one-shot block size $B(\beta)$ from \eqref{eq:h60_Bbeta_def}.
Fix once and for all a target value $u_{\mathrm{ren}}\in(0,u_\ast]$ in the weak window and let
\[
a_n:=\ell_0\,2^{-n},\qquad \beta_n:=\beta(a_n;u_{\mathrm{ren}}),\qquad B_n:=B(\beta_n)
\]
be the associated tuned dyadic trajectory from Proposition~\ref{prop:h59_uv_tuning_exists}.
Corollary~\ref{cor:route1_bounded_physical_entrance} gives a uniform constant $C_{\mathrm{ent}}<\infty$ and an entry index $N_{\mathrm{ent}}(u_{\mathrm{ren}})$ such that
\begin{equation}
 a_n\,B_n\le C_{\mathrm{ent}}\,\ell_0\qquad(n\ge N_{\mathrm{ent}}(u_{\mathrm{ren}})).
 \label{eq:bounded_phys_entrance_main}
\end{equation}
Whenever a generic entrance-scale symbol appears in earlier comparison lemmas, the final theorem chain instantiates it with this same one-shot scale $B(\beta)$; no second entrance mechanism is used in the paper-level proof.

\subsubsection{QE1/QE2$_{\mathrm{lat}}$/QE3 chain used in Theorem~\ref{thm:main}}

\begin{proposition}[QE1: one-shot entrance and tuned-sequence physical scale]\label{prop:route1_QE1}
Within the standing choices above there exist $\beta_\ast<\infty$, $q_\ast\in(0,1)$, and $C_{\mathrm{ent}}<\infty$ such that:
\begin{enumerate}[label=\textup{(\roman*)},leftmargin=2.5em]
\item for every $\beta\ge \beta_\ast$ the exact one-shot blocked specification at the block size $B(\beta)$ satisfies
\begin{equation}
\delta\!\bigl(\gamma^{\langle B(\beta)\rangle}\bigr)\le q_\ast<1;
\label{eq:bounded_phys_entrance_main_repack_a}
\end{equation}
\item for every fixed target value $u_{\mathrm{ren}}\in(0,u_\ast]$, along the associated tuned dyadic trajectory $a_n=\ell_0 2^{-n}$, $\beta_n=\beta(a_n;u_{\mathrm{ren}})$, $B_n=B(\beta_n)$, there exists an entry index $N_{\mathrm{ent}}(u_{\mathrm{ren}})$ such that
\begin{equation}
 a_n B_n\le C_{\mathrm{ent}}\ell_0
 \qquad (n\ge N_{\mathrm{ent}}(u_{\mathrm{ren}})).
 \label{eq:bounded_phys_entrance_main_repack}
\end{equation}
\end{enumerate}
In particular the exact one-shot blocked state is unique and exponentially mixing in entrance-block distance, and Proposition~\ref{prop:route1_QE1} supplies the scale data used later in the exact block-to-microscopic transfer and non-collapse chain. Only the entry index, not the constant $C_{\mathrm{ent}}$, depends on the chosen target value.
\end{proposition}

\begin{proof}
This is exactly Theorem~\ref{thm:h60_one_shot_entrance} together with Corollary~\ref{cor:route1_bounded_physical_entrance}.
\end{proof}

\begin{proposition}[QE2$_{\mathrm{lat}}$: unconditional tuned-sequence full fixed-lattice gap]\label{prop:route1_QE2}\label{cor:route1_QE2_full_lattice_gap}
Within the standing choices above, let
\[
a_n:=\ell_0 2^{-n},\qquad \beta_n:=\beta(a_n;u_{\mathrm{ren}}),\qquad B_n:=B(\beta_n).
\]
Then there exist $n_\ast<\infty$ and $\mathfrak m_\ast>0$ such that for all $n\ge n_\ast$ the fixed-lattice OS Hamiltonian for the infinite-volume Wilson Gibbs state at bare coupling $\beta_n$ has full vacuum gap $\Delta_{\mathrm{lat}}(\beta_n)>0$ and
\begin{equation}
 \frac{\Delta_{\mathrm{lat}}(\beta_n)}{a_n}\ge \mathfrak m_\ast.
 \label{eq:route1_QE2_lat_main}
\end{equation}
One may take
\begin{equation}
 \mathfrak m_\ast=\frac{m_{\mathrm{blk}}}{2C_{\mathrm{ent}}\ell_0},
 \qquad m_{\mathrm{blk}}:=-\log q_\ast.
 \label{eq:mstar_explicit}
\end{equation}
\end{proposition}

\begin{proof}
This is exactly Theorem~\ref{thm:physical_gap_from_gluing}.
\end{proof}

\begin{corollary}[QE3: sharp-local continuum vacuum gap on the Route~1 spine]\label{cor:route1_QE3_gap}
Assume Proposition~\ref{prop:route1_QE2} and let $\omega^{(u_{\mathrm{ren}})}$ be the full sharp-local state furnished by Corollary~\ref{cor:LaneA_full_scope} over the same tuned dyadic bounded base state.
Let $(\mathcal H,\Omega,H)$ be its OS reconstruction.
Then
\[
\mathrm{Spec}(H)\subset\{0\}\cup[\mathfrak m_\ast,\infty),
\qquad
\ker(H)=\mathbb C\Omega,
\qquad
P_{(0,\mathfrak m_\ast)}=0.
\]
Equivalently,
\[
\|e^{-tH}\psi\|\le e^{-\mathfrak m_\ast t}\,\|\psi\|
\qquad (t\ge 0,\ \psi\perp\Omega).
\]
This is the continuum sharp-local spectral consequence used in Theorem~\ref{thm:main}(D).
\end{corollary}

\begin{proof}
Proposition~\ref{prop:route1_QE2} yields
\[
\liminf_{n\to\infty}\frac{\Delta_{\mathrm{lat}}(\beta_n)}{a_n}\ge \mathfrak m_\ast.
\]
Apply Theorem~\ref{thm:phys_gap_noncollapse} with $m_*=\mathfrak m_\ast$.
The semigroup estimate is then the spectral-theorem reformulation of the same gap statement on $\Omega^\perp$.
\end{proof}

\begin{theorem}[Appendix~F primary Route~1 proof spine]\label{thm:route1_laneB_package}
Within the standing choices above, QE1, QE2$_{\mathrm{lat}}$, and QE3 hold on the tuned dyadic Route~1 spine.
The dense-core Route~1/Lane~B bridge retained later is therefore a downstream reformulation rather than an additional load-bearing interface.
\end{theorem}

\begin{proof}
Proposition~\ref{prop:route1_QE1} gives QE1, Proposition~\ref{prop:route1_QE2} gives QE2$_{\mathrm{lat}}$, and Corollary~\ref{cor:route1_QE3_gap} gives QE3.
\end{proof}

\begin{remark}[Reading guide for a first audit pass]
A referee auditing only the primary Route~1 package may start with \Cref{rem:route1_packet38_transport_certification}, which freezes the fixed-lattice-to-continuum theorem chain, the only live weak-window input, the first Lane~A theorem-level dependency, and the fact that Lane~B is downstream only.
Then verify QE1 from Theorem~\ref{thm:h60_one_shot_entrance} together with the tuned-sequence physical-scale corollary, Corollary~\ref{cor:route1_bounded_physical_entrance}.
Then read Lemma~\ref{lem:block_decay_to_micro}, Proposition~\ref{prop:lattice_OS_cyclicity_local_algebra}, Corollary~\ref{cor:uniform_phys_gap_from_bounded_entrance}, and Theorem~\ref{thm:physical_gap_from_gluing}, which give the unconditional block-to-microscopic transfer and the tuned-sequence full fixed-lattice gap lower bound for QE2$_{\mathrm{lat}}$.
Then read Theorem~\ref{thm:Wilson_patch_transport}, Theorem~\ref{thm:continuum_tightness}, Theorem~\ref{thm:continuum_time_OS_upgrade}, Theorem~\ref{thm:phys_gap_noncollapse}, and Corollary~\ref{cor:route1_QE3_gap}, which close the continuum spectral-gap step.
The dense-core route through Theorem~\ref{thm:route1_E2_internal} and Lane~B is secondary in this organization.
\end{remark}

\begin{corollary}[Optional diagnostic: weak-window step scaling has $b_0>0$ (not used)]\label{cor:b0_optional}
The gradient-flow step-scaling function admits a quadratic expansion with coefficient $b_0>0$
(Proposition~\ref{prop:continuum_stepscaling_b0}).
This diagnostic is not used anywhere in the mass-gap proof.
\end{corollary}

\paragraph{Organization note.}
The remainder of this appendix proves QE1 in full detail, closes QE2$_{\mathrm{lat}}$ by exact entrance-scale covariance transfer together with the fixed-lattice OS gap theorem, and then closes QE3 by transporting dyadic bounded-state limits through the Lane~A continuum-time OS reconstruction.
The dense-core Euclidean-clustering route is retained later as a secondary bridge to Lane~B.
Corollary~\ref{cor:b0_optional} records an optional weak-window diagnostic and is not used elsewhere.
Any later harmonization of notation or constant names across modules is purely expository and does not affect the stated inequalities or conclusions.


\subsection{Normalization, activity, and the basic envelope \texorpdfstring{$\gb(f)$}{g\_b(f)}}

Fix lattice dimension $d\ge 2$, compact Lie group $G$, block size $b\ge 2$, and polymer parameter $\theta>0$.
Set
\begin{equation}
\Nb:=b^2.
\end{equation}

Let $dg$ be Haar \emph{probability} measure on $G$.
For each irrep $R$ with character $\chi_R$ and dimension $d_R$, define the normalized character
\begin{equation}
\varphi_R(g):=\frac{\chi_R(g)}{d_R}
\qquad\Rightarrow\qquad |\varphi_R(g)|\le 1\ \ \forall g\in G.
\end{equation}

A (canonically normalized) plaquette weight is expanded as
\begin{equation}\label{eq:W_expand}
W(g)=1+\sum_{R\ne\mathbf 1} a_R\,\varphi_R(g),
\end{equation}
with scalar activity
\begin{equation}\label{eq:f_def}
 f:=\sum_{R\ne\mathbf 1}|a_R|.
\end{equation}

\begin{lemma}[Tile-mass envelope]\label{lem:Cint_one}
For all $g\in G$, $|W(g)|\le 1+f$.
Consequently the tile-mass majorant
\begin{equation}\label{eq:gb_def}
\gb(f):=(1+\Cint f)^{\Nb}-1
\end{equation}
holds with $\boxed{\Cint=1}$.
\end{lemma}
\begin{proof}
Immediate from $|\sum a_R\varphi_R(g)|\le \sum |a_R|=f$.
\end{proof}

\paragraph*{Short-distance products and OPE-style bounds (optional refinement)}
\label{subsubsec:ope_bounds}

The mass gap objective (existence of an OS-positive continuum Yang--Mills theory with a spectral gap) does not require an operator product
expansion at coincident points. However, if one wishes to promote the inductive-limit algebra $\mathfrak A_{\mathrm{loc}}$ to a
\emph{fully local} field algebra closed under pointwise products (as distributions), then one needs quantitative short-distance control.
In the present paper this is treated as an \emph{additional refinement module} built on the same multiscale inputs already used for
dyadic Cauchy convergence.

\begin{lemma}[Uniform moment bounds for renormalized multiplets at fixed truncation]
\label{lem:uniform_Lp_moments_trunc}
Fix $\Delta_{\max}$ and let $\mathbf{\mathcal O}^{(\le \Delta_{\max})}_{\mathrm{ren}}$ be the renormalized multiplet from
Theorem~\ref{thm:field_algebra_truncation}.
For each $p<\infty$ and each test function $\phi\in C_c^{\infty}(\mathbb R^4)$,
\[
\sup_{j\ge 0}\ \|X^{(\le\Delta_{\max})}_j(\phi)\|_{L^p(\mathbb R^{M(\Delta_{\max})})}\ <\ \infty,
\]
and the same bound holds uniformly over actions in the universality class of Definition~\ref{def:action_universality_class}
for $\beta$ sufficiently large.
\end{lemma}

\begin{proof}
On the small-curvature event $\Omega_{\varepsilon_0}$ each flowed local observable is a finite sum of local curvature polynomials with
uniformly bounded coefficients (heat-kernel norms) at the relevant scale.
Curvature Brascamp--Lieb domination (Proposition~\ref{prop:curvature_BL}) implies sub-Gaussian bounds for linear curvature functionals and
Hanson--Wright-type bounds for quadratic polynomials; by standard hypercontractivity for polynomial functionals of sub-Gaussian fields,
all moments of these finite-degree local polynomials are controlled uniformly.
The complement $\Omega_{\varepsilon_0}^c$ contributes an exponentially small tail by Lemma~\ref{lem:weak_localization}.
Finally, $Z^{(\le\Delta_{\max})}(t_j)=I+O(u_j)$ is uniformly bounded in operator norm along dyadic scales (since $\sum u_j^2<\infty$),
so multiplying by $Z^{(\le\Delta_{\max})}(t_j)$ preserves moment bounds.
\end{proof}

\begin{proposition}[Renormalized normal products via multiscale subtraction]
\label{prop:normal_products_multiscale}
Let $A,B\in\mathfrak A_{\mathrm{loc}}$ be two renormalized local fields obtained from the dyadic construction.
There exists a renormalized product (``normal product'') $A\diamond B\in\mathfrak A_{\mathrm{loc}}$ defined as follows:

\smallskip
\noindent
For each dyadic scale $j$, define the bounded approximants $A_j,B_j$ as the renormalized flowed fields at time $t_j$.
Set
\[
(A\diamond B)_j\ :=\ A_j B_j\ -\ \Pi_{<\Delta(A)+\Delta(B)}\bigl(A_j B_j\bigr),
\]
where $\Pi_{<D}$ denotes the projection (in $L^2$) onto the finite-dimensional span of renormalized fields of engineering dimension
$<D$ (in a fixed graded basis), with coefficients chosen to keep a fixed reference Gram matrix of two-point functions scale-independent.
Then $(A\diamond B)_j$ converges in $L^2$ as $j\to\infty$, and its limit is independent of the choice of dyadic sequence.
Moreover, if the supports of the smearing test functions of $A$ and $B$ are disjoint, then $A\diamond B$ coincides with the ordinary
product $AB$.
\end{proposition}

\begin{proof}
The key point is that $A_jB_j$ is itself a bounded local observable at flow time $t_j$ and therefore admits the same one-shell
decomposition as in Lemma~\ref{lem:one_shell_increment}, with coefficients controlled by curvature domination.
The only potentially non-summable part of the shell drift is the projection of $A_jB_j$ onto local operators of dimension
$\le \Delta(A)+\Delta(B)$.
By subtracting the projection onto strictly lower dimensions and fixing the remaining dimension-$\Delta(A)+\Delta(B)$ component by
normalization (exactly as in the scalar and finite-multiplet cases), the renormalized shell increment becomes $O(u_j^2)$ in $L^2$.
Summability follows from $\sum_j u_j^2<\infty$.
If the smearings are disjoint at physical distance $\gg \sqrt{t_j}$, locality of the flow implies no coincident-point divergence occurs,
so the projection term vanishes for $j$ large and the renormalized product equals the ordinary product.
\end{proof}

\begin{remark}[Associativity and higher products]
Higher normal products are defined recursively (or via the standard Zimmermann forest formula) by iterating the dyadic subtraction.
Because the subtraction is organized by the engineering-dimension filtration and uses orthogonal projections onto finite-dimensional
subspaces, one obtains an associative $*$--algebra structure on $\mathfrak A_{\mathrm{loc}}$ compatible with locality.
\end{remark}

\begin{proposition}[OPE-style remainder bound at short distance (finite truncation)]
\label{prop:ope_remainder_bound}
Fix $\Delta_{\max}$ and work inside the truncated algebra $\mathfrak A_{\mathrm{loc}}^{(\le\Delta_{\max})}$.
Let $A,B$ be two renormalized local fields of engineering dimensions $\Delta_A,\Delta_B\le\Delta_{\max}$.
Then there exists a finite set of renormalized fields $\{\mathcal O_k\}$ with dimensions $\Delta_k\le \Delta_A+\Delta_B$ and
coefficient distributions $C_{AB}^k\in \mathcal D'(\mathbb R^4)$ with \emph{singular support} contained in $\{0\}$ (allowing contact terms) such that for any test function $\phi$ supported in the unit ball and
any scale $r\in(0,1]$,
\[
\biggl\|\,A\bigl(\phi_r\bigr)\,B\bigl(\phi_r\bigr)
-\sum_{\Delta_k\le \Delta_A+\Delta_B} \bigl\langle C_{AB}^k,\phi_r\bigr\rangle\,\mathcal O_k(\phi_r)\,\biggr\|_{L^2}
\ \le\ C_{A,B,\phi}\, r^{\delta},
\]
for some $\delta>0$ depending only on the first omitted dimension.
Here $\phi_r(x):=r^{-4}\phi(x/r)$.
\end{proposition}

\begin{proof}
We follow the multiscale power counting argument.

Choose the dyadic scale $j$ so that $\sqrt{t_j}\asymp r$.
Write the flowed product $A_j(\phi_r)B_j(\phi_r)$ as a sum over shells using the finite-range decomposition
\eqref{eq:finite_range_decomp}.
At scales $<r$ (higher momenta) the two insertions are effectively coincident and produce local counterterms; at scales $>r$ the
insertions are separated and contribute only regular terms.
Projecting onto the local operator basis up to dimension $\Delta_A+\Delta_B$ yields the coefficients $C_{AB}^k$.
The first omitted dimension produces an $r^{\delta}$ gain by standard power counting.
Curvature domination gives the needed uniform $L^2$ control on each shell contribution; asymptotic freedom (Lemma~\ref{lem:dyadic_coupling_decay})
makes the logarithmic corrections summable.
\end{proof}



\paragraph*{Wilson coefficients: existence, derivative bounds, and scheme covariance (optional refinement)}
\label{subsubsec:wilson_coefficients}

This subsection upgrades Proposition~\ref{prop:ope_remainder_bound} by (i) defining the coefficient distributions
$C_{AB}^k$ as \emph{limits of dyadic coefficient functionals} extracted from the multiscale subtraction,
and (ii) recording quantitative regularity bounds away from the coincidence point.
Nothing here is needed for the lattice mass-gap or for the continuum existence statements on bounded macroscopic algebras; it is
only required if one wants a fully local field algebra equipped with quantitative OPE control.

\paragraph{Dyadic coefficient extraction.}
Fix a truncation level $\Delta_{\max}$ and a graded basis $\{\mathcal O_k\}$ of renormalized local fields with
$\Delta_k\le \Delta_{\max}$.
For a dyadic scale $j$ let $t_j:=2^{-2j}t_0$ and let $\mathcal O_{k,j}$ denote the bounded flowed approximant at time $t_j$
after applying the renormalization matrix $Z^{(\le \Delta_{\max})}(t_j)$ from
Theorem~\ref{thm:field_algebra_truncation}.
Let $A_j,B_j$ be the corresponding approximants for $A,B$.

Fix once and for all a reference test function $\phi\in C_c^\infty(\mathbb R^4)$ supported in the unit ball
and define $\phi_r(x)=r^{-4}\phi(x/r)$.
For $r\asymp 2^{-j}$ define the coefficient vector $c_{AB,j}(\phi)\in\mathbb R^{m(\Delta_{\max})}$ by the orthogonal-projection identity
\begin{equation}\label{eq:dyadic_ope_coeff_def}
\Pi^{(\le \Delta_A+\Delta_B)}_j\bigl(A_j(\phi_r)B_j(\phi_r)\bigr)
=\sum_{\Delta_k\le \Delta_A+\Delta_B} c_{AB,j}^k(\phi)\,\mathcal O_{k,j}(\phi_r),
\end{equation}
where $\Pi^{(\le D)}_j$ is the $L^2$-orthogonal projection onto the span of $\{\mathcal O_{\ell,j}(\phi_r):\Delta_\ell\le D\}$.
Because the Gram matrix of this basis is held scale-independent by the normalization defining $Z^{(\le \Delta_{\max})}(t_j)$,
the coefficients $c_{AB,j}^k(\phi)$ are well-defined and stable under small changes of $r$ within a fixed dyadic shell.

\paragraph{One-shell control of coefficient increments.}
Write $d_k:=\Delta_A+\Delta_B-\Delta_k\ge 0$ for the engineering dimension deficit.
For multiindices $\alpha$ we write $\partial^\alpha$ for derivatives in the relative-coordinate variable.

\begin{lemma}[Shell increment bound for dyadic OPE coefficients]\label{lem:ope_shell_increment}
Assume weak-coupling curvature control (Lemma~\ref{lem:weak_localization} and Proposition~\ref{prop:curvature_BL})
and dyadic coupling decay (Lemma~\ref{lem:dyadic_coupling_decay}).
Then for each fixed truncation $\Delta_{\max}$ and each fixed test function $\phi$ there exist constants
$C_{\phi,\Delta_{\max}}<\infty$ and $j_0$ such that for all $j\ge j_0$ and all $k$ with $\Delta_k\le \Delta_A+\Delta_B$,
\begin{equation}\label{eq:ope_shell_increment_bound}
|c_{AB,j+1}^k(\phi)-c_{AB,j}^k(\phi)|
\ \le\ C_{\phi,\Delta_{\max}}\; 2^{j d_k}\,\bigl(u_j+u_j^2\bigr)\ +\ C_{\phi,\Delta_{\max}}e^{-c_{\mathrm{loc}}\beta}.
\end{equation}
Moreover, if $\phi$ is replaced by $\partial^\alpha\phi$ then the same bound holds with an extra factor $2^{j|\alpha|}$.
\end{lemma}

\begin{proof}

Fix $r\asymp 2^{-j}$. The projected coefficient $c_{AB,j}^k(\phi)$ is a finite linear functional of mixed moments of
the local bounded observables $A_j(\phi_r)B_j(\phi_r)$ against the basis fields $\mathcal O_{k,j}(\phi_r)$.
Replacing scale $j$ by $j+1$ replaces the flow time $t_j$ by $t_{j+1}=t_j/4$ and inserts one additional ``shell''
in the finite-range decomposition \eqref{eq:finite_range_decomp}.
Exactly as in Lemma~\ref{lem:one_shell_increment}, the shell contribution is a finite sum of local curvature polynomials whose coefficients
obey the dimensional scaling $2^{jd_k}$ and (for derivatives) $2^{j|\alpha|}$ by power counting.
Curvature Brascamp--Lieb domination bounds the $L^2$ size of each shell polynomial by $C(u_j^{1/2})^{\deg}$, and the only potentially
non-summable part is the order-$u_j$ drift in the dimension-$\Delta_A+\Delta_B$ sector.
That drift is precisely absorbed into the renormalization/mixing matrices $Z^{(\le\Delta_{\max})}(t_j)$ defining the basis
$\{\mathcal O_{k,j}\}$, leaving a remainder of order $u_j^2$; we record the bound in the safe form $(u_j+u_j^2)$.
The tail outside $\Omega_{\varepsilon_0}$ is exponentially small in $\beta$ by Lemma~\ref{lem:weak_localization}.
\end{proof}

\paragraph{Existence and regularity of Wilson coefficients.}
The coefficients $c_{AB,j}^k(\phi)$ determine a distribution $C_{AB}^k$ by dyadic Littlewood--Paley reconstruction:
one defines the action of $C_{AB}^k$ on a general test function $\psi$ by decomposing $\psi=\sum_{j\ge 0}\psi_j$ into dyadic shells
of relative scale and setting $\langle C_{AB}^k,\psi_j\rangle$ equal to the coefficient extracted at the matching dyadic scale.
The shell increment bound \eqref{eq:ope_shell_increment_bound} implies convergence and yields explicit derivative bounds.

\begin{theorem}[Wilson coefficient limits and derivative bounds]\label{thm:wilson_coeff_limits}
Fix $\Delta_{\max}$ and let $A,B\in\mathfrak A_{\mathrm{loc}}^{(\le\Delta_{\max})}$ with dimensions $\Delta_A,\Delta_B$.
Assume Lemma~\ref{lem:ope_shell_increment}.
Then for each $k$ with $\Delta_k\le \Delta_A+\Delta_B$ there exists a distribution
$C_{AB}^k\in\mathcal D'(\mathbb R^4)$ with $\mathrm{sing\,supp}(C_{AB}^k)\subset\{0\}$ such that:

\begin{enumerate}[leftmargin=*]
\item (\emph{OPE remainder}) The bound of Proposition~\ref{prop:ope_remainder_bound} holds with these coefficients.
\item (\emph{Derivative control on scaled tests})
For every multiindex $\alpha$ there exist constants $C_{\alpha,A,B,k,\phi}$ and an integer $M\ge 0$ such that for all $r\in(0,1]$,
\begin{equation}\label{eq:wilson_scaled_derivative_bound}
\bigl|\langle C_{AB}^k,\partial^\alpha\phi_r\rangle\bigr|
\ \le\ C_{\alpha,A,B,k,\phi}\;
r^{-d_k-|\alpha|}\,\bigl(1+\log(1/r)\bigr)^{M}.
\end{equation}
\item (\emph{Smoothness away from coincidence})
There exists a smooth function $c_{AB}^k\in C^\infty(\mathbb R^4\setminus\{0\})$ such that
$C_{AB}^k|_{\mathbb R^4\setminus\{0\}}$ equals $c_{AB}^k(x)\,dx$, and for every multiindex $\alpha$,
\begin{equation}\label{eq:wilson_pointwise_derivative_bound}
|\partial^\alpha c_{AB}^k(x)|
\ \le\ C_{\alpha,A,B,k}\;
|x|^{-d_k-|\alpha|}\,\bigl(1+\log(1/|x|)\bigr)^{M},
\qquad 0<|x|\le \tfrac12.
\end{equation}
\end{enumerate}
\end{theorem}

\begin{proof}

Decompose a test function $\psi$ into dyadic annuli $\psi=\sum_j \psi_j$ with $\psi_j$ supported where $|x|\asymp 2^{-j}$.
Define $\langle C_{AB}^k,\psi_j\rangle$ by transporting $\psi_j$ to the unit scale and applying the coefficient functional
$c_{AB,j}^k(\cdot)$.
The shell increment estimate \eqref{eq:ope_shell_increment_bound} and dyadic coupling decay $u_j\lesssim (1+j)^{-1}$
imply the series defining $\langle C_{AB}^k,\psi\rangle$ converges absolutely, with at most a polynomial loss in $j$
giving the logarithmic factor in \eqref{eq:wilson_scaled_derivative_bound}.
Smoothness on $\mathbb R^4\setminus\{0\}$ follows because only finitely many shells interact with any compact set bounded away from $0$,
and each shell coefficient is given by a smooth kernel (finite-range heat-kernel coefficients).
Derivative bounds are obtained by repeating the preceding estimate for $\partial^\alpha\psi$ and tracking the extra factor $2^{j|\alpha|}$.
\end{proof}

\begin{corollary}[Scheme covariance of Wilson coefficients]\label{cor:wilson_scheme_covariance}
Let $\mathsf S,\mathsf S'$ be two admissible renormalization schemes in the sense of Definition~\ref{def:admissible_scheme},
with scheme-change matrices $R^{(\le\Delta_{\max})}$ from Theorem~\ref{thm:scheme_change_isomorphism}.
Write the (truncated) renormalized bases as $\{\mathcal O_i^{\mathsf S}\}$ and $\{\mathcal O_i^{\mathsf S'}\}$ and assume they are related by
\[
\mathcal O_i^{\mathsf S'}=\sum_{a} R_{ia}\,\mathcal O_a^{\mathsf S}.
\]
Then the Wilson coefficients transform by the standard change-of-basis rule: at the level of coefficient distributions,
\begin{equation}\label{eq:wilson_covariance_rule}
C^{k,\mathsf S'}_{ij}\;=\;\sum_{a,b,c} R_{ia}\,R_{jb}\,(R^{-1})_{ck}\; C^{c,\mathsf S}_{ab},
\end{equation}
where $C^{k}_{ij}$ denotes the coefficient in the OPE of $\mathcal O_i\cdot \mathcal O_j$ onto $\mathcal O_k$.
In particular, the singular orders and derivative bounds of Theorem~\ref{thm:wilson_coeff_limits} are invariant within the admissible class.
\end{corollary}


\subsection{KP majorant and admissibility}
\label{sec:kp_majorant}

\subsubsection{Connected-set counting constant \texorpdfstring{$\Ccount$}{C\_count}}

Let $\mathcal G_p$ be the plaquette adjacency graph: plaquettes adjacent if they share a link.
A plaquette has 4 links.
A link in $\mathbb Z^d$ belongs to $2(d-1)$ plaquettes, so a plaquette shares a link with at most
\begin{equation}\label{eq:Delta_p}
\Delta_p(d)\le 4\,(2(d-1)-1)=8d-12
\end{equation}
other plaquettes.
The standard tree-count bound implies
\begin{equation}\label{eq:connected_count}
\#\{\text{connected plaquette sets of size $n$ containing a root}\}\le (e\Delta_p)^{n-1}.
\end{equation}
We therefore take
\begin{equation}\label{eq:Ccount_def}
\Ccount:=\lceil e\Delta_p(d)\rceil.
\end{equation}

\subsubsection{Majorant function and ceiling}

Define the KP majorant
\begin{equation}\label{eq:q_def}
\qkp(f):=\Ccount\,e^{\theta}\,\Kzero\,\gb(f).
\end{equation}
Admissibility is $\qkp(f)<1$.
The maximal admissible activity is the solution of $\qkp(f)=1$:
\begin{equation}\label{eq:fadm_def}
 f_{\mathrm{adm}}
 =
 \frac{\Bigl(1+\frac{1}{\Ccount e^{\theta}\Kzero}\Bigr)^{1/\Nb}-1}{\Cint}.
\end{equation}

\begin{lemma}[Explicit inverse threshold for the KP majorant]\label{lem:qkp_inverse_explicit}
For any target $\bar q\in(0,1)$ define
\[
f_{\mathrm{KP}}(\bar q)
:=
\frac{\Bigl(1+\frac{\bar q}{\Ccount e^{\theta}\Kzero}\Bigr)^{1/\Nb}-1}{\Cint}.
\]
Then $f\le f_{\mathrm{KP}}(\bar q)$ implies $\qkp(f)\le \bar q$.
In particular, $\bar q=1$ recovers the maximal admissible activity $f_{\mathrm{adm}}=f_{\mathrm{KP}}(1)$.
\end{lemma}
\begin{proof}
By definition $\qkp(f)=\Ccount e^{\theta}\Kzero\,\gb(f)$ and $\gb(f)=(1+\Cint f)^{\Nb}-1$; invert this monotone map.
\end{proof}


\subsection{Projection remainder, drift, and the ratio audit}\label{sec:proj_remainder}

Stages 1--10 treated the following as explicit inputs:
\begin{enumerate}[leftmargin=*]
\item polymer coefficient bound $|J_P|\le \Ktwo^{|P|} f^{|P|}$;
\item activity Lipschitz bound $|f_1^{\mathrm{exact}}-f_1^{\square}|\le \Cact\,\|\Delta(1)\|_{\theta}$;
\item deterministic projected recursion $f_1^{\square}\le \Kact\, f^{\Nb}$.
\end{enumerate}

\subsubsection{Canonical choices \texorpdfstring{$\Kzero=\Krep=1$}{K0=Krep=1}, \texorpdfstring{$\Ktwo\approx 1$}{K2\ approx\ 1}, \texorpdfstring{$\Cact=1$}{Cact=1}}

\begin{lemma}[Contractive Haar averaging]\label{lem:K0_one}
Working with Haar \emph{probability} measure and $|\varphi_R|\le 1$, one may take
\[\boxed{\Kzero=1\qquad\text{and}\qquad \Krep(G)=1.}\]
\end{lemma}

\begin{lemma}[Raw coefficient majorant]\label{lem:K2_raw}
Expanding products of \eqref{eq:W_expand} and taking absolute values yields
\[
|J_P|\le f^{|P|}.
\]
So one may take $\boxed{\Ktwo=1}$ (raw).
For cumulant-safe bookkeeping one may take $\boxed{\Ktwo=e}$.
\end{lemma}

\begin{lemma}[Activity Lipschitz constant]\label{lem:Cact_one}
In the $\ell^1$ activity norm \eqref{eq:f_def}, coefficient extraction is $1$-Lipschitz.
Thus $\boxed{\Cact=1}$.
\end{lemma}

\subsubsection{Projection constant and geometric-series bound}

Assume
\begin{equation}\label{eq:K2_assumption}
|J_P|\le \Ktwo^{|P|}\,f^{|P|}.
\end{equation}
Define
\begin{equation}\label{eq:A_def}
A:=\Ccount\,\Ktwo\,e^{\theta}.
\end{equation}
Then for $Af<1$, the polymer remainder obeys the standard geometric-series bound
\begin{equation}\label{eq:Kproj_def}
\|\Delta(1)\|_{\theta}\le \Kproj(f)\,f^2,
\qquad
\Kproj(f):=\frac{\Ccount\,\Ktwo^2 e^{2\theta}}{1-Af}.
\end{equation}

\subsubsection{Stage-8 ratio audit}

Bernoulli gives $\gb(f)\ge \Nb\Cint f$ for $f\ge 0$.
Hence $\qkp(f)<1$ implies $\Ccount e^{\theta}\Kzero\Nb\Cint f<1$.
Multiplying by $\Ktwo/(\Nb\Cint\Kzero)$ yields
\begin{equation}\label{eq:R0_def}
R_0:=\frac{\Ktwo}{\Nb\Cint\Kzero}
\quad\Rightarrow\quad
\qkp(f)<1\ \Longrightarrow\ Af\le R_0.
\end{equation}
If $R_0<1$, define $\eta_0:=1-R_0>0$.
Then admissibility forces the uniform margin $1-Af\ge \eta_0$.

\subsubsection{Exact one-step contraction inside the KP basin}\label{sec:one_step_contraction}

Inside the KP admissibility window, the projected convolution-power decimation contracts extremely strongly.
We record an explicit ``once in, stay in'' one-step inequality for the \emph{exact} step.

\begin{lemma}[Exact step dominated by the projected step]\label{lem:exact_step_2u}
Assume the projected step satisfies $f_1^{\square}\le \Kact f^{\Nb}$ with $\Kact=1$ (true for the convolution-power map in the positive-type sector), and assume the drift bound
\[
|f_1^{\mathrm{exact}}-f_1^{\square}|\le \Cact\,\|\Delta(1)\|_{\theta},
\qquad
\|\Delta(1)\|_{\theta}\le \Kproj(u)\,u^2,
\quad u:=f_1^{\square}.
\]
If $u$ obeys $\Cact\,\Kproj(u)\,u\le 1$, then
\begin{equation}\label{eq:exact_step_2u}
 f_1^{\mathrm{exact}}\le 2u\le 2 f^{\Nb}.
\end{equation}
\end{lemma}

\begin{proof}
By the drift bound and the geometric-series estimate,
\[
f_1^{\mathrm{exact}}
\le
u+\Cact\,\Kproj(u)\,u^2
=
u\bigl(1+\Cact\,\Kproj(u)\,u\bigr)
\le 2u.
\]
The last inequality uses $\Cact\,\Kproj(u)\,u\le 1$.
Finally $u=f_1^{\square}\le f^{\Nb}$ when $\Kact=1$.
\end{proof}

\begin{remark}[Baseline sanity check]
For $(d,b,\theta)=(4,2,1)$ one has $\Nb=4$ and $\Ccount=55$.
At $f=f_{\mathrm{adm}}\approx 1.668\times 10^{-3}$, the projected value $u\le f^{4}\approx 7.7\times 10^{-12}$ is so small that $\Cact\Kproj(u)u\ll 1$ holds by many orders of magnitude.
Thus \eqref{eq:exact_step_2u} is automatic once inside the KP window.
\end{remark}


\subsection{Exponential clustering in the KP regime}

For finite plaquette sets $X,Y$, let $\dist(X,Y)$ be the distance in $\mathcal G_p$.
For bounded observables $A,B$ define the truncated correlation $\langle A;B\rangle:=\langle AB\rangle-\langle A\rangle\langle B\rangle$.

\begin{theorem}[Uniform exponential clustering]\label{thm:clustering}
Assume $Af<1$ and \eqref{eq:K2_assumption}.
Then for bounded observables supported on plaquette sets $X,Y$,
\begin{equation}\label{eq:cluster_bound}
|\langle A;B\rangle|
\le
2\,\|A\|_{\infty}\,\|B\|_{\infty}\,\frac{(Af)^{\dist(X,Y)}}{\Ccount\,(1-Af)}.
\end{equation}
\end{theorem}

\begin{corollary}[Decay rate in the admissible window]\label{cor:mu_from_R0}
If $R_0<1$ and $f$ is admissible, then $Af\le R_0$.
Hence $\mu:=-\log(Af)\ge -\log(R_0)=\log\bigl(\tfrac{\Nb\Cint\Kzero}{\Ktwo}\bigr)>0$.
\end{corollary}

\subsection{Reflection positivity and transfer matrix}

\begin{definition}[Reflection positivity]
Fix a time reflection $\Theta$.
A measure $\mu$ on link configurations is reflection positive if for all $F$ supported in positive times,
\begin{equation}\label{eq:RP_def}
\langle \Theta F\cdot F\rangle_{\mu}\ge 0.
\end{equation}
\end{definition}

\begin{lemma}[Positive type $\Rightarrow$ positive kernel]\label{lem:pos_kernel}
If $W$ is of positive type (Definition~\ref{def:positive_type} below), then
\[
K(g,h):=W(gh^{-1})
\]
defines a positive semidefinite kernel on $G$.
\end{lemma}

\begin{definition}[Positive type]\label{def:positive_type}
A central function $W\ge 0$ on $G$ is \emph{of positive type} if
\[
W(g)=\sum_R c_R\,\chi_R(g)
\qquad\text{with }c_R\ge 0.
\]
\end{definition}

\begin{definition}[Conditionally negative definite (cnd)]
A central function $\psi:G\to\mathbb R$ is \emph{conditionally negative definite} if for all $n\in\mathbb N$, all $g_1,\dots,g_n\in G$ and all $c_1,\dots,c_n\in\mathbb C$ with $\sum_{i=1}^n c_i=0$,
\[
\sum_{i,j=1}^n \overline c_i\,c_j\,\psi(g_i^{-1}g_j)\le 0.
\]
\end{definition}

\begin{lemma}[Normalized trace defect is cnd]\label{lem:trace_defect_cnd}
Let $\pi:G\to U(d)$ be a unitary representation and set
\[
\phi_\pi(g):=\frac{1}{d}\Re\mathrm{Tr}\,\pi(g),
\qquad
A_\pi(g):=1-\phi_\pi(g).
\]
Then $A_\pi$ is conditionally negative definite.
\end{lemma}
\begin{proof}
Fix an orthonormal basis $(e_k)_{k=1}^d$ of $\mathbb C^d$.
For each $k$ define
\[
\psi_k(g):=\frac12\|e_k-\pi(g)e_k\|^2 = 1-\Re\langle e_k,\pi(g)e_k\rangle.
\]
The kernels $(g,h)\mapsto \psi_k(g^{-1}h)$ are conditionally negative definite because $\psi_k$ is a squared Hilbert-space distance under a unitary representation.
Since $A_\pi=\frac1d\sum_{k=1}^d \psi_k$, it is a convex combination of cnd functions, hence cnd.
\end{proof}

\begin{lemma}[Schoenberg]\label{lem:Schoenberg_cnd}
If $\psi$ is cnd and $t\ge0$, then the function $g\mapsto e^{-t\psi(g)}$ is positive definite.
If in addition $\psi$ is central, then $e^{-t\psi}$ is of positive type (Definition~\ref{def:positive_type}).
\end{lemma}
\begin{proof}

This is the classical theorem of Schoenberg: the kernels $K_t(g,h):=\exp(-t\psi(g^{-1}h))$ are positive semidefinite for cnd $\psi$.
For central $e^{-t\psi}$, Peter--Weyl then gives a character expansion with nonnegative coefficients.
\end{proof}

\begin{corollary}[Wilson-class weights are positive type]\label{cor:Wilson_positive_type}
Let $A_\pi$ be as in Lemma~\ref{lem:trace_defect_cnd}. Then for every $\beta\ge0$,
\[
\widehat W_{\beta,\pi}(g):=\exp(-\beta A_\pi(g))
\]
is a central positive-type function. Moreover, for every $\alpha\in(0,1]$ the power $\widehat W_{\beta,\pi}^\alpha=\exp(-\alpha\beta A_\pi)$ is also positive type.
\end{corollary}
\begin{proof}
By Lemma~\ref{lem:trace_defect_cnd}, $A_\pi$ is cnd; by Lemma~\ref{lem:Schoenberg_cnd}, $\exp(-\beta A_\pi)$ is positive definite.
Since it is central, its character expansion has nonnegative coefficients, i.e.\ it is positive type. The identical character-expansion calculation applies to $\alpha\beta$.
\end{proof}

\begin{corollary}[Reflection positivity for the $\alpha$--reference family]\label{cor:alpha_reference_RP}
The anisotropic $\alpha$--reference measures $\nu_L^{(\alpha)}$ of \eqref{eq:alpha_reference_block}
are reflection positive with respect to reflections through coordinate hyperplanes orthogonal to any axis, for every $\alpha\in(0,1]$.
\end{corollary}
\begin{proof}
The in-plane plaquette weight is $\widehat W_{\beta,\pi}$ and the out-of-plane weight is $\widehat W_{\beta,\pi}^\alpha$, both positive type by Corollary~\ref{cor:Wilson_positive_type}.
Apply Proposition~\ref{prop:pos_type_RP} plaquette-wise.
\end{proof}



\begin{proposition}[Plaquette-wise positive type $\Rightarrow$ reflection positivity]\label{prop:pos_type_RP}
Consider a lattice gauge theory with density
\[
d\mu(U)\propto \Bigl(\prod_{p} W_p(U_p)\Bigr)\,dU,
\]
where each $W_p$ is a central function of positive type (Definition~\ref{def:positive_type}) and the family is reflection-covariant in the sense
$W_{\theta p}=W_p$ for the reflection $\theta$ under consideration.
Then $\mu$ is reflection positive with respect to $\theta$ in the sense of \eqref{eq:RP_def}.
Consequently OS reconstruction yields a Hilbert space $\mathcal H$, a vacuum $\Omega$, and a positive self-adjoint transfer matrix $T$.
\end{proposition}
\begin{proof}

Expand each $W_p$ in characters with nonnegative coefficients. Plaquettes crossing the reflection hyperplane produce positive kernels by Lemma~\ref{lem:pos_kernel};
after integrating the reflected half-lattice variables the expectation $\langle \Theta F\cdot F\rangle$ becomes a sum of squares with nonnegative coefficients.
\end{proof}


\subsection{Uniform decay \texorpdfstring{$\Rightarrow$}{implies} spectral gap}

\begin{theorem}[Uniform decay of all bounded local observables implies a full fixed-lattice gap]\label{thm:F_decay_implies_gap}
Assume reflection positivity (so $T$ and $\Omega$ exist) and $\Omega$ is cyclic for bounded gauge-invariant local observables.
If there exists $m>0$ such that for every bounded gauge-invariant local observable $\mathcal O$ with $\langle\mathcal O\rangle=0$,
\begin{equation}\label{eq:exp_decay_assumption}
0\le \langle \mathcal O(t)\mathcal O(0)\rangle\le C_{\mathcal O}e^{-mt}\qquad (t\in\mathbb N),
\end{equation}
then $H:=-\log T$ has a vacuum gap at least $m$.
\end{theorem}

\begin{proof}
For every centered bounded gauge-invariant local observable $\mathcal O$, OS reconstruction gives
\[
\langle \mathcal O(t)\mathcal O(0)\rangle
=
\langle [\mathcal O],e^{-tH}[\mathcal O]\rangle
=\int_{[0,\infty)} e^{-t\lambda}\,d\mu_{\mathcal O}(\lambda),
\]
where $\mu_{\mathcal O}$ is the spectral measure of $H$ associated to $[\mathcal O]$.
If $\mu_{\mathcal O}([0,m-\varepsilon])>0$ for some $\varepsilon\in(0,m)$, then the spectral theorem yields the lower bound
\[
\langle [\mathcal O],e^{-tH}[\mathcal O]\rangle\ge e^{-t(m-\varepsilon)}\mu_{\mathcal O}([0,m-\varepsilon]),
\]
which contradicts \eqref{eq:exp_decay_assumption} as $t\to\infty$.
Hence $\mu_{\mathcal O}([0,m))=0$ for every centered local class.
Because $\Omega$ is cyclic for bounded gauge-invariant local observables, centered local classes are dense in $\Omega^\perp$.
Therefore the spectral projector $P_{(0,m)}$ vanishes on a dense subspace of $\Omega^\perp$ and hence on all of $\Omega^\perp$.
Thus
\[
\mathrm{Spec}(H)\subset\{0\}\cup[m,\infty),
\]
as claimed.
\end{proof}

\begin{proposition}[Fixed-lattice positive-time local algebra is cyclic]\label{prop:lattice_OS_cyclicity_local_algebra}
Let $\mu_\beta$ be a translation-invariant reflection-positive infinite-volume Wilson Gibbs state and let $(\mathcal H_\beta,\Omega_\beta,H_\beta)$ be its OS reconstruction on the bounded gauge-invariant positive-time cylinder algebra $\mathcal A_{+,\beta}^{\mathrm{loc}}$.
Then $\Omega_\beta$ is cyclic for bounded gauge-invariant local cylinder observables.
\end{proposition}

\begin{proof}
By construction, $\mathcal H_\beta$ is the completion of $\mathcal A_{+,\beta}^{\mathrm{loc}}/\mathcal N_\beta$, where $\mathcal N_\beta$ is the OS null ideal.
Hence the classes of bounded gauge-invariant local positive-time cylinder observables are dense in $\mathcal H_\beta$, and $\Omega_\beta=[1]$ is cyclic for that local algebra.
\end{proof}

\begin{corollary}[Gap bound in the admissible window]\label{cor:gap_from_R0}
Assume $W$ is positive type and $R_0<1$.
If $f$ is admissible then by Corollary~\ref{cor:mu_from_R0} and Theorem~\ref{thm:F_decay_implies_gap},
\begin{equation}\label{eq:gap_bound_R0}
\boxed{\Delta_{\mathrm{lat}}(\beta)\ge -\log(R_0)=\log\Bigl(\frac{\Nb\Cint\Kzero}{\Ktwo}\Bigr)>0.}
\end{equation}
\end{corollary}

\subsection{Local non-toy majorant: exact tile integration vs convolution power}\label{sec:majorant}

This is the hard ``non-toy'' step on the current critical path: a rigorous one-block comparison between an \emph{actual} $b\times b$ block integration in the full $d$-dimensional theory and an explicit positivity-preserving convolution/power map.

\subsubsection{Sup normalization}

\begin{lemma}[Positive definite functions are maximized at the identity]\label{lem:pd_max_identity}
If $W$ is positive definite then $|W(g)|\le W(e)$ for all $g$.
If also $W\ge 0$ then $0\le W(g)\le W(e)$.
\end{lemma}

\begin{definition}[Sup-normalized weight]\label{def:sup_norm}
Assume $W\ge 0$ is positive type and $W(e)>0$.
Define $\widehat W(g):=W(g)/W(e)$.
Then $\widehat W$ is positive type and $0\le \widehat W\le 1$.
\end{definition}

\subsubsection{Tile factors}

Fix a $b\times b$ tile $T$ of plaquettes in some $(\mu,\nu)$ plane.
Let $U_{\partial T}$ be boundary links and $U_{\mathrm{int}(T)}$ internal links.
Define the \emph{full} internal block factor
\begin{equation}\label{eq:block_factor_full}
\mathcal Z_T^{\mathrm{full}}(U_{\partial T})
:=
\int
\Bigl(\prod_{p\supset \ell,\ \ell\in\mathrm{int}(T)} \widehat W(U_p)\Bigr)
\,dU_{\mathrm{int}(T)},
\end{equation}
where the product runs over all plaquettes containing at least one internal link.
Define the 2D-reduced factor (drop all out-of-plane plaquettes)
\begin{equation}\label{eq:block_factor_2d}
\mathcal Z_T^{(2D)}(U_{\partial T})
:=
\int
\Bigl(\prod_{p\subset T} \widehat W(U_p)\Bigr)
\,dU_{\mathrm{int}(T)}.
\end{equation}

\begin{lemma}[Dropping out-of-plane plaquettes majorizes pointwise]\label{lem:drop_majorant}
For every boundary configuration $U_{\partial T}$,
\[
\mathcal Z_T^{\mathrm{full}}(U_{\partial T})\le \mathcal Z_T^{(2D)}(U_{\partial T}).
\]
\end{lemma}

\subsubsection{2D tile factor equals a convolution power}

For integrable $f,h$ on $G$, define convolution $(f*h)(x):=\int f(xy^{-1})h(y)\,dy$ and $f^{*n}$ the $n$-fold convolution power.

\begin{lemma}[2D tile factor is a convolution power]\label{lem:tile_is_convolution}
Assume $\widehat W$ is central and positive type so
\[
\widehat W(g)=\sum_R d_R\,\widehat F_R\,\chi_R(g),\qquad \widehat F_R\ge 0.
\]
Let $\mathrm{Hol}_{\partial T}(U_{\partial T})$ be the boundary holonomy.
Then
\[
\mathcal Z_T^{(2D)}(U_{\partial T})=\widehat W^{*\Nb}\!\bigl(\mathrm{Hol}_{\partial T}(U_{\partial T})\bigr).
\]
\end{lemma}

\begin{theorem}[Tile majorant by convolution/power map]\label{thm:block_majorant_convolution}
Let $W$ be nonnegative, central, and positive type, and let $\widehat W=W/W(e)$.
Then for every tile $T$,
\[
\mathcal Z_T^{\mathrm{full}}(U_{\partial T})
\le
\widehat W^{*\Nb}\!\bigl(\mathrm{Hol}_{\partial T}(U_{\partial T})\bigr).
\]
\end{theorem}
\begin{proof}
Lemma~\ref{lem:drop_majorant} and Lemma~\ref{lem:tile_is_convolution}.
\end{proof}

\subsection{Global positivity-preserving coarse object and the \texorpdfstring{$\gb$}{g\_b} window}\label{sec:global}

\subsubsection{A global majorant partition function (planar reduction)}

Let $\Lambda\subset\mathbb Z^d$ be finite.
Write
\[
Z_{\Lambda}(W)=W(e)^{|P(\Lambda)|}\,Z_{\Lambda}(\widehat W),
\qquad
Z_{\Lambda}(\widehat W)=\int\prod_{p\in P(\Lambda)} \widehat W(U_p)\,dU.
\]
Fix an orientation $(1,2)$ and let $P_{12}(\Lambda)$ be the set of $(1,2)$ plaquettes.
Since $0\le \widehat W\le 1$,
\begin{equation}\label{eq:drop_to_planes}
Z_{\Lambda}(\widehat W)
\le
\int\prod_{p\in P_{12}(\Lambda)} \widehat W(U_p)\,dU.
\end{equation}
The RHS factorizes into 2D slices over fixed transverse coordinates.
Exact $b\times b$ blocking on each slice yields the coarse 2D model with plaquette weight $\widehat W^{*\Nb}$.
Thus the explicit positivity-preserving majorant partition function
\begin{equation}\label{eq:Z_majorant_def}
Z^{\mathrm{maj}}_{\Lambda,b}(W)
:=W(e)^{|P(\Lambda)|}\,\prod_{x_{\perp}} Z^{(2D)}_{\Lambda_b}\!\bigl(\widehat W^{*\Nb}\bigr)
\end{equation}
satisfies
\begin{equation}\label{eq:Z_majorant_ineq}
Z_{\Lambda}(W)\le Z^{\mathrm{maj}}_{\Lambda,b}(W).
\end{equation}

\subsubsection{Coefficient power map and activity contraction}

If $W$ is central positive type in the canonical form \eqref{eq:W_expand} with $a_R\ge 0$, then convolution powers give
\begin{equation}\label{eq:a_power_map}
 a_R^{(b)}=\frac{a_R^{\Nb}}{d_R^{2\Nb-2}}\qquad(R\ne\mathbf 1).
\end{equation}
Hence
\begin{equation}\label{eq:fb_le_fNb}
 f^{(b)}\le f^{\Nb}\le (1+f)^{\Nb}-1=\gb(f).
\end{equation}
This is the precise link between the convolution-power coarse object and the KP admissibility envelope $f\mapsto \gb(f)$.

\subsubsection{A full-dimensional assembly tool: generalized H\"older/Finner}

\begin{lemma}[Generalized H\"older (Finner) inequality]\label{lem:finner}
Let $(\Omega,\mu)$ be a product probability space with coordinates $(x_v)_{v\in V}$.
Let $\{F_i\}_{i\in I}$ be nonnegative functions where $F_i$ depends only on the coordinate subset $S_i\subset V$.
If exponents $p_i\in[1,\infty]$ satisfy, for each coordinate $v\in V$,
\begin{equation}\label{eq:finner_condition}
\sum_{i:\ v\in S_i}\frac{1}{p_i}\le 1,
\end{equation}
then
\begin{equation}\label{eq:finner_ineq}
\int_\Omega \prod_{i\in I} F_i \,d\mu
\le
\prod_{i\in I}\|F_i\|_{L^{p_i}(\mu)}.
\end{equation}
\end{lemma}

\begin{remark}[Why this matters]
Theorem~\ref{thm:block_majorant_convolution} is \emph{local}.
Lemma~\ref{lem:finner} is the clean analytic device to glue such local majorants into a \emph{global} $d$-dimensional positivity-preserving decimation, with a single geometric price: the maximum overlap number $m$ (how many chosen blocks share a link variable).
Choosing a block cover and optimizing $m(d,b)$ is the next concrete step for a Tier-3 entrance attempt.
\end{remark}

\subsubsection{A symmetric all-orientation majorant via Finner}\label{sec:symm_majorant}

The planar reduction \eqref{eq:drop_to_planes} is extremely lossy (it discards most plaquettes).
A less wasteful alternative uses Lemma~\ref{lem:finner} across all plaquette orientations.

For each orientation $(\mu,\nu)$, set
\[
F_{\mu\nu}(U):=\prod_{p\in P_{\mu\nu}(\Lambda)} \widehat W(U_p),
\qquad
Z_{\mu\nu}:=\int F_{\mu\nu}(U)\,dU.
\]
Then
\[
Z_\Lambda(\widehat W)=\int \prod_{1\le \mu<\nu\le d} F_{\mu\nu}(U)\,dU.
\]
Fix $p_{\mu\nu}=d-1$ for all $(\mu,\nu)$.
Each link variable in direction $\mu$ appears only in the factors $F_{\mu\nu}$ with $\nu\ne \mu$, i.e.\ in exactly $(d-1)$ such factors.
Hence the Finner condition \eqref{eq:finner_condition} holds with equality, and Lemma~\ref{lem:finner} gives
\begin{equation}\label{eq:symm_finner_step}
Z_\Lambda(\widehat W)\le \prod_{\mu<\nu}\|F_{\mu\nu}\|_{L^{d-1}}.
\end{equation}
Since $0\le \widehat W\le 1$, we have $F_{\mu\nu}^{\,d-1}\le F_{\mu\nu}$, hence
\[
\|F_{\mu\nu}\|_{L^{d-1}}
=\Bigl(\int F_{\mu\nu}^{\,d-1}\Bigr)^{1/(d-1)}
\le Z_{\mu\nu}^{1/(d-1)}.
\]
Therefore
\begin{equation}\label{eq:symm_majorant}
Z_\Lambda(\widehat W)
\le
\prod_{\mu<\nu} Z_{\mu\nu}^{1/(d-1)}.
\end{equation}
Each $Z_{\mu\nu}$ factorizes into 2D slices and exact $b\times b$ blocking on each slice yields the coarse 2D model with plaquette weight $\widehat W^{*\Nb}$ (as in \eqref{eq:Z_majorant_def}).
Thus \eqref{eq:symm_majorant} gives a global positivity-preserving majorant using \emph{all} orientations, at the cost of fractional exponents.

\subsubsection{Explicit \texorpdfstring{$2\times 2$}{2x2} tiling and overlap constant \texorpdfstring{$m(4,2)=6$}{m(4,2)=6}}\label{sec:overlap_42}

For $b=2$, tile each oriented plane $(\mu,\nu)$ by disjoint $2\times 2$ plaquette blocks whose lower-left corners have even coordinates in directions $\mu,\nu$.
Let $\mathcal T_{\mu\nu}$ be the family of such tiles in the $(\mu,\nu)$ planes, and let $\mathcal T:=\bigcup_{\mu<\nu}\mathcal T_{\mu\nu}$.

\begin{lemma}[Overlap constant for the $2\times 2$ tiling in $d=4$]\label{lem:m42}
In $d=4$ with the above disjoint $2\times 2$ tiling, each fine link variable belongs to the support of at most
\[
m(4,2)=6
\]
tile factors in $\{\,\mathcal Z_T^{\mathrm{full}}\,:\,T\in\mathcal T\}$.
\end{lemma}

\begin{proof}
Fix a link $\ell$ oriented in direction $\mu$.
It lies in plaquettes only in planes $(\mu,\nu)$ with $\nu\neq \mu$ (there are $d-1=3$ such planes in $d=4$).
For each such plane $(\mu,\nu)$, the link $\ell$ is contained in exactly two plaquettes of that plane (one on each side of $\ell$), and the $2\times 2$ tiling is disjoint within the plane, so those two plaquettes lie in at most two tiles of $\mathcal T_{\mu\nu}$.
Hence $\ell$ meets at most $2(d-1)=6$ tiles in total.
\end{proof}

\begin{remark}
The overlap constant $m(d,b)$ is the only geometry that enters when gluing local tile bounds via Lemma~\ref{lem:finner}: taking $p_T=m(d,b)$ for all tile factors makes the Finner condition \eqref{eq:finner_condition} hold.
\end{remark}


\subsection{Invariant observable algebra and gauge-covariant blocking}\label{sec:inv}

Let $\mathcal U_\Lambda:=G^{E(\Lambda)}$ and let gauge transformations act by $(g\cdot U)_{x\to y}=g(x)U_{x\to y}g(y)^{-1}$.
Let $\mathfrak A_\Lambda$ be the $*$-algebra of bounded gauge-invariant cylinder functions on $\mathcal U_\Lambda$.

\begin{proposition}[Wilson loop generators]\label{prop:loop_generators}
$\mathfrak A_\Lambda$ is generated (as a $*$-algebra) by Wilson loop characters $W_{R,\gamma}(U)=\chi_R(U(\gamma))$.
\end{proposition}

\begin{proposition}[Blocking pullback preserves the invariant algebra]\label{prop:pullback_invariant_algebra}
Let $\pi:\mathcal U_\Lambda\to\mathcal U_{\Lambda_b}$ be a gauge-covariant blocking map built from parallel transport (coarse links are products of fine links along fixed paths).
Then $\pi^*: \mathfrak A_{\Lambda_b}\to\mathfrak A_\Lambda$ is a $*$-algebra homomorphism.
In particular, the pullback of a coarse Wilson loop is a gauge-invariant function of fine holonomies and hence lies in the loop-generated algebra.
\end{proposition}

\begin{remark}[Uniqueness mechanism]
The entrance/non-collapse and UV matching hypotheses imply uniqueness of macroscopic laws on $\mathcal O_{\ell}$ by KP/Dobrushin contraction;
see Theorem~\ref{thm:uniqueness_on_Oell}.
\end{remark}

\begin{remark}[Structural ``parameter collapse'']
Any coupling parameter that changes only the presentation of variables but leaves the induced functional on $\mathfrak A_\Lambda$ unchanged is redundant.
This is the natural lattice formalization of a \emph{redundant reparameterization}: different parameter choices that induce the same functional on $\mathfrak A_\Lambda$ are identified and carry no physical content.
\end{remark}

\subsection{A first quantitative remainder variable (KP-basin stability)}\label{sec:remainder}

To compare exact block integration to the convolution-power local reference, define the defect for the sup-normalized weight $\widehat W$:
\[
\delta(g):=1-\widehat W(g)\in[0,1].
\]
If $W$ is positive type in the canonical form \eqref{eq:W_expand} with $a_R\ge 0$, then $W(e)=1+f$ and
\begin{equation}\label{eq:defect_bound}
\sup_{g\in G}\delta(g)\le \frac{2f}{1+f}\le 2f.
\end{equation}

For a $b\times b$ tile $T$, let $P_{\perp}(T)$ be the set of out-of-plane plaquettes touching internal links.
Then
\begin{equation}\label{eq:Pperp_count}
|P_{\perp}(T)|\le 4(d-2)b(b-1)
\end{equation}
(because $T$ has $2b(b-1)$ internal links and each belongs to $2(d-2)$ out-of-plane plaquettes).
A crude bound gives
\begin{equation}\label{eq:tile_polymer_bound}
\Bigl|1-\prod_{p\in P_{\perp}(T)} \widehat W(U_p)\Bigr|
\le |P_{\perp}(T)|\,\sup\delta
\le 8(d-2)b(b-1)\,f\qquad(f\le 1).
\end{equation}
This is the first explicit input for bounding the polymer remainder when upgrading the local majorant to a full KP-norm comparison.


\subsection{From partition-function majorant to measure-level control of coarse observables}\label{sec:measure_level}

The local non-toy majorant (Section~\ref{sec:majorant}) and the global partition-function majorant (Section~\ref{sec:global}) support an alternative measure-level domination approach for Tier-3 comparisons of coarse observables; we record it for completeness but do not use it in the primary proof.
To compare blocked expectations, one needs either
\begin{enumerate}[leftmargin=*]
\item a genuine \emph{one-sided} domination of the blocked \emph{measure} by a positivity-preserving coarse measure (rare and very strong), or
\item a representation of the exact blocked measure as a \emph{local} coarse reference measure times a \emph{KP-small polymer correction}, yielding quantitative stability bounds for local observables.
\end{enumerate}
The second approach is the natural ``KP-small correction'' upgrade and is what we formalize here.

\subsubsection{Local coarse reference measure}

Let $\Lambda_b$ be the blocked lattice (spacing $b$).
Let $U^b\in G^{E(\Lambda_b)}$ denote coarse link variables and $U^b_p$ the coarse plaquette holonomy.
Define the local (positivity-preserving) coarse plaquette weight
\begin{equation}\label{eq:Wb_def}
W_b(g):=\widehat W^{*\Nb}(g),
\end{equation}
and the corresponding local coarse Gibbs measure
\begin{equation}\label{eq:mu_loc_def}
d\mu^{\mathrm{loc}}_{\Lambda_b}(U^b)
:=
\frac{1}{Z^{\mathrm{loc}}_{\Lambda_b}}
\Bigl(\prod_{p\in P(\Lambda_b)} W_b(U^b_p)\Bigr)\,dU^b,
\end{equation}
with $dU^b$ Haar probability measure on coarse links.
If $\widehat W$ is positive type, then so is $W_b$, hence $\mu^{\mathrm{loc}}_{\Lambda_b}$ is reflection positive (Proposition~\ref{prop:pos_type_RP}).

\subsubsection{Exact blocked measure and polymer remainder ansatz}

Let $\pi:\mathcal U_\Lambda\to\mathcal U_{\Lambda_b}$ be a gauge-covariant blocking map (Section~\ref{sec:inv}).
Let $\mu_\Lambda$ be the fine-lattice gauge measure with plaquette weight $\widehat W$.
Define the induced (pushforward) coarse measure
\begin{equation}
\mu^{\mathrm{ex}}_{\Lambda_b}:=\pi_{\#}\mu_\Lambda.
\end{equation}
In general $\mu^{\mathrm{ex}}_{\Lambda_b}$ is \emph{nonlocal} on $\Lambda_b$.
We package that nonlocality as a polymer remainder around the local reference \eqref{eq:mu_loc_def}.

\begin{definition}[Polymer remainder (coarse)]\label{def:poly_remainder}
We say that $\mu^{\mathrm{ex}}_{\Lambda_b}$ admits a polymer remainder representation around $\mu^{\mathrm{loc}}_{\Lambda_b}$ if there exist bounded functions $\{R_P\}$ indexed by finite connected plaquette sets $P\subset P(\Lambda_b)$ such that
\begin{equation}\label{eq:mu_ex_as_loc_times_poly}
d\mu^{\mathrm{ex}}_{\Lambda_b}
=
\frac{1}{Z_{\mathrm{rem}}}\,
\Bigl(\prod_{p\in P(\Lambda_b)} W_b(U^b_p)\Bigr)
\Bigl(\prod_{P\in\mathcal P(\Lambda_b)}(1+R_P(U^b_P))\Bigr)
\,dU^b,
\end{equation}
where $\mathcal P(\Lambda_b)$ is a chosen polymer family (e.g.\ all connected plaquette sets) and $Z_{\mathrm{rem}}$ is the normalizing constant.
Define the polymer activity bound
\begin{equation}\label{eq:rho_def}
\rho(P):=\sup_{U^b_P}|R_P(U^b_P)|.
\end{equation}
\end{definition}

\begin{remark}[What is non-toy here]
The existence of \eqref{eq:mu_ex_as_loc_times_poly} is a standard consequence of rewriting a nonlocal effective action as a polymer gas.
The hard content is to control $\rho(P)$ in a KP norm.
The defect bounds in Section~\ref{sec:remainder} are precisely the first quantitative handle needed for that control.
\end{remark}

\subsubsection{KP majorant for the remainder and observable stability}

Given $\rho(P)$, define a KP majorant for the remainder
\begin{equation}\label{eq:qrem_def}
\qkp_{\mathrm{rem}}
:=
\sup_{p_0\in P(\Lambda_b)}
\sum_{P\ni p_0}\rho(P)\,e^{\theta|P|}.
\end{equation}
For a finite plaquette set $S\subset P(\Lambda_b)$ define the localized majorant
\begin{equation}\label{eq:qrem_S_def}
\qkp_{\mathrm{rem}}(S)
:=
\sup_{p_0\in S}
\sum_{P\ni p_0}\rho(P)\,e^{\theta|P|}.
\end{equation}

\begin{theorem}[KP-small remainder controls coarse observables]\label{thm:obs_stability}
Assume $\qkp_{\mathrm{rem}}<1$ in \eqref{eq:qrem_def} and that \eqref{eq:mu_ex_as_loc_times_poly} holds.
Let $\mathcal O$ be any bounded gauge-invariant observable on $\Lambda_b$ supported on a finite plaquette set $S$.
Then
\begin{equation}\label{eq:obs_stability_bound}
\bigl|\langle \mathcal O\rangle_{\mu^{\mathrm{ex}}_{\Lambda_b}}-\langle \mathcal O\rangle_{\mu^{\mathrm{loc}}_{\Lambda_b}}\bigr|
\le
2\,\|\mathcal O\|_{\infty}\,
\frac{\qkp_{\mathrm{rem}}(S)}{1-\qkp_{\mathrm{rem}}}.
\end{equation}
In particular, in a translation-invariant setting one has $\qkp_{\mathrm{rem}}(S)\le |S|\,\qkp_{\mathrm{rem}}$, yielding
\[
\bigl|\langle \mathcal O\rangle_{\mathrm{ex}}-\langle \mathcal O\rangle_{\mathrm{loc}}\bigr|
\le
2\,\|\mathcal O\|_{\infty}\,
\frac{|S|\,\qkp_{\mathrm{rem}}}{1-\qkp_{\mathrm{rem}}}.
\]
\end{theorem}

\begin{proof}

Expand both the numerator and denominator of $\langle\mathcal O\rangle_{\mu^{\mathrm{ex}}}$ using the polymer product in \eqref{eq:mu_ex_as_loc_times_poly}.
By inclusion--exclusion/Mayer expansion, contributions to the difference from $\langle\mathcal O\rangle_{\mu^{\mathrm{loc}}}$ come from polymer clusters that intersect $S$.
The KP condition $\qkp_{\mathrm{rem}}<1$ implies absolute convergence and bounds the sum of connected clusters touching $S$ by a geometric series with ratio $\qkp_{\mathrm{rem}}$.
This yields \eqref{eq:obs_stability_bound}.
\end{proof}

\subsubsection{A crude universal bound from the defect \texorpdfstring{$\delta$}{delta}}

The out-of-plane plaquette sector enters through factors $\widehat W(U_p)=1-\delta(U_p)$ with $\delta\in[0,2f]$ (Equation~\eqref{eq:defect_bound}).
Treating each such defect insertion as an elementary polymer and bounding by $\sup\delta\le 2f$ yields the crude activity estimate
\begin{equation}\label{eq:rho_crude}
\rho(P)\le (2f)^{|P|},
\end{equation}
and hence, using the connected-set count \eqref{eq:connected_count},
\begin{equation}\label{eq:qrem_crude}
\qkp_{\mathrm{rem}}
\le
\sum_{n\ge 1} \Ccount^{n-1}\,(2f e^{\theta})^{n}
=
\frac{2f e^{\theta}}{1-2f e^{\theta}\Ccount}.
\end{equation}
In particular, a sufficient condition for $\qkp_{\mathrm{rem}}<1$ is
\begin{equation}\label{eq:qrem_simple_sufficient}
2f e^{\theta}\Ccount<\tfrac12.
\end{equation}
This bound is intentionally conservative.


\paragraph*{Eliminating the one-plaquette sector: a quadratic remainder majorant}

If the local reference $\mu^{\mathrm{loc}}_{\Lambda_b}$ is chosen so that all \emph{single-plaquette} (one-body) terms of the exact blocked weight are absorbed into $W_b$ (equivalently, we take $\rho(P)=0$ for $|P|=1$ in the polymerization of the Radon--Nikod\'ym derivative), then the KP majorant starts at $n=2$ and becomes manifestly quadratic:
\begin{equation}\label{eq:qrem_ge2}
\qkp_{\mathrm{rem}}
\le
\sum_{n\ge 2}\Ccount^{\,n-1}\,(2f e^{\theta})^{n}
=
\frac{\Ccount\,(2f e^{\theta})^{2}}{1-\Ccount\,2f e^{\theta}}.
\end{equation}
In particular, $\qkp_{\mathrm{rem}}<1$ follows from $\Ccount\,2f e^{\theta}<\tfrac12$.

\paragraph*{Flux/foam admissibility and a dimension-suppression refinement}

Write the positive-type expansion in normalized characters $\varphi_R=\chi_R/d_R$:
\[
\widehat W(g)=1+\sum_{R\ne \mathbf 1} a_R\,\varphi_R(g),\qquad a_R\ge 0,\qquad f=\sum_{R\ne\mathbf 1} a_R.
\]
Let $d_{\min}:=\min_{R\ne\mathbf 1} d_R$.
Haar orthogonality of matrix coefficients implies, schematically,
\[
\int_G \varphi_R(g h)\,\varphi_R(g^{-1}k)\,dg=\frac{1}{d_R^2}\,\varphi_R(hk)
\]
(and vanishes unless representation labels match), so each internal link integration that glues two nontrivial plaquette insertions produces at most a factor $1/d_{\min}^2$ in absolute value.

For a connected plaquette polymer $P$ with $|P|=n\ge 2$, choose a spanning tree in the adjacency graph of $P$ (edges correspond to shared links).
Integrating sequentially along the tree links yields at least $(n-1)$ such orthogonality gains.
Consequently one has the conservative bound
\begin{equation}\label{eq:rho_dim_supp}
\rho(P)\ \lesssim\ d_{\min}^{2}\,\Bigl(\frac{2f}{d_{\min}^{2}}\Bigr)^{|P|}\qquad(|P|\ge 2),
\end{equation}
and hence, after summing as in \eqref{eq:qrem_ge2},
\begin{equation}\label{eq:qrem_dim}
\qkp_{\mathrm{rem}}
\ \lesssim\
d_{\min}^{2}\,
\frac{\Ccount\,\bigl(2f e^{\theta}/d_{\min}^{2}\bigr)^{2}}{1-\Ccount\,2f e^{\theta}/d_{\min}^{2}}.
\end{equation}
For an arbitrary compact connected simple $G$ one still has a strictly positive lower bound $d_{\min}(G)\ge 2$ for the dimension of the smallest nontrivial irreducible representation. Thus the refinement \eqref{eq:qrem_dim} yields at least as strong a remainder majorant as \eqref{eq:qrem_ge2}, and typically a strictly stronger one. For the classical series the explicit values are recorded in the audited appendices.



\begin{remark}[Interpretation]
Theorem~\ref{thm:obs_stability} is the promised upgrade from a partition-function majorant to \emph{measure-level} control: it compares \emph{expectations of coarse observables} under the exact blocked measure and the positivity-preserving local reference, provided the multi-loop remainder is KP-small.
Closing the remainder bound sharply (improving \eqref{eq:rho_crude}) is the next technical battleground for a true Tier-3 entrance theorem.
\end{remark}


% [v5 hardening] Optional numerical illustration section removed (not used in proofs).

\subsection{Multiscale iteration inside the KP basin: stability of \texorpdfstring{$(\qkp,\qkp_{\mathrm{rem}})$}{(q,q\_rem)}}
\label{sec:iteration}

This section packages the iteration-friendly consequence of the preceding modules:
\begin{quote}
\emph{Once an exact blocked specification enters the KP admissibility window for the single-plaquette activity (and the one-plaquette sector is absorbed into the local reference), both (i) the KP majorant $\qkp$ and (ii) the measure-level remainder majorant $\qkp_{\mathrm{rem}}$ remain valid at all later coarse scales, and in fact collapse extremely fast.}
\end{quote}
This is the lattice-analytic form of a \emph{certified-domain} viewpoint: the admissibility majorant and its remainder majorant are fixed (regime-locked), and the exact RG step is only invoked inside the domain where these quantitative bounds hold.

\subsubsection{Two majorantes at scale \texorpdfstring{$n$}{n}}

Fix a blocking factor $b$ and write $\Nb=b^2$.
At scale $n$ (lattice spacing $a_n=b^n a$) suppose the effective single-plaquette weight is positive type and (after canonical normalization) has activity $f_n$.
Define
\[
q_n:=\qkp(f_n).
\]
After one exact block step, choose the local coarse reference to be the convolution-power decimation $W_{n+1}^{\mathrm{loc}}:=\widehat W_n^{*\Nb}$ (Section~\ref{sec:measure_level}).
Denote by $\mu^{\mathrm{ex}}_{n+1}$ the exact blocked measure and $\mu^{\mathrm{loc}}_{n+1}$ the local reference measure.
Assume the Radon--Nikod\'ym derivative admits a polymerization in the sense of Definition~\ref{def:poly_remainder}, and define
\[
q^{\mathrm{rem}}_{n+1}:=\qkp_{\mathrm{rem}}(\mu^{\mathrm{ex}}_{n+1}\,\|\,\mu^{\mathrm{loc}}_{n+1}).
\]

\subsubsection{One-step propagation bounds}

\begin{proposition}[One-step two-majorant propagation]\label{prop:two_majorant_step}
Assume $\Kact=1$ (convolution-power projection) and the drift hypothesis of Lemma~\ref{lem:exact_step_2u}.
Let $f_n\in(0,1)$.
Then after one exact step:
\begin{enumerate}[leftmargin=*]
\item (\textbf{KP majorant contracts}) If $\Cact\,\Kproj(u_n)u_n\le 1$ for $u_n:=f_n^{\Nb}$, then
\begin{equation}\label{eq:fnplus1_bound}
f_{n+1}\le 2 f_n^{\Nb},
\qquad\text{hence}\qquad
q_{n+1}=\qkp(f_{n+1})\le \qkp(2 f_n^{\Nb})<1.
\end{equation}
In particular, if $f_n\le f_{\mathrm{adm}}$ then $f_{n+1}\ll f_{\mathrm{adm}}$ and the KP regime is preserved.

\item (\textbf{Remainder majorant is $O(f_n^2)$}) If the one-plaquette sector is absorbed into the local reference (so the polymer remainder starts at $|P|\ge 2$), then the crude quadratic majorant bound gives
\begin{equation}\label{eq:qrem_step_bound}
q^{\mathrm{rem}}_{n+1}\le
\Phi(f_n)
:=\frac{\Ccount\,(2 f_n e^{\theta})^{2}}{1-\Ccount\,2 f_n e^{\theta}},
\end{equation}
and with the flux/foam (dimension-suppressed) refinement,
\begin{equation}\label{eq:qrem_step_bound_dim}
q^{\mathrm{rem}}_{n+1}\ \lesssim\
d_{\min}^{2}
\frac{\Ccount\,\bigl(2 f_n e^{\theta}/d_{\min}^{2}\bigr)^{2}}{1-\Ccount\,2 f_n e^{\theta}/d_{\min}^{2}}.
\end{equation}
In particular, $f_n\le f_{\mathrm{adm}}$ implies $q^{\mathrm{rem}}_{n+1}<1$, so Theorem~\ref{thm:obs_stability} applies to local coarse observables at scale $n+1$.
\end{enumerate}
\end{proposition}

\begin{proof}
Item (1) is Lemma~\ref{lem:exact_step_2u} with $u_n=f_n^{\Nb}$.
Item (2) is \eqref{eq:qrem_ge2} (and \eqref{eq:qrem_dim}) applied at activity $f_n$.
\end{proof}

\subsubsection{Induction: once in, you stay in (and rapidly deepen)}

\begin{corollary}[Iteration inside the KP basin]\label{cor:two_majorant_induction}
Assume there exists a scale $n_0$ such that $f_{n_0}\le f_{\mathrm{adm}}$ and the hypotheses of Proposition~\ref{prop:two_majorant_step} hold for all subsequent steps.
Then for all $n\ge n_0$:
\begin{enumerate}[leftmargin=*]
\item $q_n<1$ (KP admissibility persists) and in fact $f_{n+1}\le 2 f_n^{\Nb}$, so $f_n\to 0$ super-exponentially.
\item The measure-level remainder majorant $q^{\mathrm{rem}}_{n+1}$ is well-defined and bounded by $\Phi(f_n)=O(f_n^2)$, hence $q^{\mathrm{rem}}_{n}\to 0$.
\item For any bounded gauge-invariant local coarse observable $\mathcal O$ supported on a finite plaquette set $S$ at scale $n+1$, one has the explicit comparison bound
\[
\bigl|\langle \mathcal O\rangle_{\mu^{\mathrm{ex}}_{n+1}}-\langle \mathcal O\rangle_{\mu^{\mathrm{loc}}_{n+1}}\bigr|
\le
2\,\|\mathcal O\|_{\infty}\,\frac{|S|\,q^{\mathrm{rem}}_{n+1}}{1-q^{\mathrm{rem}}_{n+1}}.
\]
\end{enumerate}
\end{corollary}

% [v5 hardening] Removed purely numerical baseline illustration.



\subsection{Deterministic global entry for the convolution-power majorant}\label{sec:det_entry}

This section is deliberately \emph{deterministic} and isolates a key simplification:
the positivity-preserving convolution-power map $W\mapsto W^{*\Nb}$ has a
global ``mixing-to-uniform'' mechanism on compact groups that does \emph{not} rely on small-$f$.
It therefore provides a clean target for a finite-depth entrance transfer
argument (Section~\ref{sec:finite_transfer}).

\subsubsection{Setup: central probability density and its spectral radius}

Assume in this section that the plaquette weight is \emph{nonnegative, central, and positive type} and is
normalized so that
\begin{equation}
\int_G W(g)\,dg = 1.
\end{equation}
Equivalently, $W$ is a probability density on $G$.
Write its (scalar) Fourier--Peter--Weyl eigenvalues as
\begin{equation}
\lambda_R := \int_G W(g)\,\overline{\varphi_R(g)}\,dg \in [0,1],
\qquad (R\ \text{irrep}),
\end{equation}
so that
\begin{equation}
W^{*N}(g)=1+\sum_{R\ne\mathbf 1} d_R^2\,\lambda_R^{N}\,\varphi_R(g).
\end{equation}
In particular, for positive type,
\begin{equation}\label{eq:f_as_value_at_id}
f_N := W^{*N}(e)-1 = \sum_{R\ne\mathbf 1} d_R^2\,\lambda_R^{N}.
\end{equation}
Define the spectral radius
\begin{equation}
\rho(W):=\sup_{R\ne\mathbf 1}\lambda_R \in [0,1].
\end{equation}
For any genuinely mixing density (i.e.\ not supported in a proper closed subgroup and not Haar itself),
$\rho(W)<1$ and $W^{*N}\to 1$ in $L^2$.

\subsubsection{A sharp \texorpdfstring{$L^2$}{L2}-mixing bound for \texorpdfstring{$f_N$}{f\_N}}

\begin{proposition}[Mixing bound]\label{prop:mixing_bound}
Let $W$ be as above and assume $\rho(W)<1$.
Then for all integers $N\ge 2$,
\begin{equation}\label{eq:mixing_bound}
f_N \le \rho(W)^{N-2}\,\bigl(W^{*2}(e)-1\bigr)
= \rho(W)^{N-2}\,\|W-1\|_2^2.
\end{equation}
\end{proposition}
\begin{proof}
From~\eqref{eq:f_as_value_at_id},
\[
f_N=\sum_{R\ne\mathbf 1} d_R^2\,\lambda_R^{N}
=\sum_{R\ne\mathbf 1} d_R^2\,\lambda_R^{2}\,\lambda_R^{N-2}
\le \rho(W)^{N-2}\sum_{R\ne\mathbf 1} d_R^2\,\lambda_R^{2}.
\]
By Parseval,
$\sum_{R\ne\mathbf 1} d_R^2\lambda_R^2 = \|W-1\|_2^2$,
and also $\|W-1\|_2^2=W^{*2}(e)-1$ for central real $W$.
\end{proof}

\begin{corollary}[Finite-depth deterministic entry for $b=2$]\label{cor:det_entry_b2}
Fix $b=2$ so $\Nb=4$.
Let $f_{\mathrm{adm}}$ be the KP admissibility ceiling defined by \eqref{eq:fadm_def}.
If $\rho(W)<1$, then there exists a finite $k_{\mathrm{enter}}$ such that
\[
f_{4^{k_{\mathrm{enter}}}}\le f_{\mathrm{adm}}.
\]
A sufficient explicit condition is
\begin{equation}\label{eq:k_enter_condition}
4^{k_{\mathrm{enter}}}-2\ \ge\
\frac{\log\!\bigl(f_{\mathrm{adm}}/\|W-1\|_2^2\bigr)}{\log\rho(W)}.
\end{equation}
\end{corollary}

\subsubsection{Heat-kernel majorant mixing via the Casimir gap (explicit \texorpdfstring{$k(\beta)$}{k(beta)} and \texorpdfstring{$\epsilon_n(\beta)$}{epsilon\_n(beta)})}\label{subsec:hardened50_casimir_gap_mixing}

This subsection pins an explicit finite-depth schedule for the \emph{majorant mixing} step used in the shortcut viability test:
given a fixed target $\delta\in(0,1)$ (the KP/Plan-A threshold), we construct an explicit integer depth $k(\beta)$ so that
\[
D_{\chi^2}\!\bigl(\rho_\beta^{(*L^{2k(\beta)})}\,\|\,dg\bigr)\ \le\ \delta,
\]
for the heat-kernel majorant single-plaquette density $\rho_\beta$.
Unlike the polynomial bound based on the crude representation-growth envelope $f(\beta)\le K_G\beta^{m_G+1}$,
the present schedule uses the \emph{Casimir gap} and yields a $\log(1/\delta)$ dependence.

\paragraph{Heat-kernel normalization.}
Let $(K_t)_{t>0}$ denote the (normalized) heat kernel on $G$ for the bi-invariant Laplacian with spectrum $\{C_2(R)\}$.
For $\beta>0$ define the central probability density
\begin{equation}\label{eq:h50_rho_beta_def}
\rho_\beta(g):=K_{1/\beta}(g),\qquad \int_G \rho_\beta\,dg=1.
\end{equation}
By the semigroup property, for each integer $m\ge 1$,
\begin{equation}\label{eq:h50_conv_power_semigroup}
\rho_\beta^{(*m)} \;=\; K_{m/\beta}.
\end{equation}

\paragraph{Casimir gap constant.}
Define the (group-dependent) Casimir gap
\begin{equation}\label{eq:h50_casimir_gap_def}
C_{\min}:=\min\{C_2(R):R\neq \mathbf 1\}\ >\ 0.
\end{equation}

\begin{lemma}[Exact $\chi^2$ formula for the heat-kernel majorant]\label{lem:h50_chi2_heat_kernel_exact}
For $m\in\mathbb N$,
\begin{equation}\label{eq:h50_exact_chi2_series}
D_{\chi^2}\!\bigl(\rho_\beta^{(*m)}\,\|\,dg\bigr)
=\chi^2\!\bigl(\rho_\beta^{(*m)}\bigr)
=\sum_{R\neq\mathbf 1} d_R^2\exp\!\Bigl(-\frac{2m}{\beta}\,C_2(R)\Bigr).
\end{equation}
Equivalently,
\begin{equation}\label{eq:h50_exact_chi2_f}
D_{\chi^2}\!\bigl(\rho_\beta^{(*m)}\,\|\,dg\bigr)
=f\!\Bigl(\frac{\beta}{2m}\Bigr),
\qquad
f(x):=\sum_{R\neq\mathbf 1} d_R^2\,e^{-C_2(R)/x}.
\end{equation}
\end{lemma}

\begin{proof}
By~\eqref{eq:h50_conv_power_semigroup}, $\rho_\beta^{(*m)}=K_{m/\beta}$.
In Peter--Weyl coordinates, for central $K_t$ the (scalar) eigenvalues are $\lambda_R(K_t)=e^{-tC_2(R)}$.
Thus
\[
\chi^2(K_t)=\|K_t-1\|_2^2=\sum_{R\ne\mathbf 1} d_R^2\,\lambda_R(K_t)^2
=\sum_{R\ne\mathbf 1} d_R^2\,e^{-2tC_2(R)}.
\]
Set $t=m/\beta$ to obtain~\eqref{eq:h50_exact_chi2_series}. The rewriting~\eqref{eq:h50_exact_chi2_f} is the identity
$e^{-2tC_2}=e^{-C_2/(1/(2t))}$.
\end{proof}

\begin{lemma}[$L^2$ contraction from the Casimir gap]\label{lem:h50_L2_contraction_gap}
For all $t\ge s>0$,
\begin{equation}\label{eq:h50_L2_gap_contraction}
\|K_t-1\|_2 \le e^{-C_{\min}(t-s)}\,\|K_s-1\|_2.
\end{equation}
\end{lemma}

\begin{proof}
Expand $K_t-1=\sum_{R\ne\mathbf 1} d_R\,e^{-tC_2(R)}\,\chi_R$ in $L^2(G)$.
Then
\[
\|K_t-1\|_2^2=\sum_{R\ne\mathbf 1} d_R^2\,e^{-2tC_2(R)}
=\sum_{R\ne\mathbf 1} d_R^2\,e^{-2(t-s)C_2(R)}\,e^{-2sC_2(R)}
\le e^{-2C_{\min}(t-s)}\sum_{R\ne\mathbf 1} d_R^2\,e^{-2sC_2(R)},
\]
which is exactly~\eqref{eq:h50_L2_gap_contraction} after taking square roots.
\end{proof}

\begin{corollary}[Exponential $\chi^2$ decay]\label{cor:h50_chi2_decay_gap}
Fix $s=\tfrac12$ and define the group constant
\begin{equation}\label{eq:h50_Cgap_def}
C_{1/2}:=\|K_{1/2}-1\|_2^2,\qquad C_{\mathrm{gap}}:=e^{C_{\min}}\,C_{1/2}.
\end{equation}
Then for all integers $m\ge \beta/2$,
\begin{equation}\label{eq:h50_gap_decay}
D_{\chi^2}\!\bigl(\rho_\beta^{(*m)}\,\|\,dg\bigr)
=\|K_{m/\beta}-1\|_2^2
\le C_{\mathrm{gap}}\exp\!\Bigl(-\frac{2C_{\min}}{\beta}\,m\Bigr).
\end{equation}
Moreover $C_{1/2}=f(1)$ in~\eqref{eq:h50_exact_chi2_f} and hence may be bounded using the bound $f(1)\le K_G$ as
$C_{1/2}\le K_G$ (so $C_{\mathrm{gap}}\le e^{C_{\min}}K_G$) when an explicit constant is desired.
\end{corollary}

\begin{proof}
Assume $m\ge \beta/2$ so that $t:=m/\beta\ge \tfrac12$.
Apply Lemma~\ref{lem:h50_L2_contraction_gap} with $s=\tfrac12$ and square the result:
\[
\|K_t-1\|_2^2\le \|K_{1/2}-1\|_2^2\,e^{-2C_{\min}(t-1/2)}
=C_{1/2}\,e^{C_{\min}}\,e^{-2C_{\min}t}.
\]
Substitute $t=m/\beta$ and use~\eqref{eq:h50_conv_power_semigroup}.
The final inequality $C_{1/2}=f(1)\le K_G$ uses the definition of $f$ in~\eqref{eq:h50_exact_chi2_f}
and the bound $f(1)\le K_G$ recorded above.
\end{proof}

\begin{proposition}[Explicit mixing depth $k(\beta)$ for block size $L$]\label{prop:h50_kbeta}
Fix a block factor $L\ge 2$ and a target $\delta\in(0,1)$.
Define
\begin{equation}\label{eq:h50_kbeta_def}
m(\beta):=\max\Bigl\{\frac{\beta}{2},\ \frac{\beta}{2C_{\min}}\log\!\Bigl(\frac{C_{\mathrm{gap}}}{\delta}\Bigr)\Bigr\},
\qquad
k(\beta):=\left\lceil \frac{\log m(\beta)}{2\log L}\right\rceil.
\end{equation}
Then
\begin{equation}\label{eq:h50_majorant_mixing_delta}
D_{\chi^2}\!\bigl(\rho_\beta^{(*L^{2k(\beta)})}\,\|\,dg\bigr)\ \le\ \delta.
\end{equation}
\end{proposition}

\begin{proof}
Let $m_\beta:=L^{2k(\beta)}$. By definition, $m_\beta\ge m(\beta)\ge \beta/2$, so Corollary~\ref{cor:h50_chi2_decay_gap} applies:
\[
D_{\chi^2}\!\bigl(\rho_\beta^{(*m_\beta)}\|dg\bigr)\le C_{\mathrm{gap}}\,e^{-2C_{\min}m_\beta/\beta}.
\]
Also $m_\beta\ge \frac{\beta}{2C_{\min}}\log(\frac{C_{\mathrm{gap}}}{\delta})$ implies
$e^{-2C_{\min}m_\beta/\beta}\le \delta/C_{\mathrm{gap}}$. Multiply to obtain~\eqref{eq:h50_majorant_mixing_delta}.
\end{proof}

\begin{definition}[Equal-split finite-depth transfer budget]\label{def:h50_eps_budget}
With $k(\beta)$ as in~\eqref{eq:h50_kbeta_def}, define
\begin{equation}\label{eq:h50_eps_def}
\epsilon_n(\beta):=\frac{\delta}{k(\beta)},\qquad n=0,1,\dots,k(\beta)-1.
\end{equation}
Then $\sum_{n<k(\beta)}\epsilon_n(\beta)=\delta$.
\end{definition}

\begin{remark}[How this plugs into the finite-depth transfer criterion]
The finite-depth transfer criterion~\eqref{eq:transfer_criterion} requires explicit bounds
$\mathrm{Err}_{n+1}\le \epsilon_{n+1}$ for the first $k_{\mathrm{enter}}$ steps.
Corollary~\ref{cor:h50_chi2_decay_gap} provides a concrete \emph{mixing} depth $k(\beta)$ and a canonical slack allocation~\eqref{eq:h50_eps_def}
satisfying $\sum_{n<k(\beta)}\epsilon_n(\beta)=\delta$.
To complete the finite-depth transfer argument one must still bound the actual change-of-measure errors
$\mathrm{Err}_{n+1}$ by this budget, e.g.\ by proving per-step $L^2$ tilt bounds
of the form $\mathrm{Err}_{n+1}\lesssim \|\Delta_{n+1}(1)\|_\theta$
and using the already-established shared-edge/outside-factor controls to enforce
$\|\Delta_{n+1}(1)\|_\theta \le \epsilon_{n+1}(\beta)$.
\end{remark}



\subsubsection{Optimizing the Casimir-gap schedule (best reference time \texorpdfstring{$s_\delta$}{s\_delta})}\label{subsec:hardened51_optimal_time}

The Casimir-gap bound in Corollary~\ref{cor:h50_chi2_decay_gap} fixed a reference time $s=\tfrac12$ for convenience.
Here we record the sharpest version of the same idea for the heat-kernel majorant:
choose the reference time so that the target $\chi^2$ level $\delta$ is met exactly at that time.
This removes the prefactor $C_{\mathrm{gap}}$ from the required convolution power.

\begin{lemma}[Monotonicity and level-set time for $\chi^2(K_t)$]\label{lem:h51_chi2_monotone}
Let
\[
\Phi(t):=\chi^2(K_t)=\|K_t-1\|_2^2=\sum_{R\ne \mathbf 1} d_R^2\,e^{-2tC_2(R)}.
\]
Then $\Phi(t)$ is finite for all $t>0$, continuous, and strictly decreasing in $t$, with
$\Phi(t)\to\infty$ as $t\downarrow 0$ and $\Phi(t)\to 0$ as $t\to\infty$.
Consequently, for every $\delta>0$ there exists a unique $s_\delta>0$ such that $\Phi(s_\delta)=\delta$.
\end{lemma}

\begin{proof}
Finiteness for $t>0$ follows from Lemma~\ref{lem:h50_chi2_heat_kernel_exact} and the finiteness of
$f(x)$ in~\eqref{eq:h50_exact_chi2_f} at $x=\frac{1}{2t}$:
\[
\Phi(t)=f\!\Bigl(\frac{1}{2t}\Bigr)<\infty.
\]
Termwise continuity of the series is immediate. Strict decrease holds termwise since $C_2(R)>0$ for $R\ne\mathbf 1$
and all coefficients are nonnegative.
The limit $\Phi(t)\to 0$ as $t\to\infty$ follows by dominated convergence from the series representation, because each term decays exponentially
and the sum is finite for every fixed $t>0$.
For the $t\downarrow 0$ limit, rewrite using~\eqref{eq:h50_exact_chi2_f}:
$\Phi(t)=f(\frac{1}{2t})$, and invoke the recorded growth bound on $f(x)$ as $x\to\infty$ to conclude $\Phi(t)\to\infty$.
Since $\Phi$ is strictly decreasing and continuous with range $(0,\infty)$, the level-set time $s_\delta$ exists and is unique.
\end{proof}

\begin{proposition}[Optimal mixing depth via $s_\delta$]\label{prop:h51_kbeta_opt}
Fix $L\ge 2$ and $\delta\in(0,1)$.
Let $s_\delta$ be as in Lemma~\ref{lem:h51_chi2_monotone}.
Define
\begin{equation}\label{eq:h51_kbeta_opt_def}
m_{\mathrm{opt}}(\beta):=\max\{1,\beta\,s_\delta\},
\qquad
k_{\mathrm{opt}}(\beta):=\left\lceil \frac{\log m_{\mathrm{opt}}(\beta)}{2\log L}\right\rceil.
\end{equation}
Then
\begin{equation}\label{eq:h51_majorant_mixing_delta}
D_{\chi^2}\!\bigl(\rho_\beta^{(*L^{2k_{\mathrm{opt}}(\beta)})}\,\|\,dg\bigr)\le \delta.
\end{equation}
Moreover, for the heat-kernel majorant one has
\[
D_{\chi^2}\!\bigl(\rho_\beta^{(*m)}\,\|\,dg\bigr)\le \delta
\quad\Longleftrightarrow\quad
\frac{m}{\beta}\ge s_\delta.
\]
\end{proposition}

\begin{proof}
By the semigroup identity~\eqref{eq:h50_conv_power_semigroup}, $\rho_\beta^{(*m)}=K_{m/\beta}$.
Let $m_\beta:=L^{2k_{\mathrm{opt}}(\beta)}$. By definition, $m_\beta\ge m_{\mathrm{opt}}(\beta)\ge \beta s_\delta$, hence $m_\beta/\beta\ge s_\delta$.
By monotonicity of $\Phi(t)=\chi^2(K_t)$ (Lemma~\ref{lem:h51_chi2_monotone}),
\[
D_{\chi^2}\!\bigl(\rho_\beta^{(*m_\beta)}\|dg\bigr)
=\chi^2(K_{m_\beta/\beta})
\le \chi^2(K_{s_\delta})
=\delta,
\]
proving~\eqref{eq:h51_majorant_mixing_delta}.
The final equivalence is the same strict monotonicity of $\Phi$.
\end{proof}

\begin{definition}[Equal-split budget based on $k_{\mathrm{opt}}(\beta)$]\label{def:h51_eps_budget}
With $k_{\mathrm{opt}}(\beta)$ as in~\eqref{eq:h51_kbeta_opt_def}, define
\[
\epsilon_n^{\mathrm{opt}}(\beta):=\frac{\delta}{k_{\mathrm{opt}}(\beta)},\qquad n=0,1,\dots,k_{\mathrm{opt}}(\beta)-1.
\]
Then $\sum_{n<k_{\mathrm{opt}}(\beta)}\epsilon_n^{\mathrm{opt}}(\beta)=\delta$.
\end{definition}

% [v5 hardening] Removed group-specific numerical illustration (not used).




\subsection{Finite-depth entrance transfer: the exact comparison target}\label{sec:finite_transfer}

This section states (cleanly) what must be proved to transfer the deterministic entry in
Section~\ref{sec:det_entry} to the \emph{exact} blocked Yang--Mills action.

\subsubsection{The finite-depth reduction}

Let $S_0$ be the microscopic (fine-lattice) action and let
\[
S_{n+1}:=B_{\mathrm{exact}}(S_n)
\]
denote exact blocking on a fixed scale factor $b$.
Let $S_{n+1}^{\square}$ denote the local positive-type reference produced by the
convolution-power decimation ($W\mapsto W^{*\Nb}$) applied to the local plaquette weight extracted from $S_n$.

Section~\ref{sec:iteration} already proves the following:
\begin{quote}
If for some $n_0$ the scale-$n_0$ local activity obeys $f_{n_0}\le f_{\mathrm{adm}}$, then for all $n\ge n_0$ the
pair of majorantes $(\qkp,\qkp_{\mathrm{rem}})$ remains admissible, contracts rapidly, and coarse observables are
controlled in a KP-small sense relative to the local reference measure.
\end{quote}
Thus, the only missing step is to show that such an $n_0$ exists on the physically relevant trajectory.

Deterministic majorant entry (Section~\ref{sec:det_entry}) produces a small integer $k_{\mathrm{enter}}$
such that the purely local decimation would have entered:
\[
f^{\square}_{k_{\mathrm{enter}}}\le f_{\mathrm{adm}}.
\]
Therefore the entrance problem reduces to the following finite-depth \emph{transfer} statement:
prove that the exact blocked action does not deviate too far from the local decimation during the first
$k_{\mathrm{enter}}$ steps.

\subsubsection{A clean transfer criterion (what must be bounded)}

A sufficient (and natural) finite-depth criterion is:
\begin{equation}\label{eq:transfer_criterion}
f_{n+1}\ \le\ f_{n+1}^{\square} + \mathrm{Err}_{n+1},
\qquad
\mathrm{Err}_{n+1}\le \epsilon_{n+1}\ \text{for}\ n=0,1,\dots,k_{\mathrm{enter}}-1,
\end{equation}
with a small, explicit error budget $\epsilon_{n+1}$ that ensures
$f_{k_{\mathrm{enter}}}\le f_{\mathrm{adm}}$.
In the finite-depth certification viewpoint, $\epsilon_{n+1}$ is the finite-depth \emph{change-of-measure slack}:
it is the maximal permitted deviation before the trajectory loses the admissible envelope.

The work already done in Section~\ref{sec:proj_remainder} provides a template:
the error is controlled by a remainder norm at the \emph{incoming} scale,
\begin{equation}\label{eq:error_template}
\mathrm{Err}_{n+1}\ \le\ \Cact\,\|\Delta_{n+1}(1)\|_{\theta}
\ \lesssim\ \Cact\,\Kproj\!\bigl(f_{n+1}^{\square}\bigr)\,\bigl(f_{n+1}^{\square}\bigr)^2,
\end{equation}
once the remainder admits a KP-style polymer bound.

\begin{remark}[Key technical point (handled below)]
Outside the KP basin, the remainder does not admit a small-activity expansion \emph{a priori}.
The deterministic entry theorem localizes this issue to a \emph{finite} set of early steps.
The remainder of this appendix resolves these early steps by replacing the small-activity majorant with
the mixing majorant introduced in Section~\ref{sec:mix_majorant} and by proving the associated transfer bounds,
thereby verifying the finite-depth transfer criterion \eqref{eq:transfer_criterion} with explicit constants.
\end{remark}


\subsection{An early-step mixing majorant: spectral radius and \texorpdfstring{$\chi^2$}{chi2}-distance to Haar}\label{sec:mix_majorant}

The KP majorant $\qkp$ is powerful once $f$ is small, but it is not a good \emph{early-step} majorant
in the weak-coupling regime where the plaquette density is sharply peaked at the identity and $f$ is not small.
A more structural majorant---suggested by the finite-depth certification viewpoint---is a \emph{mixing majorant}
attached to the single-plaquette kernel viewed as a class Markov kernel on $G$.
This section formalizes two such majorantes and records the precise inequalities that would transfer
deterministic entry to the exact RG.


\paragraph{Early-step mixing criterion.}
Under the standing choices above, the \emph{early-step} entrance criterion is formulated in terms of the spectral radius $\rho(W)$
(and the associated conditional $\chi^2$ transfer inequality used to compare exact and reference kernels).
Any alternative early-step diagnostic is admissible only if it either (i) is proved to imply this quantitative input with explicit constants,
or (ii) is declared as a genuine scope change with its own transport theorem .

\subsubsection{Central Markov kernels and the spectral-radius majorant}

Let $W\ge 0$ be a central (class) function on $G$ with $\int_G W\,dg=1$.
Define the convolution operator on $L^2(G)$,
\[
(T_W f)(x):=\int_G W(y)\,f(y^{-1}x)\,dy.
\]
If $W$ is positive type, then $T_W$ is self-adjoint, positivity-preserving, and
$\|T_W\|_{2\to 2}=1$ with constants as the top eigenspace.

In Peter--Weyl coordinates, for central $W$ the matrix Fourier coefficient is scalar:
\[
\widehat W(R)=\int_G W(g)\,\rho_R(g^{-1})\,dg = \lambda_R(W)\,I_{d_R},
\qquad \lambda_R(W)\in\mathbb R.
\]
Equivalently, with our normalization $W(g)=\sum_R d_R\,\lambda_R(W)\,\chi_R(g)$ and
$\lambda_{\mathbf 1}(W)=1$.
For positive type, $\lambda_R(W)\ge 0$ for all $R$.

\begin{definition}[Spectral-radius majorant]\label{def:rho_majorant}
For a central probability density $W$ define
\[
\rho(W):=\sup_{R\neq \mathbf 1}\lambda_R(W)\in[0,1].
\]
\end{definition}

\begin{lemma}[Convolution powers contract $\rho$ exactly]\label{lem:rho_convolution}
For central probability densities $W$,
\[
\lambda_R(W^{*N})=\lambda_R(W)^N,\qquad
\rho(W^{*N})=\rho(W)^N.
\]
\end{lemma}

\begin{proof}
The matrix Fourier transform turns convolution into matrix multiplication.
For central $W$, the Fourier coefficients are scalar multiples of the identity, hence multiply as scalars.
\end{proof}

\subsubsection{\texorpdfstring{$\chi^2$}{chi2}-distance to Haar and its relation to \texorpdfstring{$\rho$}{rho}}

Define the $\chi^2$-distance to Haar (the $L^2$ distance to the constant density $1$):
\[
\chi^2(W):=\|W-1\|_2^2 = \int_G (W(g)-1)^2\,dg.
\]
In the character basis,
\begin{equation}\label{eq:chi2_pw}
\chi^2(W)=\sum_{R\neq\mathbf 1} d_R^2\,\lambda_R(W)^2.
\end{equation}

\begin{lemma}[Dimension lower bound]\label{lem:rho_vs_chi2}
Let $d_{\min}:=\min\{d_R: R\neq \mathbf 1\}$.
Then for central probability densities $W$,
\[
d_{\min}^2\,\rho(W)^2 \le \chi^2(W),
\qquad\text{hence}\qquad
\rho(W)\le \frac{\sqrt{\chi^2(W)}}{d_{\min}}.
\]
\end{lemma}

\begin{proof}
By definition $\rho(W)=\sup_{R\neq 1}\lambda_R(W)$.
In \eqref{eq:chi2_pw}, choose an $R$ attaining (or approximating) the supremum.
Then $\chi^2(W)\ge d_R^2\rho(W)^2\ge d_{\min}^2\rho(W)^2$.
\end{proof}

\begin{proposition}[$L^2$-mixing bound for convolution powers]\label{prop:chi2_convolution}
For $N\ge 2$,
\[
\chi^2(W^{*N})
=\sum_{R\neq\mathbf 1} d_R^2\,\lambda_R(W)^{2N}
\le \rho(W)^{2N-4}\,\chi^2(W^{*2}).
\]
\end{proposition}

\begin{proof}
Factor $\lambda_R(W)^{2N}=\lambda_R(W)^{2N-4}\lambda_R(W)^4\le \rho(W)^{2N-4}\lambda_R(W)^4$
and sum over $R\neq\mathbf 1$.
\end{proof}

\subsubsection{From \texorpdfstring{$L^2$}{L2} to total variation: a mixing-time bound}

For probability densities $\nu,\pi$ on a finite state space, the total variation distance is
$d_{\mathrm{TV}}(\nu,\pi)=\tfrac12\sum_x|\nu(x)-\pi(x)|$.
In our compact-group setting with Haar probability $dg$, if $W$ is a density then
\begin{equation}
 d_{\mathrm{TV}}(W\,dg,\,dg)
 := \tfrac12\,\|W-1\|_{L^1(dg)}.
\end{equation}
The Cauchy--Schwarz inequality gives the clean $L^2\to L^1$ conversion
\begin{equation}\label{eq:tv_vs_chi2_group}
 d_{\mathrm{TV}}(W\,dg,\,dg)
\le \tfrac12\,\|W-1\|_2
=\tfrac12\,\sqrt{\chi^2(W)}.
\end{equation}

\begin{proposition}[Spectral gap $\Rightarrow$ mixing time on compact groups]\label{prop:mixing_time_group}
Let $W$ be a central probability density on a compact group $G$ and assume $\rho(W)<1$.
Then for all integers $t\ge 1$,
\begin{equation}\label{eq:tv_mixing_bound_group}
 d_{\mathrm{TV}}(W^{*t}\,dg,\,dg)
\le \tfrac12\,\rho(W)^{t-1}\,\|W-1\|_2.
\end{equation}
Consequently, for any $\epsilon\in(0,\tfrac12)$ the $\epsilon$-mixing time
\[
\tau(\epsilon):=\min\{t\ge 1: d_{\mathrm{TV}}(W^{*t}\,dg,\,dg)\le \epsilon\}
\]
satisfies
\begin{equation}\label{eq:mixing_time_bound_group}
 \tau(\epsilon)
\le 1+\frac{\log\bigl(\|W-1\|_2/(2\epsilon)\bigr)}{-\log\rho(W)}
\le 1+\frac{1}{1-\rho(W)}\,\log\Bigl(\frac{\|W-1\|_2}{2\epsilon}\Bigr).
\end{equation}
\end{proposition}

\begin{proof}
The Peter--Weyl expansion gives
$\chi^2(W^{*t})=\sum_{R\ne\mathbf 1} d_R^2\,\lambda_R(W)^{2t}$.
For $t\ge 1$, factor $\lambda_R(W)^{2t}=\lambda_R(W)^2\,\lambda_R(W)^{2(t-1)}\le \lambda_R(W)^2\,\rho(W)^{2(t-1)}$ and sum:
\[
\chi^2(W^{*t})\le \rho(W)^{2(t-1)}\sum_{R\ne\mathbf 1} d_R^2\,\lambda_R(W)^2
=\rho(W)^{2(t-1)}\,\chi^2(W).
\]
Take square roots and use~\eqref{eq:tv_vs_chi2_group}.
For the final inequality note that $-\log x\ge 1-x$ for $x\in(0,1)$.
\end{proof}

\begin{remark}[Relation to the finite-state bound]
For a finite, reversible Markov chain with stationary distribution $\pi$ and spectral gap
$\gamma=1-\lambda_2$ (or absolute gap $\gamma^*=1-\max(|\lambda_2|,|\lambda_n|)$), a standard bound is
\[
\tau(\epsilon)\ \lesssim\ \frac{1}{\gamma^*}\Bigl(\log\frac{1}{\epsilon}+\tfrac12\log\frac{1}{\pi_*}\Bigr),
\qquad \pi_*:=\min_x\pi(x).
\]
The factor involving $\pi_*$ arises when converting $L^2(\pi)$ control to total variation in the finite-state setting.
In our compact-group setting with Haar probability $dg$, the conversion~\eqref{eq:tv_vs_chi2_group} is immediate,
so the role of the ``initial condition'' is played by $\|W-1\|_2$.
\end{remark}

\subsubsection{Why \texorpdfstring{$\rho$}{rho} is the right early-step majorant}

In a finite-depth certification approach, each admissible RG step is required to preserve a chosen envelope majorant.
Inside the KP basin the natural majorant is $\qkp$; outside the KP basin $\qkp$ is too crude.
The majorant $\rho$ is structural: it measures the maximum nontrivial representation content of the
single-plaquette kernel, is bounded in $[0,1]$, and contracts \emph{globally} under convolution.
It therefore provides a concrete, conventional-mathematics target for early-step control.

\subsubsection{Finite-depth entrance transfer in terms of \texorpdfstring{$\rho$}{rho} (exact RG target)}

Let $W_n^{\mathrm{ex}}$ denote the \emph{single coarse plaquette} density extracted from the exact blocked
measure at scale $n$ (by integrating out everything except one coarse plaquette holonomy, normalized to a density).
Let $W_n^{\square}$ denote the deterministic convolution-power decimation at the same scale.
Then $W_{n+1}^{\square}=(W_n^{\square})^{*\Nb}$ and $\rho(W_{n+1}^{\square})=\rho(W_n^{\square})^{\Nb}$.

\begin{proposition}[A clean $\rho$-transfer criterion]\label{prop:rho_transfer}
Fix a depth $k$.
Assume that for each $n=0,1,\dots,k-1$ one has a comparison
\[
\rho(W_{n+1}^{\mathrm{ex}})\le \rho(W_{n+1}^{\square}) + \delta_{n+1}
\qquad\text{with}\qquad
\delta_{n+1}\le \tfrac12\bigl(\rho(W_n^{\square})^{\Nb}-\rho(W_{n+1}^{\square})\bigr).
\]
Then
\[
\rho(W_n^{\mathrm{ex}})\le 2\,\rho(W_n^{\square})\qquad\text{for all }n\le k.
\]
Consequently, by Lemma~\ref{lem:rho_vs_chi2} and Proposition~\ref{prop:chi2_convolution}, the deterministic
entry bound of Section~\ref{sec:det_entry} transfers (with explicit constants) to the exact lineage.
\end{proposition}

\begin{proof}
Induct on $n$.
At each step, the hypothesis implies
$\rho(W_{n+1}^{\mathrm{ex}})\le \rho(W_{n+1}^{\square})+\frac12(\rho(W_n^{\square})^{\Nb}-\rho(W_{n+1}^{\square}))
\le \frac12\rho(W_n^{\square})^{\Nb}+\frac12\rho(W_{n+1}^{\square})$.
Using $\rho(W_{n+1}^{\square})=\rho(W_n^{\square})^{\Nb}$ gives
$\rho(W_{n+1}^{\mathrm{ex}})\le \rho(W_{n+1}^{\square})$ up to a factor $2$,
and the claim follows.
\end{proof}

\begin{remark}[Finite-depth entrance inequality]
A key step is to prove, for the first few steps before KP-smallness applies, a bound of the form
\[
\rho(W_{n+1}^{\mathrm{ex}})\ \le\ \rho(W_{n+1}^{\square}) + \delta_{n+1},
\]
or a comparable $L^2$-mixing bound on $\chi^2(W_{n+1}^{\mathrm{ex}})$.
This is a \emph{finite-depth} problem for Wilson-type inputs because
Section~\ref{sec:det_entry} shows that $k_{\mathrm{enter}}$ is finite and admits an explicit bound.
(For illustration only, in baseline Wilson examples at $b=2$ one typically finds $k_{\mathrm{enter}}\le 5$, but no numerical value is used.)
\end{remark}




\subsection{From measure-level remainder to mixing: an \texorpdfstring{$L^2$}{L2} transfer module}\label{sec:L2_transfer}

Section~\ref{sec:mix_majorant} isolates the key entrance goal as a finite-depth control of a
mixing majorant ($\rho$ or $\chi^2$) for the \emph{exact} blocked single-plaquette marginal.
This section records a concrete, conventional-mathematics bridge:
\begin{quote}
\emph{If one can bound an $L^2$ change-of-measure between the exact blocked measure and the
positivity-preserving local reference, then one gets explicit control of the nontrivial
representation eigenvalues, hence of $\rho$ and $\chi^2$.}
\end{quote}

The point of this module is to make the early-step target inequality completely explicit and local.
In this submission the required non-collapse lower bound on the resulting change-of-measure slack is supplied by the
pivot and mixing analysis of Section~\ref{sec:slack_collapse} and the one-plaquette coefficient discharge (A1)
(Proposition~\ref{prop:TBmix_one_plaq_coeff_crossover}, implied by Theorem~\ref{thm:A1_one_plaq_classical}).

\subsubsection{\texorpdfstring{$\chi^2$}{chi2} divergence and basic transfer inequalities}

Let $(X,\mathcal F)$ be a measurable space and let $\nu$ be a probability measure.
If $\mu\ll\nu$ with Radon--Nikod\'ym derivative $r:=\frac{d\mu}{d\nu}$, define the $\chi^2$ divergence
\begin{equation}\label{eq:chi2_div}
D_{\chi^2}(\mu\|\nu):=\int_X (r-1)^2\,d\nu=\|r-1\|_{L^2(\nu)}^2.
\end{equation}

\begin{lemma}[$L^2$ change-of-measure control]\label{lem:L2_change_measure}
Let $\mu\ll\nu$ with density $r=d\mu/d\nu$.
Then for any $F\in L^2(\nu)$,
\begin{equation}\label{eq:change_measure_bound}
\bigl|\mathbb E_{\mu}[F]-\mathbb E_{\nu}[F]\bigr|
\le
\|F-\mathbb E_{\nu}[F]\|_{L^2(\nu)}\,\sqrt{D_{\chi^2}(\mu\|\nu)}.
\end{equation}
\end{lemma}

\begin{proof}
Write $\mathbb E_{\mu}[F]-\mathbb E_{\nu}[F]=\int (F-\mathbb E_{\nu}[F])(r-1)\,d\nu$ and apply Cauchy--Schwarz.
\end{proof}

\begin{lemma}[$\chi^2$ transfer of second moments and variances]\label{lem:chi2_transfer_variance}
Let $\mu\ll\nu$ with density $r=d\mu/d\nu$ and set $\varepsilon:=D_{\chi^2}(\mu\|\nu)$.
Then for any $F\in L^4(\nu)$,
\begin{align}
\bigl|\mathbb E_{\mu}[F]-\mathbb E_{\nu}[F]\bigr|
&\le
\|F-\mathbb E_{\nu}[F]\|_{L^2(\nu)}\,\sqrt{\varepsilon},
\label{eq:chi2_transfer_mean}
\\
\bigl|\mathbb E_{\mu}[F^2]-\mathbb E_{\nu}[F^2]\bigr|
&\le
\|F^2-\mathbb E_{\nu}[F^2]\|_{L^2(\nu)}\,\sqrt{\varepsilon}
\ \le\ 
\|F\|_{L^4(\nu)}^2\,\sqrt{\varepsilon}.
\label{eq:chi2_transfer_second_moment}
\end{align}
Consequently, if $\varepsilon\le 1$ then
\begin{equation}\label{eq:chi2_variance_transfer}
\bigl|\mathrm{Var}_{\mu}(F)-\mathrm{Var}_{\nu}(F)\bigr|
\ \le\
\bigl(\|F\|_{L^4(\nu)}^2 + 4\|F\|_{L^2(\nu)}^2\bigr)\,\sqrt{\varepsilon}.
\end{equation}
\end{lemma}

\begin{proof}
The mean bound \eqref{eq:chi2_transfer_mean} is Lemma~\ref{lem:L2_change_measure}.
Apply Lemma~\ref{lem:L2_change_measure} to $F^2$ to get the first inequality in \eqref{eq:chi2_transfer_second_moment}, and then use
$\|F^2-\mathbb E_{\nu}[F^2]\|_2\le \|F^2\|_2=\|F\|_4^2$.

For the variance, write
\[
\mathrm{Var}_{\mu}(F)-\mathrm{Var}_{\nu}(F)
=
(\mathbb E_{\mu}[F^2]-\mathbb E_{\nu}[F^2])
-(\mathbb E_{\mu}[F]^2-\mathbb E_{\nu}[F]^2).
\]
The first difference is bounded by \eqref{eq:chi2_transfer_second_moment}.
For the mean term,
\[
|\mathbb E_{\mu}[F]^2-\mathbb E_{\nu}[F]^2|
=
|\mathbb E_{\mu}[F]-\mathbb E_{\nu}[F]|\cdot|\mathbb E_{\mu}[F]+\mathbb E_{\nu}[F]|.
\]
By Cauchy--Schwarz,
$|\mathbb E_{\mu}[F]|\le \|F\|_{2}\,\|r\|_{2}=\|F\|_{2}\sqrt{1+\varepsilon}$,
and similarly $|\mathbb E_{\nu}[F]|\le \|F\|_{2}$.
When $\varepsilon\le 1$, $\sqrt{1+\varepsilon}\le \sqrt 2\le 2$, so
$|\mathbb E_{\mu}[F]+\mathbb E_{\nu}[F]|\le 3\|F\|_{2}\le 4\|F\|_{2}$.
Combine with \eqref{eq:chi2_transfer_mean} to get \eqref{eq:chi2_variance_transfer}.
\end{proof}

\begin{lemma}[$\chi^2$ divergence under conditioning]\label{lem:chi2_conditioning}
Let $\mu\ll\nu$ on $(X,\mathcal F)$ with density $r=d\mu/d\nu$.
Let $\Omega\in\mathcal F$ satisfy $\mu(\Omega),\nu(\Omega)>0$ and write $\mu^{\Omega}:=\mu(\,\cdot\,|\Omega)$ and $\nu^{\Omega}:=\nu(\,\cdot\,|\Omega)$.
Then $\mu^{\Omega}\ll\nu^{\Omega}$ with density
\[
\frac{d\mu^{\Omega}}{d\nu^{\Omega}}
=
\frac{\nu(\Omega)}{\mu(\Omega)}\,r
\qquad\text{on }\Omega,
\]
and
\begin{equation}\label{eq:chi2_conditioning}
D_{\chi^2}(\mu^{\Omega}\|\nu^{\Omega})
=
\frac{\nu(\Omega)}{\mu(\Omega)^2}\,\mathbb E_{\nu}\!\big[r^2\,\mathbf 1_{\Omega}\big]-1
\ \le\
\frac{\nu(\Omega)}{\mu(\Omega)^2}\,\bigl(1+D_{\chi^2}(\mu\|\nu)\bigr)-1.
\end{equation}
\end{lemma}

\begin{proof}
The conditional density formula is immediate from Bayes' rule.
For the divergence, compute
\[
D_{\chi^2}(\mu^{\Omega}\|\nu^{\Omega})
=
\mathbb E_{\nu^{\Omega}}\!\left[\left(\frac{\nu(\Omega)}{\mu(\Omega)}\,r-1\right)^2\right]
=
\frac{1}{\nu(\Omega)}\int_{\Omega}\left(\frac{\nu(\Omega)}{\mu(\Omega)}\,r-1\right)^2\,d\nu
=
\frac{\nu(\Omega)}{\mu(\Omega)^2}\,\mathbb E_{\nu}\!\big[r^2\,\mathbf 1_{\Omega}\big]-1.
\]
The inequality follows from $\mathbb E_{\nu}[r^2\,\mathbf 1_{\Omega}]\le \mathbb E_{\nu}[r^2]=1+D_{\chi^2}(\mu\|\nu)$.
\end{proof}



\subsubsection{Application to central single-plaquette marginals}

Let $w^{\mathrm{ex}}$ and $w^{\mathrm{loc}}$ be \emph{central probability densities} on $G$
(e.g.\ single-plaquette marginals at a given RG scale).
Let $\lambda_R(w):=\int_G w(g)\,\varphi_R(g)\,dg$ be the normalized character coefficient.
Define the spectral-radius majorant $\rho(w)=\sup_{R\neq \mathbf 1}\lambda_R(w)$ (Definition~\ref{def:rho_majorant}).

\begin{lemma}[Coefficient transfer from Haar-$L^2$ density control]\label{lem:coeff_transfer_L2}
For each irrep $R\neq\mathbf 1$,
\begin{equation}\label{eq:coeff_transfer_L2}
\bigl|\lambda_R(w^{\mathrm{ex}})-\lambda_R(w^{\mathrm{loc}})\bigr|
\le
\frac{1}{d_R}\,\|w^{\mathrm{ex}}-w^{\mathrm{loc}}\|_{L^2(dg)}.
\end{equation}
Consequently,
\begin{equation}\label{eq:rho_transfer_L2}
\rho(w^{\mathrm{ex}})
\le
\rho(w^{\mathrm{loc}})
+
\frac{1}{d_{\min}}\,\|w^{\mathrm{ex}}-w^{\mathrm{loc}}\|_{2}.
\end{equation}
\end{lemma}

\begin{proof}
By definition,
$\lambda_R(w^{\mathrm{ex}})-\lambda_R(w^{\mathrm{loc}})=\int_G \varphi_R(g)\,(w^{\mathrm{ex}}-w^{\mathrm{loc}})(g)\,dg$.
Apply Cauchy--Schwarz and $\|\varphi_R\|_2^2=\int \chi_R\overline{\chi_R}/d_R^2=1/d_R^2$.
Take the supremum over $R\neq \mathbf 1$.
\end{proof}

The $\chi^2$ divergence relative to $w^{\mathrm{loc}}$ controls the Haar-$L^2$ difference by a trivial bound.

\begin{lemma}[$\chi^2$ divergence implies Haar-$L^2$ density control]\label{lem:chi2_to_L2}
Assume $w^{\mathrm{ex}}\ll w^{\mathrm{loc}}$ and write $r:=w^{\mathrm{ex}}/w^{\mathrm{loc}}$.
Then
\begin{equation}\label{eq:L2_from_chi2}
\|w^{\mathrm{ex}}-w^{\mathrm{loc}}\|_2^2
=\int_G (r-1)^2\,(w^{\mathrm{loc}})^2\,dg
\le \|w^{\mathrm{loc}}\|_{\infty}\,D_{\chi^2}(w^{\mathrm{ex}}\|w^{\mathrm{loc}}).
\end{equation}
\end{lemma}

\begin{proof}
Bound $(w^{\mathrm{loc}})^2\le \|w^{\mathrm{loc}}\|_{\infty}\,w^{\mathrm{loc}}$.
\end{proof}

Combining Lemmas~\ref{lem:coeff_transfer_L2} and \ref{lem:chi2_to_L2} yields the explicit bound
\begin{equation}\label{eq:rho_transfer_chi2}
\rho(w^{\mathrm{ex}})
\le
\rho(w^{\mathrm{loc}})
+
\frac{\sqrt{\|w^{\mathrm{loc}}\|_{\infty}}}{d_{\min}}\,
\sqrt{D_{\chi^2}(w^{\mathrm{ex}}\|w^{\mathrm{loc}})}.
\end{equation}

\subsubsection{Bounding \texorpdfstring{$\chi^2$}{chi2} divergence when the RN factor is bounded by \texorpdfstring{$1$}{1}}

A particularly clean situation occurs when the change of measure is given by a factor $R$ with
$0\le R\le 1$ (as happens naturally when ``dropping'' plaquette factors after sup-normalization).

\begin{lemma}[If $0\le R\le 1$, then $\chi^2$ divergence is controlled by $\mathbb E_\nu(R)$]
\label{lem:R_le_1_chi2}
Let $(X,\nu)$ be a probability space and let $R:X\to[0,1]$ be measurable with $\mathbb E_{\nu}[R]>0$.
Define the tilted probability measure
\[
d\mu:=\frac{R}{\mathbb E_{\nu}[R]}\,d\nu.
\]
Then
\begin{equation}\label{eq:chi2_bound_ER}
D_{\chi^2}(\mu\|\nu)
=\mathbb E_{\nu}\Bigl[\Bigl(\frac{R}{\mathbb E_{\nu}[R]}-1\Bigr)^2\Bigr]
\le \frac{1}{\mathbb E_{\nu}[R]}-1.
\end{equation}
\end{lemma}

\begin{proof}
Compute
$D_{\chi^2}(\mu\|\nu)=\mathbb E_{\nu}[R^2]/\mathbb E_{\nu}[R]^2-1$.
Since $0\le R\le 1$ implies $R^2\le R$, we have $\mathbb E_{\nu}[R^2]\le \mathbb E_{\nu}[R]$.
\end{proof}

\subsubsection{Finite-depth entrance transfer rewritten as a local integral bound}

Combine \eqref{eq:rho_transfer_chi2} and Lemma~\ref{lem:R_le_1_chi2}.
Suppose that at an early RG step $n\mapsto n+1$, the exact blocked single-plaquette marginal
$w_{n+1}^{\mathrm{ex}}$ can be expressed as a tilt of the local reference marginal $w_{n+1}^{\mathrm{loc}}$ by
a factor $R_{n+1}\in[0,1]$:
\begin{equation}\label{eq:tilt_assumption}
 w_{n+1}^{\mathrm{ex}}(g)\,dg
=
\frac{R_{n+1}(g)}{\mathbb E_{w_{n+1}^{\mathrm{loc}}}[R_{n+1}]}
\,w_{n+1}^{\mathrm{loc}}(g)\,dg.
\end{equation}
Then
\begin{equation}\label{eq:rho_transfer_ER}
\rho\bigl(w_{n+1}^{\mathrm{ex}}\bigr)
\le
\rho\bigl(w_{n+1}^{\mathrm{loc}}\bigr)
+
\frac{\sqrt{\|w_{n+1}^{\mathrm{loc}}\|_{\infty}}}{d_{\min}}\,
\sqrt{\frac{1}{\mathbb E_{w_{n+1}^{\mathrm{loc}}}[R_{n+1}]}-1}.
\end{equation}

\begin{remark}[Why this is useful]
The right-hand side is now expressed in terms of:
(i) the deterministic/local contraction $\rho(w_{n+1}^{\mathrm{loc}})$ (which is known exactly for convolution powers),
(ii) a mild shape constant $\|w_{n+1}^{\mathrm{loc}}\|_{\infty}$, and
(iii) a single local integral $\mathbb E_{w_{n+1}^{\mathrm{loc}}}[R_{n+1}]$.

For Wilson-type inputs, deterministic contraction drives $\rho(w_{n+1}^{\mathrm{loc}})$ to be extremely small after
only a few $b=2$ steps (Section~\ref{sec:det_entry}).
Thus, to transfer entrance it suffices to prove that the local expectation
$\mathbb E[R_{n+1}]$ is not absurdly tiny for those few early steps.
This is a finite-dimensional problem on a fixed block.
In the load-bearing primary chain below it is handled by an analytic representation-theoretic mixing bound;
alternatively one may compute the relevant integrals explicitly by finite-dimensional recoupling identities.
\end{remark}

\subsection{Local tilt construction for a single coarse plaquette and reduction to a finite block integral}\label{sec:local_tilt}

The $L^2$ transfer module of Section~\ref{sec:L2_transfer} reduces early-step entrance transfer to verifying,
for a \emph{finite depth} of RG steps, a lower bound on a local expectation
$\mathbb E[R]$ for a factor $R\in[0,1]$.
This section makes that reduction concrete by constructing such a factor for a
\emph{single coarse plaquette} and expressing its expectation as a \emph{finite-dimensional block integral}.

\subsubsection{Exact vs reduced local measures on a block}

Fix an oriented plane $(\mu,\nu)$ and let $T_L$ be the $L\times L$ plaquette tile in that plane (with $L=b^{n+1}$ at RG depth $n\mapsto n+1$).

\begin{definition}[Thin block neighborhood of a coarse plaquette]\label{def:block_BL}
Let $\mathcal B_L$ be the set of fine links needed to support:
\begin{enumerate}[label=(\roman*)]
\item all in-plane plaquettes of $T_L$, and
\item all out-of-plane plaquettes that share at least one link with an in-plane plaquette of $T_L$.
\end{enumerate}
Let $P_{\parallel}(\mathcal B_L)$ be the set of in-plane plaquettes and $P_{\perp}(\mathcal B_L)$ the set of out-of-plane plaquettes in $\mathcal B_L$.
\end{definition}

Work throughout with a sup-normalized plaquette weight $\widehat W$ so that $0\le \widehat W\le 1$.
(As noted earlier, multiplying each plaquette weight by a constant does not change the Gibbs measure.)

Define two probability measures on the block-link space $G^{E(\mathcal B_L)}$:
\begin{align}
d\nu_L(U) &:= \frac{1}{Z_L^{\parallel}}\Bigl(\prod_{p\in P_{\parallel}(\mathcal B_L)} \widehat W(U_p)\Bigr)\,dU,
\\
d\mu_L(U) &:= \frac{1}{Z_L^{\mathrm{full}}}\Bigl(\prod_{p\in P_{\parallel}(\mathcal B_L)\cup P_{\perp}(\mathcal B_L)} \widehat W(U_p)\Bigr)\,dU,
\end{align}
where $dU$ is Haar product measure on the block links.

\begin{lemma}[Local change-of-measure factor is in $[0,1$]]\label{lem:block_tilt_R}
Let
\[
R_L(U):=\prod_{p\in P_{\perp}(\mathcal B_L)} \widehat W(U_p)\in[0,1].
\]
Then $\mu_L$ is a tilt of $\nu_L$:
\begin{equation}\label{eq:block_tilt}
d\mu_L = \frac{R_L}{\mathbb E_{\nu_L}[R_L]}\,d\nu_L,
\qquad\text{and}\qquad
\mathbb E_{\nu_L}[R_L]=\frac{Z_L^{\mathrm{full}}}{Z_L^{\parallel}}\in(0,1].
\end{equation}
\end{lemma}

\begin{proof}
Immediate from the definitions: the only difference between $d\mu_L$ and $d\nu_L$ is the extra factor
$R_L$ and the normalizing constants.
Since $0\le\widehat W\le 1$, one has $0\le R_L\le 1$ and $Z_L^{\mathrm{full}}\le Z_L^{\parallel}$.
\end{proof}

\subsubsection{Pushforward to a single coarse plaquette holonomy}

Let $\Phi_L:G^{E(\mathcal B_L)}\to G$ be the map that assigns to a block configuration $U$
the holonomy around the \emph{outer boundary} of the tile $T_L$ in the $(\mu,\nu)$ plane:
\[
\Phi_L(U):=\mathrm{Hol}_{\partial T_L}(U).
\]
Let $w_L^{\mathrm{ex}}$ and $w_L^{\parallel}$ be the pushforward (single-plaquette) densities:
\[
w_L^{\mathrm{ex}}\,dg := (\Phi_L)_{\#}\mu_L,
\qquad
w_L^{\parallel}\,dg := (\Phi_L)_{\#}\nu_L.
\]

\begin{lemma}[Marginal tilt factor is a conditional expectation]\label{lem:marginal_tilt}
Define
\[
\overline R_L(g):=\mathbb E_{\nu_L}\!\bigl[R_L \,\big|\, \Phi_L=g\bigr]\in[0,1].
\]
Then $w_L^{\mathrm{ex}}$ is a tilt of $w_L^{\parallel}$:
\begin{equation}\label{eq:marginal_tilt}
w_L^{\mathrm{ex}}(g)=\frac{\overline R_L(g)}{\mathbb E_{w_L^{\parallel}}[\overline R_L]}\,w_L^{\parallel}(g),
\qquad
\mathbb E_{w_L^{\parallel}}[\overline R_L]=\mathbb E_{\nu_L}[R_L]=\frac{Z_L^{\mathrm{full}}}{Z_L^{\parallel}}.
\end{equation}
\end{lemma}

\begin{proof}
For any bounded measurable $F$ on $G$,
\[
\int_G F(g)\,w_L^{\mathrm{ex}}(g)\,dg
=\mathbb E_{\mu_L}[F(\Phi_L)]
=\frac{\mathbb E_{\nu_L}[F(\Phi_L)R_L]}{\mathbb E_{\nu_L}[R_L]}.
\]
Disintegrate $\nu_L$ over $\Phi_L$; the numerator becomes
$\int F(g)\,\overline R_L(g)\,w_L^{\parallel}(g)\,dg$.
This identifies the density ratio as stated.
\end{proof}

\subsubsection{A data-processing bound for \texorpdfstring{$\chi^2$}{chi2} divergence}

\begin{lemma}[Data processing for $\chi^2$ divergence]\label{lem:chi2_data_processing}
Let $\mu\ll\nu$ on $(X,\mathcal F)$ and let $\Phi:X\to Y$ be measurable.
Then
\[
D_{\chi^2}(\Phi_{\#}\mu\|\Phi_{\#}\nu)\le D_{\chi^2}(\mu\|\nu).
\]
\end{lemma}

\begin{proof}
This is a standard property of $f$-divergences (here $f(t)=(t-1)^2$) under Markov kernels.
A direct proof is obtained by writing the density ratio of $\Phi_{\#}\mu$ w.r.t.\ $\Phi_{\#}\nu$ as the conditional
expectation of $d\mu/d\nu$ given $\Phi$, and applying Jensen's inequality to the convex function $(t-1)^2$.
\end{proof}

Combining Lemmas~\ref{lem:block_tilt_R}, \ref{lem:marginal_tilt}, \ref{lem:R_le_1_chi2}, and \ref{lem:chi2_data_processing} yields:

\begin{proposition}[Early-step mixing transfer reduced to a block integral]\label{prop:rho_transfer_block_integral}
With notation as above,
\[
D_{\chi^2}\!\bigl(w_L^{\mathrm{ex}}\|w_L^{\parallel}\bigr)
\le D_{\chi^2}(\mu_L\|\nu_L)
\le \frac{1}{\mathbb E_{\nu_L}[R_L]}-1
=\frac{Z_L^{\parallel}}{Z_L^{\mathrm{full}}}-1.
\]
Consequently,
\begin{equation}\label{eq:rho_transfer_block_final}
\rho\bigl(w_L^{\mathrm{ex}}\bigr)
\le
\rho\bigl(w_L^{\parallel}\bigr)
+
\frac{\sqrt{\|w_L^{\parallel}\|_{\infty}}}{d_{\min}}\,
\sqrt{\frac{Z_L^{\parallel}}{Z_L^{\mathrm{full}}}-1}.
\end{equation}
\end{proposition}

\begin{remark}[Uniqueness mechanism]
The entrance/non-collapse and UV matching hypotheses imply uniqueness of macroscopic laws on $\mathcal O_{\ell}$ by KP/Dobrushin contraction;
see Theorem~\ref{thm:uniqueness_on_Oell}.
\end{remark}

\begin{remark}[What is genuinely left to bound]
Equation~\eqref{eq:rho_transfer_block_final} reduces the early-step entrance-transfer inequality to
lower-bounding the \emph{finite-dimensional} ratio $Z_L^{\mathrm{full}}/Z_L^{\parallel}$ for the
block region $\mathcal B_L$ at depths $n=0,1,\dots,k_{\mathrm{enter}}-1$.
For Wilson-type inputs, $k_{\mathrm{enter}}$ is finite and admits an explicit bound (e.g.\ Corollary~\ref{cor:det_entry_b2}).
	For illustration (not used in the proof), a numerical evaluation for a baseline Wilson one-plaquette density at block size $b=2$
	typically yields $k_{\mathrm{enter}}\le 5$ even at large $\beta$, so this reduces concretely to a \emph{finite list of finite integrals}.
	In this submission the required \emph{non-collapse} of the early-step slack ratio is proved analytically by the structural pivot of
	Section~\ref{sec:slack_collapse} and the subsequent coefficient-level mixing analysis, culminating in the one-plaquette coefficient discharge (A1)
	(Proposition~\ref{prop:TBmix_one_plaq_coeff_crossover}, implied by Theorem~\ref{thm:A1_one_plaq_classical}).
	In particular, no numerical verification is used in the proof.
\end{remark}

\subsection{Remark: the planar-reference slack collapses at weak coupling (and why we must pivot)}\label{sec:slack_collapse}

Section~\ref{sec:local_tilt} reduced early-step entrance transfer (in the mixing majorant $\rho$) to a
\emph{lower bound} on the block ratio $Z_L^{\mathrm{full}}/Z_L^{\parallel}$.
This is a perfectly sharp reduction, but it hides a crucial issue:
if the reference $\nu_L$ is \emph{purely planar} (i.e.\ contains no out-of-plane plaquettes),
then $Z_L^{\mathrm{full}}/Z_L^{\parallel}$ typically collapses rapidly in weak coupling.
This section records an explicit upper bound which already shows that one cannot expect a
\emph{uniform-in-$\beta$} lower bound for Wilson-type inputs.
The correct response is a structural pivot: the reference measure used for the tilt must
include out-of-plane plaquettes with \emph{positive exponent}, so that the tilt factor does not lose all mass.

\subsubsection{A Finner/H\"older upper bound for the planar slack ratio}

For $p\ge 1$ define
\[
\|\widehat W\|_{p}:=\Bigl(\int_G \widehat W(g)^{p}\,dg\Bigr)^{1/p}.
\]

\begin{lemma}[Planar slack has an explicit dimension-free upper bound]\label{lem:planar_slack_upper}
Fix a plane $(\mu,\nu)$ and tile side $L$.
Let $\nu_L$ be the planar reference measure of Section~\ref{sec:local_tilt} and
$R_L=\prod_{p\in P_{\perp}(\mathcal B_L)}\widehat W(U_p)$ the out-of-plane product.
Then
\begin{equation}\label{eq:planar_slack_upper}
\mathbb E_{\nu_L}[R_L] \le \|\widehat W\|_{4}^{|P_{\perp}(\mathcal B_L)|}
=
\Bigl(\int_G \widehat W(g)^{4}\,dg\Bigr)^{|P_{\perp}(\mathcal B_L)|/4}.
\end{equation}
In particular, for fixed $d\ge 3$,
\[
|P_{\perp}(\mathcal B_L)| = 4(d-2)L(L+1),
\]
so the right-hand side decays like a \emph{stretched power} of $\|\widehat W\|_4$ with exponent $\asymp L^2$.
\end{lemma}

\begin{proof}
Under $\nu_L$, the out-of-plane link variables appearing in $R_L$ are integrated with Haar product measure.
Each out-of-plane plaquette factor depends on at most four such link variables and is bounded by~$1$.
Moreover, each out-of-plane link appears in at most four plaquette factors from $P_\perp(\mathcal B_L)$
(two with the $\mu$-direction and two with the $\nu$-direction, corresponding to the two ``sides'' of the link).
Thus the incidence hypergraph has maximum degree~$\le 4$.
Applying Finner's inequality (generalized H\"older) with uniform exponents $p_e=4$ gives
\[
\int \prod_{p\in P_{\perp}} \widehat W(U_p)\,dU_{\perp}
\ \le\
\prod_{p\in P_{\perp}}\Bigl(\int \widehat W(U_p)^{4}\,dU_{\perp}\Bigr)^{1/4}.
\]
For a fixed plaquette $p$, the holonomy map from the four incident link variables (with Haar product measure)
pushes forward to Haar on $G$, so the $L^4$ norm equals $\|\widehat W\|_4$.
This proves \eqref{eq:planar_slack_upper}.
The count $|P_\perp|=4(d-2)L(L+1)$ follows from:
there are $2L(L+1)$ in-plane links in the $L\times L$ tile, and each such link lies in $2(d-2)$ out-of-plane plaquettes.
\end{proof}

\subsubsection{Implication for Wilson-type weak coupling}

For a Wilson-class one-plaquette weight of the form
\[
\widehat W_\beta(U):=\exp\big(-\beta\,A(U)\big),\qquad A(\1)=0,
\]
one has (by Laplace asymptotics at the identity; see Remark~\ref{rem:zbeta_scaling_G})
\begin{equation}\label{eq:W4_asymptotic}
\int_{G} \widehat W_\beta(U)^4\,dU
=\int_{G} \widehat W_{4\beta}(U)\,dU
\asymp (4\beta)^{-\dim(G)/2}\qquad(\beta\to\infty).
\end{equation}

Combining \eqref{eq:planar_slack_upper}--\eqref{eq:W4_asymptotic} yields, for fixed $L$,
\[
\mathbb E_{\nu_L}[R_L]
\ \lesssim\
\beta^{-\frac{\dim(G)}{8}|P_\perp(\mathcal B_L)|}\qquad(\beta\to\infty),
\]
so the planar slack ratio collapses at weak coupling.
In particular, there cannot exist a lower bound on $Z_L^{\mathrm{full}}/Z_L^{\parallel}$ that is uniform in $\beta$
for Wilson-type inputs.

\begin{remark}[Structural consequence: planar reference is degenerate at weak coupling]
The quantitative ``slack'' quantity that must remain nondegenerate to transfer entrance is precisely
\[
\mathbb E_{\nu_L}[R_L]=\frac{Z_L^{\mathrm{full}}}{Z_L^{\parallel}}.
\]
Lemma~\ref{lem:planar_slack_upper} shows that the purely planar reference $\nu_L$ forces this slack to collapse in weak coupling.
Thus the planar reference is \emph{not} a usable comparison reference in the UV scaling regime: it loses too much mass to carry
a stable change-of-measure comparison.
The appropriate pivot is to use a reference measure that includes out-of-plane plaquettes with positive exponent
so that the change-of-measure factor is not catastrophically small.
\end{remark}

\subsubsection{Pivot target: a weighted reference that includes out-of-plane plaquettes}

Introduce an interpolation parameter $\alpha\in(0,1)$ and define a \emph{weighted reference} on the same block:
\begin{equation}\label{eq:alpha_reference_block}
d\nu_L^{(\alpha)}(U)
:=
\frac{1}{Z_L^{(\alpha)}}\Bigl(\prod_{p\in P_{\parallel}(\mathcal B_L)}\widehat W(U_p)\Bigr)
\Bigl(\prod_{p\in P_{\perp}(\mathcal B_L)}\widehat W(U_p)^{\alpha}\Bigr)\,dU.
\end{equation}
Then the full block measure is the tilt of $\nu_L^{(\alpha)}$ by
\[
R_L^{(\alpha)}(U):=\prod_{p\in P_{\perp}(\mathcal B_L)}\widehat W(U_p)^{1-\alpha}\in[0,1],
\qquad
d\mu_L=\frac{R_L^{(\alpha)}}{\mathbb E_{\nu_L^{(\alpha)}}[R_L^{(\alpha)}]}\,d\nu_L^{(\alpha)},
\]
with
\[
\mathbb E_{\nu_L^{(\alpha)}}[R_L^{(\alpha)}]=\frac{Z_L^{\mathrm{full}}}{Z_L^{(\alpha)}}.
\]
Taking $\alpha$ closer to $1$ reduces the exponent $1-\alpha$ and therefore mitigates slack collapse.
The remaining difficulty is to obtain a mixing estimate for the boundary holonomy marginal of $\nu_L^{(\alpha)}$
that is strong enough, over a \emph{finite depth} of steps, to enter the KP basin.
This is the next genuinely non-toy entrance-transfer target.



\paragraph{Thin-block leaf factorization: exact slack formula and the right $\alpha$-scaling}

A useful (and slightly sobering) structural fact about the specific thin block $\mathcal B_L$ defined in
Section~\ref{sec:local_tilt} is that every out-of-plane plaquette in $P_\perp(\mathcal B_L)$ carries an
\emph{opposite in-plane edge} that is not shared by any other plaquette in the block.
Consequently the out-of-plane sector can be integrated out \emph{exactly} as a product of
one-plaquette integrals.

For a general nonnegative plaquette weight $\widehat W:G\to(0,1]$ and exponent $t>0$ write
\begin{equation}\label{eq:z_t_def}
z(t):=\int_G \widehat W(g)^{t}\,dg .
\end{equation}

\begin{lemma}[Leaf factorization on the thin block]\label{lem:leaf_factorization_thin_block}
Let $\mathcal B_L$ be the thin block built from the in-plane tile $T_L$ together with the out-of-plane plaquettes
that share at least one in-plane link with $T_L$ (Definition~\ref{def:block_BL}).
Let $\nu_L^{(\alpha)}$ be the weighted reference \eqref{eq:alpha_reference_block} with $\alpha\in(0,1)$.
Then the corresponding block partition functions satisfy the exact identities
\begin{equation}\label{eq:Zalpha_factorizes}
Z_L^{(\alpha)} = Z_L^{\parallel}\,z(\alpha)^{|P_\perp(\mathcal B_L)|},
\qquad
Z_L^{\mathrm{full}} = Z_L^{\parallel}\,z(1)^{|P_\perp(\mathcal B_L)|},
\end{equation}
and therefore the change-of-measure slack expectation is
\begin{equation}\label{eq:ER_alpha_exact_thin}
\mathbb E_{\nu_L^{(\alpha)}}\!\bigl[R_L^{(\alpha)}\bigr]
=\frac{Z_L^{\mathrm{full}}}{Z_L^{(\alpha)}}
=\Bigl(\frac{z(1)}{z(\alpha)}\Bigr)^{|P_\perp(\mathcal B_L)|}.
\end{equation}
In particular, for the purely planar reference ($\alpha=0$, so $z(0)=1$),
\begin{equation}\label{eq:ER_planar_exact_thin}
\mathbb E_{\nu_L}[R_L]=z(1)^{|P_\perp(\mathcal B_L)|}.
\end{equation}
\end{lemma}

\begin{proof}
Each $p\in P_\perp(\mathcal B_L)$ is a plaquette in a plane $(\mu,\rho)$ with $\rho\neq \mu,\nu$
that shares a unique in-plane link $\ell$ of the tile $T_L$ (at perpendicular coordinate $0$).
Its opposite edge is the translate $\ell^{(\rho,\pm)}$ of that link at perpendicular coordinate $\pm 1$
(in the same in-plane direction as $\ell$).
By construction of the \emph{thin} block $\mathcal B_L$, no plaquettes lying entirely in the shifted
in-plane slices are included, hence each translated edge $\ell^{(\rho,\pm)}$ appears in \emph{exactly one}
plaquette factor, namely the corresponding $p$.

Fix all other link variables and integrate first over $\ell^{(\rho,\pm)}$ with Haar measure.
Writing the plaquette holonomy as $U_p=A\,\ell^{(\rho,\pm)}\,B$ (with $A,B$ depending on the fixed links),
Haar invariance gives
\[
\int_G \widehat W(U_p)^{\alpha}\,d\ell^{(\rho,\pm)}
=
\int_G \widehat W(g)^{\alpha}\,dg
=z(\alpha),
\]
independently of $A,B$.
Thus one may integrate out the unique opposite edge for each $p\in P_\perp(\mathcal B_L)$ sequentially,
factoring out $z(\alpha)^{|P_\perp|}$ and leaving precisely the planar block integral $Z_L^{\parallel}$.
Setting $\alpha=1$ in the preceding computation gives $Z_L^{\mathrm{full}}$.
Taking the ratio yields \eqref{eq:ER_alpha_exact_thin} and \eqref{eq:ER_planar_exact_thin}.
\end{proof}

\begin{remark}[Why this matters: it pins down the correct $\alpha$ scaling]\label{rem:alpha_scaling}
Equation~\eqref{eq:ER_alpha_exact_thin} shows that on the thin block the slack depends on $\alpha$
\emph{only through} the one-plaquette integral ratio $z(1)/z(\alpha)$, raised to the power $|P_\perp(\mathcal B_L)|\asymp L^2$.
Therefore, to keep the $\chi^2$ correction term in \eqref{eq:rho_transfer_ER} bounded over a finite depth of RG steps with growing $L$,
one needs
\[
|P_\perp(\mathcal B_L)|\,\log\!\Bigl(\frac{z(\alpha)}{z(1)}\Bigr)=O(1),
\qquad\text{i.e.}\qquad
1-\alpha \ \lesssim\ \frac{c}{|P_\perp(\mathcal B_L)|}\asymp \frac{c}{L^2}.
\]
This is the \emph{reference-stabilizing} choice: keep the slack nondegenerate by letting
$\alpha$ approach $1$ as the block grows.

A convenient explicit schedule (for $d=4$) is
\begin{equation}\label{eq:alpha_schedule_L}
\alpha(L):=1-\frac{\kappa}{|P_\perp(\mathcal B_L)|}
=1-\frac{\kappa}{8L(L+1)},
\qquad \kappa\in(0,1)\ \text{fixed}.
\end{equation}

\end{remark}

\begin{remark}[Weak-coupling scaling of Wilson one-plaquette integrals]\label{rem:zbeta_scaling_G}
Let $G$ be a fixed compact connected Lie group and let $A:G\to[0,\infty)$ be the Wilson one-plaquette action density associated to a fixed unitary representation $\pi$.
By \Cref{rem:effective_group_convention,def:effective_group} we may and do assume $\pi$ is faithful by replacing $G$ with $G_{\mathrm{eff}}:=G/\ker(\pi)$.
\[
A(g):=1-\frac{1}{\dim(\pi)}\Re\mathrm{Tr}\,\pi(g),
\qquad A(\1)=0.
\]
Define the (sup-normalized) one-plaquette weight and its integral
\[
\widehat W_\beta(U):=\exp\big(-\beta\,A(U)\big),
\qquad
z_\beta:=\int_G \widehat W_\beta(U)\,dU.
\]
By Laplace's method at the unique nondegenerate minimum of $A$ (at $U=\1$, under the effective-group convention), one has as $\beta\to\infty$
\begin{equation}\label{eq:zbeta_laplace_asymp}
 z_\beta = c_G\,\beta^{-\dim(G)/2}\,\bigl(1+O(\beta^{-1})\bigr),
\end{equation}
for an explicit constant $c_G>0$ depending only on $G$, $\pi$, and the chosen inner product on $\mathfrak g$ (equivalently on the Hessian of $A$ at the identity).
Consequently, for fixed $\alpha\in(0,1]$,
\begin{equation}\label{eq:zalpha_ratio}
 \frac{z_{\alpha\beta}}{z_\beta} = \alpha^{-\dim(G)/2}\,\bigl(1+o(1)\bigr)\qquad (\beta\to\infty).
\end{equation}
Inserting \eqref{eq:zalpha_ratio} into the exact thin-block slack identity \eqref{eq:ER_alpha_exact_thin} shows that the schedule
\eqref{eq:alpha_schedule_L} yields
\[
\mathbb E_{\nu_L^{(\alpha(L))}}\!\big[R_L^{(\alpha(L))}\big]
\approx
\exp\!\Big(-\frac{\dim(G)}{2}\,\kappa\Big)
\qquad (\beta\ \text{large}),
\]
uniformly in $L$.

The asymptotic statements \eqref{eq:zbeta_laplace_asymp}--\eqref{eq:zalpha_ratio} are included only for orientation and are not used in the proofs below; all subsequent bounds are derived from the explicit inequalities proved in this appendix.

\smallskip
\noindent\emph{Orientation only (not used in the proof).} Write $U=\exp(X)$ with $X\in\mathfrak g$ near $0$, expand $A(\exp X)=\frac12\langle X, H X\rangle+O(\|X\|^3)$ with $H$ positive definite, and use the standard Laplace rescaling $X\mapsto X/\sqrt\beta$.
\end{remark}




\paragraph{Weak-coupling $\chi^2$--decay for the $\alpha$--pivot on genuine 4D blocks}\label{sec:chi2_decay_alpha}

The thin-block leaf factorization (Lemma~\ref{lem:leaf_factorization_thin_block}) is special to $\mathcal B_L$.
For a genuine four-dimensional block, the out-of-plane plaquette sector shares edges and does not factor.
Nevertheless, in weak coupling one can still prove the required $\chi^2$ decay under the schedule
\eqref{eq:alpha_schedule_L} by combining:
(i) \emph{curvature} (rather than edge) Brascamp--Lieb domination, and
(ii) a quadratic-form concentration bound (Hanson--Wright) applied to the linearized curvature field.

\medskip
\noindent\textbf{Transverse energy and RN factor.}
Define the transverse energy observable
\begin{equation}\label{eq:X_def_alpha}
X(U):=\sum_{p\in P_\perp(\mathcal B_L)}A(U_p),
\qquad
A(g):=1-\tfrac{1}{\dim(\pi)}\Re\mathrm{Tr}\,\pi(g),
\end{equation}
for a fixed unitary representation $\pi$ (without loss of generality faithful under the effective-group convention).
For the sup-normalized Wilson class weights
\[
\widehat W_\beta(g):=\exp\Big(\beta\Big(\tfrac{1}{\dim(\pi)}\Re\mathrm{Tr}\,\pi(g)-1\Big)\Big)=e^{-\beta A(g)},
\]
the $\alpha$--reference measure \eqref{eq:alpha_reference_block} satisfies
\begin{equation}\label{eq:RN_factor_beta}
R_L^{(\alpha)}(U)
:=\prod_{p\in P_\perp(\mathcal B_L)}\widehat W_\beta(U_p)^{1-\alpha}
=\exp\!\big(-\beta(1-\alpha)\,X(U)\big)\in[0,1].
\end{equation}
Recall the boundary holonomy map $\Phi_L:G^{E(\mathcal B_L)}\to G$ from Section~\ref{sec:local_tilt} and define the corresponding
$\alpha$--reference single-plaquette marginal by
\[
w_L^{(\alpha)}\,dg := (\Phi_L)_{\#}\nu_L^{(\alpha)}.
\]
Since $d\mu_L = (R_L^{(\alpha)}/\mathbb E_{\nu_L^{(\alpha)}}[R_L^{(\alpha)}])\,d\nu_L^{(\alpha)}$, the block-level $\chi^2$ divergence is
\[
D_{\chi^2}(\mu_L\|\nu_L^{(\alpha)})
=\frac{\mathbb E_{\nu_L^{(\alpha)}}\!\big[e^{-2tX}\big]}{\mathbb E_{\nu_L^{(\alpha)}}\!\big[e^{-tX}\big]^2}-1,
\qquad t:=\beta(1-\alpha).
\]
By data processing for $\chi^2$ divergence (Lemma~\ref{lem:chi2_data_processing}) under $\Phi_L$,
\begin{equation}\label{eq:chi2_mgf_upper}
D_{\chi^2}\!\bigl(w_L^{\mathrm{ex}}\|w_L^{(\alpha)}\bigr)
\le D_{\chi^2}(\mu_L\|\nu_L^{(\alpha)})
=\frac{\mathbb E_{\nu_L^{(\alpha)}}\!\big[e^{-2tX}\big]}{\mathbb E_{\nu_L^{(\alpha)}}\!\big[e^{-tX}\big]^2}-1.
\end{equation}
Thus it suffices to control exponential moments of $X$ under $\nu_L^{(\alpha)}$.

\paragraph{Gauge fixing and curvature variables.}
Fix a spanning tree $T$ on $E(\mathcal B_L)$ and impose tree gauge.
Let $A=(A_e)_{e\in E_f}$ be exponential coordinates for free edges and let
\[
F:=dA\in \mathfrak g^{P(\mathcal B_L)}
\]
denote the \emph{linearized curvature} 2-form (cochain differential $d$).
Let $\Omega_{\varepsilon_0}$ be the small-curvature event of Lemma~\ref{lem:weak_localization}.
On $\Omega_{\varepsilon_0}$, Proposition~\ref{prop:curvature_BL} implies that all centered linear functionals
$\langle J,F\rangle$ are sub-Gaussian with variance proxy $(\beta m_-)^{-1}$, uniformly in $L$.

\medskip
The key point is that, although the edge Laplacian $\Delta_1=d^*d$ has small eigenvalues at large $L$, the curvature covariance involves
the projection $\Pi_{\mathrm{ex}}=d\Delta_1^{-1}d^*$ (Lemma~\ref{lem:projection_identity}), hence has operator norm $1$.
This removes the long-mode degeneracy and yields uniform control for curvature observables.

\begin{lemma}[Variance scaling for the transverse energy]\label{lem:VarX_beta2}
Fix $d=4$ and compact connected $G$.
There exist $\beta_0=\beta_0(G)$ and $C=C(G)$ such that for all $\beta\ge\beta_0$, all $L$, and all $\alpha\in[1/2,1]$,
\begin{equation}\label{eq:VarX_bound_beta2}
\mathrm{Var}_{\nu_L^{(\alpha)}}(X)\ \le\ \frac{C}{\beta^2}\,|P_\perp(\mathcal B_L)|.
\end{equation}
\end{lemma}

\begin{proof}
Work conditionally on $\Omega_{\varepsilon_0}$ and then remove the conditioning using Lemma~\ref{lem:weak_localization}.
On $\Omega_{\varepsilon_0}$ each plaquette holonomy is within a fixed normal neighborhood of the identity, and the class function
$A(g)$ admits the Taylor expansion
\begin{equation}\label{eq:A_Taylor}
A(\exp \xi)=\tfrac12\langle \xi,\mathcal Q \xi\rangle + R(\xi),
\qquad |R(\xi)|\le C_R\,\|\xi\|^3
\quad (\|\xi\|\le \varepsilon_0),
\end{equation}
with $\mathcal Q\succeq c_Q I$ depending only on $(G,\pi)$.
Writing the plaquette logarithm as $\xi_p=\log U_p$, one has $\xi_p=F_p+O(\|A\|^2)$ on $\Omega_{\varepsilon_0}$, hence
\[
X
=\frac12\sum_{p\in P_\perp}\langle F_p,\mathcal Q F_p\rangle
+ \mathrm{Rem},
\qquad
|\mathrm{Rem}|\ \le\ C\sum_{p\in P_\perp}\|F_p\|^3,
\]
after enlarging $C$.
The leading term is a quadratic form in the curvature vector $F$:
\[
X_{\mathrm{quad}}:=\frac12\sum_{p\in P_\perp}\langle F_p,\mathcal Q F_p\rangle
=\langle F, Q F\rangle,
\]
where $Q$ is block-diagonal on plaquettes with $\|Q\|_{\mathrm{op}}\le C$ and
$\|Q\|_{\mathrm{HS}}^2\asymp |P_\perp|$.

By Proposition~\ref{prop:curvature_BL}, $F$ is a sub-Gaussian vector (in the exact 2-form subspace) with parameter
$\sigma^2\asymp (\beta m_-)^{-1}$.
The Hanson--Wright inequality for sub-Gaussian vectors therefore yields
\[
\mathrm{Var}(X_{\mathrm{quad}}\mid \Omega_{\varepsilon_0})
\ \le\ C_{HW}\,\sigma^4\,\|Q\|_{\mathrm{HS}}^2
\ \le\ \frac{C}{\beta^2}\,|P_\perp|.
\]
For the remainder, sub-Gaussian moment bounds imply
$\mathbb E[\|F_p\|^6\mid\Omega_{\varepsilon_0}] \le C/\beta^3$ uniformly in $p$,
hence by Cauchy--Schwarz,
\[
\mathbb E[\mathrm{Rem}^2\mid\Omega_{\varepsilon_0}]
\ \le\ C\sum_{p\in P_\perp}\mathbb E[\|F_p\|^6\mid\Omega_{\varepsilon_0}]
\ \le\ \frac{C}{\beta^3}\,|P_\perp|.
\]
Thus $\mathrm{Var}(X\mid\Omega_{\varepsilon_0})\le (C/\beta^2)|P_\perp|$ after enlarging $C$.
Finally, Lemma~\ref{lem:weak_localization} implies $\nu_L^{(\alpha)}(\Omega_{\varepsilon_0}^c)\le |P(\mathcal B_L)|e^{-c\beta}$,
so removing the conditioning changes $\mathrm{Var}(X)$ by at most an exponentially small term in $\beta$ that can be absorbed into $C$ for
$\beta\ge\beta_0$.
\end{proof}

\begin{lemma}[Centered mgf bound for the transverse energy]\label{lem:mgf_centered_X}
Under the assumptions of Lemma~\ref{lem:VarX_beta2}, there exist constants $\beta_1=\beta_1(G)\ge\beta_0(G)$ and $C=C(G)$ such that for all
$\beta\ge\beta_1$, all $L$, all $\alpha\in[1/2,1]$, and all $s$ with $|s|\le c_1\,\beta$,
\begin{equation}\label{eq:mgf_centered_X}
\log \mathbb E_{\nu_L^{(\alpha)}}\!\Big[\exp\big(s(X-\mathbb E X)\big)\Big]
\ \le\ \frac{C}{\beta^2}\,s^2\,|P_\perp(\mathcal B_L)|.
\end{equation}
\end{lemma}

\begin{proof}
Condition on $\Omega_{\varepsilon_0}$ and use the decomposition $X=X_{\mathrm{quad}}+\mathrm{Rem}$ from the proof of Lemma~\ref{lem:VarX_beta2}.
The Hanson--Wright mgf form for sub-Gaussian vectors implies that, for $|s|\le c_1\beta$,
\[
\log \mathbb E\!\big[\exp\big(s(X_{\mathrm{quad}}-\mathbb E X_{\mathrm{quad}})\big)\mid\Omega_{\varepsilon_0}\big]
\ \le\ \frac{C}{\beta^2}\,s^2\,|P_\perp|.
\]
For the remainder, sub-Gaussian moment bounds on $F$ on $\Omega_{\varepsilon_0}$ imply $\|\mathrm{Rem}-\mathbb E\mathrm{Rem}\|_{L^2}\le
C|P_\perp|^{1/2}\beta^{-3/2}$, hence by Cauchy--Schwarz
\[
\log \mathbb E\!\big[e^{s(\mathrm{Rem}-\mathbb E\mathrm{Rem})}\mid\Omega_{\varepsilon_0}\big]
\le \frac{C}{\beta^3}\,s^2\,|P_\perp|
\]
after restricting to $|s|\le c_1\beta$ and absorbing constants.
Combining these and removing the conditioning as in Lemma~\ref{lem:VarX_beta2} (using Lemma~\ref{lem:weak_localization}) yields
\eqref{eq:mgf_centered_X}.
\end{proof}


\begin{remark}[Compatibility of the schedule with the mgf radius]\label{rem:kappa_mgf_radius}
Lemma~\ref{lem:mgf_centered_X} holds for $|s|\le c_1\,\beta$.
In the $\chi^2$--decay application below we take $s=-2t_L$ with $t_L=\beta(1-\alpha(L))=\beta\kappa/|P_\perp(\mathcal B_L)|$.
In $d=4$ one has $|P_\perp(\mathcal B_L)|=8L(L+1)\ge 16$ for $L\ge 1$, hence $|s|\le \beta\kappa/8$.
Thus it suffices (and we henceforth assume) that $\kappa\le \kappa_0(G):=8c_1(G)$ so that $|s|\le c_1\beta$ for all $L$.
(If one prefers to allow arbitrary $\kappa\in(0,1)$, it is enough to check finitely many small $L$ separately and then use the argument below for large $L$.)
\end{remark}



\begin{proposition}[$\chi^2$ decay under the schedule \eqref{eq:alpha_schedule_L}]\label{prop:chi2_decay_alpha_schedule}
Fix $\kappa\in(0,\kappa_0(G)]$ (cf.\ Remark~\ref{rem:kappa_mgf_radius}) and take the schedule \eqref{eq:alpha_schedule_L} so that $1-\alpha(L)=\kappa/|P_\perp(\mathcal B_L)|$.
Set $t_L:=\beta(1-\alpha(L))=\beta\kappa/|P_\perp(\mathcal B_L)|$.
Then for $\beta\ge\beta_0(G)$,
\begin{equation}\label{eq:chi2_decay_alpha_v2}
D_{\chi^2}\!\bigl(w_L^{\mathrm{ex}}\|w_L^{(\alpha(L))}\bigr)
\ \le\ \frac{C(d,G)\,\kappa^2}{|P_\perp(\mathcal B_L)|},
\end{equation}
and in $d=4$ this is $O(L^{-2})$ since $|P_\perp(\mathcal B_L)|=8L(L+1)$.
\end{proposition}

\begin{proof}
By \eqref{eq:chi2_mgf_upper} with $\alpha=\alpha(L)$ it suffices to bound the block-level ratio
\[
D_{\chi^2}(\mu_L\|\nu_L^{(\alpha(L))})
=\frac{\mathbb E[e^{-2t_L X}]}{\mathbb E[e^{-t_L X}]^2}-1.
\]
Write $Y:=X-\mathbb E_{\nu_L^{(\alpha(L))}}X$.
Then
\[
\log\frac{\mathbb E[e^{-2t_L X}]}{\mathbb E[e^{-t_L X}]^2}
=
\log \mathbb E\big[e^{-2t_L Y}\big] - 2\log \mathbb E\big[e^{-t_L Y}\big]
\le \log \mathbb E\big[e^{-2t_L Y}\big],
\]
since $\log \mathbb E[e^{-t_L Y}]\ge 0$ by Jensen (recall $\mathbb E Y=0$ and $e^{-t_L y}$ is convex).
Apply Lemma~\ref{lem:mgf_centered_X} with $s=-2t_L$ (per Remark~\ref{rem:kappa_mgf_radius}) to obtain
\[
\log\frac{\mathbb E[e^{-2t_L X}]}{\mathbb E[e^{-t_L X}]^2}
\le \frac{C}{\beta^2}\,(2t_L)^2\,|P_\perp|
= \frac{4C\,\kappa^2}{|P_\perp|}.
\]
Let $u:=4C\kappa^2/|P_\perp|$.
Then
\[
D_{\chi^2}(\mu_L\|\nu_L^{(\alpha(L))})
\le e^{u}-1 \le u\,e^{u}.
\]
In $d=4$ one has $|P_\perp|\ge 16$ and $\kappa\le 1$, hence $u\le C/4$ and so $e^{u}\le e^{C/4}$.
Therefore
\[
D_{\chi^2}(\mu_L\|\nu_L^{(\alpha(L))})
\le C'(d,G)\,\frac{\kappa^2}{|P_\perp|}.
\]
Combine with \eqref{eq:chi2_mgf_upper} to conclude \eqref{eq:chi2_decay_alpha_v2}.

\end{proof}

\medskip
The bound \eqref{eq:chi2_decay_alpha_v2} is exactly the local non-toy input required by the $\chi^2$ transfer module in
Section~\ref{sec:L2_transfer} to compare the exact blocked single-plaquette marginal to the $\alpha$--reference marginal, with a correction of
order $|P_\perp|^{-1/2}$ in the spectral-radius majorant $\rho$ and order $|P_\perp|^{-1}$ in $\chi^2$.



\paragraph{KP-budgeted slack: converting a majorant entry depth into an exact entry depth}

The deterministic entry mechanism (Proposition~\ref{prop:mixing_bound} and Corollary~\ref{cor:det_entry_b2})
produces, for each microscopic input $W$, a finite depth $k$ such that the convolution-power reference
$W^{*4^k}$ has small activity at the identity,
\[
f^{\mathrm{ref}}_{4^k}:=W^{*4^k}(e)-1 \le f_{\mathrm{adm}}.
\]
To transfer this to an exact blocked plaquette marginal, one must control the \emph{normalization loss}
coming from the change of measure. The $\chi^2$ transfer bounds in
Section~\ref{sec:L2_transfer} are robust but can be overly conservative for controlling the single scalar
$w(e)-1$ that defines the KP entrance condition in the positive-type sector.
A sharper pointwise budget is available because the tilt factor is bounded by $1$.

\begin{proposition}[Tilt budget for the activity at the identity]\label{prop:tilt_budget_identity}
Let $w^{\mathrm{ref}}$ be a probability density on $G$ and let $\overline R:G\to[0,1]$ be measurable.
Define the tilted density
\[
w^{\mathrm{ex}}(g):=\frac{\overline R(g)}{\mathbb E_{w^{\mathrm{ref}}}[\overline R]}\,w^{\mathrm{ref}}(g),
\qquad
\mathbb E_{w^{\mathrm{ref}}}[\overline R]:=\int_G \overline R(g)\,w^{\mathrm{ref}}(g)\,dg.
\]
Then
\begin{equation}\label{eq:tilt_budget_identity}
w^{\mathrm{ex}}(e)-1
\le
\frac{w^{\mathrm{ref}}(e)-1}{\mathbb E_{w^{\mathrm{ref}}}[\overline R]}
+\Bigl(\frac{1}{\mathbb E_{w^{\mathrm{ref}}}[\overline R]}-1\Bigr).
\end{equation}
Equivalently, with $\delta_0:=\mathbb E_{w^{\mathrm{ref}}}[\overline R]^{-1}-1\ge 0$,
\begin{equation}\label{eq:tilt_budget_identity_delta}
w^{\mathrm{ex}}(e)-1 \le \delta_0 + (1+\delta_0)\,(w^{\mathrm{ref}}(e)-1).
\end{equation}
\end{proposition}

\begin{proof}
Since $\overline R(e)\le 1$ we have
$w^{\mathrm{ex}}(e)=\overline R(e)\,w^{\mathrm{ref}}(e)/\mathbb E_{w^{\mathrm{ref}}}[\overline R]\le w^{\mathrm{ref}}(e)/\mathbb E_{w^{\mathrm{ref}}}[\overline R]$.
Subtract $1$ and split $w^{\mathrm{ref}}(e)=1+(w^{\mathrm{ref}}(e)-1)$.
\end{proof}

\begin{corollary}[A two-line sufficient condition for KP entrance via a reference]\label{cor:tilt_budget_KP}
Let $f_{\mathrm{adm}}$ be the KP ceiling \eqref{eq:fadm_def}.
Assume $w^{\mathrm{ref}}(e)-1\le f_{\mathrm{ref}}$.
If the slack obeys
\begin{equation}\label{eq:slack_budget_condition}
\delta_0:=\mathbb E_{w^{\mathrm{ref}}}[\overline R]^{-1}-1
\le \frac{f_{\mathrm{adm}}-f_{\mathrm{ref}}}{1+f_{\mathrm{ref}}},
\end{equation}
then $w^{\mathrm{ex}}(e)-1\le f_{\mathrm{adm}}$, hence (in the positive-type sector) the exact step lies in the KP admissible window.
A simpler (slightly stronger) sufficient condition is
\[
w^{\mathrm{ref}}(e)-1\le \frac{f_{\mathrm{adm}}}{2}
\qquad\text{and}\qquad
\delta_0\le \frac{f_{\mathrm{adm}}}{2}.
\]
\end{corollary}

\begin{remark}[What this changes about the $\alpha$ schedule]\label{rem:kappa_budgeted}
Remark~\ref{rem:alpha_scaling} focused on keeping the slack \emph{nondegenerate} ($\mathbb E[R]=\Theta(1)$) as $L$ grows.
Corollary~\ref{cor:tilt_budget_KP} shows that for an actual entrance transfer into the KP window one needs a much tighter requirement:
the normalization loss must be on the same order as the KP ceiling, i.e.\ $\delta_0=O(f_{\mathrm{adm}})$.
On the thin block, using \eqref{eq:ER_alpha_exact_thin} and the weak-coupling scaling
\eqref{eq:zalpha_ratio} from Remark~\ref{rem:zbeta_scaling_G}, one has
$z(\alpha)/z(1)\sim \alpha^{-\dim(G)/2}$.
Thus the slack-stabilizing schedule $\alpha(L)=1-\kappa/|P_\perp|$ gives
\[
\delta_0
=
\Bigl(\frac{z(\alpha)}{z(1)}\Bigr)^{|P_\perp|}-1
\approx
e^{\frac{\dim(G)}{2}\,\kappa}-1.
\]
Therefore the KP-budgeted constraint \eqref{eq:slack_budget_condition} forces
\begin{equation}\label{eq:kappa_budget}
\kappa\ \lesssim\ \frac{2}{\dim(G)}\,\log\bigl(1+O(f_{\mathrm{adm}})\bigr)\ \approx\ O\bigl(f_{\mathrm{adm}}\bigr).
\end{equation}
In particular, in the baseline window $f_{\mathrm{adm}}\approx 1.668\times 10^{-3}$ one is in the regime $\kappa\sim 10^{-3}$.
(For larger $\dim(G)$ the constraint tightens further through the factor $\dim(G)$.)
\end{remark}








\begin{proposition}[Finite-depth slack/mixing bound for a slack-stabilizing reference (FD)]\label{assump:FD_slack_mixing}
Fix the $\alpha$--pivot reference family $\nu_L^{(\alpha)}$ and tilt $R_L^{(\alpha)}$ of \eqref{eq:alpha_reference_block},
with the $\alpha$--schedule $\alpha=\alpha(L)$ declared in the manuscript (e.g.\ $\alpha(L)=1-\kappa/L^2$ with $\kappa$ fixed).
Assume there exist constants $\beta_{\mathrm{FD}}<\infty$ and a budget $\delta_{\mathrm{budget}}>0$ such that for all $\beta\ge\beta_{\mathrm{FD}}$,
for all admissible boundary conditions, and for each block size $L$ that occurs in the finite-depth entrance comparison up to
depth $k_{\mathrm{enter}}$, the slack denominator is uniformly nondegenerate in the sense that
\begin{equation}\label{eq:FD_slack_bound}
\delta_{\mathrm{slack}}(L,\beta):=\Big(\mathbb E_{\nu_L^{(\alpha(L))}}\!\big[R_L^{(\alpha(L))}\big]\Big)^{-1}-1
=\Big(\frac{Z_L^{\mathrm{full}}}{Z_L^{(\alpha(L))}}\Big)^{-1}-1
\ \le\ \delta_{\mathrm{budget}},
\end{equation}
with $\delta_{\mathrm{budget}}$ small enough to close the tilt budget inequality \eqref{eq:tilt_budget_identity_delta} at the KP threshold $f_{\mathrm{adm}}$.

In addition, assume that the corresponding \emph{full $4$D} $\alpha$--reference boundary marginal satisfies the early-step mixing estimate
needed by the transfer module of Section~\ref{sec:mix_majorant} (spectral-radius / $\chi^2$ control) with constants uniform in $\beta$,
volume, and boundary conditions; a concrete sufficient form is Proposition~\ref{prop:TBmix_coeff_decay_reference} below.
\end{proposition}

\noindent\emph{Comment.} This assumption is the ``thick-block'' analogue of the thin-block surrogate used above: it isolates the remaining mixing input needed for the finite-depth entrance comparison. In the present manuscript, the coefficient-decay form Proposition~\ref{prop:TBmix_coeff_decay_reference} is discharged on the load-bearing dyadic path from two theorem-level ingredients: (i) the convolution-dominance/chessboard comparison proved in Proposition~\ref{prop:TBmix_chessboard_proved} by reflection positivity and bond-moving bookkeeping, and (ii) the one-plaquette coefficient bound recorded in Corollary~\ref{cor:TBmix_one_plaq_closed}. All remaining steps in the mass-gap module are independent of $N$ once these coefficient bounds are in place.



\begin{proposition}[Thick-block $\alpha$--reference mixing (crossover form)]\label{prop:TBmix_coeff_decay_reference}
Fix $G$ and fix $\kappa\in(0,\kappa_0(G)]$ and the schedule \eqref{eq:alpha_schedule_L}.
For each $L$ let $w_L^{(\alpha(L))}\,dg = (\Phi_L)_{\#}\nu_L^{(\alpha(L))}$ be the boundary holonomy marginal of the
(thick-block) $\alpha$--reference family.

Assume there exists $\beta_{\mathrm{mix}}\ge 1$ such that for all $\beta\ge \beta_{\mathrm{mix}}$, all $L$ in the finite-depth entrance window,
and every nontrivial irrep $R\neq \mathbf 1$,
\begin{equation}\label{eq:TBmix_coeff_decay}
\begin{aligned}
\lambda_R\!\bigl(w_L^{(\alpha(L))}\bigr)
&:=\int_{G} w_L^{(\alpha(L))}(g)\,\varphi_R(g)\,dg,\\
\bigl|\lambda_R\!\bigl(w_L^{(\alpha(L))}\bigr)\bigr|
&\le
\exp\!\Big(-|T_L|\,c_{\mathrm{A1}}(G,\rho)
\frac{C_2(R)}{\alpha(L)\beta+\sqrt{C_2(R)}}\Big),
\qquad |T_L|=L^2.
\end{aligned}
\end{equation}

\noindent\emph{Remark.} The bound \eqref{eq:TBmix_coeff_decay} interpolates between: (i) a heat-kernel regime $C_2(R)\le \beta^2$, where
$|\lambda_R|\le \exp(-\tfrac{c_{\mathrm{A1}}(G,\rho)}{2}\,\tfrac{|T_L|}{\beta}\,C_2(R))$, and (ii) a large-representation regime $C_2(R)\ge \beta^2$, where
$|\lambda_R|\le \exp(-\tfrac{c_{\mathrm{A1}}(G,\rho)}{2}\,|T_L|\,\sqrt{C_2(R)})$.
\end{proposition}




\begin{remark}[Roadmap for proving Proposition~\ref{prop:TBmix_coeff_decay_reference}]\label{rem:TBmix_roadmap}
The coefficient bound \eqref{eq:TBmix_coeff_decay} is equivalent to a quantitative area-law / heat-kernel type upper bound for a Wilson loop
in the thick $\alpha$--reference block, since
\[
\lambda_R\!\bigl(w_L^{(\alpha(L))}\bigr)=\mathbb E_{\nu_L^{(\alpha(L))}}\!\big[\varphi_R(\Phi_L)\big].
\]
In particular, it is \emph{strictly stronger} than what one gets from the purely abelian / Gaussianized $(dA)$--approximation:
in that linearized model a discrete Stokes identity collapses the surface flux to a boundary line integral, leading at most to
perimeter-scale fluctuations (Lemma~\ref{lem:TBmix_stokes_perimeter} below).

Two complementary closure approaches suggest themselves.
\begin{enumerate}[leftmargin=*]
\item[(A)] \emph{Nonabelian Gaussianization/CLT along the tile.}
Work on $\Omega_{\varepsilon_0}^{\parallel}$, write $\Phi_L$ as a surface-ordered product of plaquette holonomies based at the root in tree gauge,
and prove that the resulting random product behaves like the group heat kernel at time $\tau\asymp |T_L|/\beta$.
Analytically, one needs a quantitative \emph{area-scale nondegeneracy} in Lie-algebra coordinates that does not collapse to the boundary,
together with probabilistic control of the commutator/BCH corrections.

\item[(B)] \emph{Positivity-preserving decimation (bond-moving/Migdal--Kadanoff type comparison).}
Use reflection positivity / positivity of the character expansion to compare
$\mathbb E_{\nu_L^{(\alpha(L))}}[\varphi_R(\Phi_L)]$ to the same expectation in an analytically solvable
decimated reference model in which the boundary marginal is an explicit convolution power (hence exhibits the desired
$\exp\{-c\,C_2(R)|T_L|/\beta\}$ coefficient decay).
\end{enumerate}

The remainder of Section~\ref{sec:TBmix_reduction} records one explicit lemma chain of type (A) and isolates the precise analytic inputs
that would be required if one pursued Alternative~A in the thick-block geometry (the primary proof does not require these inputs). It also records an explicit sufficient closure approach of type (B)
in Section~\ref{sec:TBmix_chessboard} (chessboard/tiling estimate + one-plaquette coefficient control).
\end{remark}


\subsubsection{Thick-block mixing: explicit reduction to a variance bound and commutator control}\label{sec:TBmix_reduction}


\noindent\textbf{Signpost.} The variance/CLT reduction recorded in Proposition~\ref{prop:TBmix_var_lower} and Proposition~\ref{prop:B9_planA_reduction} is an optional side-branch. The load-bearing inputs used in the primary chain are proved in Subsubsections~\ref{sec:TBmix_chessboard}--\ref{sec:TBmix_routeC} (Lemmas~\ref{lem:BM_channel_decomp_G}, \ref{lem:BM_channel_chessboard_G}, \ref{lem:BM_tilt_form_G}, \ref{lem:BM_plaq_defect_L2} and Proposition~\ref{prop:BM_partition_correction}); these discharge $H3\to H6\to H2$ and feed Theorem~\ref{thm:h60_one_shot_entrance}.\\

\begin{definition}[Tile small-field event]\label{def:Omega_parallel}
For the planar $L\times L$ tile $T_L$ in the $(1,2)$--plane define the \emph{tile small-field event}
\begin{equation}\label{eq:Omega_parallel_def}
\Omega_{\varepsilon_0}^{\parallel}
:=
\bigl\{\,\|F_p\|\le \varepsilon_0\ \text{for all}\ p\in T_L\,\bigr\}.
\end{equation}
This is the minimal event required to make $F_p=\log U_p$ well-defined for all plaquettes appearing in the surface observable
$S_L^{\mathrm{plaq}}=\sum_{p\in T_L}F_p$.
\end{definition}


Proposition~\ref{prop:TBmix_coeff_decay_reference} is formulated directly in terms of normalized character coefficients of the thick-block
$\alpha$--reference boundary marginal. For completeness and referee navigation we spell out an explicit lemma chain
showing \emph{exactly} which analytic ingredients suffice to prove \eqref{eq:TBmix_coeff_decay}.

\smallskip
\noindent
Throughout this subsection we work conditionally on the small-curvature event $\Omega_{\varepsilon_0}^{\parallel}$ of Lemma~\ref{lem:weak_localization},
so that exponential coordinates are valid and curvature variables are uniformly bounded.

\begin{definition}[Linearized flux and plaquette-log surface sum]\label{def:surface_curvature_sum}
Let $T_L$ be the oriented $L\times L$ tile in the $(\mu,\nu)$--plane spanning the boundary loop $\partial T_L$ used in the definition
of $\Phi_L$ in Section~\ref{sec:local_tilt}. 

For each plaquette $p$ let $U_p\in G$ be its holonomy (based at the tree-gauge root). On the event $\Omega_{\varepsilon_0}^{\parallel}$,
the plaquette logarithm is well-defined and we write
\[
U_p=\exp(F_p),\qquad F_p\in\mathfrak g,\quad \|F_p\|\le \varepsilon_0.
\]
In addition, write the edge variables in exponential coordinates as $U_e=\exp(A_e)$ and let $d$ denote the cochain differential so that
$(dA)_p$ is the linearized curvature 2-form on plaquettes.

Define the $\mathfrak g$--valued sums
\[
S_L^{\mathrm{plaq}}:=\sum_{p\in T_L}F_p,\qquad
S_L^{\mathrm{lin}}:=\sum_{p\in T_L}(dA)_p,
\]
and for $v\in\mathfrak g$ set $S_L^{\mathrm{plaq}}(v):=\langle v,S_L^{\mathrm{plaq}}\rangle$ and
$S_L^{\mathrm{lin}}(v):=\langle v,S_L^{\mathrm{lin}}\rangle$.
\end{definition}


\begin{lemma}[Sub-Gaussian upper bound for the linearized flux $S_L^{\mathrm{lin}}$ from curvature BL]\label{lem:TBmix_subgaussian}
Assume Proposition~\ref{prop:curvature_BL}. There exists a constant $C_{\mathrm{BL}}=C_{\mathrm{BL}}(G)$ such that for every $v\in\mathfrak g$
and every $s\in\mathbb R$,
\begin{equation}\label{eq:TBmix_mgf_upper}
\log\,\mathbb E_{\nu_L^{(\alpha(L))}}\Big[\exp\big(s\,S_L^{\mathrm{lin}}(v)\big)\ \big|\ \Omega_{\varepsilon_0}^{\parallel}\Big]
\ \le\ C_{\mathrm{BL}}\frac{s^2}{\beta}\,|T_L|\,\|v\|^2.
\end{equation}
In particular, $S_L^{\mathrm{lin}}(v)$ is sub-Gaussian with variance proxy $\lesssim (\beta^{-1}|T_L|)\|v\|^2$, uniformly in volume and boundary conditions.
\end{lemma}

\begin{proof}
Take the finitely supported $\mathfrak g$--valued 2-form $J^{(v)}$ defined by $J^{(v)}_p=v$ for $p\in T_L$ and $J^{(v)}_p=0$ otherwise.
Then $\langle J^{(v)},dA\rangle=\sum_{p\in T_L}\langle v,(dA)_p\rangle=S_L^{\mathrm{lin}}(v)$.
Applying \eqref{eq:curvature_mgf} with $J=J^{(v)}$ and using $\langle J,\Pi_{\mathrm{ex}}J\rangle\le\|J\|_2^2=|T_L|\|v\|^2$
gives \eqref{eq:TBmix_mgf_upper} (absorbing $(2m_-)^{-1}$ into $C_{\mathrm{BL}}$).
\end{proof}




\begin{lemma}[Stokes collapse yields perimeter-scale variance in the abelianized Gaussian model]\label{lem:TBmix_stokes_perimeter}
Consider the Gaussian (abelianized) block measure on tree-gauge $1$--forms $A$ with density proportional to
$\exp\{-(\beta/2)\|dA\|_2^2\}$.
Then for the tile indicator $J^{(v)}$ from Lemma~\ref{lem:TBmix_subgaussian},
\[
S_L^{\mathrm{lin}}(v)=\langle J^{(v)},dA\rangle = \langle d^*J^{(v)},A\rangle,
\]
so $S_L^{\mathrm{lin}}(v)$ depends only on the boundary $1$--current $d^*J^{(v)}$ (supported on $\partial T_L$).
Moreover,
\[
\mathrm{Var}\big(S_L^{\mathrm{lin}}(v)\big)
=\frac{1}{\beta}\,\langle d^*J^{(v)},\Delta_1^{-1}d^*J^{(v)}\rangle
\ \le\ C_{\mathrm{per}}\,\frac{L}{\beta}\,\|v\|^2,
\]
with $C_{\mathrm{per}}$ depending only on the lattice geometry and the chosen norm on $\mathfrak g$.
In particular, an area-scale lower bound of the form $\mathrm{Var}\gtrsim |T_L|/\beta$ cannot hold for the linearized flux in this Gaussian model.
\end{lemma}

\begin{proof}
The identity $S_L^{\mathrm{lin}}(v)=\langle d^*J^{(v)},A\rangle$ is discrete integration by parts.
For the Gaussian measure, $\mathrm{Cov}(A)=\beta^{-1}\Delta_1^{-1}$, hence the variance formula.
Finally, standard Green function bounds for the (tree-gauge) $1$--form Laplacian in four dimensions imply
$(\Delta_1^{-1})(e,e')\le C(1+\mathrm{dist}(e,e'))^{-2}$, and since $d^*J^{(v)}$ is supported on a loop of length $\asymp L$
one obtains the stated $O(L)$ bound on $\langle d^*J^{(v)},\Delta_1^{-1}d^*J^{(v)}\rangle$ by a direct double-sum estimate.
\end{proof}


\begin{proposition}[B9: nondegenerate surface variance for the plaquette-log sum (via Plan A)]\label{prop:TBmix_var_lower}
Assume Proposition~\ref{ass:routeC_cut_local_surrogate} at scale $L$ and schedule $\alpha(L)$ (hence with parameters
$m_L$, $\Delta_X$, $c_{\Psi}$ and $\epsilon_L^{(C)}$).
Let $\Omega_{\varepsilon_0}^{\parallel}$ be the tile small-field event from Definition~\ref{def:Omega_parallel}.
Then there exist constants $\beta_{\mathrm{var}}=\beta_{\mathrm{var}}(G,\kappa,\varepsilon_0)<\infty$ and
$c_{\mathrm{var}}=c_{\mathrm{var}}(G,\kappa,\varepsilon_0)>0$ such that for all $\beta\ge \beta_{\mathrm{var}}$,
all $L$ in the finite-depth entrance window, all admissible boundary conditions, and all $v\in\mathfrak g$ with $\|v\|=1$,
\begin{equation}\label{eq:TBmix_var_lower}
\mathrm{Var}_{\nu_L^{(\alpha(L))}}\!\left(S_L^{\mathrm{plaq}}(v)\,\bigm|\,\Omega_{\varepsilon_0}^{\parallel}\right)
\ \ge\
c_{\mathrm{var}}\frac{|T_L|}{\beta}.
\end{equation}
Moreover one may take $c_{\mathrm{var}}=c_1/2$, where $c_1(G)$ is the planar one-plaquette constant from Lemma~\ref{lem:planar_one_plaq_var_lower}.
\end{proposition}

\begin{proof}
Fix $\beta$ and $L$ and abbreviate $\Omega:=\Omega_{\varepsilon_0}^{\parallel}$.
By the planar variance computation in Proposition~\ref{prop:B9_planA_reduction} (Step~1), using Lemma~\ref{lem:planar_one_plaq_var_lower}, one has
\[
\mathrm{Var}_{\nu_L^{\parallel}}\!\left(S_L^{\mathrm{plaq}}(v)\,\bigm|\,\Omega\right)\ \ge\ c_1\frac{|T_L|}{\beta}.
\]
Let $\nu_{\alpha}:=\nu_L^{(\alpha(L))}|_{\mathscr F_{\parallel}}$ and $\nu_{\parallel}:=\nu_L^{\parallel}|_{\mathscr F_{\parallel}}$.
Proposition~\ref{ass:routeC_cut_local_surrogate} and Lemma~\ref{lem:routeC_surrogate_implies_chi2_flatness} give the (unconditioned) $L^2$-flatness
\[
\delta:=\bigl\|\widetilde\Psi_L^{(\alpha)}-1\bigr\|_{L^2(\nu_{\parallel})}
\le \epsilon_L^{(C)}+\frac{t\,\Delta_X}{\beta\sqrt2\,c_{\Psi}}\sqrt{m_L}.
\]
Since $\Omega$ depends only on the in-plane plaquettes, weak-localization (union bound with the one-plaquette tail from Proposition~\ref{prop:curvature_BL})
yields $\nu_{\parallel}(\Omega)\ge 1-|T_L|e^{-c\beta}$, hence $\nu_{\parallel}(\Omega)\ge 1/2$ for $\beta$ large.
Applying Lemma~\ref{lem:chi2_from_L2_on_event} with $\nu=\nu_{\parallel}$, $\mu=\nu_{\alpha}$ and event $\Omega$ gives
\[
D_{\chi^2}\!\left(\nu_{\alpha}^{\Omega}\Big\|\nu_{\parallel}^{\Omega}\right)
\ \le\ \frac{4}{\nu_{\parallel}(\Omega)}\,\delta^2
\ \le\ 8\,\delta^2,
\]
provided $\delta\le \tfrac12\sqrt{\nu_{\parallel}(\Omega)}$, which holds for $\beta$ large since $\delta\to 0$ in the entrance window.
Choosing $\beta_{\mathrm{var}}$ large enough so that this $\chi^2$ divergence is at most the threshold required in
Proposition~\ref{prop:B9_planA_reduction}, the conclusion \eqref{eq:TBmix_var_lower} follows from Proposition~\ref{prop:B9_planA_reduction}
with $c_{\mathrm{var}}=c_1/2$.
\end{proof}


\begin{corollary}[B9 via the tilt influence bound (Plan~A without a cut surrogate)]\label{cor:TBmix_var_lower_influence}
Assume Proposition~\ref{ass:routeC_tilt_influence} at scale $L$ and schedule $\alpha(L)$.
Then the conclusion \eqref{eq:TBmix_var_lower} of Proposition~\ref{prop:TBmix_var_lower} holds
(with possibly different constants $\beta_{\mathrm{var}}$ and $c_{\mathrm{var}}$ depending only on $G,\kappa,\varepsilon_0$).
\end{corollary}

\begin{proof}
Fix $\beta$ and $L$ and abbreviate $\Omega:=\Omega_{\varepsilon_0}^{\parallel}$.
As in the proof of Proposition~\ref{prop:TBmix_var_lower}, one has the planar conditional lower bound
\[
\mathrm{Var}_{\nu_L^{\parallel}}\!\left(S_L^{\mathrm{plaq}}(v)\,\bigm|\,\Omega\right)\ \ge\ c_1\frac{|T_L|}{\beta}.
\]
Let $\nu_{\alpha}:=\nu_L^{(\alpha(L))}|_{\mathscr F_{\parallel}}$ and $\nu_{\parallel}:=\nu_L^{\parallel}|_{\mathscr F_{\parallel}}$.
By Corollary~\ref{cor:routeC_chi2_from_tilt_influences}, Proposition~\ref{ass:routeC_tilt_influence} implies
\[
\delta^2:=\bigl\|\widetilde\Psi_L^{(\alpha)}-1\bigr\|_{L^2(\nu_{\parallel})}^2
=
D_{\chi^2}\!\left(\nu_{\alpha}\Big\|\nu_{\parallel}\right)
\ \le\ \frac{|T_L|\,\Delta_{\Psi}^2}{2\beta^2}.
\]
As above, $\nu_{\parallel}(\Omega)\ge 1-|T_L|e^{-c\beta}$ and hence $\nu_{\parallel}(\Omega)\ge 1/2$ for $\beta$ large.
Applying Lemma~\ref{lem:chi2_from_L2_on_event} with $\nu=\nu_{\parallel}$, $\mu=\nu_{\alpha}$ and event $\Omega$ gives
\[
D_{\chi^2}\!\left(\nu_{\alpha}^{\Omega}\Big\|\nu_{\parallel}^{\Omega}\right)
\ \le\ \frac{4}{\nu_{\parallel}(\Omega)}\,\delta^2
\ \le\ 8\,\delta^2.
\]
In the entrance window $|T_L|\lesssim \beta$, this $\chi^2$ divergence tends to $0$ as $\beta\to\infty$, so for $\beta$ large it is below the threshold required in
Proposition~\ref{prop:B9_planA_reduction}. The conclusion \eqref{eq:TBmix_var_lower} then follows from Proposition~\ref{prop:B9_planA_reduction}
with $c_{\mathrm{var}}=c_1/2$.
\end{proof}



\begin{lemma}[Planar one-plaquette variance lower bound]\label{lem:planar_one_plaq_var_lower}
Fix $\varepsilon_0>0$ within the logarithm chart as in Proposition~\ref{prop:curvature_BL}.
There exist $\beta_1=\beta_1(G)<\infty$ and $c_1=c_1(G)>0$ such that for all $\beta\ge \beta_1$ and all $v\in\mathfrak g$ with $\|v\|=1$,
\begin{equation}\label{eq:planar_one_plaq_var_lower}
\mathrm{Var}_{\nu_1^{\parallel}}\!\Big(\langle v,F_{p_0}\rangle\ \Big|\ \|F_{p_0}\|\le \varepsilon_0\Big)
\ \ge\ \frac{c_1}{\beta}.
\end{equation}
\end{lemma}

\begin{proof}
Work in exponential coordinates $g=\exp(X)$ on the ball $\|X\|\le \varepsilon_0$, where $F=\log g=X$.
In these coordinates Haar measure has the form $dg=J(X)\,dX$ with a smooth Jacobian $J$ and $J(0)=1$.
By Taylor expansion of the Wilson action at the identity and compactness, there exist constants $m_-,m_+>0$ and $0<j_-\le j_+<\infty$ (depending only on $G$ and the chosen normalization) such that for all $\|X\|\le \varepsilon_0$,
\begin{equation}\label{eq:one_plaq_quadratic_sandwich}
\frac{m_-}{2}\|X\|^2 \ \le\ A(\exp X)\ \le\ \frac{m_+}{2}\|X\|^2,
\qquad
j_- \ \le\ J(X)\ \le\ j_+.
\end{equation}
Since $A(g)=A(g^{-1})$ and Haar is inversion-invariant, the conditional law on $\{\|X\|\le \varepsilon_0\}$ is symmetric under $X\mapsto -X$, hence
$\mathbb E[\langle v,X\rangle\mid \|X\|\le \varepsilon_0]=0$ and the conditional variance equals the conditional second moment.

Let $R:=1$ and take $\beta\ge \beta_1:=(R/\varepsilon_0)^2$, so that $\{\|X\|\le R/\sqrt{\beta}\}\subseteq \{\|X\|\le \varepsilon_0\}$.
Lower bound the numerator and upper bound the denominator using \eqref{eq:one_plaq_quadratic_sandwich}:
\[
\int_{\|X\|\le \varepsilon_0} \langle v,X\rangle^2 e^{-\beta A(\exp X)}J(X)\,dX
\ \ge\
j_- \int_{\|X\|\le R/\sqrt{\beta}} \langle v,X\rangle^2 e^{-(\beta m_+/2)\|X\|^2}\,dX,
\]
\[
\int_{\|X\|\le \varepsilon_0} e^{-\beta A(\exp X)}J(X)\,dX
\ \le\
j_+ \int_{\mathbb R^{\dim G}} e^{-(\beta m_-/2)\|X\|^2}\,dX.
\]
In the numerator substitute $X=Y/\sqrt{\beta}$ to obtain
\[
\begin{aligned}
\int_{\|X\|\le R/\sqrt{\beta}} \langle v,X\rangle^2 e^{-(\beta m_+/2)\|X\|^2}\,dX
&=
\beta^{-\frac{\dim G}{2}-1}
\int_{\|Y\|\le R} \langle v,Y\rangle^2 e^{-(m_+/2)\|Y\|^2}\,dY\\
&=:\beta^{-\frac{\dim G}{2}-1} C_{\mathrm{num}}(G).
\end{aligned}
\]
where $C_{\mathrm{num}}(G)>0$ by positivity of the integrand.
The Gaussian integral in the denominator equals $\beta^{-\dim G/2}C_{\mathrm{den}}(G)$ with $C_{\mathrm{den}}(G)<\infty$.
Therefore the conditional second moment (and hence variance) is at least
\[
\frac{j_- C_{\mathrm{num}}(G)}{j_+ C_{\mathrm{den}}(G)}\cdot \frac{1}{\beta}.
\]
Set $c_1:=j_- C_{\mathrm{num}}(G)/(j_+ C_{\mathrm{den}}(G))>0$.
\end{proof}

\begin{proposition}[Plan A reduction of Proposition~\ref{prop:TBmix_var_lower}]\label{prop:B9_planA_reduction}
Let $\mathscr F_{\parallel}:=\sigma(\{U_p\}_{p\in T_L})$ be the $\sigma$--algebra generated by the in-plane plaquette variables on the tile $T_L$.
Write $\nu_{\parallel}$ for the purely planar reference on $T_L$ (Lemma~\ref{lem:planar_exact_convolution}) and
$\nu_{\alpha}$ for the thick-block $\alpha(L)$--reference $\nu_L^{(\alpha(L))}$ restricted to $\mathscr F_{\parallel}$.
Let $\nu_{\parallel}^{\Omega}$ and $\nu_{\alpha}^{\Omega}$ denote these laws further conditioned on $\Omega_{\varepsilon_0}^{\parallel}$.

Assume that for some $\varepsilon_{\ast}\in(0,1]$ one has the uniform in-plane $\chi^2$ control
\begin{equation}\label{eq:B9_planA_chi2_assump}
D_{\chi^2}\big(\nu_{\alpha}^{\Omega}\,\big\|\,\nu_{\parallel}^{\Omega}\big)\le \varepsilon_{\ast},
\end{equation}
uniformly in admissible boundary conditions.
Then there exists a constant $C_{\mathrm{A}}=C_{\mathrm{A}}(G)$ such that for every $v\in\mathfrak g$ with $\|v\|=1$,
\begin{equation}\label{eq:B9_planA_var_transfer}
\mathrm{Var}_{\nu_L^{(\alpha(L))}}\!\Big(S_L^{\mathrm{plaq}}(v)\ \Big|\ \Omega_{\varepsilon_0}^{\parallel}\Big)
\ \ge\
\mathrm{Var}_{\nu_{\parallel}}\!\Big(S_L^{\mathrm{plaq}}(v)\ \Big|\ \Omega_{\varepsilon_0}^{\parallel}\Big)
\ -\ 
C_{\mathrm{A}}\,\frac{|T_L|}{\beta}\,\sqrt{\varepsilon_{\ast}}.
\end{equation}

In particular, Lemma~\ref{lem:planar_one_plaq_var_lower} provides constants $\beta_1(G)$ and $c_1(G)>0$ such that for all $\beta\ge \beta_1(G)$,
\[
\mathrm{Var}_{\nu_{\parallel}}\!\Big(S_L^{\mathrm{plaq}}(v)\ \Big|\ \Omega_{\varepsilon_0}^{\parallel}\Big)\ge \frac{c_1(G)}{\beta}\,|T_L|.
\]
Hence if additionally $\sqrt{\varepsilon_{\ast}}\le c_1(G)/(2C_{\mathrm{A}}(G))$ then Proposition~\ref{prop:TBmix_var_lower} holds with $c_{\mathrm{var}}=c_1(G)/2$.
\end{proposition}

\begin{proof}
We work under the conditional measures $\nu_{\alpha}^{\Omega}$ and $\nu_{\parallel}^{\Omega}$ on $\mathscr F_{\parallel}$.

\smallskip
\noindent\emph{Step 1: planar variance is exactly extensive.}
By Lemma~\ref{lem:planar_exact_convolution}, in spanning-tree gauge the plaquette variables $\{U_p\}_{p\in T_L}$ are independent under $\nu_{\parallel}$.
The event $\Omega_{\varepsilon_0}^{\parallel}$ is a product event over plaquettes in this gauge, hence independence is preserved under conditioning.
Moreover $A(g)=A(g^{-1})$ and Haar is inversion-invariant, so each $F_p=\log U_p$ is centered under $\nu_{\parallel}^{\Omega}$ and
\[
\mathrm{Var}_{\nu_{\parallel}}\!\Big(S_L^{\mathrm{plaq}}(v)\ \Big|\ \Omega_{\varepsilon_0}^{\parallel}\Big)
=
\sum_{p\in T_L}\mathrm{Var}_{\nu_1^{\parallel}}\!\Big(\langle v,F_{p}\rangle\ \Big|\ \|F_{p}\|\le \varepsilon_0\Big)
=
|T_L|\cdot \mathrm{Var}_{\nu_1^{\parallel}}\!\Big(\langle v,F_{p_0}\rangle\ \Big|\ \|F_{p_0}\|\le \varepsilon_0\Big).
\]
In particular, Lemma~\ref{lem:planar_one_plaq_var_lower} implies
\begin{equation}\label{eq:B9_planA_planar_var_lb}
\mathrm{Var}_{\nu_{\parallel}}\!\Big(S_L^{\mathrm{plaq}}(v)\ \Big|\ \Omega_{\varepsilon_0}^{\parallel}\Big)\ \ge\ \frac{c_1(G)}{\beta}\,|T_L|.
\end{equation}

\smallskip
\noindent\emph{Step 2: $L^4$ size of the planar surface sum.}
On $\Omega_{\varepsilon_0}^{\parallel}$ one has $|\langle v,F_p\rangle|\le \varepsilon_0$ for each $p$, hence for $X_p:=\langle v,F_p\rangle$,
$\mathbb E[(X_p)^4]\le \varepsilon_0^2\,\mathbb E[(X_p)^2]$.
By independence and the standard fourth-moment expansion for sums of centered independent variables,
\[
\mathbb E_{\nu_{\parallel}^{\Omega}}\!\Big[\big(S_L^{\mathrm{plaq}}(v)\big)^4\Big]
\ \le\
C\Big(|T_L|\,\mathbb E[X_{p_0}^4]+|T_L|^2\,\mathbb E[X_{p_0}^2]^2\Big)
\ \le\
C'(G)\Big(\frac{|T_L|}{\beta^2}+\frac{|T_L|^2}{\beta^2}\Big),
\]
using the curvature-scale upper bound $\mathbb E[X_{p_0}^2]\lesssim 1/\beta$ from Proposition~\ref{prop:curvature_BL}.
Thus
\begin{equation}\label{eq:B9_planA_L4_bound}
\|S_L^{\mathrm{plaq}}(v)\|_{L^4(\nu_{\parallel}^{\Omega})}^2\ \le\ C''(G)\,\frac{|T_L|}{\beta}.
\end{equation}
Similarly $\|S_L^{\mathrm{plaq}}(v)\|_{L^2(\nu_{\parallel}^{\Omega})}^2=\mathrm{Var}_{\nu_{\parallel}^{\Omega}}(S_L^{\mathrm{plaq}}(v))\lesssim |T_L|/\beta$.

\smallskip
\noindent\emph{Step 3: $\chi^2$ transfer of variance.}
Apply Lemma~\ref{lem:chi2_transfer_variance} with $\mu=\nu_{\alpha}^{\Omega}$, $\nu=\nu_{\parallel}^{\Omega}$, $F=S_L^{\mathrm{plaq}}(v)$ and
$\varepsilon=D_{\chi^2}(\nu_{\alpha}^{\Omega}\|\nu_{\parallel}^{\Omega})\le \varepsilon_{\ast}$.
Using \eqref{eq:B9_planA_L4_bound} and the corresponding $L^2$ bound gives
\[
\mathrm{Var}_{\nu_{\alpha}^{\Omega}}\big(S_L^{\mathrm{plaq}}(v)\big)
\ \ge\
\mathrm{Var}_{\nu_{\parallel}^{\Omega}}\big(S_L^{\mathrm{plaq}}(v)\big)
-
C_{\mathrm{A}}(G)\,\frac{|T_L|}{\beta}\,\sqrt{\varepsilon_{\ast}},
\]
which is \eqref{eq:B9_planA_var_transfer}.

Finally combine with \eqref{eq:B9_planA_planar_var_lb}.
If $\sqrt{\varepsilon_{\ast}}\le c_1(G)/(2C_{\mathrm{A}}(G))$ and $\beta\ge \beta_1(G)$, then
\[
\mathrm{Var}_{\nu_L^{(\alpha(L))}}\!\Big(S_L^{\mathrm{plaq}}(v)\ \Big|\ \Omega_{\varepsilon_0}^{\parallel}\Big)
\ge \frac{c_1(G)}{2\beta}\,|T_L|,
\]
i.e.\ Proposition~\ref{prop:TBmix_var_lower} holds with $c_{\mathrm{var}}=c_1(G)/2$.
\end{proof}

\begin{remark}[What the $\chi^2$ hypothesis is]\label{rem:B9_planA_chi2_is_tilt}
By Lemma~\ref{lem:routeC_effective_planar_tilt}, the in-plane marginal of the thick-block $\alpha$--reference is a tilt of the planar reference:
$\nu_{\alpha}|_{\mathscr F_{\parallel}}=\widetilde\Psi_L^{(\alpha)}\cdot \nu_{\parallel}|_{\mathscr F_{\parallel}}$.
Therefore the hypothesis \eqref{eq:B9_planA_chi2_assump} is equivalently the $L^2$ flatness condition
\[
\|\widetilde\Psi_L^{(\alpha)}-1\|_{L^2(\nu_{\parallel}^{\Omega})}^2\le \varepsilon_{\ast}.
\]
In the thin-block leaf geometry one has $\widetilde\Psi_L^{(\alpha)}\equiv 1$ exactly (Lemma~\ref{lem:thin_block_tilt_constant}),
whereas in genuine thick blocks this is precisely the ``effective tilt control'' problem of Main chain.
\end{remark}




\begin{lemma}[Surface representation of $\Phi_L$ (lattice nonabelian Stokes; deterministic form)]\label{lem:TBmix_surface_log}
In tree gauge, there exists an ordering of plaquettes $p_1,\dots,p_{|T_L|}$ in the tile $T_L$ such that the boundary holonomy satisfies the
exact surface-ordered product identity
\begin{equation}\label{eq:TBmix_surface_product}
\Phi_L\ =\ \overrightarrow{\prod_{i=1}^{|T_L|}}\,U_{p_i},
\end{equation}
where $U_{p_i}\in G$ is the plaquette holonomy around $p_i$ (with the induced orientation from $T_L$) and the arrow indicates the chosen surface order.

On the event $\Omega_{\varepsilon_0}^{\parallel}$ one may write $U_{p_i}=\exp(F_{p_i})$ with $\|F_{p_i}\|\le\varepsilon_0$.
Consequently, there exists a $\mathfrak g$--valued error term $E_L$ such that
\begin{equation}\label{eq:TBmix_surface_log}
\Phi_L\ =\ \exp\!\big(S_L^{\mathrm{plaq}}+E_L\big),
\qquad S_L^{\mathrm{plaq}}=\sum_{p\in T_L}F_p,
\end{equation}
with the crude deterministic BCH bound
\begin{equation}\label{eq:TBmix_surface_log_error}
\|E_L\|\ \le\ C_{\mathrm{stokes}}\,|T_L|^2\,\varepsilon_0^2,
\end{equation}
for a constant $C_{\mathrm{stokes}}=C_{\mathrm{stokes}}(G)$ depending only on the compact group $G$ and the chosen norm on $\mathfrak g$.
\end{lemma}

\begin{proof}
The exact product identity \eqref{eq:TBmix_surface_product} is a standard lattice nonabelian Stokes statement:
enumerate the plaquettes of $T_L$ in an order compatible with a spanning tree of the dual tile so that when multiplying the plaquette boundary words,
interior edges cancel adjacently (tree gauge is used only to make the resulting surface order canonical; the identity itself is gauge invariant).

On $\Omega_{\varepsilon_0}^{\parallel}$, exponential coordinates are valid and each plaquette holonomy can be written as $U_{p_i}=\exp(F_{p_i})$
with $\|F_{p_i}\|\le\varepsilon_0$.
Write $X_i:=F_{p_i}$ and consider the partial products $P_k:=\overrightarrow{\prod_{i=1}^k}\exp(X_i)$.
For $\varepsilon_0$ fixed below the injectivity radius, the Baker--Campbell--Hausdorff formula implies a uniform estimate
\[
\big\|\log(\exp(X)\exp(Y))-(X+Y)\big\|\ \le\ C_{\mathrm{BCH}}\,\|X\|\,\|Y\|
\qquad(\|X\|,\|Y\|\le\varepsilon_0),
\]
with $C_{\mathrm{BCH}}=C_{\mathrm{BCH}}(G)$.
Iterating this bound over $k=1,\dots,|T_L|$ and using $\|X_i\|\le\varepsilon_0$ yields the crude estimate
$\|\log P_{|T_L|}-\sum_{i}X_i\|\le C_{\mathrm{BCH}}|T_L|^2\varepsilon_0^2$,
which is \eqref{eq:TBmix_surface_log_error} (absorbing constants into $C_{\mathrm{stokes}}$).
\end{proof}


\begin{remark}
Lemma~\ref{lem:TBmix_surface_log} is the precise form of the “surface representation’’ step in Remark~\ref{rem:TBmix_roadmap}.
Its proof is purely deterministic (surface ordering + BCH commutator bookkeeping) once one works in a normal chart
and uses $\|F_p\|\le \varepsilon_0$ on $\Omega_{\varepsilon_0}^{\parallel}$.
\end{remark}

\begin{lemma}[Heat-kernel coefficient decay from Gaussian surface fluctuations]\label{lem:TBmix_gaussian_to_coeff}
Assume Lemma~\ref{lem:TBmix_surface_log} and Proposition~\ref{prop:TBmix_var_lower}.
Assume in addition that the one-dimensional marginals $S_L^{\mathrm{plaq}}(v)$ are \emph{uniformly strongly log-concave} on $\Omega_{\varepsilon_0}^{\parallel}$ in the sense that
their conditional densities (with respect to Lebesgue measure on $\mathbb R$) have second derivative bounded below by
$-C(\beta/|T_L|)$ and above by $-c(\beta/|T_L|)$, with constants uniform in $v$ on the unit sphere.\footnote{This hypothesis is a quantitative
“Gaussian comparison’’ statement for the surface sum; it is expected to follow from the curvature domination mechanism plus a nondegeneracy
estimate such as \eqref{eq:TBmix_var_lower}, but it is recorded separately here to make the required analytic input explicit.}
Then there exist constants $\beta_{\mathrm{hk}}(G)$ and $c_{\mathrm{mix}}(G)>0$ such that for all $\beta\ge\beta_{\mathrm{hk}}$ and all $L$
in the finite-depth entrance window,
\[
\lambda_R\!\bigl(w_L^{(\alpha(L))}\bigr)
=\mathbb E_{\nu_L^{(\alpha(L))}}\!\big[\varphi_R(\Phi_L)\big]
\ \le\
\exp\!\Big(-c_{\mathrm{mix}}\frac{C_2(R)}{\beta}\,|T_L|\Big)
\qquad\forall R\neq\mathbf 1,
\]
i.e. Proposition~\ref{prop:TBmix_coeff_decay_reference} holds.
\end{lemma}

\begin{remark}[What is genuinely new]
Proposition~\ref{prop:curvature_BL} supplies the \emph{upper} sub-Gaussian control \eqref{eq:TBmix_mgf_upper} for the \emph{linearized} flux
$S_L^{\mathrm{lin}}(v)$ at essentially no additional cost.
However, Lemma~\ref{lem:TBmix_stokes_perimeter} shows that in the corresponding Gaussian/abelian model this linearized flux has only
perimeter-scale variance, so an area-law coefficient decay cannot come from the linear sector alone.

The genuinely new quantitative content needed for thick-block mixing is therefore twofold:
(i) an area-scale variance noncollapse for the \emph{plaquette-log} sum $S_L^{\mathrm{plaq}}(v)$ (Proposition~\ref{prop:TBmix_var_lower}),
and (ii) probabilistic control of the commutator/BCH error $E_L$ in Lemma~\ref{lem:TBmix_surface_log} strong enough to transfer this
nondegeneracy to the group-valued holonomy $\Phi_L$.
The additional strong log-concavity hypothesis in Lemma~\ref{lem:TBmix_gaussian_to_coeff} is a convenient way to package the needed
Gaussian comparison input.
\end{remark}

\paragraph{Alternative approach (B): convolution dominance (bond-moving) and a one-plaquette coefficient bound}\label{sec:TBmix_chessboard}

\noindent\textbf{Load-bearing (Main chain).} The inputs proved in Subsubsections~\ref{sec:TBmix_chessboard}--\ref{sec:TBmix_routeC} are used to discharge Proposition~\ref{prop:BM_dyadic}, yielding $H3$ for $G$ and feeding the main chain $H3\to H6\to H2\to$ Theorem~\ref{thm:h60_one_shot_entrance}.


The variance/CLT approach above isolates two genuinely nontrivial analytic inputs (area-scale noncollapse and commutator control).
There is, however, a qualitatively different approach that aims to bypass Lie-algebra Gaussianization:
compare the thick-block boundary marginal to an analytically solvable \emph{convolution power}.
The structural point is that for central densities the normalized character coefficients $\lambda_R(\cdot)$ are precisely the Fourier coefficients,
and convolution powers multiply coefficients exactly.

\begin{lemma}[Convolution powers multiply coefficients]\label{lem:conv_coeff_mult}
Let $w\,dg$ be a central probability density on $G$ and let $w^{*n}$ denote its $n$-fold convolution.
Then for every irrep $R$ one has
\[
\lambda_R(w^{*n})=\lambda_R(w)^n,
\qquad\text{and hence}\qquad
\bigl|\lambda_R(w^{*n})\bigr|=\bigl|\lambda_R(w)\bigr|^n.
\]
\end{lemma}
\begin{proof}
This is immediate from Peter--Weyl: convolution corresponds to multiplication of Fourier coefficients.
\end{proof}


\begin{definition}[Coefficient domination order on central densities]\label{def:coeff_domination}
Let $u\,dg$ and $v\,dg$ be central probability densities on $G$ and write their normalized character coefficients as
$\lambda_R(u)=\int_G u(g)\,\overline{\varphi_R(g)}\,dg$ as in Section~\ref{sec:mix_majorant}.
We write
\[
u \preceq_{\mathrm{coeff}} v
\]
if for every irrep $R$ one has the coefficient domination (in modulus)
\[
\big|\lambda_R(u)\big|\le \big|\lambda_R(v)\big|.
\]
If $u$ and $v$ are of positive type, then $\lambda_R(u),\lambda_R(v)\ge 0$ for all $R$ and the modulus may be dropped.
\end{definition}

\begin{lemma}[Positive type $\Rightarrow$ nonnegative coefficients]\label{lem:pos_type_nonneg_coeff}
Let $u\,dg$ be a central probability density of positive type. Then $\lambda_R(u)\ge 0$ for all irreps $R$ and
$0\le \lambda_R(u)\le 1$.
\end{lemma}
\begin{proof}
For a positive-type $u$, the convolution operator $T_u$ on $L^2(G)$ is positive semidefinite.
In Peter--Weyl coordinates one has $\widehat u(R)=\lambda_R(u)I_{d_R}$ for central $u$ (Section~\ref{sec:mix_majorant}),
so positivity of $\widehat u(R)$ forces $\lambda_R(u)\ge 0$.
Moreover, since $u\,dg$ is a probability measure and $|\varphi_R(g)|\le 1$ for all $g\in G$,
\[
0\le \lambda_R(u)=\Big|\int_G \varphi_R(g)\,u(g)\,dg\Big|\le \int_G |\varphi_R(g)|\,u(g)\,dg\le 1.
\]
\end{proof}

\begin{lemma}[Heat kernel coefficients and semigroup]\label{lem:heat_kernel_coeff}
Let $k_t$ denote the (bi-invariant) heat kernel on $G$ at time $t>0$ corresponding to the Casimir/Laplacian normalization used for $C_2$,
i.e.\ the unique central probability density whose normalized character coefficients satisfy
\[
\lambda_R(k_t)=e^{-tC_2(R)}
\qquad\text{for every irrep }R,
\]
equivalently
$
k_t(g)=\sum_R d_R\,e^{-tC_2(R)}\,\chi_R(g)
$
in Peter--Weyl form.
Then $\{k_t\}_{t\ge 0}$ is a convolution semigroup: $k_t*k_s=k_{t+s}$ and hence $k_t^{*n}=k_{nt}$.
\end{lemma}
\begin{proof}
In Peter--Weyl, convolution corresponds to multiplication of Fourier coefficients.
Thus for each irrep $R$,
\[
\lambda_R(k_t*k_s)=\lambda_R(k_t)\lambda_R(k_s)=e^{-tC_2(R)}e^{-sC_2(R)}=e^{-(t+s)C_2(R)}=\lambda_R(k_{t+s}).
\]
Since the characters form an orthonormal basis for central $L^2(G)$ functions and $\lambda_{\mathbf 1}(k_t)=1$,
this identifies $k_t*k_s$ with $k_{t+s}$ and shows $\int_G k_t\,dg=1$.
Iterating gives $k_t^{*n}=k_{nt}$.
\end{proof}


\begin{remark}[One-plaquette bounds and heat-kernel domination]\label{rem:plaq_as_heatkernel_dom}
For a compact Lie group $G$, a Gaussian-in-$C_2/\beta$ one-plaquette coefficient bound implies coefficient domination by a heat kernel:
if $|\lambda_R(w_1^{(\alpha)})|\le e^{-c\,C_2(R)/\beta}$ uniformly in $\alpha\in[1/2,1]$, then
$w_1^{(\alpha)} \preceq_{\mathrm{coeff}} k_{c/\beta}$ in the sense of Definition~\ref{def:coeff_domination}.

In the present manuscript we use instead the crossover bound \eqref{eq:TBmix_one_plaq_coeff}.
It still yields a heat-kernel style decay in the low-representation regime $C_2(R)\le \beta^2$ via \eqref{eq:TBmix_one_plaq_lowrep_heatkernel},
and a stronger $\sqrt{C_2}$ decay outside that regime.
\end{remark}



\begin{lemma}[Planar (2D) exact convolution power]\label{lem:planar_exact_convolution}
Consider the purely planar reference on a simply connected $L\times L$ tile $T_L$ with plaquette weight $W$ and no out-of-plane plaquettes.
Let $w_L^{\parallel}\,dg=(\Phi_L)_{\#}\nu_L$ be the planar boundary marginal around $T_L$, and let $w_1^{\parallel}$ be the corresponding $1\times 1$ marginal.
Then one may fix a spanning-tree gauge in which the plaquette holonomies $\{U_p\}_{p\in T_L}$ become independent with marginal density proportional to $W$,
and the boundary holonomy is the ordered product of these plaquette variables.
Consequently,
\[
w_L^{\parallel}=(w_1^{\parallel})^{*|T_L|},
\qquad\text{and in particular}\qquad
\lambda_R(w_L^{\parallel})=\lambda_R(w_1^{\parallel})^{|T_L|}\quad\text{for all }R.
\]
\end{lemma}
\begin{proof}
Choose a rooted spanning tree of the tile that contains all boundary edges and all but one edge of each plaquette.
Gauge-fix all tree edges to the identity. The remaining (non-tree) edge variables can be chosen one-per-plaquette, and in this gauge each plaquette holonomy $U_p$
is exactly its unique non-tree edge variable.
Since the Haar measure factorizes over the remaining edges, the planar reference density $\prod_{p\in T_L}W(U_p)$ becomes a product measure over plaquettes.
Moreover, because the boundary edges lie in the tree, the boundary holonomy $\Phi_L$ is the ordered product of the plaquette holonomies.
Thus $\Phi_L$ is the product of $|T_L|$ i.i.d.\ $G$-valued variables with density $w_1^{\parallel}$, and its law is the convolution power $(w_1^{\parallel})^{*|T_L|}$.
\end{proof}


\begin{lemma}[Schur insertion bound (uniform deterministic boundary data)]\label{lem:schur_insertion_bound}
Let $w\,dg$ be a central probability density on $G$ and let $U_1,\dots,U_N$ be independent $G$--valued random variables with law $w\,dg$.
Fix deterministic elements $g_0,\dots,g_N\in G$ and define the interleaved product
\[
\Phi:= g_0\,U_1\,g_1\,U_2\cdots g_{N-1}\,U_N\,g_N.
\]
Then for every irreducible unitary representation $R$,
\begin{equation}\label{eq:schur_insertion_expect}
\mathbb E\!\big[\varphi_R(\Phi)\big]
=
\lambda_R(w)^N\cdot \frac{1}{d_R}\,\mathrm{Tr}\!\Big(\prod_{k=0}^N \pi_R(g_k)\Big),
\end{equation}
and in particular
\begin{equation}\label{eq:schur_insertion_bound}
\big|\mathbb E[\varphi_R(\Phi)]\big|\le \lambda_R(w)^N.
\end{equation}
\end{lemma}

\begin{proof}
Let
\[
M_R:=\int_G \pi_R(u)\,w(u)\,du.
\]
Since $w$ is central and Haar is conjugation-invariant,
\[
\pi_R(h)\,M_R\,\pi_R(h)^{-1}=\int_G \pi_R(huh^{-1})\,w(u)\,du=M_R
\qquad\text{for all }h\in G,
\]
so $M_R$ commutes with $\pi_R(h)$ for all $h\in G$.
By Schur's lemma $M_R=c_R\,\mathrm{Id}$ for some scalar $c_R$.
Taking traces yields
\[
c_R=\frac{1}{d_R}\,\mathrm{Tr}(M_R)=\int_G \varphi_R(u)\,w(u)\,du=\lambda_R(w).
\]
By independence,
\[
\mathbb E[\pi_R(\Phi)]
=
\pi_R(g_0)\,M_R\,\pi_R(g_1)\,M_R\cdots M_R\,\pi_R(g_N)
=
\lambda_R(w)^N\,\prod_{k=0}^N \pi_R(g_k).
\]
Taking normalized traces gives \eqref{eq:schur_insertion_expect}.
Finally, each $\pi_R(g_k)$ is unitary, hence $|\mathrm{Tr}(U)|\le d_R$ for unitary $U$, yielding \eqref{eq:schur_insertion_bound}.
\end{proof}

\begin{remark}[Planar power bounds are uniform in boundary conditions]\label{rem:planar_power_uniform_eta}
In a planar reference model on a simply connected tile $T_L$, a spanning-tree gauge expresses the boundary holonomy $\Phi_L$
as an interleaved product of the independent plaquette variables (each with marginal $w_1^{\parallel}$) and deterministic
group insertions coming from the exterior boundary data.
Lemma~\ref{lem:schur_insertion_bound} then yields the uniform estimate
\[
\big|\lambda_R\big((\Phi_L)_\#\nu_L^{\parallel}(\cdot\mid\eta)\big)\big|
\le \lambda_R(w_1^{\parallel})^{|T_L|}
\qquad\text{for all boundary conditions }\eta.
\]
Thus boundary data can modify only the trace factor in \eqref{eq:schur_insertion_expect}; it cannot increase the coefficient magnitude.
\end{remark}


\begin{proposition}[Dyadic bond-moving / decimation inequality (local form)]\label{prop:BM_dyadic_local}
There exists $\beta_{\mathrm{BM}}=\beta_{\mathrm{BM}}(G)$ such that for all $\beta\ge\beta_{\mathrm{BM}}$ and all $\alpha\in[1/2,1]$ the following holds.
For each dyadic $L=2^n$ let $w_L^{(\alpha)}\,dg=(\Phi_L)_{\#}\nu_{L}^{(\alpha)}$ denote the thick-block $\alpha$--reference boundary marginal around the $L\times L$ tile $T_L$.
Assume that for every dyadic $L$ and every nontrivial irrep $R\neq\mathbf 1$,
\begin{equation}\label{eq:BM_dyadic_step_template}
\big|\lambda_R\!\bigl(w_{2L}^{(\alpha)}\bigr)\big|
\ \le\
\big|\lambda_R\!\bigl(w_{L}^{(\alpha)}\bigr)\big|^{4}.
\end{equation}
\noindent Equivalently, by Definition~\ref{def:coeff_domination} and Lemma~\ref{lem:conv_coeff_mult},
\[
w_{2L}^{(\alpha)} \preceq_{\mathrm{coeff}} (w_L^{(\alpha)})^{*4}.
\]
\end{proposition}

\paragraph*{Hardening note: a direct $\chi^2$ approach via influences of $\widetilde{\Psi}$}

For Plan~A in Section~\ref{sec:TBmix_reduction}, it is enough to control the $\chi^2$ divergence
\[
D_{\chi^2}\!\left(\nu_L^{(\alpha)}\big|_{\mathscr F_{\parallel}}\Big\| \nu_L^{\parallel}\right)
=
\big\|\widetilde{\Psi}_L^{(\alpha)}-1\big\|_{L^2(\nu_L^{\parallel})}^2,
\]
and the cut-local surrogate mechanism in Proposition~\ref{ass:routeC_cut_local_surrogate}
is one way to do so.
Independently, one can bound the same quantity by an Efron--Stein influence estimate on the
\emph{in-plane plaquette coordinates} under the planar reference (tree gauge / i.i.d.\ plaquettes).

\begin{lemma}[Conditional Efron--Stein bound for localization]\label{lem:routeC_conditional_efronstein}
Let $(\sigma_a)_{a=1}^{n}$ be independent coordinates under a product probability space, and let
$\mathscr F_{\mathrm{cut}}=\sigma(\sigma_1,\dots,\sigma_m)$ with $m\le n$.
For any $Y=Y(\sigma_1,\dots,\sigma_n)\in L^2$ and independent resamplings $\sigma_a'\sim\sigma_a$,
let $Y^{(a)}$ denote $Y$ with $\sigma_a$ replaced by $\sigma_a'$.
Then
\begin{equation}\label{eq:routeC_conditional_es}
\big\|Y-E[Y\mid \mathscr F_{\mathrm{cut}}]\big\|_{L^2}^2
=
E\!\left[\mathrm{Var}(Y\mid \mathscr F_{\mathrm{cut}})\right]
\;\le\;
\frac12\sum_{a=m+1}^{n}\big\|Y-Y^{(a)}\big\|_{L^2}^2.
\end{equation}

\begin{proof}
Fix the cut coordinates $(\sigma_1,\dots,\sigma_m)$.
Conditionally on $\mathscr F_{\mathrm{cut}}$, the remaining coordinates $(\sigma_{m+1},\dots,\sigma_n)$ remain independent.
Applying the standard Efron--Stein inequality to the conditional law gives
\[
\mathrm{Var}(Y\mid \mathscr F_{\mathrm{cut}})
\;\le\;
\frac12\sum_{a=m+1}^{n}E\!\left[(Y-Y^{(a)})^2\mid \mathscr F_{\mathrm{cut}}\right].
\]
Taking expectation of both sides yields \eqref{eq:routeC_conditional_es}.
\end{proof}

\end{lemma}

\begin{proposition}[Main-chain per-plaquette influence bound for the normalized effective tilt]\label{ass:routeC_tilt_influence}
Under the planar reference $\nu_L^{\parallel}$ in tree gauge, realize the in-plane plaquette variables as independent coordinates $(\sigma_a)_{a=1}^{|T_L|}$.
The normalized effective tilt $\widetilde{\Psi}_L^{(\alpha)}=\widetilde{\Psi}_L^{(\alpha)}(\sigma_1,\dots,\sigma_{|T_L|})$ satisfies the per-plaquette influence bound
\begin{equation}\label{eq:routeC_tilt_influence_bound}
\big\|\widetilde{\Psi}_L^{(\alpha)}-\big(\widetilde{\Psi}_L^{(\alpha)}\big)^{(a)}\big\|_{L^2(\nu_L^{\parallel})}
\;\le\;
\frac{\Delta_{\Psi}}{\beta}
\qquad\text{for all }a\in\{1,\dots,|T_L|\},
\end{equation}
with a constant $\Delta_{\Psi}$ independent of $L$ and uniform in all admissible boundary conditions.
\end{proposition}


\begin{corollary}[Sufficient influence bound for Plan~A $\chi^2$-flatness]\label{cor:routeC_chi2_from_tilt_influences}
Assume Proposition~\ref{ass:routeC_tilt_influence}.
Then
\begin{equation}\label{eq:routeC_chi2_flatness_from_influence}
\big\|\widetilde{\Psi}_L^{(\alpha)}-1\big\|_{L^2(\nu_L^{\parallel})}^2
\;=\;
D_{\chi^2}\!\left(\nu_L^{(\alpha)}\big|_{\mathscr F_{\parallel}}\Big\| \nu_L^{\parallel}\right)
\;\le\;
\frac{|T_L|\,\Delta_{\Psi}^2}{2\beta^2}
\;=\;
\frac{L^2\,\Delta_{\Psi}^2}{2\beta^2}.
\end{equation}
In particular, in the entrance window $L\lesssim\sqrt\beta$ this yields $D_{\chi^2}\lesssim \beta^{-1}$.

\begin{proof}
Since $\widetilde{\Psi}_L^{(\alpha)}$ is a Radon--Nikod\'ym derivative, $E_{\nu_L^{\parallel}}[\widetilde{\Psi}_L^{(\alpha)}]=1$ and hence
$\|\widetilde{\Psi}_L^{(\alpha)}-1\|_2^2=\mathrm{Var}_{\nu_L^{\parallel}}(\widetilde{\Psi}_L^{(\alpha)})$.
Apply Lemma~\ref{lem:routeC_conditional_efronstein} with $m=0$ (trivial $\mathscr F_{\mathrm{cut}}$) and $Y=\widetilde{\Psi}_L^{(\alpha)}$ to obtain
\[
\mathrm{Var}_{\nu_L^{\parallel}}(\widetilde{\Psi}_L^{(\alpha)})
\;\le\;
\frac12\sum_{a=1}^{|T_L|}\big\|\widetilde{\Psi}_L^{(\alpha)}-(\widetilde{\Psi}_L^{(\alpha)})^{(a)}\big\|_2^2
\;\le\;
\frac{|T_L|\,\Delta_{\Psi}^2}{2\beta^2},
\]
which is \eqref{eq:routeC_chi2_flatness_from_influence}.
\end{proof}

\end{corollary}

\begin{remark}[Thin-block case]\label{rem:routeC_tilt_influence_thin_block}
In the thin-block model of Lemma~\ref{lem:thin_block_tilt_constant}, $\widetilde{\Psi}_L^{(\alpha)}\equiv 1$,
so \eqref{eq:routeC_tilt_influence_bound} holds with $\Delta_{\Psi}=0$ and the $\chi^2$-flatness is exact.
\end{remark}

\begin{lemma}[Surrogate approximation $\Rightarrow$ a tilt influence bound]\label{lem:routeC_tilt_influence_from_surrogate}
Assume Proposition~\ref{ass:routeC_cut_local_surrogate}.
Realize the in-plane plaquette field under $\nu_L^{\parallel}$ in tree gauge as independent coordinates
$(\sigma_a)_{a=1}^{|T_L|}$, and assume without loss that the cut $\sigma$-algebra
$\mathscr F_{\mathrm{cut}}$ is generated by $(\sigma_1,\dots,\sigma_{m_L})$.
Then for each coordinate $a\in\{1,\dots,|T_L|\}$ and the corresponding independent resampling $\sigma_a'$, one has
\begin{equation}\label{eq:routeC_tilt_influence_from_surrogate}
\big\|\widetilde{\Psi}_L^{(\alpha)}-\big(\widetilde{\Psi}_L^{(\alpha)}\big)^{(a)}\big\|_{L^2(\nu_L^{\parallel})}
\ \le\
2\,\epsilon_L^{(C)}\ +\ \frac{t\,\Delta_X}{\beta\,c_{\Psi}}\,\mathbf 1_{\{a\le m_L\}}.
\end{equation}
In particular, if $\epsilon_L^{(C)}\le \Delta'/\beta$ for some constant $\Delta'$ independent of $L$, then the global influence target
\eqref{eq:routeC_tilt_influence_bound} holds with
\begin{equation}\label{eq:routeC_tilt_influence_constant_from_surrogate}
\Delta_{\Psi}=2\Delta'+\frac{t\,\Delta_X}{c_{\Psi}}.
\end{equation}
\end{lemma}

\begin{proof}
Write $\widetilde\Psi:=\widetilde{\Psi}_L^{(\alpha)}$ and $\widetilde\Psi_{\mathrm{cut}}$ for the surrogate from \eqref{eq:routeC_surrogate_def}.
By the triangle inequality and invariance of the $L^2(\nu_L^{\parallel})$ norm under a single-coordinate resampling,
\[
\begin{aligned}
\|\widetilde\Psi-\widetilde\Psi^{(a)}\|_2
\le{}&\|\widetilde\Psi-\widetilde\Psi_{\mathrm{cut}}\|_2
+\|\widetilde\Psi_{\mathrm{cut}}-(\widetilde\Psi_{\mathrm{cut}})^{(a)}\|_2
+\|(\widetilde\Psi_{\mathrm{cut}})^{(a)}-\widetilde\Psi^{(a)}\|_2\\
\le{}&2\,\|\widetilde\Psi-\widetilde\Psi_{\mathrm{cut}}\|_2
+\|\widetilde\Psi_{\mathrm{cut}}-(\widetilde\Psi_{\mathrm{cut}})^{(a)}\|_2.
\end{aligned}
\]
The first term is bounded by $2\epsilon_L^{(C)}$ by \eqref{eq:routeC_surrogate_L2}.
For the second term, note that $\widetilde\Psi_{\mathrm{cut}}=e^{-tX_{\mathrm{cut}}}/E[e^{-tX_{\mathrm{cut}}}]$ and the denominator is nonrandom.
If $a>m_L$ then $X_{\mathrm{cut}}$ does not depend on $\sigma_a$ and the difference vanishes.
If $a\le m_L$, then using $|e^{-t x}-e^{-t y}|\le t|x-y|$ and \eqref{eq:routeC_surrogate_mean_lower},
\[
\|\widetilde\Psi_{\mathrm{cut}}-(\widetilde\Psi_{\mathrm{cut}})^{(a)}\|_2
=\frac{\|e^{-tX_{\mathrm{cut}}}-e^{-tX_{\mathrm{cut}}^{(a)}}\|_2}{E[e^{-tX_{\mathrm{cut}}}]}
\le \frac{t}{c_{\Psi}}\,\|X_{\mathrm{cut}}-X_{\mathrm{cut}}^{(a)}\|_2
\le \frac{t\,\Delta_X}{\beta\,c_{\Psi}}
\]
by \eqref{eq:routeC_cut_influence_L2}.
Combining these bounds gives \eqref{eq:routeC_tilt_influence_from_surrogate}; the final statement is immediate.
\end{proof}

\begin{corollary}[Influences outside the cut $\Rightarrow$ $L^2$ cut-localization]\label{cor:routeC_L2_localization_from_outside_influences}
In the setting of Lemma~\ref{lem:routeC_conditional_efronstein}, assume
\begin{equation}\label{eq:routeC_outside_influence_bound}
\|Y-Y^{(a)}\|_{L^2}\le \frac{\Delta}{\beta}\qquad\text{for all }a\in\{m+1,\dots,n\}.
\end{equation}
Then
\begin{equation}\label{eq:routeC_L2_localization_bound}
\big\|Y-E[Y\mid\mathscr F_{\mathrm{cut}}]\big\|_{L^2}
\le
\frac{\Delta}{\beta\sqrt2}\,\sqrt{n-m}.
\end{equation}
In particular, if $Y=\widetilde\Psi_L^{(\alpha)}$ and $n=|T_L|$, this gives a fully explicit sufficient condition for the existence of an $\mathscr F_{\mathrm{cut}}$-measurable $L^2$ approximation
\emph{without} committing to the exponential form \eqref{eq:routeC_surrogate_def}.
\end{corollary}

\begin{proof}
By Lemma~\ref{lem:routeC_conditional_efronstein} and \eqref{eq:routeC_outside_influence_bound},
\[
\big\|Y-E[Y\mid\mathscr F_{\mathrm{cut}}]\big\|_{2}^2
\le \frac12\sum_{a=m+1}^{n}\|Y-Y^{(a)}\|_2^2
\le \frac12\,(n-m)\,\frac{\Delta^2}{\beta^2}.
\]
Taking square roots yields \eqref{eq:routeC_L2_localization_bound}.
\end{proof}

\medskip
\noindent
The influence bound \eqref{eq:routeC_tilt_influence_bound} is not an additional hypothesis in the main proof chain.
It is recorded only as a particularly sharp \emph{sufficient} target for discharging the Plan~A $\chi^2$ hypothesis
without invoking the cut-local exponential surrogate \eqref{eq:routeC_surrogate_def}.
Lemma~\ref{lem:routeC_tilt_influence_from_surrogate} shows how it can be derived from Proposition~\ref{ass:routeC_cut_local_surrogate}
when the surrogate error $\epsilon_L^{(C)}$ improves to $O(\beta^{-1})$.


\begin{proposition}[Dyadic bond-moving inequality on the theorem-level dyadic path]\label{prop:BM_dyadic}
Fix $\beta\ge\beta_{\mathrm{BM}}$ and $\alpha\in[1/2,1]$.
Then the dyadic bond-moving inequality holds with the explicit partition-correction factor
\begin{equation}\label{eq:BM_dyadic_with_correction}
\lambda_R\!\bigl(w^{(\alpha)}_{2L}\bigr)\;\le\;\exp\!\Big(\frac{C}{L\beta}\Big)\,\lambda_R\!\bigl(w^{(\alpha)}_{L}\bigr)^4,
\end{equation}
for every dyadic $L$ and every nontrivial irrep $R\neq\mathbf 1$, where $C=C(G)$ is the partition-correction constant of Proposition~\ref{prop:BM_partition_correction}.
In particular, after absorbing the finite number of initial dyadic correction factors into the universal weak-window constants,
\begin{equation}\label{eq:BM_dyadic_step}
\big|\lambda_R\!\bigl(w^{(\alpha)}_{2L}\bigr)\big|
\;\le\;
\big|\lambda_R\!\bigl(w^{(\alpha)}_{L}\bigr)\big|^{4}
\end{equation}
for the dyadic scales used on the Route~1 load-bearing path.

\begin{center}
\fbox{\begin{minipage}{0.94\textwidth}
\textbf{Theorem-level ingredients used in the proof.}
\begin{itemize}[leftmargin=*]
\item Lemma~\ref{lem:BM_channel_decomp_G}: channel decomposition on the cut links.
\item Lemma~\ref{lem:BM_channel_chessboard_G}: channel chessboard inequality (reflection positivity).
\item Proposition~\ref{prop:BM_partition_correction}: partition correction under the $\alpha$-schedule.
\end{itemize}
\end{minipage}}
\end{center}
\end{proposition}

\begin{proof}
By Lemma~\ref{lem:BM_channel_decomp_G} and Lemma~\ref{lem:BM_channel_chessboard_G},
\[
Z^{(2L)}_{R}=\sum_c Z^{(2L)}_{R,c}\le \sum_c \bigl(Z^{(L)}_{R,c}\bigr)^4.
\]
Since $x\mapsto x^4$ is convex on $\mathbb R_{\ge 0}$,
\begin{equation}\label{eq:BM_recombine}
\sum_c \bigl(Z^{(L)}_{R,c}\bigr)^4 \;\le\;\Big(\sum_c Z^{(L)}_{R,c}\Big)^4=\bigl(Z^{(L)}_{R}\bigr)^4.
\end{equation}
Divide by the partition functions and use Proposition~\ref{prop:BM_partition_correction}:
\[
\lambda_R(w^{(\alpha)}_{2L})=\frac{Z^{(2L)}_{R}}{Z^{(2L)}}\le 
\frac{(Z^{(L)}_{R})^4}{(Z^{(L)})^4}\,\exp\!\Big(\frac{C}{L\beta}\Big)
=\exp\!\Big(\frac{C}{L\beta}\Big)\,\lambda_R(w^{(\alpha)}_{L})^4,
\]
which is \eqref{eq:BM_dyadic_with_correction}.
The finite dyadic overhead remark following Proposition~\ref{prop:BM_partition_correction} then yields the clean step form \eqref{eq:BM_dyadic_step} on the scales used later in Route~1.
\end{proof}

\begin{lemma}[Corrected dyadic iteration with explicit overhead]\label{lem:BM_dyadic_iter_corrected}
Assume the one-step dyadic bond-moving inequality in the explicit form
\begin{equation}\label{eq:BM_dyadic_step_corrected_local}
\big|\lambda_R(w_{2L}^{(\alpha)})\big|
\le
\exp\!\Big(\frac{C_{\mathrm{pc}}(G)\,\kappa_{\mathrm{pc}}}{L\beta}\Big)
\big|\lambda_R(w_L^{(\alpha)})\big|^{4}
\end{equation}
for every dyadic $L=2^m$ and every nontrivial irrep $R\neq\mathbf 1$.
Then for every dyadic $L=2^n$,
\begin{equation}\label{eq:BM_dyadic_power_corrected}
\big|\lambda_R(w_{L}^{(\alpha)})\big|
\le
\exp\!\Big(\frac{2C_{\mathrm{pc}}(G)\,\kappa_{\mathrm{pc}}}{7\beta}\,(L^2-L^{-1})\Big)
\big|\lambda_R(w_1^{(\alpha)})\big|^{L^2}.
\end{equation}
\end{lemma}

\begin{proof}
Write $E_n$ for the minimal nonnegative number such that
\[
\big|\lambda_R(w_{2^n}^{(\alpha)})\big|
\le
\exp(E_n/\beta)\,\big|\lambda_R(w_1^{(\alpha)})\big|^{4^n}.
\]
Then $E_0=0$, and \eqref{eq:BM_dyadic_step_corrected_local} gives the recurrence
\[
E_{n+1}=4E_n+\frac{C_{\mathrm{pc}}(G)\,\kappa_{\mathrm{pc}}}{2^n}.
\]
Dividing by $4^{n+1}$ and summing the resulting geometric series yields
\[
\frac{E_n}{4^n}
=\sum_{m=0}^{n-1}\frac{C_{\mathrm{pc}}(G)\,\kappa_{\mathrm{pc}}}{2^m4^{m+1}}
=\frac{C_{\mathrm{pc}}(G)\,\kappa_{\mathrm{pc}}}{4}\sum_{m=0}^{n-1}8^{-m}
=\frac{2C_{\mathrm{pc}}(G)\,\kappa_{\mathrm{pc}}}{7}\,(1-8^{-n}).
\]
Hence
\[
E_n=\frac{2C_{\mathrm{pc}}(G)\,\kappa_{\mathrm{pc}}}{7}\,4^n(1-8^{-n})
=\frac{2C_{\mathrm{pc}}(G)\,\kappa_{\mathrm{pc}}}{7}\,(4^n-2^{-n}),
\]
which is exactly \eqref{eq:BM_dyadic_power_corrected} because $L=2^n$ and $|T_L|=L^2=4^n$.
\end{proof}

% [v123 edit: removed status remark; replaced by discharge summary box in Proposition~\ref{prop:BM_dyadic}]

\paragraph{Channel decomposition and channel-wise chessboard for a compact group}\label{sec:BM_channel_closure_G}

We now discharge items (1)--(2) in Proposition~\ref{prop:BM_dyadic} for a compact group $G$, using standard reflection positivity.
Throughout this subsection, fix a dyadic $L$ and an irreducible unitary representation $R$ of $G$.
\begin{definition}[Dyadic cut graph and channel basis]\label{def:BM_cut_graph_channels}
For the dyadic split of the $2L\times 2L$ tile into four $L\times L$ quadrants, let $\Gamma_{\mathrm{cut}}=\Gamma_{\mathrm{cut}}(L)$ denote
the internal cut graph consisting of the in-plane links lying on the two midlines (together with their cut vertices).
Write $E_{\mathrm{cut}}$ and $V_{\mathrm{cut}}$ for its edge and vertex sets. The gauge group $G^{V_{\mathrm{cut}}}$ acts on
$G^{E_{\mathrm{cut}}}$ by the usual lattice gauge transformations.

Define the cut Hilbert space
\[
\mathcal H_{\mathrm{cut}}:=L^2\!\big(G^{E_{\mathrm{cut}}}\big)^{G^{V_{\mathrm{cut}}}},
\]
the closed subspace of gauge-invariant square-integrable functions of the cut-link variables.
A \emph{channel} $c$ is a choice of:
\begin{itemize}
\item an irrep label $\lambda_e\in\mathrm{Irr}(G)$ for each cut edge $e\in E_{\mathrm{cut}}$;
\item an intertwiner label $\iota_v$ at each cut vertex $v\in V_{\mathrm{cut}}$ choosing an element of an orthonormal basis of the invariant space
\[
\mathrm{Inv}\!\Big(\bigotimes_{e\ni v} V_{\lambda_e}^{\sigma(v,e)}\Big),
\qquad
V_{\lambda}^{+}:=V_{\lambda},\ \ V_{\lambda}^{-}:=V_{\lambda}^{\ast},
\]
where $\sigma(v,e)\in\{+,-\}$ is determined by the orientation of $e$ relative to $v$ (outgoing $=+$, incoming $=-$).
\end{itemize}
The resulting spin-network functions $\psi_c\in\mathcal H_{\mathrm{cut}}$ form a countable orthonormal basis of $\mathcal H_{\mathrm{cut}}$. We write $\mathcal C_L$ for the set of all such channels.
\end{definition}

\begin{lemma}[Positive-type factorization of normalized characters]\label{lem:BM_char_pos_type_G}
Let $\rho$ be a unitary irrep of a compact group $G$ with dimension $d_\rho$ and normalized character
$\varphi_\rho(g)=d_\rho^{-1}\chi_\rho(g)$.
Then for all $g,h\in G$,
\begin{equation}\label{eq:BM_char_pos_type}
\varphi_\rho(g^{-1}h)=\frac{1}{d_\rho}\sum_{a,b=1}^{d_\rho}\overline{\rho(g)_{ab}}\,\rho(h)_{ab}.
\end{equation}
In particular, the kernel $K_\rho(g,h):=\varphi_\rho(g^{-1}h)$ is positive semidefinite.
\end{lemma}

\begin{proof}
Since $\rho$ is unitary, $\rho(g^{-1})=\rho(g)^\ast$.
Thus $\rho(g^{-1}h)=\rho(g)^\ast\rho(h)$ and
\[
\chi_\rho(g^{-1}h)=\mathrm{Tr}\big(\rho(g)^\ast\rho(h)\big)=\sum_{a,b}\overline{\rho(g)_{ab}}\,\rho(h)_{ab}.
\]
Divide by $d_\rho$ to obtain \eqref{eq:BM_char_pos_type}. Positivity follows because \eqref{eq:BM_char_pos_type} is a Gram representation.
\end{proof}

\begin{proposition}[Reflection positivity for the thick-block reference]\label{ass:BM_RP_G}
The thick-block $\alpha$--reference measure $\nu_{2L}^{(\alpha)}$ is reflection positive with respect to reflections in the two midlines
used in the dyadic split, in the standard Osterwalder--Seiler sense.
\end{proposition}

\begin{lemma}[Channel decomposition at the cut links (Item (1) in Proposition~\ref{prop:BM_dyadic})]\label{lem:BM_channel_decomp_G}
Assume \ref{ass:BM_RP_G}. Then there exist nonnegative numbers $\{Z^{(2L)}_{R,c}\}_{c\in\mathcal C_L}$ and
$\{Z^{(L)}_{R,c}\}_{c\in\mathcal C_L}$ such that
\[
Z^{(2L)}_{R}=\sum_{c\in\mathcal C_L} Z^{(2L)}_{R,c},\qquad
Z^{(L)}_{R}=\sum_{c\in\mathcal C_L} Z^{(L)}_{R,c},
\]
and $Z^{(2L)}_{R,c}\ge 0$, $Z^{(L)}_{R,c}\ge 0$ for all $c$.
\end{lemma}

\begin{proof}[Proof (explicit Hilbert-space expansion)]
Fix the dyadic split and write $\Theta$ for the composition of the two midline reflections (so $\Theta$ maps the lower-left quadrant to the upper-right).
Reflection positivity implies that the sesquilinear form
\[
\langle F,G\rangle_{\mathrm{OS}}:=\int (\Theta F)\,G\,d\nu^{(\alpha)}_{2L}
\]
is positive semidefinite on the algebra $\mathcal A_{++}$ of observables supported in the lower-left quadrant together with its boundary cut links.

By Lemma~\ref{lem:BM_char_pos_type_G}, the normalized character insertion on the $2L$ boundary can be written as a finite sum of
matrix-coefficient products across the reflection interface.
Concretely, there exists a finite family of cut-dependent functions $F_{R,ab}\in\mathcal A_{++}$ (built from quadrant plaquette weights and a boundary
matrix coefficient of $\rho_R$) such that
\begin{equation}\label{eq:BM_Z_as_sum_of_norms}
Z^{(2L)}_{R}=\sum_{a,b=1}^{d_R}\langle F_{R,ab},F_{R,ab}\rangle_{\mathrm{OS}}.
\end{equation}
Each inner product is nonnegative by reflection positivity, so the right-hand side is a sum of nonnegative terms.

Now expand each $F_{R,ab}$ in the cut Hilbert space basis of Definition~\ref{def:BM_cut_graph_channels}.
More precisely: the restriction of $F_{R,ab}$ to the cut link variables defines an element of $\mathcal H_{\mathrm{cut}}$, hence admits an expansion
$F_{R,ab}=\sum_{c\in\mathcal C_L} f_{R,ab}(c)\,\psi_c$ with $\sum_c |f_{R,ab}(c)|^2<\infty$.
By Parseval,
\[
\langle F_{R,ab},F_{R,ab}\rangle_{\mathrm{OS}}
=\sum_{c\in\mathcal C_L} |f_{R,ab}(c)|^2.
\]
Define
\[
Z^{(2L)}_{R,c}:=\sum_{a,b=1}^{d_R}|f_{R,ab}(c)|^2\ \ge\ 0.
\]
Summing over $c$ and using \eqref{eq:BM_Z_as_sum_of_norms} yields $Z^{(2L)}_{R}=\sum_c Z^{(2L)}_{R,c}$.

The same construction applied at scale $L$ (with the corresponding dyadic cut graph in the $L$-tile reference geometry) yields the decomposition
$Z^{(L)}_{R}=\sum_c Z^{(L)}_{R,c}$ with $Z^{(L)}_{R,c}\ge 0$.
\end{proof}

\begin{lemma}[Channel-wise chessboard estimate (Item (2) in Proposition~\ref{prop:BM_dyadic})]\label{lem:BM_channel_chessboard_G}
Assume \ref{ass:BM_RP_G}. With the channel pieces defined in Lemma~\ref{lem:BM_channel_decomp_G}, for every channel $c\in\mathcal C_L$,
\[
Z^{(2L)}_{R,c}\ \le\ \bigl(Z^{(L)}_{R,c}\bigr)^4.
\]
\end{lemma}

\begin{proof}[Proof (two-step RP / Cauchy--Schwarz)]
Fix $c$ and fix the basis element $\psi_c$ on the cut graph. By construction in Lemma~\ref{lem:BM_channel_decomp_G},
$Z^{(2L)}_{R,c}$ is a finite sum of Osterwalder--Schrader norms of observables supported in a fundamental quadrant with boundary channel $\psi_c$.
Apply reflection positivity across the vertical midline to double the quadrant to a half-tile, and then across the horizontal midline to double
to the full $2L$ tile. At each doubling step, Cauchy--Schwarz bounds the doubled norm by the product of the norms on the two halves.
Because the channel insertion $\psi_c$ is reflection-even (spin-network basis elements are real class functions on each cut edge and gauge-invariant at vertices),
the same channel label is preserved under reflections.
After two doublings one obtains precisely the fourth power bound.
\end{proof}

\begin{corollary}[Discharge of items (1)--(2) and closure of \texttt{BM\_dyadic}]\label{cor:BM_dyadic_G_closed}
Assume \ref{ass:BM_RP_G}. Then items (1)--(2) of Proposition~\ref{prop:BM_dyadic} hold for $G$.
Since item (3) is discharged in the strengthened branch (Lemma~\ref{lem:BM_tilt_form_G} + Proposition~\ref{prop:BM_partition_correction}),
Proposition~\ref{prop:BM_dyadic_local} holds for $G$ up to the explicit correction of \eqref{eq:BM_dyadic_with_correction}.
\end{corollary}







\paragraph{Partition correction inequality: an explicit attack}\label{sec:BM_partition_correction}

The analytic input (3) in Proposition~\ref{prop:BM_dyadic} is a statement about how the thick-block partition function scales under dyadic
gluing. The purpose of this subsection is to reduce \eqref{eq:BM_partition_correction_proved} to a cut-local finite-dimensional tilt bound and then
to bound that tilt explicitly in weak coupling.

\begin{definition}[Dyadic cut set and cut-adjacent plaquettes]\label{def:BM_cut_set}
Let $T_{2L}$ be the $2L\times 2L$ tile in the $(\mu,\nu)$--plane and split it into four $L\times L$ quadrants by the coordinate midlines.
Let $E_{\mathrm{cut}}=E_{\mathrm{cut}}(L)$ be the set of in-plane links lying on the two midlines (the internal cuts).
Let $P_{\mathrm{cut}}$ denote the set of \emph{out-of-plane} plaquettes in the thick block incident to at least one link in $E_{\mathrm{cut}}$.
In $d=4$ one has $|E_{\mathrm{cut}}|\asymp L$ and
\begin{equation}\label{eq:BM_Pcut_count}
|P_{\mathrm{cut}}|\ \le\ C_{\mathrm{cut}}\,L,
\end{equation}
since each cut link lies in at most $2(d-2)=4$ out-of-plane plaquettes.
\end{definition}

\begin{lemma}[Gluing without extra cross weights is supermultiplicative]\label{lem:BM_glue_supermult}
Let $Z_{\mathrm{glue}}^{(2L)}$ be the partition function of the geometry obtained by gluing four $L$--quadrant blocks along $E_{\mathrm{cut}}$
(identifying the cut links) \emph{without introducing any additional plaquette weights beyond those already present in the quadrants}.
Then
\begin{equation}\label{eq:BM_glue_supermult}
Z_{\mathrm{glue}}^{(2L)} \ \ge\ (Z^{(L)})^4.
\end{equation}
\end{lemma}
\begin{proof}
Integrate out interior variables of each quadrant, leaving a nonnegative effective weight $F(U_{E_{\mathrm{cut}}})$ on the cut links.
By symmetry the same $F$ appears for each quadrant, and with Haar-product measure on $U_{E_{\mathrm{cut}}}$ normalized to $1$ one has
\[
Z^{(L)}=\int F\,dU_{E_{\mathrm{cut}}},
\qquad
Z_{\mathrm{glue}}^{(2L)}=\int F^4\,dU_{E_{\mathrm{cut}}}.
\]
Since $x\mapsto x^4$ is convex on $\mathbb R_{\ge 0}$, Jensen gives $\int F^4\ge (\int F)^4$, which is \eqref{eq:BM_glue_supermult}.
\end{proof}


\begin{lemma}[Verification of the cut-local tilt form for the dyadic reference]\label{lem:BM_tilt_form_G}
Fix $G$ and dyadic $L$.
In the strengthened dyadic construction, the $2L$ thick-block reference partition function is obtained from the glued-quadrant geometry
of Lemma~\ref{lem:BM_glue_supermult} by adding \emph{only} the cut-adjacent plaquette weights on $P_{\mathrm{cut}}$
(with exponent $\theta\in(0,1]$), i.e. the density ratio is exactly
\begin{equation}\label{eq:BM_tilt_form_G}
\mathsf D^{(2L)}(U)
=\mathsf D^{(2L)}_{\mathrm{glue}}(U)\cdot \prod_{p\in P_{\mathrm{cut}}}\widehat W_\beta(U_p)^{\theta}.
\end{equation}
Consequently the partition functions satisfy the exact identity
\begin{equation}\label{eq:BM_partition_tilt_form_G}
Z^{(2L)} \;=\; Z_{\mathrm{glue}}^{(2L)}\,
\mathbb E_{\nu_{\mathrm{glue}}}\Big[\prod_{p\in P_{\mathrm{cut}}}\widehat W_\beta(U_p)^{\theta}\Big],
\end{equation}
which is the hypothesis \eqref{eq:BM_partition_tilt_form} of Proposition~\ref{prop:BM_partition_correction}.
Moreover, under the strengthened schedule of Remark~\ref{rem:BM_theta_schedule} (taking $\theta=1-\alpha_{\mathrm{pc}}(2L,\beta)$),
one has $\beta\theta\lesssim 1/L^2$ as required in \eqref{eq:BM_theta_small}.
\end{lemma}

\begin{proof}
Equation \eqref{eq:BM_tilt_form_G} is the defining relationship between the glued-quadrant dyadic reference and the true $2L$ reference in the strengthened construction: the only additional factors beyond those present in the quadrants are the cut-adjacent plaquette weights on
$P_{\mathrm{cut}}$, inserted with exponent $\theta$.
Integrating \eqref{eq:BM_tilt_form_G} against Haar measure and dividing by $Z_{\mathrm{glue}}^{(2L)}$ yields
\eqref{eq:BM_partition_tilt_form_G}.
The schedule statement follows from Remark~\ref{rem:BM_theta_schedule}.
\end{proof}

\begin{lemma}[Weak-coupling plaquette defect bound]\label{lem:BM_plaq_defect_mean}
Fix a unitary representation $\pi$ and set $A(g):=1-\frac{1}{\dim(\pi)}\Re\mathrm{Tr}\,\pi(g)$ as in \eqref{eq:RN_factor_beta}.
As above, we may assume $\pi$ is faithful after replacing $G$ by $G_{\mathrm{eff}}:=G/\ker(\pi)$.
There exist $\beta_{\mathrm{plaq}}>0$ and $C_{\mathrm{plaq}}<\infty$ (depending only on $G$ and $\pi$) such that for every weak-coupling
$\alpha$--reference measure of Wilson class on a finite block (including the glued geometry above) and every plaquette $p$ in that block,
\begin{equation}\label{eq:BM_plaq_defect_mean}
\mathbb E[A(U_p)] \ \le\ \frac{C_{\mathrm{plaq}}}{\beta}
\qquad\text{for all }\beta\ge \beta_{\mathrm{plaq}},\ \alpha\in[1/2,1].
\end{equation}
\end{lemma}
\begin{proof}
Work conditionally on the small-curvature event $\Omega_{\varepsilon_0}$ of Lemma~\ref{lem:weak_localization}.
On $\Omega_{\varepsilon_0}$ one may write $U_p=\exp(\xi_p)$ with $\|\xi_p\|\le \varepsilon_0$, and the Taylor expansion \eqref{eq:A_Taylor} yields
$A(U_p)\le C\,\|\xi_p\|^2$ after shrinking $\varepsilon_0$ and enlarging $C$.
Moreover, on $\Omega_{\varepsilon_0}$ the plaquette logarithm is controlled by the curvature 2-form, so $\|\xi_p\|\le C\|F_p\|$.
By Proposition~\ref{prop:curvature_BL} (curvature Brascamp--Lieb domination), $F_p$ is sub-Gaussian with parameter $\sigma^2\asymp 1/\beta$,
hence $\mathbb E[\|F_p\|^2\mid\Omega_{\varepsilon_0}]\le C/\beta$ uniformly in $p$.
Therefore $\mathbb E[A(U_p)\mid\Omega_{\varepsilon_0}]\le C/\beta$.

Finally, remove the conditioning using Lemma~\ref{lem:weak_localization}, which gives
$\mathbb P(\Omega_{\varepsilon_0}^c)\le |P(\mathcal B_L)|e^{-c\beta}$ and $0\le A\le 2$.
Absorbing the exponentially small contribution into $C_{\mathrm{plaq}}$ for $\beta\ge\beta_{\mathrm{plaq}}$ yields \eqref{eq:BM_plaq_defect_mean}.
\end{proof}

\begin{lemma}[Weak-coupling plaquette defect $L^2$ bound]\label{lem:BM_plaq_defect_L2}
Under the hypotheses of Lemma~\ref{lem:BM_plaq_defect_mean}, there exist $\beta_{\mathrm{plaq},2}>0$ and $C_{\mathrm{plaq},2}<\infty$
(depending only on $G$ and $\pi$) such that for every weak-coupling $\alpha$--reference measure of Wilson class on a finite block and every plaquette $p$,
\begin{equation}\label{eq:BM_plaq_defect_L2}
\mathbb E\!\big[A(U_p)^2\big] \ \le\ \frac{C_{\mathrm{plaq},2}}{\beta^2}
\qquad\text{for all }\beta\ge \beta_{\mathrm{plaq},2},\ \alpha\in[1/2,1].
\end{equation}
In particular, $\|A(U_p)\|_{L^2}\le \sqrt{C_{\mathrm{plaq},2}}/\beta$.
\end{lemma}
\begin{proof}
Repeat the argument of Lemma~\ref{lem:BM_plaq_defect_mean} but keep second moments.
On $\Omega_{\varepsilon_0}$, the Taylor bound gives $A(U_p)\le C\|\xi_p\|^2\le C\|F_p\|^2$, hence $A(U_p)^2\le C\|F_p\|^4$.
By Proposition~\ref{prop:curvature_BL}, $F_p$ is sub-Gaussian with parameter $\sigma^2\asymp 1/\beta$, so $\mathbb E[\|F_p\|^4\mid\Omega_{\varepsilon_0}]\le C/\beta^2$
uniformly in $p$.
Thus $\mathbb E[A(U_p)^2\mid\Omega_{\varepsilon_0}]\le C/\beta^2$.

Remove the conditioning as before: $\mathbb P(\Omega_{\varepsilon_0}^c)\le |P(\mathcal B_L)|e^{-c\beta}$ and $0\le A\le 2$ imply
$\mathbb E[A(U_p)^2\mathbf 1_{\Omega_{\varepsilon_0}^c}]\le 4\,\mathbb P(\Omega_{\varepsilon_0}^c)$, which is exponentially small in $\beta$ and can be absorbed into $C_{\mathrm{plaq},2}$
for $\beta\ge \beta_{\mathrm{plaq},2}$.
\end{proof}



\begin{proposition}[Partition correction from a cut-local tilt]\label{prop:BM_partition_correction}
Assume that, relative to the glued-quadrant geometry of Lemma~\ref{lem:BM_glue_supermult}, the true $2L$ thick-block partition function differs
only by the inclusion of additional \emph{cut-adjacent} plaquette weights on $P_{\mathrm{cut}}$ with a (possibly small) exponent $\theta\in(0,1]$:
\begin{equation}\label{eq:BM_partition_tilt_form}
Z^{(2L)} \;=\; Z_{\mathrm{glue}}^{(2L)}\,
\mathbb E_{\nu_{\mathrm{glue}}}\Big[\prod_{p\in P_{\mathrm{cut}}}\widehat W_\beta(U_p)^{\theta}\Big],
\end{equation}
where $\nu_{\mathrm{glue}}$ is the normalized glued measure (so $\mathbb E_{\nu_{\mathrm{glue}}}[\cdot]$ is expectation under the density defining
$Z_{\mathrm{glue}}^{(2L)}$), and $\widehat W_\beta(g)=e^{-\beta A(g)}$ is the sup-normalized Wilson class weight.

If the exponent satisfies the quantitative smallness condition
\begin{equation}\label{eq:BM_theta_small}
\beta\,\theta \ \le\ \frac{\kappa_{\mathrm{pc}}}{L^2}
\end{equation}
for some fixed $\kappa_{\mathrm{pc}}>0$, then with
\begin{equation}\label{eq:BM_Cpc_def}
C_{\mathrm{pc}}(G):=C_{\mathrm{cut}}\,C_{\mathrm{plaq}}
\end{equation}
one has the explicit lower bound
\begin{equation}\label{eq:BM_partition_correction_proved}
Z^{(2L)}\ \ge\ (Z^{(L)})^4\,\exp\!\Big(-\frac{C_{\mathrm{pc}}(G)\,\kappa_{\mathrm{pc}}}{L\beta}\Big).
\end{equation}
In particular, \eqref{eq:BM_partition_correction_proved} holds with $C=C_{\mathrm{pc}}(G)\kappa_{\mathrm{pc}}$ whenever \eqref{eq:BM_partition_tilt_form} and
\eqref{eq:BM_theta_small} hold.
\end{proposition}
\begin{proof}
Combine \eqref{eq:BM_partition_tilt_form} with Lemma~\ref{lem:BM_glue_supermult}:
\[
Z^{(2L)} \ \ge\ (Z^{(L)})^4\,
\mathbb E_{\nu_{\mathrm{glue}}}\Big[\prod_{p\in P_{\mathrm{cut}}}\widehat W_\beta(U_p)^{\theta}\Big].
\]
Since $\log$ is concave and $\widehat W_\beta(U_p)^{\theta}\in(0,1]$, Jensen gives
\[
\log \mathbb E_{\nu_{\mathrm{glue}}}\Big[\prod_{p\in P_{\mathrm{cut}}}\widehat W_\beta(U_p)^{\theta}\Big]
\ \ge\
\mathbb E_{\nu_{\mathrm{glue}}}\Big[\log \prod_{p\in P_{\mathrm{cut}}}\widehat W_\beta(U_p)^{\theta}\Big]
=
-\beta\theta\sum_{p\in P_{\mathrm{cut}}}\mathbb E_{\nu_{\mathrm{glue}}}[A(U_p)].
\]
Apply Lemma~\ref{lem:BM_plaq_defect_mean} to bound each expectation by $C_{\mathrm{plaq}}/\beta$ and use
the counting bound \eqref{eq:BM_Pcut_count}:
\[
\log \mathbb E_{\nu_{\mathrm{glue}}}\Big[\prod_{p\in P_{\mathrm{cut}}}\widehat W_\beta(U_p)^{\theta}\Big]
\ \ge\
-\beta\theta\cdot |P_{\mathrm{cut}}|\cdot \frac{C_{\mathrm{plaq}}}{\beta}
\ \ge\
- C_{\mathrm{cut}}C_{\mathrm{plaq}}\,\theta\,L.
\]
Under \eqref{eq:BM_theta_small}, $\theta L\le \kappa_{\mathrm{pc}}/(L\beta)$, hence the right-hand side is at least
$-C_{\mathrm{pc}}(G)\kappa_{\mathrm{pc}}/(L\beta)$ by \eqref{eq:BM_Cpc_def}.
Exponentiating yields \eqref{eq:BM_partition_correction_proved}.
\end{proof}

\begin{remark}[How the $\alpha$--schedule can enforce \eqref{eq:BM_theta_small}]\label{rem:BM_theta_schedule}
Proposition~\ref{prop:BM_partition_correction} is designed for the situation where the only genuinely ``cross'' contributions are
cut-local weights with a \emph{small} exponent $\theta$.
If these cut-local factors arise from a tilt sector with $\theta=1-\alpha$, then \eqref{eq:BM_theta_small} is equivalent to choosing
\[
1-\alpha \ \le\ \frac{\kappa_{\mathrm{pc}}}{\beta L^2}.
\]
One convenient sufficient form is the strengthened schedule
\[
\alpha_{\mathrm{pc}}(L,\beta)
:=1-\frac{\kappa_{\mathrm{pc}}}{\beta\,|P_\perp(\mathcal B_L)|}
\asymp 1-\frac{\kappa_{\mathrm{pc}}}{\beta L^2}\qquad(d=4),
\]
which is admissible since it only \emph{reduces} the tilt exponent $\beta(1-\alpha)$ (hence improves the $\chi^2$ and slack budgets recorded
earlier).
\end{remark}


\begin{lemma}[Iteration of the dyadic bond-moving inequality]\label{lem:BM_dyadic_iter}
Assume Proposition~\ref{prop:BM_dyadic_local}. Then for every dyadic $L=2^n$ and every irrep $R$,
\begin{equation}\label{eq:BM_dyadic_power}
\big|\lambda_R\!\bigl(w_{L}^{(\alpha)}\bigr)\big|
\ \le\
\big|\lambda_R\!\bigl(w_{1}^{(\alpha)}\bigr)\big|^{|T_L|}
\qquad(|T_L|=L^2).
\end{equation}
\noindent Equivalently,
\[
w_L^{(\alpha)} \preceq_{\mathrm{coeff}} (w_1^{(\alpha)})^{*|T_L|}.
\]
\end{lemma}
\begin{proof}
\noindent\emph{Remark on the partition-correction factor.} The $\exp(C/(L\beta))$ correction produced in Proposition~\ref{prop:BM_dyadic} is absorbed into a finite dyadic overhead (at most a constant factor $2^{K_0}$ as discussed there), so it does not affect the scale-by-scale iteration of the pure step inequality \eqref{eq:BM_dyadic_step}.

For $L=1$ there is nothing to prove.
Assume \eqref{eq:BM_dyadic_power} holds at scale $L$ and apply \eqref{eq:BM_dyadic_step}:
\[
\big|\lambda_R(w_{2L}^{(\alpha)})\big|
\le \big|\lambda_R(w_{L}^{(\alpha)})\big|^{4}
\le \Big(\big|\lambda_R(w_{1}^{(\alpha)})\big|^{|T_L|}\Big)^{4}
= \big|\lambda_R(w_{1}^{(\alpha)})\big|^{4|T_L|}.
\]
Since $|T_{2L}|=(2L)^2=4L^2=4|T_L|$, this is exactly \eqref{eq:BM_dyadic_power} at scale $2L$.
\end{proof}

\begin{remark}[Local dyadic decimation as a sufficient form of Proposition~\ref{prop:TBmix_chessboard_domination}]\label{rem:BM_to_chessboard}
Lemma~\ref{lem:BM_dyadic_iter} shows that the global power bound \eqref{eq:TBmix_chessboard} follows (for dyadic $L$) from the local
dyadic step \eqref{eq:BM_dyadic_step}.
Thus, in implementations where the blocking factor $L$ is taken dyadic (cf.\ the $N=4^k$ conventions in Section~\ref{sec:mix_majorant}),
it suffices to establish Proposition~\ref{prop:BM_dyadic_local} uniformly in $\alpha\in[1/2,1]$ in the required weak-coupling window.
\end{remark}

\begin{proposition}[Convolution dominance for the thick-block $\alpha$--reference Wilson loop]\label{prop:TBmix_chessboard_domination}
Let $w_L^{(\alpha(L))}\,dg=(\Phi_L)_{\#}\nu_L^{(\alpha(L))}$ be the thick-block boundary marginal around the $L\times L$ tile $T_L$, and let
$w_1^{(\alpha(L))}$ be the corresponding $1\times 1$ marginal.
Assume that for every nontrivial irrep $R\neq\mathbf 1$ one has the coefficient bound (in modulus)
\begin{equation}\label{eq:TBmix_chessboard}
\bigl|\lambda_R\!\bigl(w_L^{(\alpha(L))}\bigr)\bigr|
=\Bigl|\mathbb E_{\nu_L^{(\alpha(L))}}\!\big[\varphi_R(\Phi_L)\big]\Bigr|
\ \le\
\Bigl|\mathbb E_{\nu_1^{(\alpha(L))}}\!\big[\varphi_R(\Phi_1)\big]\Bigr|^{|T_L|}.
\end{equation}
\noindent Equivalently, by Definition~\ref{def:coeff_domination} and Lemma~\ref{lem:conv_coeff_mult},
\[
w_L^{(\alpha(L))} \preceq_{\mathrm{coeff}} (w_1^{(\alpha(L))})^{*|T_L|}.
\]

\end{proposition}

\noindent\emph{Comment.} For the planar reference the inequality is an equality by Lemma~\ref{lem:planar_exact_convolution}. Even when exterior boundary data cannot be gauged away, the same \emph{planar power upper bound} remains valid: deterministic boundary insertions can only affect the trace factor, so by Lemma~\ref{lem:schur_insertion_bound} (Remark~\ref{rem:planar_power_uniform_eta}) one still has $|\lambda_R|\le \lambda_R(w_1^{\parallel})^{|T_L|}$ uniformly in boundary conditions.
In the thick-block $\alpha$--reference the out-of-plane plaquettes couple the tile degrees of freedom and destroy exact convolution structure.
Proposition~\ref{prop:TBmix_chessboard_domination} packages the desired bond-moving / domination comparison: it asserts that this coupling can only \emph{reduce}
the magnitude of the boundary character coefficient relative to the solvable convolution-power model.
Reflection positivity of the $\alpha$--reference family (Corollary~\ref{cor:alpha_reference_RP}) is a natural structural prerequisite for such comparison inequalities,
but does not by itself imply \eqref{eq:TBmix_chessboard}; one needs an additional monotone decimation or domination argument.

\begin{proposition}[Thick-block convolution dominance is proved on the dyadic load-bearing path]\label{prop:TBmix_chessboard_proved}
Assume the dyadic bond-moving inequality of Proposition~\ref{prop:BM_dyadic} on the weak-coupling window under consideration.
Then for every dyadic block size $L=2^n$ and every nontrivial irrep $R\neq \mathbf 1$ one has
\begin{equation}\label{eq:TBmix_chessboard_proved}
\bigl|\lambda_R\!\bigl(w_L^{(\alpha(L))}\bigr)\bigr|
\le
\Bigl|\mathbb E_{\nu_1^{(\alpha(L))}}\!\big[\varphi_R(\Phi_1)\big]\Bigr|^{|T_L|}.
\end{equation}
Equivalently,
\[
 w_L^{(\alpha(L))} \preceq_{\mathrm{coeff}} (w_1^{(\alpha(L))})^{*|T_L|}.
\]
In particular, Proposition~\ref{prop:TBmix_chessboard_domination} is no longer an external hypothesis on the load-bearing Route~1 chain.
\end{proposition}

\begin{proof}
Proposition~\ref{prop:BM_dyadic} yields the clean dyadic step inequality \eqref{eq:BM_dyadic_step} after absorbing the finite number of initial correction factors into the universal weak-window constants.
Iterating that step by Lemma~\ref{lem:BM_dyadic_iter} gives
\[
\big|\lambda_R(w_L^{(\alpha(L))})\big|
\le
\big|\lambda_R(w_1^{(\alpha(L))})\big|^{|T_L|}
\]
for every dyadic $L$.
By the definition of the boundary marginal and \eqref{eq:TBmix_chessboard}, this is exactly the desired coefficient bound.
The coefficient-domination reformulation follows from Definition~\ref{def:coeff_domination} and Lemma~\ref{lem:conv_coeff_mult}.
\end{proof}



\begin{proposition}[One-plaquette weak-coupling coefficient bound (crossover form)]\label{prop:TBmix_one_plaq_coeff_crossover}
Assume Theorem~\ref{thm:A1_one_plaq_classical}.
Then there exist constants $c_{\mathrm{A1}}(G,\rho)>0$ and $\beta_{\mathrm{A1}}(G,\rho)\ge 1$ such that for all $\beta\ge \beta_{\mathrm{A1}}(G,\rho)$,
all $\alpha\in[1/2,1]$, and all nontrivial irreps $R\neq \mathbf 1$,
\begin{equation}\label{eq:TBmix_one_plaq_coeff}
\Bigl|\mathbb E_{\nu_1^{(\alpha)}}\!\big[\varphi_R(\Phi_1)\big]\Bigr|
\ \le\
\exp\!\Big(-c_{\mathrm{A1}}(G,\rho)\,\frac{C_2(R)}{\alpha\beta+\sqrt{C_2(R)}}\Big).
\end{equation}
Moreover, in the \emph{low-representation heat-kernel regime} $C_2(R)\le \beta^2$ one has the simplified bound
\begin{equation}\label{eq:TBmix_one_plaq_lowrep_heatkernel}
\Bigl|\mathbb E_{\nu_1^{(\alpha)}}\!\big[\varphi_R(\Phi_1)\big]\Bigr|
\ \le\
\exp\!\Big(-\frac{c_{\mathrm{A1}}(G,\rho)}{2}\,\frac{C_2(R)}{\beta}\Big).
\end{equation}
\end{proposition}

\begin{corollary}[One-plaquette crossover bound on the Route~1 load-bearing path]\label{cor:TBmix_one_plaq_closed}
For every compact connected simple $G$ and fixed Wilson representation $\rho$, the coefficient bounds \eqref{eq:TBmix_one_plaq_coeff} and \eqref{eq:TBmix_one_plaq_lowrep_heatkernel}
hold uniformly in $\alpha\in[1/2,1]$ once $\beta\ge\beta_{\mathrm{A1}}(G,\rho)$.
In particular, the one-plaquette input used by the thick-block comparison is now available in theorem-level form, with no remaining condition wrapper on the load-bearing Route~1 spine.
\end{corollary}

\begin{proof}
This is exactly the compact-simple analytic input proved in Appendix~\ref{app:A1_compact_simple} and summarized in Theorem~\ref{thm:A1_one_plaq_compact_simple}; the classical-family appendices remain direct audited special cases.
\end{proof}

\begin{proposition}[Tile crossover from theorem-level dyadic inputs]\label{prop:TBmix_from_chessboard}
For every compact connected simple $G$ and fixed Wilson representation $\rho$, the thick-block coefficient bound \eqref{eq:TBmix_coeff_decay} holds for all dyadic block sizes in the weak-coupling window relevant to Route~1.
\end{proposition}

\begin{proof}
By Proposition~\ref{prop:TBmix_chessboard_proved}, for each nontrivial $R\neq\mathbf 1$,
\[
\bigl|\lambda_R\!\bigl(w_L^{(\alpha(L))}\bigr)\bigr|
\le
\Bigl|\mathbb E_{\nu_1^{(\alpha(L))}}\!\big[\varphi_R(\Phi_1)\big]\Bigr|^{|T_L|}.
\]
Applying Corollary~\ref{cor:TBmix_one_plaq_closed} to the one-plaquette factor yields
\[
\Big(\exp\!\big(-c_{\mathrm{A1}}(G,\rho)\,\tfrac{C_2(R)}{\alpha(L)\beta+\sqrt{C_2(R)}}\big)\Big)^{|T_L|}
=
\exp\!\Big(-|T_L|\,c_{\mathrm{A1}}(G,\rho)\,\frac{C_2(R)}{\alpha(L)\beta+\sqrt{C_2(R)}}\Big),
\]
which is exactly \eqref{eq:TBmix_coeff_decay}.
\end{proof}

\begin{corollary}[Thick-block crossover coefficient decay is closed on the load-bearing path]\label{cor:TBmix_heatkernel_closed}
For every compact connected simple $G$ and fixed Wilson representation $\rho$, the coefficient decay \eqref{eq:TBmix_coeff_decay} holds for all dyadic block sizes in the weak-coupling window relevant to Route~1.
Equivalently, the thick-block reference satisfies the mixing input packaged in Proposition~\ref{prop:TBmix_coeff_decay_reference}.
\end{corollary}

\begin{proof}
Apply Proposition~\ref{prop:TBmix_from_chessboard}.
\end{proof}

\begin{remark}
Proposition~\ref{prop:TBmix_chessboard_domination} is a global convolution-dominance statement comparing the thick-block boundary marginal to a solvable convolution power.
In dyadic implementations it is implied by the local decimation step Proposition~\ref{prop:BM_dyadic_local} via Lemma~\ref{lem:BM_dyadic_iter}
(see Remark~\ref{rem:BM_to_chessboard}).
Reflection positivity of the $\alpha$--reference family (Corollary~\ref{cor:alpha_reference_RP}) is a structural prerequisite,
but establishing the local step \eqref{eq:BM_dyadic_step} requires an additional monotone domination/decimation argument (bond moving).

Proposition~\ref{prop:TBmix_one_plaq_coeff_crossover} is a separate local weak-coupling inequality.
It is proved for arbitrary compact connected simple $G$ in Appendix~\ref{app:A1_compact_simple}; Appendices~G--I retain explicit classical-family audits and are the only places where special-function/determinantal formulas enter the manuscript.
\end{remark}

\paragraph{Alternative approach (C): convolution-power mixing input for coefficient pinning}\label{sec:TBmix_routeC}

The bond-moving/chessboard approach above packages the planar comparison into a single hard inequality (BM\_dyadic).
Main chain is a complementary mechanism: it isolates the \emph{representation-dependent} part of the needed decorrelation
into an explicit, finite-dimensional representation-theoretic lemma (Lemma~\ref{lem:routeC_matrix_coeff_decay} below).
This lemma is the missing line-by-line input behind the heuristic statement
``products of many i.i.d.\ plaquettes look Haar in high representations'', and it is the place where
the familiar $1/\sqrt{d_R}$ factor (from Schur orthogonality) meets the additional
$\lambda_R(w_1)^k$ damping (from convolution-power mixing).

\medskip
\noindent
In Main chain one combines:
\begin{enumerate}[label=(\roman*)]
\item a \emph{covariance identity} for the shift of Fourier/character coefficients under a (small) transverse tilt;
\item a cut-local Lipschitz/influence bound for the tilt observable (already proved elsewhere in the paper);
\item the convolution-power damping supplied by Lemma~\ref{lem:routeC_matrix_coeff_decay}.
\end{enumerate}

\medskip
\noindent
Item~(i) is a one-line consequence of the marginal-tilt identity (Lemma~\ref{lem:marginal_tilt}) and is recorded explicitly here.

\begin{lemma}[Coefficient shift identity under a positive tilt]\label{lem:routeC_coeff_shift_identity}
Let $w^{\mathrm{ref}}$ be a central probability density on $G$ and let $\overline R:G\to[0,\infty)$ be measurable with
$0<\mathbb E_{w^{\mathrm{ref}}}[\overline R]<\infty$.
Define the tilted density
\begin{equation}\label{eq:routeC_tilt_density_def}
w^{\mathrm{tilt}}(g)
 :=\ 
\frac{\overline R(g)}{\mathbb E_{w^{\mathrm{ref}}}[\overline R]}\,w^{\mathrm{ref}}(g).
\end{equation}
Then for every irrep $R$ with normalized character $\varphi_R$,
\begin{equation}\label{eq:routeC_coeff_shift}
\lambda_R\bigl(w^{\mathrm{tilt}}\bigr)
=\frac{\mathbb E_{w^{\mathrm{ref}}}[\overline R\,\varphi_R]}{\mathbb E_{w^{\mathrm{ref}}}[\overline R]}
=\lambda_R\bigl(w^{\mathrm{ref}}\bigr)
+\frac{\mathrm{Cov}_{w^{\mathrm{ref}}}(\varphi_R,\overline R)}{\mathbb E_{w^{\mathrm{ref}}}[\overline R]}.
\end{equation}
\end{lemma}

\begin{proof}
The first equality is immediate from \eqref{eq:routeC_tilt_density_def}:
\[
\lambda_R\bigl(w^{\mathrm{tilt}}\bigr)
=\int_G \varphi_R(g)\,w^{\mathrm{tilt}}(g)\,dg
=\frac{\int_G \overline R(g)\,\varphi_R(g)\,w^{\mathrm{ref}}(g)\,dg}{\int_G \overline R(g)\,w^{\mathrm{ref}}(g)\,dg}
=\frac{\mathbb E_{w^{\mathrm{ref}}}[\overline R\,\varphi_R]}{\mathbb E_{w^{\mathrm{ref}}}[\overline R]}.
\]
For the covariance form, expand
$\mathbb E_{w^{\mathrm{ref}}}[\overline R\,\varphi_R]=\mathbb E_{w^{\mathrm{ref}}}[\overline R]\,\mathbb E_{w^{\mathrm{ref}}}[\varphi_R]+\mathrm{Cov}_{w^{\mathrm{ref}}}(\varphi_R,\overline R)$
and note that $\mathbb E_{w^{\mathrm{ref}}}[\varphi_R]=\lambda_R(w^{\mathrm{ref}})$ by definition.
\end{proof}

\begin{remark}[How \eqref{eq:routeC_coeff_shift} is used for thick blocks]
In the thick-block setting one typically has two measures on the same block, e.g.\ a planar reference block measure $\nu_L^{\parallel}$
(out-of-plane plaquettes turned off) and an $\alpha$--reference block measure $\nu_L^{(\alpha)}$.
Let $\Phi_L$ be the induced boundary holonomy map and write $w_L^{\parallel}=\Phi_{L\,*}\nu_L^{\parallel}$,
$w_L^{(\alpha)}=\Phi_{L\,*}\nu_L^{(\alpha)}$.
If $R_L:=d\nu_L^{(\alpha)}/d\nu_L^{\parallel}$ exists, then Lemma~\ref{lem:marginal_tilt} gives
\[
w_L^{(\alpha)}(g)
=\frac{\overline R_L(g)}{\mathbb E_{w_L^{\parallel}}[\overline R_L]}\,w_L^{\parallel}(g),
\qquad \overline R_L(g):=\mathbb E_{\nu_L^{\parallel}}[R_L\mid \Phi_L=g].
\]
Applying Lemma~\ref{lem:routeC_coeff_shift_identity} with $\overline R=\overline R_L$ yields
\[
\lambda_R\bigl(w_L^{(\alpha)}\bigr)
=\lambda_R\bigl(w_L^{\parallel}\bigr)
+\frac{\mathrm{Cov}_{w_L^{\parallel}}\bigl(\varphi_R,\overline R_L\bigr)}{\mathbb E_{w_L^{\parallel}}[\overline R_L]}.
\]
Moreover, by conditional expectation,
$\mathrm{Cov}_{w_L^{\parallel}}(\varphi_R,\overline R_L)=\mathrm{Cov}_{\nu_L^{\parallel}}\bigl(\varphi_R(\Phi_L),R_L\bigr)$.
Thus, bounding the coefficient shift reduces to a covariance estimate on the planar reference probability space,
which is exactly what Lemma~\ref{lem:routeC_covariance_pinning} provides once a cut-local representation of $R_L$ (or a surrogate) is available.
\end{remark}

\begin{lemma}[Effective planar tilt after transverse integration]\label{lem:routeC_effective_planar_tilt}
Fix $L$ and $\beta>0$. Let $\nu_L^{\parallel}$ denote the planar reference block measure: the density includes the in-plane plaquette weights $\widehat W_{\beta}(U_p)$ for $p\in T_L$ and no weights for out-of-plane plaquettes $p\in P_{\perp}(\mathcal B_L)$. Let $\nu_L^{(\alpha)}$ denote the $\alpha$--reference block measure: the same in-plane weights and out-of-plane weights $\widehat W_{\beta}(U_p)^{\alpha}$.
Define the (unnormalized) transverse tilt on configurations by
\begin{equation}\label{eq:routeC_transverse_tilt_T}
T_L^{(\alpha)}(U)\;:=\;\prod_{p\in P_{\perp}(\mathcal B_L)} \widehat W_{\beta}(U_p)^{\alpha}.
\end{equation}
Then $\nu_L^{(\alpha)}$ is the tilt of $\nu_L^{\parallel}$ by $T_L^{(\alpha)}$:
\begin{equation}\label{eq:routeC_block_tilt_form}
d\nu_L^{(\alpha)} \;=\; \frac{T_L^{(\alpha)}}{E_{\nu_L^{\parallel}}[T_L^{(\alpha)}]}\, d\nu_L^{\parallel}.
\end{equation}

Let $\mathscr F_{\parallel}:=\sigma(\{U_p\}_{p\in T_L})$ be the $\sigma$--algebra generated by the in-plane plaquette variables (e.g.\ in the spanning-tree gauge of Lemma~\ref{lem:planar_exact_convolution}).
Define the transverse-averaged tilt
\begin{equation}\label{eq:routeC_transverse_averaged_tilt_Psi}
\Psi_L^{(\alpha)} \;:=\; E_{\nu_L^{\parallel}}\!\big[T_L^{(\alpha)} \,\big|\, \mathscr F_{\parallel}\big].
\end{equation}
Then for every bounded $\mathscr F_{\parallel}$--measurable observable $F$,
\begin{equation}\label{eq:routeC_inplane_tilt_expectation}
E_{\nu_L^{(\alpha)}}[F] \;=\; \frac{E_{\nu_L^{\parallel}}\!\big[F\,\Psi_L^{(\alpha)}\big]}{E_{\nu_L^{\parallel}}\!\big[\Psi_L^{(\alpha)}\big]}.
\end{equation}

In particular, for the boundary holonomy $\Phi_L$ (which is $\mathscr F_{\parallel}$--measurable), the induced boundary marginal $w_L^{(\alpha)}=\Phi_{L\ast}\nu_L^{(\alpha)}$ is a tilt of $w_L^{\parallel}=\Phi_{L\ast}\nu_L^{\parallel}$ by the function
\begin{equation}\label{eq:routeC_effective_boundary_tilt}
\overline\Psi_L^{(\alpha)}(g)\;:=\;E_{\nu_L^{\parallel}}\!\big[\Psi_L^{(\alpha)} \,\big|\, \Phi_L=g\big]
\;=\;E_{\nu_L^{\parallel}}\!\big[T_L^{(\alpha)} \,\big|\, \Phi_L=g\big],
\end{equation}
so that
\begin{equation}\label{eq:routeC_effective_boundary_tilt_density}
w_L^{(\alpha)} \;=\; \frac{\overline\Psi_L^{(\alpha)}}{E_{w_L^{\parallel}}[\overline\Psi_L^{(\alpha)}]}\, w_L^{\parallel}.
\end{equation}
Moreover, for each irrep $R$,
\begin{equation}\label{eq:routeC_coeff_shift_Psi}
\lambda_R(w_L^{(\alpha)})-\lambda_R(w_L^{\parallel})
\;=\;\frac{\mathrm{Cov}_{\nu_L^{\parallel}}\!\big(\varphi_R(\Phi_L),\, \Psi_L^{(\alpha)}\big)}{E_{\nu_L^{\parallel}}[\Psi_L^{(\alpha)}]}
\;=\;\frac{\mathrm{Cov}_{w_L^{\parallel}}\!\big(\varphi_R,\, \overline\Psi_L^{(\alpha)}\big)}{E_{w_L^{\parallel}}[\overline\Psi_L^{(\alpha)}]}.
\end{equation}
\end{lemma}

\begin{proof}
Equation~\eqref{eq:routeC_block_tilt_form} is immediate from the definitions of $\nu_L^{(\alpha)}$ and $\nu_L^{\parallel}$ (they share the same in-plane plaquette weights; only the out-of-plane plaquette weights differ, producing the Radon--Nikodym factor~\eqref{eq:routeC_transverse_tilt_T}).
For any bounded $\mathscr F_{\parallel}$--measurable $F$, we have
\[
E_{\nu_L^{(\alpha)}}[F]
=\frac{E_{\nu_L^{\parallel}}[F\,T_L^{(\alpha)}]}{E_{\nu_L^{\parallel}}[T_L^{(\alpha)}]}
=\frac{E_{\nu_L^{\parallel}}\!\big[F\,E_{\nu_L^{\parallel}}[T_L^{(\alpha)}\mid\mathscr F_{\parallel}]\big]}
 {E_{\nu_L^{\parallel}}\!\big[E_{\nu_L^{\parallel}}[T_L^{(\alpha)}\mid\mathscr F_{\parallel}]\big]}
=\frac{E_{\nu_L^{\parallel}}[F\,\Psi_L^{(\alpha)}]}{E_{\nu_L^{\parallel}}[\Psi_L^{(\alpha)}]},
\]
which is~\eqref{eq:routeC_inplane_tilt_expectation}.
The identities~\eqref{eq:routeC_effective_boundary_tilt}--\eqref{eq:routeC_effective_boundary_tilt_density} are the specialization of Lemma~\ref{lem:marginal_tilt} to the tilt~\eqref{eq:routeC_block_tilt_form}, together with the tower property
$E[T_L^{(\alpha)}\mid \Phi_L]=E[E[T_L^{(\alpha)}\mid\mathscr F_{\parallel}]\mid \Phi_L]$ (since $\Phi_L$ is $\mathscr F_{\parallel}$--measurable).
Finally,~\eqref{eq:routeC_coeff_shift_Psi} is~\eqref{eq:routeC_effective_boundary_tilt_density} tested against $\varphi_R$ and rewritten as a covariance, exactly as in Lemma~\ref{lem:routeC_coeff_shift_identity}.
\end{proof}

\begin{corollary}[Normalized effective tilt]\label{cor:routeC_normalized_tilt}
In the setting of Lemma~\ref{lem:routeC_effective_planar_tilt}, define the (in-plane) \emph{normalized} effective tilt
\begin{equation}\label{eq:routeC_normalized_Psi}
\widetilde\Psi_L^{(\alpha)}
\;:=\;
\frac{\Psi_L^{(\alpha)}}{E_{\nu_L^{\parallel}}[\Psi_L^{(\alpha)}]}
\;=\;
E_{\nu_L^{\parallel}}\!\left[\frac{T_L^{(\alpha)}}{E_{\nu_L^{\parallel}}[T_L^{(\alpha)}]}\,\Big|\,\mathscr F_{\parallel}\right].
\end{equation}
Then $E_{\nu_L^{\parallel}}[\widetilde\Psi_L^{(\alpha)}]=1$ and the coefficient shift identity admits the denominator-free form
\begin{equation}\label{eq:routeC_coeff_shift_Psi_normalized}
\lambda_R(w_L^{(\alpha)})-\lambda_R(w_L^{\parallel})
\;=\;\mathrm{Cov}_{\nu_L^{\parallel}}\!\big(\varphi_R(\Phi_L),\,\widetilde\Psi_L^{(\alpha)}\big)
\;=\;\mathrm{Cov}_{w_L^{\parallel}}\!\big(\varphi_R,\,\overline{\widetilde\Psi}_L^{(\alpha)}\big),
\end{equation}
where $\overline{\widetilde\Psi}_L^{(\alpha)}(g):=E_{\nu_L^{\parallel}}[\widetilde\Psi_L^{(\alpha)}\mid \Phi_L=g]=\overline\Psi_L^{(\alpha)}(g)/E_{w_L^{\parallel}}[\overline\Psi_L^{(\alpha)}]$.
\end{corollary}

\begin{lemma}[Thin-block degeneration of the effective tilt]\label{lem:thin_block_tilt_constant}
Assume the block is the thin neighborhood $\mathcal B_L$ of Definition~\ref{def:block_BL}.
Then, in the setting of Lemma~\ref{lem:routeC_effective_planar_tilt}, the conditional transverse partition function is \emph{exactly} constant:
\begin{equation}\label{eq:thin_block_Psi_constant}
\Psi_L^{(\alpha)}
\;\equiv\;
z(\alpha)^{|P_{\perp}(\mathcal B_L)|}.
\end{equation}
Consequently $\widetilde\Psi_L^{(\alpha)}\equiv 1$ and the boundary holonomy marginal coincides with the planar marginal:
\begin{equation}\label{eq:thin_block_w_equals_planar}
w_L^{(\alpha)} \;=\; w_L^{\parallel}.
\end{equation}
\end{lemma}

\begin{proof}
By Lemma~\ref{lem:leaf_factorization_thin_block}, each out-of-plane plaquette $p\in P_{\perp}(\mathcal B_L)$ contains an \emph{opposite in-plane edge} $\ell^{(\rho,\pm)}$
which appears in no other plaquette holonomy in the thin block.
Fix all link variables except $\ell^{(\rho,\pm)}$ and write $U_p=A\,\ell^{(\rho,\pm)}B$ for some $A,B$ depending on the fixed links.
Haar bi-invariance gives
\[
\int \widehat W_{\beta}(U_p)^{\alpha}\,d\ell^{(\rho,\pm)}
\;=\;
\int \widehat W_{\beta}(g)^{\alpha}\,dg
\;=\;
z(\alpha),
\]
independent of $A,B$.
Iterating this integration over all such leaf edges shows that, \emph{for every fixed in-plane configuration},
\[
E_{\nu_L^{\parallel}}\!\left[T_L^{(\alpha)}\,\big|\,\mathscr F_{\parallel}\right]
\;=\;
z(\alpha)^{|P_{\perp}(\mathcal B_L)|},
\]
which is \eqref{eq:thin_block_Psi_constant}. The remaining statements follow immediately from
Corollary~\ref{cor:routeC_normalized_tilt}.
\end{proof}


\begin{corollary}[Thin-block tilt constancy is uniform under exterior conditioning]\label{cor:thin_block_tilt_constant_uniform_eta}
Let $\eta$ be any fixed exterior configuration on the complement links of the thin block $\mathcal B_L$.
Let $\nu_{L,\eta}^{\parallel}$ and $\nu_{L,\eta}^{(\alpha)}$ denote the corresponding planar and $\alpha$--reference conditional measures on the links in $\mathcal B_L$
(obtained by freezing the exterior links to $\eta$ in the local block weight).
Then the conclusion of Lemma~\ref{lem:thin_block_tilt_constant} holds pointwise in $\eta$:
the effective tilt is constant and the boundary holonomy marginal is unchanged,
\[
\Psi_{L,\eta}^{(\alpha)}\equiv z(\alpha)^{|P_{\perp}(\mathcal B_L)|},
\qquad
(\Phi_L)_\#\nu_{L,\eta}^{(\alpha)}=(\Phi_L)_\#\nu_{L,\eta}^{\parallel}.
\]
In particular, for every irrep $R$,
\[
\lambda_R\big((\Phi_L)_\#\nu_{L,\eta}^{(\alpha)}\big)=\lambda_R\big((\Phi_L)_\#\nu_{L,\eta}^{\parallel}\big).
\]
\end{corollary}

\begin{proof}
The proof of Lemma~\ref{lem:thin_block_tilt_constant} integrates, one by one, the Haar leaf edges $\ell^{(\rho,\pm)}$ that appear in exactly one out-of-plane plaquette holonomy.
This integration is local and uses only Haar bi-invariance on those leaf edges.
Freezing an exterior configuration $\eta$ changes only the deterministic group elements carried by edges \emph{outside} $\mathcal B_L$;
it does not remove the Haar integration over the leaf edges \emph{inside} $\mathcal B_L$.
Therefore the same leaf integration yields \eqref{eq:thin_block_Psi_constant} for $\nu_{L,\eta}^{\parallel}$, and the remaining conclusions follow exactly as in Lemma~\ref{lem:thin_block_tilt_constant}.
\end{proof}



\begin{remark}
Lemma~\ref{lem:thin_block_tilt_constant} highlights that the ``thin'' neighborhood $\mathcal B_L$ is a toy model for transverse integration:
all out-of-plane plaquettes are Haar-leaves and hence contribute only a constant factor.
The genuine thick-block mixing wall begins when the out-of-plane plaquette sector has a nontrivial ``core'' (shared transverse edges),
so that $\Psi_L^{(\alpha)}$ acquires nontrivial dependence on the in-plane field.
\end{remark}


\medskip
\noindent
The present subsection records item (iii) in enough detail to upgrade Main chain from a heuristic to a fully checkable inequality chain.
Once item (ii) and the geometric instantiation of the effective tilt $\overline R$ are supplied, Main chain can be treated as ``proved'' rather than ``standard''.


\begin{proposition}[Main-chain cut-local surrogate for the effective planar tilt]\label{ass:routeC_cut_local_surrogate}
Fix $L$ and $\alpha$ and let $\widetilde\Psi_L^{(\alpha)}$ be the normalized effective tilt from \eqref{eq:routeC_normalized_Psi}.
There exist
\begin{itemize}
\item a cut $\sigma$-algebra $\mathscr F_{\mathrm{cut}}\subset \mathscr F_{\parallel}$ generated by at most $m_L$ in-plane plaquettes, with $m_L\lesssim L$, and
\item an $\mathscr F_{\mathrm{cut}}$-measurable nonnegative observable $X_{\mathrm{cut}}=X_{\mathrm{cut}}(\sigma_{\mathrm{cut}})$,
\end{itemize}
such that the following hold.
\begin{enumerate}[label=(\roman*)]
\item \textbf{L\textsuperscript{2}-influence bound.}
There is a constant $\Delta_X>0$ so that, for each coordinate $a\in\{1,\dots,m_L\}$ and an independent resampling of that coordinate,
\begin{equation}\label{eq:routeC_cut_influence_L2}
\big\|X_{\mathrm{cut}}-X_{\mathrm{cut}}^{(a)}\big\|_{L^2(\nu_L^{\parallel})}
\;\le\;
\frac{\Delta_X}{\beta}.
\end{equation}
\item \textbf{Exponential surrogate and L\textsuperscript{2}-approximation.}
There exists a tilt parameter $t=t_L\ge 0$ (possibly depending on $\alpha,\beta,L$) such that the \emph{normalized} cut-local surrogate
\begin{equation}\label{eq:routeC_surrogate_def}
\widetilde\Psi_{\mathrm{cut}}
\;:=\;
\frac{e^{-t X_{\mathrm{cut}}}}{E_{\nu_L^{\parallel}}\!\left[e^{-t X_{\mathrm{cut}}}\right]}
\end{equation}
approximates the true normalized effective tilt in $L^2(\nu_L^{\parallel})$:
\begin{equation}\label{eq:routeC_surrogate_L2}
\big\|\widetilde\Psi_L^{(\alpha)}-\widetilde\Psi_{\mathrm{cut}}\big\|_{L^2(\nu_L^{\parallel})}
\;\le\;
\epsilon_L^{(C)}.
\end{equation}
\item \textbf{Mean lower bound.}
There exists $c_{\Psi}>0$ (uniform in $L$) such that
\begin{equation}\label{eq:routeC_surrogate_mean_lower}
E_{\nu_L^{\parallel}}\!\left[e^{-t X_{\mathrm{cut}}}\right]\;\ge\;c_{\Psi}.
\end{equation}
\end{enumerate}
\end{proposition}


\begin{lemma}[Main-chain reduction via a cut-local surrogate]\label{lem:routeC_surrogate_reduction}
Assume Proposition~\ref{ass:routeC_cut_local_surrogate}. Then for every nontrivial irrep $R$,
\begin{align*}
\bigl|\lambda_R(w_L^{(\alpha)})-\lambda_R(w_L^{\parallel})\bigr|
&=\bigl|\mathrm{Cov}_{\nu_L^{\parallel}}\big(\varphi_R(\Phi_L),\widetilde\Psi_L^{(\alpha)}\big)\bigr|
\qquad\text{(by \eqref{eq:routeC_coeff_shift_Psi_normalized})}
\\
&\le \bigl|\mathrm{Cov}_{\nu_L^{\parallel}}\big(\varphi_R(\Phi_L),\widetilde\Psi_{\mathrm{cut}}\big)\bigr|
\;+
\bigl|\mathrm{Cov}_{\nu_L^{\parallel}}\big(\varphi_R(\Phi_L),\widetilde\Psi_L^{(\alpha)}-\widetilde\Psi_{\mathrm{cut}}\big)\bigr|
\\
&\le \frac{1}{E_{\nu_L^{\parallel}}[e^{-tX_{\mathrm{cut}}}]}
\bigl|\mathrm{Cov}_{\nu_L^{\parallel}}\big(\varphi_R(\Phi_L),e^{-tX_{\mathrm{cut}}}\big)\bigr|
\;+
\varepsilon_L^{\mathrm{(C)}}\,\|\varphi_R(\Phi_L)\|_{L^2(\nu_L^{\parallel})}.
\end{align*}
Moreover, $\|\varphi_R(\Phi_L)\|_{L^2(\nu_L^{\parallel})}\le d_R^{-1}\,\|w_L^{\parallel}\|_{\infty}^{1/2}$, and the first covariance term is bounded by Lemma~\ref{lem:routeC_covariance_pinning} applied to $X_{\mathrm{cut}}$.
\end{lemma}

\begin{remark}[The surrogate normalizer is bounded away from $0$]\label{rem:routeC_surrogate_normalizer}
Since $X_{\mathrm{cut}}\ge 0$ and $t>0$, one has $0<E_{\nu_L^{\parallel}}[e^{-tX_{\mathrm{cut}}}]\le 1$.
Moreover, Jensen's inequality gives
\begin{equation}\label{eq:routeC_surrogate_normalizer_jensen}
E_{\nu_L^{\parallel}}[e^{-tX_{\mathrm{cut}}}]\ge \exp\bigl(-t\,E_{\nu_L^{\parallel}}[X_{\mathrm{cut}}]\bigr).
\end{equation}
For the explicit strip observable of Lemma~\ref{lem:routeC_cut_instantiation}, the one-plaquette defect mean bound \eqref{eq:BM_plaq_defect_mean} yields
\begin{equation}\label{eq:routeC_Xcut_mean_bound}
E_{\nu_L^{\parallel}}[X_{\mathrm{cut}}]\le |T_{\mathrm{cut}}|\,\frac{C_{\mathrm{plaq}}}{\beta}\le \frac{C_{\mathrm{cut}}L\,C_{\mathrm{plaq}}}{\beta}.
\end{equation}
Consequently,
\begin{equation}\label{eq:routeC_surrogate_normalizer_inverse_bound}
1\le \frac{1}{E_{\nu_L^{\parallel}}[e^{-tX_{\mathrm{cut}}}]}
\le \exp\Bigl( t\,\frac{C_{\mathrm{cut}}L\,C_{\mathrm{plaq}}}{\beta}\Bigr),
\end{equation}
which is harmless provided $t\,L/\beta$ stays bounded; for instance this holds under the common choice $t=\beta(1-\alpha)$ with $1-\alpha\asymp \kappa/L^2$ at the entrance scale $L^2/\beta\asymp 1$.
\end{remark}

\begin{remark}[Diagnostic on shared-edge Haar integrations at weak coupling]\label{rem:routeC_shared_edge_diagnostic}
A basic obstruction to a ``one-plaquette-at-a-time'' leaf integration on a \emph{thick} out-of-plane sector is the appearance of \emph{shared} transverse edges.
The simplest local diagnostic is the two-plaquette shared-edge integral
\begin{equation}\label{eq:routeC_shared_edge_integral}
\mathcal I_\alpha(X)\;:=\;\int_G \widehat W_\beta(g)^\alpha\,\widehat W_\beta(gX)^\alpha\,dg,
\qquad X\in G,
\end{equation}
which arises when two transverse plaquettes share the same integrated edge-variable $g$.
Let $z_\alpha:=\int_G \widehat W_\beta(g)^\alpha\,dg$.

By Peter--Weyl and Schur orthogonality, one has the exact character expansion
\[
\mathcal I_\alpha(X)
=
\sum_{S\in\widehat G} d_S\,\lambda_S(\widehat W_\beta^\alpha)^2\,\chi_S(X)
=
z_\alpha^2
+
\sum_{S\neq \mathbf 1} d_S\,\lambda_S(\widehat W_\beta^\alpha)^2\,\chi_S(X),
\]
so using $|\chi_S(X)|\le d_S$ gives the uniform bound
\begin{equation}\label{eq:routeC_shared_edge_bound_z2a}
\bigl|\mathcal I_\alpha(X)-z_\alpha^2\bigr|
\;\le\;
\sum_{S\neq \mathbf 1} d_S^2\,\lambda_S(\widehat W_\beta^\alpha)^2
\;=\;
\|\widehat W_\beta^\alpha\|_{L^2(G)}^2-z_\alpha^2
\;=\;
z_{2\alpha}-z_\alpha^2
\;\le\;
z_{2\alpha}.
\end{equation}
In particular, each shared-edge integration produces a \emph{polynomially small} absolute error at weak coupling,
since $z_{2\alpha}\asymp (2\alpha\beta)^{-\dim(G)/2}$ as $\beta\to\infty$ (by Laplace expansion in a chart near the identity).

\medskip
\noindent\textbf{Higher degree.}
If an integrated transverse edge is incident to $k\ge 2$ plaquette-factors, the same philosophy still applies.
Indeed, for any $X_1,\dots,X_k\in G$ and $f(g):=\widehat W_{\beta}(g)^{\alpha}$ one has
\(
0\le \int_G \prod_{j=1}^k f(gX_j)\,dg\le \int_G f(g)f(gX)\,dg\le \int_G f(g)^2\,dg=z_{2\alpha}
\)
by dropping factors $\le 1$ and using Cauchy--Schwarz (here $X:=X_1^{-1}X_2$).
Moreover $z_{\alpha}^k\le z_{\alpha}^2\le z_{2\alpha}$ for $k\ge 2$.
Hence the deviation from the naive ``decoupled'' value is still uniformly bounded by $z_{2\alpha}$:
\(
\bigl|\int \prod_{j=1}^k f(gX_j)\,dg\; -\; z_{\alpha}^k\bigr|\le z_{2\alpha}.
\)
Consequently, whenever a global shared-edge correction is expressed as a product over edge-integrals,
Lemma~\ref{lem:routeC_shared_edge_telescoping} canonically yields an edge-by-edge decomposition with per-edge defect bounded by $z_{2\alpha}$.

This estimate is only a \emph{local} diagnostic: to control a thick-block effective tilt one must still (i) identify the relevant set of shared edges in the transverse core,
(ii) count how many such shared-edge integrations survive after all available leaf eliminations, and (iii) propagate the local bound \eqref{eq:routeC_shared_edge_bound_z2a}
through the remaining conditional expectations (typically via an $L^2$ tensor-network estimate rather than a pointwise union bound).
\end{remark}


\begin{lemma}[Sufficient conditions for the cut-local surrogate estimate]\label{lem:routeC_surrogate_sufficient_conditions}
Work under $\nu_L^{\parallel}$ and fix a cut $\sigma$-algebra $\mathscr F_{\mathrm{cut}}\subset\mathscr F_{\parallel}$.
Let $T:=T_L^{(\alpha)}$ and $\widetilde\Psi_L^{(\alpha)}$ be defined as in \eqref{eq:routeC_normalized_Psi}.
Assume that on $\Omega_{\varepsilon_0}$ there exist a parameter $t\ge 0$, a nonnegative $\mathscr F_{\mathrm{cut}}$-measurable random variable $X_{\mathrm{cut}}$,
a nonnegative random variable $X_{\mathrm{out}}$, and an error term $\Delta_{\mathrm{sh}}$ such that
\begin{equation}\label{eq:routeC_T_factor_cut_out}
T \;=\; e^{-t X_{\mathrm{cut}}}\,e^{-t X_{\mathrm{out}}}\,\bigl(1+\Delta_{\mathrm{sh}}\bigr).
\end{equation}
Define the $\mathscr F_{\parallel}$-conditional objects
\[
m_\parallel := \mathbb E_{\nu_L^{\parallel}}\bigl[e^{-tX_{\mathrm{out}}}\bigm|\mathscr F_\parallel\bigr],
\qquad
\mathrm{Err}_\parallel := \mathbb E_{\nu_L^{\parallel}}\bigl[e^{-tX_{\mathrm{out}}}\Delta_{\mathrm{sh}}\bigm|\mathscr F_\parallel\bigr],
\]
and set $\bar m := \mathbb E_{\nu_L^{\parallel}}[e^{-tX_{\mathrm{out}}}]$ and $Z_{\mathrm{cut}}:=\mathbb E_{\nu_L^{\parallel}}[e^{-tX_{\mathrm{cut}}}]$.
Assume the quantitative bounds
\begin{enumerate}[label=(\roman*)]
\item $\| m_\parallel - \bar m\|_{L^2(\nu_L^{\parallel}\mid\Omega_{\varepsilon_0})}\le \varepsilon_{\mathrm{out}}$,
\item $\| \mathrm{Err}_\parallel\|_{L^2(\nu_L^{\parallel}\mid\Omega_{\varepsilon_0})}\le \varepsilon_{\mathrm{sh}}$,
\item $Z_{\mathrm{cut}} \ge c_\Psi>0$ and $\bar m \ge c_m>0$,
\item $\varepsilon_{\mathrm{out}}+\varepsilon_{\mathrm{sh}} \le \tfrac12 c_\Psi c_m$.
\end{enumerate}
Then the cut-local surrogate $\widetilde\Psi_{\mathrm{cut}} := e^{-tX_{\mathrm{cut}}}/\mathbb E_{\nu_L^{\parallel}}[e^{-tX_{\mathrm{cut}}}]$ satisfies
\[
\bigl\|\widetilde\Psi_L^{(\alpha)} - \widetilde\Psi_{\mathrm{cut}}\bigr\|_{L^2(\nu_L^{\parallel}\mid\Omega_{\varepsilon_0})}
\;\le\;
\frac{8}{c_\Psi c_m}\,(\varepsilon_{\mathrm{out}}+\varepsilon_{\mathrm{sh}}).
\]
\end{lemma}

\begin{proof}
On $\Omega_{\varepsilon_0}$, the decomposition \eqref{eq:routeC_T_factor_cut_out} and the $\mathscr F_{\mathrm{cut}}$-measurability of $X_{\mathrm{cut}}$
give
\[
\Psi_L^{(\alpha)} \;=\; \mathbb E_{\nu_L^{\parallel}}\!\bigl[T\bigm|\mathscr F_\parallel\bigr]
\;=\;
e^{-tX_{\mathrm{cut}}}\,
\mathbb E_{\nu_L^{\parallel}}\!\bigl[e^{-tX_{\mathrm{out}}}(1+\Delta_{\mathrm{sh}})\bigm|\mathscr F_\parallel\bigr]
\;=\;
e^{-tX_{\mathrm{cut}}}\,(m_\parallel+\mathrm{Err}_\parallel).
\]
Hence, with $Z:=\mathbb E_{\nu_L^{\parallel}}[\Psi_L^{(\alpha)}]=\mathbb E_{\nu_L^{\parallel}}[T]$, we have
\[
\widetilde\Psi_L^{(\alpha)} \;=\; \frac{\Psi_L^{(\alpha)}}{Z}
\;=\;
\frac{e^{-tX_{\mathrm{cut}}}(m_\parallel+\mathrm{Err}_\parallel)}{Z}.
\]
By the triangle inequality and Cauchy--Schwarz,
\[
\bigl| Z - Z_{\mathrm{cut}}\bar m \bigr|
\,=\,
\Bigl|
\mathbb E_{\nu_L^{\parallel}}\!\bigl[e^{-tX_{\mathrm{cut}}}(m_\parallel-\bar m)\bigr]
+
\mathbb E_{\nu_L^{\parallel}}\!\bigl[e^{-tX_{\mathrm{cut}}}\mathrm{Err}_\parallel\bigr]
\Bigr|
\;\le\;
\varepsilon_{\mathrm{out}}+\varepsilon_{\mathrm{sh}}.
\]
Using $Z_{\mathrm{cut}}\ge c_\Psi$ and $\bar m\ge c_m$ together with (iv) yields the lower bound
\[
Z \;\ge\; Z_{\mathrm{cut}}\bar m - (\varepsilon_{\mathrm{out}}+\varepsilon_{\mathrm{sh}})
\;\ge\; \tfrac12 c_\Psi c_m.
\]
Now write
\[
\widetilde\Psi_L^{(\alpha)}-\widetilde\Psi_{\mathrm{cut}}
\;=\;
e^{-tX_{\mathrm{cut}}}\Bigl(\frac{m_\parallel+\mathrm{Err}_\parallel}{Z}-\frac{1}{Z_{\mathrm{cut}}}\Bigr),
\]
so that $e^{-tX_{\mathrm{cut}}}\le 1$ gives
\[
\bigl\|\widetilde\Psi_L^{(\alpha)}-\widetilde\Psi_{\mathrm{cut}}\bigr\|_{2}
\;\le\;
\Bigl\|\frac{m_\parallel+\mathrm{Err}_\parallel}{Z}-\frac{1}{Z_{\mathrm{cut}}}\Bigr\|_{2}.
\]
Decompose the right-hand side as
\[
\frac{m_\parallel+\mathrm{Err}_\parallel}{Z}-\frac{1}{Z_{\mathrm{cut}}}
=
\frac{(m_\parallel-\bar m)+\mathrm{Err}_\parallel}{Z}
+
\frac{\bar m}{Z}-\frac{1}{Z_{\mathrm{cut}}}
=
\frac{(m_\parallel-\bar m)+\mathrm{Err}_\parallel}{Z}
+
\frac{Z_{\mathrm{cut}}\bar m-Z}{Z\,Z_{\mathrm{cut}}}.
\]
Using the lower bounds $Z\ge \tfrac12c_\Psi c_m$ and $Z_{\mathrm{cut}}\ge c_\Psi$ and the estimates above,
\[
\Bigl\|\frac{(m_\parallel-\bar m)+\mathrm{Err}_\parallel}{Z}\Bigr\|_2
\;\le\;
\frac{2}{c_\Psi c_m}\,(\varepsilon_{\mathrm{out}}+\varepsilon_{\mathrm{sh}}),
\qquad
\Bigl|\frac{Z_{\mathrm{cut}}\bar m-Z}{Z\,Z_{\mathrm{cut}}}\Bigr|
\;\le\;
\frac{4}{c_\Psi c_m}\,(\varepsilon_{\mathrm{out}}+\varepsilon_{\mathrm{sh}}).
\]
Combining the last displays yields the stated bound.
\end{proof}

\begin{remark}[Where the hard analysis lives]
Lemma~\ref{lem:routeC_surrogate_sufficient_conditions} makes explicit that the quantitative content of
Proposition~\ref{ass:routeC_cut_local_surrogate}(ii) reduces to two concrete inputs:
\begin{itemize}
\item a \emph{geometric/probabilistic} decoupling statement that produces the factorization \eqref{eq:routeC_T_factor_cut_out} with a controlled
shared-edge error $\varepsilon_{\mathrm{sh}}$, and
\item a weak-coupling control of the \emph{in-plane fluctuations} of the outside factor, namely an $L^2$ bound on
$m_\parallel-\bar m$ where $m_\parallel=\E_{\nu_L^{\parallel}}[e^{-tX_{\mathrm{out}}}\mid\mathscr F_\parallel]$ and $\bar m=\E_{\nu_L^{\parallel}}[e^{-tX_{\mathrm{out}}}]$.
A sufficient approach is: (a) an influence estimate for $X_{\mathrm{out}}$ under single-plaquette resampling and (b) a sub-Gaussian/mgf bound for $X_{\mathrm{out}}$,
which together give stability of conditional Laplace transforms and hence the desired $L^2$ control of $m_\parallel-\bar m$.
\end{itemize}
In the thin-block model $\mathcal B_L$ the shared-edge term vanishes (Lemma~\ref{lem:thin_block_tilt_constant}),
while for genuine thick blocks one needs a quantitative version of the shared-edge diagnostic in
Remark~\ref{rem:routeC_shared_edge_diagnostic} combined with a counting argument for the number of transverse shared edges crossing the cut.
\end{remark}

\begin{lemma}[From outside-energy influences to $L^2$ flatness of $m_\parallel$]
\label{lem:routeC_mparallel_from_influences}
Work under $\nu_L^{\parallel}$ and suppose $X_{\mathrm{out}}=X_{\mathrm{out}}(\sigma_{\parallel},\omega)$ is a nonnegative random variable on a product space where
$\sigma_{\parallel}=(\sigma_1,\dots,\sigma_N)$ are the in-plane plaquette variables (independent under $\nu_L^{\parallel}$)
and $\omega$ collects all non-in-plane edge variables (Haar under $\nu_L^{\parallel}$, independent of $\sigma_{\parallel}$).
Fix $t\ge 0$ and define
\[
m_\parallel(\sigma_{\parallel}) := \mathbb E\bigl[e^{-t X_{\mathrm{out}}(\sigma_{\parallel},\omega)}\bigm| \sigma_{\parallel}\bigr],
\qquad
\bar m := \mathbb E\bigl[e^{-tX_{\mathrm{out}}}\bigr].
\]
For each $a\in\{1,\dots,N\}$ let $\sigma^{(a)}$ be obtained by resampling only $\sigma_a$ independently and set
$X_{\mathrm{out}}^{(a)} := X_{\mathrm{out}}(\sigma^{(a)},\omega)$ and $m_\parallel^{(a)} := m_\parallel(\sigma^{(a)})$.
Then
\begin{equation}\label{eq:routeC_mparallel_efronstein}
\|m_\parallel-\bar m\|_{L^2(\nu_L^{\parallel})}
\le \frac{1}{\sqrt2}\Big(\sum_{a=1}^N \|e^{-tX_{\mathrm{out}}}-e^{-tX_{\mathrm{out}}^{(a)}}\|_{L^2(\nu_L^{\parallel})}^2\Big)^{1/2}.
\end{equation}
In particular, if there is an index set $\mathcal I_{\mathrm{out}}\subset\{1,\dots,N\}$ such that
$X_{\mathrm{out}}$ does not depend on $\sigma_a$ for $a\notin\mathcal I_{\mathrm{out}}$, and
\begin{equation}\label{eq:routeC_Xout_influence}
\|X_{\mathrm{out}}-X_{\mathrm{out}}^{(a)}\|_{L^2(\nu_L^{\parallel})}\le \frac{\Delta_{\mathrm{out}}}{\beta}\qquad(a\in\mathcal I_{\mathrm{out}}),
\end{equation}
then
\begin{equation}\label{eq:routeC_mparallel_influence_bound}
\|m_\parallel-\bar m\|_{L^2(\nu_L^{\parallel})}
\le \frac{t\,\Delta_{\mathrm{out}}}{\beta\sqrt2}\,\sqrt{|\mathcal I_{\mathrm{out}}|}.
\end{equation}
\end{lemma}

\begin{proof}
Apply Lemma~\ref{lem:routeC_influence_variance} (Efron--Stein) to the function $Y:=m_\parallel(\sigma_{\parallel})$ of the independent coordinates $\sigma_1,\dots,\sigma_N$.
Then
\[
\|m_\parallel-\bar m\|_2^2=\mathrm{Var}(m_\parallel)
\le \frac12\sum_{a=1}^N\mathbb E\bigl[(m_\parallel-m_\parallel^{(a)})^2\bigr].
\]
Now fix $a$ and work on the probability space with coordinates $(\sigma_{-a},\sigma_a,\sigma_a',\omega)$ where $\sigma_a'$ is an independent copy of $\sigma_a$.
Write $Z:=e^{-tX_{\mathrm{out}}(\sigma_{-a},\sigma_a,\omega)}$ and
$Z^{(a)}:=e^{-tX_{\mathrm{out}}(\sigma_{-a},\sigma_a',\omega)}$.
Then $m_\parallel=\mathbb E_\omega[Z]$ and $m_\parallel^{(a)}=\mathbb E_\omega[Z^{(a)}]$, so Jensen gives
\[
(m_\parallel-m_\parallel^{(a)})^2
=\bigl(\mathbb E_\omega[Z-Z^{(a)}]\bigr)^2
\le \mathbb E_\omega\bigl[(Z-Z^{(a)})^2\bigr].
\]
Taking expectation over $(\sigma_{-a},\sigma_a,\sigma_a')$ yields
$\mathbb E[(m_\parallel-m_\parallel^{(a)})^2]\le \mathbb E[(Z-Z^{(a)})^2]=\|e^{-tX_{\mathrm{out}}}-e^{-tX_{\mathrm{out}}^{(a)}}\|_2^2$.
Inserting this into the Efron--Stein bound and taking square roots gives \eqref{eq:routeC_mparallel_efronstein}.
If \eqref{eq:routeC_Xout_influence} holds, then $|e^{-tx}-e^{-ty}|\le t|x-y|$ implies
$\|e^{-tX_{\mathrm{out}}}-e^{-tX_{\mathrm{out}}^{(a)}}\|_2\le t\,\|X_{\mathrm{out}}-X_{\mathrm{out}}^{(a)}\|_2\le t\Delta_{\mathrm{out}}/\beta$ for $a\in\mathcal I_{\mathrm{out}}$ and $0$ otherwise.
This yields \eqref{eq:routeC_mparallel_influence_bound}.
\end{proof}


\begin{lemma}[Bounding the shared-edge error contribution]
\label{lem:routeC_errparallel_bound}
Work under $\nu_L^{\parallel}$ and let $\Delta_{\mathrm{sh}}$ be square-integrable.
With $X_{\mathrm{out}}\ge 0$ and $t\ge 0$ define
\[
\mathrm{Err}_\parallel
:=\mathbb E_{\nu_L^{\parallel}}\bigl[e^{-tX_{\mathrm{out}}}\,\Delta_{\mathrm{sh}}\bigm|\mathscr F_\parallel\bigr].
\]
Then
\begin{equation}\label{eq:routeC_errparallel_condexp_contraction}
\|\mathrm{Err}_\parallel\|_{L^2(\nu_L^{\parallel})}
\le \|\Delta_{\mathrm{sh}}\|_{L^2(\nu_L^{\parallel})}.
\end{equation}
Moreover, if $\Delta_{\mathrm{sh}}=\sum_{e\in E_{\mathrm{sh}}}\delta_e$ with $E_{\mathrm{sh}}$ finite and $\|\delta_e\|_2\le \zeta$ for all $e$, then
\begin{equation}\label{eq:routeC_Delta_sh_counting}
\|\Delta_{\mathrm{sh}}\|_{L^2(\nu_L^{\parallel})}\le |E_{\mathrm{sh}}|\,\zeta.
\end{equation}
In particular, a sufficient pointwise control is $|\delta_e|\le z_{2\alpha}$ (cf.\ \eqref{eq:routeC_shared_edge_bound_z2a}), which implies $\|\delta_e\|_2\le z_{2\alpha}$.
\end{lemma}

\begin{proof}
Conditional expectation is a contraction in $L^2$, hence
\(
\|\mathrm{Err}_\parallel\|_2\le \|e^{-tX_{\mathrm{out}}}\,\Delta_{\mathrm{sh}}\|_2\le \|\Delta_{\mathrm{sh}}\|_2
\)
since $0\le e^{-tX_{\mathrm{out}}}\le 1$.
The counting bound \eqref{eq:routeC_Delta_sh_counting} is Minkowski's inequality:
\(\|\sum_e \delta_e\|_2\le\sum_e\|\delta_e\|_2\le |E_{\mathrm{sh}}|\,\zeta\).
\end{proof}


% =====================================================================
% : Canonical edge-by-edge telescoping for shared-edge errors
% =====================================================================

\begin{lemma}[Edge-by-edge telescoping decomposition for shared-edge products]
\label{lem:routeC_shared_edge_telescoping}
Let $E$ be a finite index set and let $(a_e)_{e\in E}$ and $(b_e)_{e\in E}$ be real-valued random variables.
Assume that $0\le a_e\le 1$ and $0\le b_e\le 1$ almost surely for every $e\in E$.
Define the product difference
\begin{equation}\label{eq:routeC_telescoping_DeltaE_def}
\Delta_E\;:=\;\prod_{e\in E} b_e\; -\; \prod_{e\in E} a_e.
\end{equation}
Then for any ordering $E=\{e_1,\dots,e_n\}$ one has the telescoping identity
\begin{equation}\label{eq:routeC_telescoping_identity}
\Delta_E
=\sum_{j=1}^n (b_{e_j}-a_{e_j})\,\prod_{i<j} b_{e_i}\,\prod_{i>j} a_{e_i}.
\end{equation}
In particular, if there is a deterministic $\zeta\ge 0$ such that $|b_e-a_e|\le \zeta$ pointwise for all $e\in E$,
then $\Delta_E$ admits a decomposition
\begin{equation}\label{eq:routeC_telescoping_decomp}
\Delta_E=\sum_{e\in E} \delta_e
\qquad\text{with}\qquad
|\delta_e|\le \zeta.
\end{equation}
\end{lemma}

\begin{proof}
The identity \eqref{eq:routeC_telescoping_identity} is the standard expansion
\(
\prod_{j=1}^n b_j-\prod_{j=1}^n a_j=\sum_{j=1}^n (b_j-a_j)\prod_{i<j} b_i\prod_{i>j} a_i
\),
applied with $b_j=b_{e_j}$ and $a_j=a_{e_j}$.
If $0\le a_e,b_e\le 1$, then each product factor in \eqref{eq:routeC_telescoping_identity} has absolute value at most $1$,
so the pointwise bound $|b_e-a_e|\le\zeta$ implies $|\delta_{e_j}|\le\zeta$ where
\(
\delta_{e_j}:=(b_{e_j}-a_{e_j})\prod_{i<j} b_{e_i}\prod_{i>j} a_{e_i}
\).
\end{proof}

\begin{corollary}[From per-edge product control to a $z_{2\alpha}$ shared-edge decomposition]
\label{cor:routeC_shared_edge_decomp_from_telescoping}
Let $E$ be a finite set and suppose that the shared-edge correction satisfies the product form
\(
\Delta_{\mathrm{sh}}=\prod_{e\in E} b_e-\prod_{e\in E} a_e
\)
with $0\le a_e,b_e\le 1$ and $|b_e-a_e|\le z_{2\alpha}$ pointwise for all $e\in E$.
Then $\Delta_{\mathrm{sh}}$ admits a decomposition
\(
\Delta_{\mathrm{sh}}=\sum_{e\in E}\delta_e
\)
with $|\delta_e|\le z_{2\alpha}$.
Consequently, Lemma~\ref{lem:routeC_errparallel_bound} applies with $E_{\mathrm{sh}}=E$.
\end{corollary}

\begin{proof}
Apply Lemma~\ref{lem:routeC_shared_edge_telescoping} with $\zeta=z_{2\alpha}$.
\end{proof}


\begin{corollary}[Quantitative surrogate error from outside influences and shared-edge counting]
\label{cor:routeC_surrogate_error_from_influences}
In the setting of Lemma~\ref{lem:routeC_surrogate_sufficient_conditions}, assume additionally that:
\begin{enumerate}[label=(\alph*)]
\item there is an index set $\mathcal I_{\mathrm{out}}$ such that \eqref{eq:routeC_Xout_influence} holds and
$X_{\mathrm{out}}$ does not depend on the remaining in-plane variables, and
\item $\Delta_{\mathrm{sh}}=\sum_{e\in E_{\mathrm{sh}}}\delta_e$ with $|E_{\mathrm{sh}}|<\infty$ and $|\delta_e|\le z_{2\alpha}$.
\end{enumerate}
\noindent\emph{Midline instantiation.} For the canonical outside choice \eqref{eq:routeC_Xout_def} associated to the strip $T_{\mathrm{cut}}$ (Lemma~\ref{lem:routeC_cut_instantiation}), Lemma~\ref{lem:routeC_Xout_influence_instantiation} supplies (a) with $\mathcal I_{\mathrm{out}}=I_{\mathrm{out}}=T_{\mathrm{out}}$ and $\Delta_{\mathrm{out}}=2\sqrt{C_{\mathrm{plaq},2}}$.
Then the surrogate approximation error in Proposition~\ref{ass:routeC_cut_local_surrogate}(ii) can be taken as
\begin{equation}\label{eq:routeC_epsC_explicit}
\epsilon_L^{(C)}
\le
\frac{6}{c_\Psi c_m}\Big(\frac{t\,\Delta_{\mathrm{out}}}{\beta\sqrt2}\,\sqrt{|\mathcal I_{\mathrm{out}}|}
\;+\;
|E_{\mathrm{sh}}|\,z_{2\alpha}\Big).
\end{equation}
In particular, if $|\mathcal I_{\mathrm{out}}|\lesssim L^2$ and $|E_{\mathrm{sh}}|\lesssim L^2$, then at the entrance scaling $L^2\asymp \beta$ one has
\(\epsilon_L^{(C)}=O(\beta^{1-\dim(G)/2})\) as $\beta\to\infty$.
\end{corollary}

\begin{proof}
In Lemma~\ref{lem:routeC_surrogate_sufficient_conditions}, the surrogate $L^2$ error is bounded by
\(\tfrac{6}{c_\Psi c_m}(\varepsilon_{\mathrm{out}}+\varepsilon_{\mathrm{sh}})\)
with $\varepsilon_{\mathrm{out}}:=\|m_\parallel-\bar m\|_2$ and $\varepsilon_{\mathrm{sh}}:=\|\mathrm{Err}_\parallel\|_2$.
Apply Lemma~\ref{lem:routeC_mparallel_from_influences} to bound $\varepsilon_{\mathrm{out}}$ by the first term in \eqref{eq:routeC_epsC_explicit}.
Apply Lemma~\ref{lem:routeC_errparallel_bound} to bound $\varepsilon_{\mathrm{sh}}\le \|\Delta_{\mathrm{sh}}\|_2\le |E_{\mathrm{sh}}|z_{2\alpha}$, giving \eqref{eq:routeC_epsC_explicit}.
\end{proof}





% =====================================================================
% : Combinatorial bookkeeping for the midline strip cut (Main chain)
% =====================================================================

\paragraph*{Concrete size bounds for the midline cut parameters}

\begin{lemma}[Midline cut bookkeeping in $d=4$]\label{lem:routeC_midline_bookkeeping}
Work in $d=4$ and let $T_L$ be the $L\times L$ planar tile in the $(1,2)$--plane used in Main chain.
Let $T_{\mathrm{cut}}$ be the unit-width midline strip from Lemma~\ref{lem:routeC_cut_instantiation}, and write
$m_L := |T_{\mathrm{cut}}|\le 2L$.

Let $\mathcal E_{\mathrm{cut}}$ be the set of in-plane links that appear in at least one plaquette of $T_{\mathrm{cut}}$.
Let $P_{\perp,\mathrm{cut}}$ be the set of out-of-plane plaquettes (in the thick block) whose boundary contains at least one link of $\mathcal E_{\mathrm{cut}}$.
Let $\mathcal E_{\mathrm{sh}}$ be the set of non-in-plane links that are contained in at least two plaquettes of $P_{\perp,\mathrm{cut}}$.

Then there exist absolute constants $C_{\mathcal E},C_{\perp},C_{\mathrm{sh}}<\infty$ (depending only on the dimension $d=4$) such that
\begin{equation}\label{eq:routeC_midline_size_bounds}
|\mathcal E_{\mathrm{cut}}|\le C_{\mathcal E}\, L,\qquad
|P_{\perp,\mathrm{cut}}|\le C_{\perp}\, L,\qquad
|\mathcal E_{\mathrm{sh}}|\le C_{\mathrm{sh}}\, L.
\end{equation}
Moreover, in any shared-edge decomposition of $\Delta_{\mathrm{sh}}$ supported on shared non-in-plane links inside $P_{\perp,\mathrm{cut}}$,
one may take the index set $E_{\mathrm{sh}}$ in Corollary~\ref{cor:routeC_surrogate_error_from_influences} to satisfy $|E_{\mathrm{sh}}|\le C_{\mathrm{sh}}L$.
\end{lemma}

\begin{proof}
The strip $T_{\mathrm{cut}}$ contains $m_L\le 2L$ plaquettes, each plaquette has $4$ boundary links, and each in-plane link can be incident
to at most $2(d-2)=4$ out-of-plane plaquettes in $d=4$ (one for each choice of transverse direction and orientation).
Thus,
\[
|\mathcal E_{\mathrm{cut}}|
\le 4 m_L
\le 8L,
\qquad\text{and}\qquad
|P_{\perp,\mathrm{cut}}|
\le 4\,|\mathcal E_{\mathrm{cut}}|
\le 32L,
\]
so one may take $C_{\mathcal E}=8$ and $C_{\perp}=32$.

Each plaquette in $P_{\perp,\mathrm{cut}}$ has $4$ boundary links, hence the total link--plaquette incidence count satisfies
$\sum_{e} \deg(e) \le 4\,|P_{\perp,\mathrm{cut}}|$.
Every $e\in\mathcal E_{\mathrm{sh}}$ has $\deg(e)\ge 2$ by definition, so
\[
|\mathcal E_{\mathrm{sh}}|
\le \frac12\sum_{e\in\mathcal E_{\mathrm{sh}}}\deg(e)
\le \frac12\sum_{e}\deg(e)
\le 2\,|P_{\perp,\mathrm{cut}}|
\le 64L.
\]
Thus one may take $C_{\mathrm{sh}}=64$.
The final statement follows by choosing the shared-edge index set $E_{\mathrm{sh}}$ to be any subset of $\mathcal E_{\mathrm{sh}}$
that supports the decomposition $\Delta_{\mathrm{sh}}=\sum_{e\in E_{\mathrm{sh}}}\delta_e$.

% =====================================================================
% : Explicit product-over-edge realization of the shared-edge correction (Main chain, midline strip cut)
% =====================================================================

\begin{lemma}[Explicit shared-edge product for the canonical midline strip cut]\label{lem:routeC_shared_edge_product_midline}
Work in the midline strip cut geometry of Lemma~\ref{lem:routeC_midline_bookkeeping} inside the thick block.
Let $P_{\perp,\mathrm{cut}}$ and $\mathcal E_{\mathrm{sh}}$ be as there, and set $E_{\mathrm{sh}}:=\mathcal E_{\mathrm{sh}}$.

Let $f(g):=\widehat W_{\beta}(g)^{\alpha}$ and $z_\alpha:=\int_G f(g)\,dg$.
For each $e\in E_{\mathrm{sh}}$, define the cut-adjacent transverse incidence number
\[
\deg_{\perp}(e):=\bigl|\{p\in P_{\perp,\mathrm{cut}}:\ e\subset \partial p\}\bigr|,
\]
so that $\deg_{\perp}(e)\ge 2$ by definition of $E_{\mathrm{sh}}$.

After performing the canonical leaf eliminations on all transverse links in $P_{\perp,\mathrm{cut}}$ that are not in $E_{\mathrm{sh}}$,
the remaining cut-adjacent transverse contribution admits the edge-factor form
\begin{equation}\label{eq:routeC_shared_edge_product_form_midline}
1+\Delta_{\mathrm{sh}}
=
\prod_{e\in E_{\mathrm{sh}}} b_e,
\qquad
b_e = \int_G \prod_{j=1}^{\deg_\perp(e)} f(g X_{e,j})\,dg,
\end{equation}
for some $X_{e,1},\dots,X_{e,\deg_\perp(e)}\in G$ measurable with respect to the in-plane field (and any already-integrated transverse leaf data).

Define the decoupled baseline factors $a_e:=z_\alpha^{\deg_\perp(e)}$ and the product difference
\begin{equation}\label{eq:routeC_Delta_sh_product_difference_midline}
\Delta_{\mathrm{sh}}
=
\prod_{e\in E_{\mathrm{sh}}} b_e - \prod_{e\in E_{\mathrm{sh}}} a_e.
\end{equation}
Moreover, for every $e\in E_{\mathrm{sh}}$ one has $0\le a_e,b_e\le 1$ and the pointwise defect bound
\begin{equation}\label{eq:routeC_shared_edge_pointwise_defect_midline}
|b_e-a_e|\le z_{2\alpha}.
\end{equation}
\end{lemma}

\begin{proof}
\emph{(Edge-factor form.)}
By construction, $P_{\perp,\mathrm{cut}}$ consists of out-of-plane plaquettes that are incident to the in-plane cut-link set $\mathcal E_{\mathrm{cut}}$
(Lemma~\ref{lem:routeC_midline_bookkeeping}).
The canonical leaf eliminations integrate, one by one, transverse link variables that appear in exactly one remaining plaquette factor; by Haar bi-invariance,
each such integration contributes the constant $z_\alpha$ and removes that plaquette-variable dependence.
After all such eliminations, the only transverse link variables that remain in $P_{\perp,\mathrm{cut}}$ are precisely the non-in-plane links incident
to at least two cut-adjacent transverse plaquettes, i.e. $\mathcal E_{\mathrm{sh}}$.
In the midline strip geometry, the remaining transverse factors group by shared link $e\in \mathcal E_{\mathrm{sh}}$, producing
\eqref{eq:routeC_shared_edge_product_form_midline} with $\deg_\perp(e)$ shifted copies of $f$; the shifts $X_{e,j}$ are the group-word remainders
of the incident plaquette holonomies with the $e$-variable factored out.

\smallskip
\emph{(Pointwise bounds.)}
Since $0\le f\le 1$, we have $0\le a_e=z_\alpha^{\deg_\perp(e)}\le 1$.
Also $0\le b_e\le 1$ because the integrand is a product of factors in $[0,1]$ and Haar has total mass $1$.
Finally, \eqref{eq:routeC_shared_edge_pointwise_defect_midline} is exactly the higher-degree extension stated in
Remark~\ref{rem:routeC_shared_edge_diagnostic} (take $k=\deg_\perp(e)\ge 2$).
\end{proof}

\begin{corollary}[Shared-edge $L^2$ bound for the midline strip cut]\label{cor:routeC_shared_edge_L2_midline}
In the setting of Lemma~\ref{lem:routeC_shared_edge_product_midline}, the shared-edge correction admits a decomposition
\(
\Delta_{\mathrm{sh}}=\sum_{e\in E_{\mathrm{sh}}}\delta_e
\)
with $|\delta_e|\le z_{2\alpha}$ pointwise. Consequently,
\[
\|\Delta_{\mathrm{sh}}\|_{L^2(\nu_L^{\parallel})}
\le |E_{\mathrm{sh}}|\,z_{2\alpha}
\le C_{\mathrm{sh}}\,L\,z_{2\alpha},
\]
where $C_{\mathrm{sh}}$ is the geometric constant from Lemma~\ref{lem:routeC_midline_bookkeeping} (one may take $C_{\mathrm{sh}}=64$).
\end{corollary}

\begin{proof}
Combine \eqref{eq:routeC_Delta_sh_product_difference_midline} with Corollary~\ref{cor:routeC_shared_edge_decomp_from_telescoping}
(using \eqref{eq:routeC_shared_edge_pointwise_defect_midline}) and then apply Lemma~\ref{lem:routeC_errparallel_bound}.
The final inequality uses \eqref{eq:routeC_midline_size_bounds}.
\end{proof}


\end{proof}

\begin{remark}[Scaling consequence for $\varepsilon_L^{(C)}$]\label{rem:routeC_midline_scaling_epsilonC}
Assume the hypotheses of Corollary~\ref{cor:routeC_surrogate_error_from_influences}.
Using \eqref{eq:routeC_midline_size_bounds} and the trivial bound $|I_{\mathrm{out}}|\le |T_L|=L^2$, we obtain the explicit estimate
\begin{equation}\label{eq:routeC_epsilonC_midline_scaling}
\varepsilon_L^{(C)}
\;\le\;
\frac{t\,\Delta_{\mathrm{out}}}{\beta\,c_{\Psi}}\; L
\;+\;
\frac{C_{\mathrm{sh}}}{c_{\Psi}}\; L\; z_{2\alpha}.
\end{equation}
In the finite-depth entrance window $L^2\simeq \beta$ and under the weak-coupling asymptotic $z_{2\alpha}\lesssim (\alpha\beta)^{-\frac12\dim G}$
(from stationary phase at the identity), the right-hand side satisfies
\[
\varepsilon_L^{(C)}
\;\lesssim\;
\frac{t}{\sqrt{\beta}}
\;+\;
\beta^{\frac12-\frac12\dim G},
\]
and in particular $\varepsilon_L^{(C)}=o_{\beta\to\infty}(1)$ for every compact connected simple Lie group $G$ (where $\dim G\ge 3$).
\end{remark}

\begin{lemma}[Convolution-power damping of matrix coefficients]\label{lem:routeC_matrix_coeff_decay}
Let $G$ be a compact group with normalized Haar measure $dg$, let $w\ge 0$ be a \emph{central} probability density on $G$
($\int w\,dg=1$ and $w(hgh^{-1})=w(g)$), and let $R$ be an irreducible unitary representation
$\pi_R:G\to U(d_R)$ with character $\chi_R$ and normalized character $\varphi_R:=\chi_R/d_R$.
For indices $1\le i,j\le d_R$, define the matrix coefficient function
\[
u_{ij}(g):=\bigl(\pi_R(g)\bigr)_{ij}\in\mathbb C,
\qquad g\in G.
\]
Let $T_w$ be the left-convolution operator on $L^2(G)$,
\[
(T_w f)(g):=\int_G f(hg)\,w(h)\,dh.
\]
Then the following hold.
\begin{enumerate}[label=(\alph*)]
\item \textbf{(Schur orthogonality / $L^2$ norm)} One has
\begin{equation}\label{eq:routeC_schur_uij_norm}
\int_G |u_{ij}(g)|^2\,dg=\frac{1}{d_R},
\qquad\text{equivalently}\qquad
\|u_{ij}\|_{L^2(G)}=\frac{1}{\sqrt{d_R}}.
\end{equation}
\item \textbf{(Convolution coefficient identity)} The coefficient $u_{ij}$ is an eigenfunction of $T_w$:
\begin{equation}\label{eq:routeC_convolution_eigen}
T_w u_{ij}=\lambda_R(w)\,u_{ij},
\qquad\text{where}\qquad
\lambda_R(w):=\int_G \varphi_R(h)\,w(h)\,dh.
\end{equation}
\item \textbf{(Convolution powers and variance bound)} For every integer $k\ge 1$,
\begin{equation}\label{eq:routeC_convolution_power_norm}
T_w^k u_{ij}=\lambda_R(w)^k\,u_{ij}
\quad\text{and}\quad
\|T_w^k u_{ij}\|_{L^2(G)}^2=\frac{|\lambda_R(w)|^{2k}}{d_R}.
\end{equation}
Equivalently, if $A=X_1\cdots X_k$ with $X_m$ i.i.d.\ having density $w$ (so $A\sim w^{*k}$), then
the averaged matrix coefficient
\[
F_{ij}(g):=\mathbb E\bigl[u_{ij}(Ag)\bigr]=\int_G u_{ij}(hg)\,w^{*k}(h)\,dh
\]
satisfies $F_{ij}(g)=\lambda_R(w)^k\,u_{ij}(g)$ and therefore has Haar $L^2$--variance
$\int |F_{ij}(g)|^2dg=|\lambda_R(w)|^{2k}/d_R$.
\end{enumerate}
\end{lemma}

\begin{proof}
We prove the three parts in order.

\medskip
\noindent\textbf{(a) Schur orthogonality and \eqref{eq:routeC_schur_uij_norm}.}
Fix indices $j,\ell$ and let $E_{j\ell}$ be the $d_R\times d_R$ matrix unit.
Define the operator
\[
A_{j\ell}:=\int_G \pi_R(g)\,E_{j\ell}\,\pi_R(g)^{-1}\,dg.
\]
By left--right invariance of Haar measure, for every $h\in G$ we have
\[
\pi_R(h)\,A_{j\ell}\,\pi_R(h)^{-1}
=\int_G \pi_R(hg)\,E_{j\ell}\,\pi_R(hg)^{-1}\,dg
=\int_G \pi_R(g)\,E_{j\ell}\,\pi_R(g)^{-1}\,dg
=A_{j\ell},
\]
so $A_{j\ell}$ commutes with $\pi_R(h)$ for all $h\in G$.
Since $R$ is irreducible, Schur's lemma implies $A_{j\ell}=c_{j\ell}\,I$ for some scalar $c_{j\ell}$.
Taking traces and using cyclicity gives
\[
c_{j\ell}\,d_R=\mathrm{tr}(A_{j\ell})
=\int_G \mathrm{tr}\!\bigl(\pi_R(g)\,E_{j\ell}\,\pi_R(g)^{-1}\bigr)\,dg
=\int_G \mathrm{tr}(E_{j\ell})\,dg
=\delta_{j\ell}.
\]
Hence $c_{j\ell}=\delta_{j\ell}/d_R$.
Now take the $(i,k)$ entry of $A_{j\ell}$:
\[
(A_{j\ell})_{ik}
=\int_G \bigl(\pi_R(g)\bigr)_{ij}\,\bigl(\pi_R(g)^{-1}\bigr)_{\ell k}\,dg.
\]
Because $\pi_R(g)$ is unitary, $\pi_R(g)^{-1}=\pi_R(g)^{\ast}$ and therefore
$\bigl(\pi_R(g)^{-1}\bigr)_{\ell k}=\overline{\bigl(\pi_R(g)\bigr)_{k\ell}}$.
Thus
\[
\int_G \bigl(\pi_R(g)\bigr)_{ij}\,\overline{\bigl(\pi_R(g)\bigr)_{k\ell}}\,dg
=(A_{j\ell})_{ik}
=\bigl(c_{j\ell}I\bigr)_{ik}
=\frac{1}{d_R}\,\delta_{j\ell}\,\delta_{ik}.
\]
Setting $k=i$ and $\ell=j$ yields $\int |u_{ij}(g)|^2dg=1/d_R$, i.e.\ \eqref{eq:routeC_schur_uij_norm}.

\medskip
\noindent\textbf{(b) Convolution coefficient identity \eqref{eq:routeC_convolution_eigen}.}
Define the $d_R\times d_R$ matrix
\[
M:=\int_G \pi_R(h)\,w(h)\,dh.
\]
Using that $w$ is central and Haar is conjugation-invariant, for any $g\in G$,
\[
\pi_R(g)\,M\,\pi_R(g)^{-1}
=\int_G \pi_R(ghg^{-1})\,w(h)\,dh
=\int_G \pi_R(h)\,w(h)\,dh
=M.
\]
Thus $M$ commutes with $\pi_R(g)$ for all $g$, so by Schur's lemma $M=c\,I$ for some scalar $c$.
Taking trace gives
\[
\begin{aligned}
c\,d_R&=\mathrm{tr}(M)=\int_G \chi_R(h)\,w(h)\,dh,\\
\text{so}\qquad
c&=\frac{1}{d_R}\int_G \chi_R(h)\,w(h)\,dh
=\int_G \varphi_R(h)\,w(h)\,dh
=\lambda_R(w).
\end{aligned}
\]
Now compute, for any $g\in G$,
\begin{align*}
(T_w u_{ij})(g)
&=\int_G u_{ij}(hg)\,w(h)\,dh
=\int_G \bigl(\pi_R(hg)\bigr)_{ij}\,w(h)\,dh
\\
&=\int_G \sum_{a=1}^{d_R}\bigl(\pi_R(h)\bigr)_{ia}\,\bigl(\pi_R(g)\bigr)_{aj}\,w(h)\,dh
=\sum_{a=1}^{d_R} M_{ia}\,\bigl(\pi_R(g)\bigr)_{aj}
\\
&=\sum_{a=1}^{d_R} (c\,\delta_{ia})\,\bigl(\pi_R(g)\bigr)_{aj}
=c\,\bigl(\pi_R(g)\bigr)_{ij}
=\lambda_R(w)\,u_{ij}(g),
\end{align*}
which is exactly \eqref{eq:routeC_convolution_eigen}.

\medskip
\noindent\textbf{(c) Convolution powers and \eqref{eq:routeC_convolution_power_norm}.}
Applying $T_w$ repeatedly and using \eqref{eq:routeC_convolution_eigen} gives $T_w^k u_{ij}=\lambda_R(w)^k u_{ij}$.
Taking $L^2(G)$ norms and using \eqref{eq:routeC_schur_uij_norm} yields
\[
\|T_w^k u_{ij}\|_{L^2(G)}^2
=|\lambda_R(w)|^{2k}\,\|u_{ij}\|_{L^2(G)}^2
=\frac{|\lambda_R(w)|^{2k}}{d_R}.
\]
For the probabilistic restatement, note that if $A\sim w^{*k}$ then
$g\mapsto\mathbb E[u_{ij}(Ag)]$ equals $(T_{w^{*k}}u_{ij})(g)$, and $T_{w^{*k}}=T_w^k$.
\end{proof}

\begin{lemma}[Efron--Stein variance bound from $L^2$ influences]\label{lem:routeC_influence_variance}
Let $\xi=(\xi_1,\dots,\xi_m)$ be independent random variables on a product probability space, and let $Y=Y(\xi)$ be square-integrable.
For each $a\in\{1,\dots,m\}$, let $\xi^{(a)}$ be obtained by resampling only $\xi_a$ (independently of everything else), and set $Y^{(a)}:=Y(\xi^{(a)})$.
Then
\begin{equation}\label{eq:routeC_efron_stein}
\mathrm{Var}(Y)\le \frac12\sum_{a=1}^m\mathbf E\bigl[(Y-Y^{(a)})^2\bigr].
\end{equation}
In particular, if $X=X(\xi)$ satisfies $\|X-X^{(a)}\|_2\le \Delta/\beta$ for all $a$, then for every $t\ge 0$,
\begin{equation}\label{eq:routeC_var_exp_influence}
\sqrt{\mathrm{Var}(e^{-tX})}
\le \frac{t\,\Delta}{\beta\sqrt2}\,\sqrt m.
\end{equation}
\end{lemma}

\begin{proof}
Eq.~\eqref{eq:routeC_efron_stein} is the standard Efron--Stein inequality.
For completeness, one may write $\mathrm{Var}(Y)=\frac12\mathbf E[(Y-Y')^2]$ where $Y'$ is an independent copy built from an independent copy $\xi'$ of $\xi$,
then reveal the coordinates one by one and compare $Y(\xi)$ to $Y(\xi_1,\dots,\xi_{a-1},\xi_a',\dots,\xi_m')$ to obtain the telescoping bound.
The consequence~\eqref{eq:routeC_var_exp_influence} follows by applying~\eqref{eq:routeC_efron_stein} to $Y=e^{-tX}$ and using
$|e^{-tX}-e^{-tX^{(a)}}|\le t|X-X^{(a)}|$.
\end{proof}

\begin{lemma}[Cut-local surrogate $\Rightarrow$ $\chi^2$ flatness of the effective tilt]\label{lem:routeC_surrogate_implies_chi2_flatness}
Assume Proposition~\ref{ass:routeC_cut_local_surrogate}.
Then the normalized cut surrogate $\widetilde\Psi_{\mathrm{cut}}$ from \eqref{eq:routeC_surrogate_def} satisfies
\begin{equation}\label{eq:routeC_surrogate_L2_flatness}
\bigl\|\widetilde\Psi_{\mathrm{cut}}-1\bigr\|_{L^2(\nu_L^{\parallel})}
\ \le\
\frac{t\,\Delta_X}{\beta\sqrt2\,c_{\Psi}}\sqrt{m_L}.
\end{equation}
Consequently,
\begin{equation}\label{eq:routeC_effective_tilt_L2_flatness}
\bigl\|\widetilde\Psi_L^{(\alpha)}-1\bigr\|_{L^2(\nu_L^{\parallel})}
\ \le\
\epsilon_L^{(C)}+\frac{t\,\Delta_X}{\beta\sqrt2\,c_{\Psi}}\sqrt{m_L}.
\end{equation}
In particular, the (unconditioned) in-plane $\chi^2$ divergence obeys
\begin{equation}\label{eq:routeC_inplane_chi2_from_surrogate}
D_{\chi^2}\!\left(\nu_L^{(\alpha)}|_{\mathscr F_{\parallel}}\Big\| \nu_L^{\parallel}|_{\mathscr F_{\parallel}}\right)
=
\bigl\|\widetilde\Psi_L^{(\alpha)}-1\bigr\|_{L^2(\nu_L^{\parallel})}^2
\ \le\
\left(\epsilon_L^{(C)}+\frac{t\,\Delta_X}{\beta\sqrt2\,c_{\Psi}}\sqrt{m_L}\right)^2.
\end{equation}
\end{lemma}

\begin{proof}
By definition $\widetilde\Psi_{\mathrm{cut}}=Y/E_{\nu_L^{\parallel}}[Y]$ with $Y=e^{-tX_{\mathrm{cut}}}$.
Since $E_{\nu_L^{\parallel}}[Y]\ge c_{\Psi}$ by \eqref{eq:routeC_surrogate_mean_lower},
\[
\bigl\|\widetilde\Psi_{\mathrm{cut}}-1\bigr\|_{2}^2
=\mathrm{Var}\!\left(\frac{Y}{E[Y]}\right)
=\frac{\mathrm{Var}(Y)}{E[Y]^2}
\le \frac{\mathrm{Var}(Y)}{c_{\Psi}^2}.
\]
Applying Lemma~\ref{lem:routeC_influence_variance} with the influence bound \eqref{eq:routeC_cut_influence_L2} yields
$\sqrt{\mathrm{Var}(Y)}\le (t\Delta_X/(\beta\sqrt2))\sqrt{m_L}$,
hence \eqref{eq:routeC_surrogate_L2_flatness}.
The bound \eqref{eq:routeC_effective_tilt_L2_flatness} follows from the triangle inequality and \eqref{eq:routeC_surrogate_L2}.
Finally, since $\widetilde\Psi_L^{(\alpha)}$ is the Radon--Nikodym derivative of $\nu_L^{(\alpha)}|_{\mathscr F_{\parallel}}$ with respect to
$\nu_L^{\parallel}|_{\mathscr F_{\parallel}}$ (Lemma~\ref{lem:routeC_effective_planar_tilt}),
one has $D_{\chi^2}=\|\widetilde\Psi_L^{(\alpha)}-1\|_2^2$, giving \eqref{eq:routeC_inplane_chi2_from_surrogate}.
\end{proof}

\begin{lemma}[$L^2$ flatness on an event $\Rightarrow$ conditional $\chi^2$ control]\label{lem:chi2_from_L2_on_event}
Let $\nu$ be a probability measure on $(X,\mathcal F)$, let $\Omega\in\mathcal F$ satisfy $\nu(\Omega)>0$, and let $r\ge 0$ satisfy $E_{\nu}[r]=1$.
Define $\mu:=r\cdot\nu$.
Set $\delta_{\Omega}:=\|r-1\|_{L^2(\nu^{\Omega})}$, where $\nu^{\Omega}:=\nu(\,\cdot\,|\Omega)$.
If $\delta_{\Omega}<1$, then
\begin{equation}\label{eq:chi2_from_L2_on_event}
D_{\chi^2}(\mu^{\Omega}\|\nu^{\Omega})
\ \le\
\frac{\delta_{\Omega}^2}{(1-\delta_{\Omega})^2}.
\end{equation}
In particular, $\delta_{\Omega}\le \|r-1\|_{L^2(\nu)}/\sqrt{\nu(\Omega)}$, so if $\|r-1\|_{L^2(\nu)}\le \frac12\sqrt{\nu(\Omega)}$ then
\begin{equation}\label{eq:chi2_from_L2_on_event_simple}
D_{\chi^2}(\mu^{\Omega}\|\nu^{\Omega})
\ \le\
\frac{4}{\nu(\Omega)}\,\|r-1\|_{L^2(\nu)}^2.
\end{equation}
\end{lemma}

\begin{proof}
Write $\nu^{\Omega}(\,\cdot\,)=\nu(\,\cdot\,\cap\Omega)/\nu(\Omega)$.
Since $d\mu^{\Omega}/d\nu^{\Omega}=r/E_{\nu^{\Omega}}[r]$ on $\Omega$,
\[
D_{\chi^2}(\mu^{\Omega}\|\nu^{\Omega})
=\mathrm{Var}_{\nu^{\Omega}}\!\left(\frac{r}{E_{\nu^{\Omega}}[r]}\right)
=\frac{\mathrm{Var}_{\nu^{\Omega}}(r)}{E_{\nu^{\Omega}}[r]^2}.
\]
Now $\mathrm{Var}_{\nu^{\Omega}}(r)\le E_{\nu^{\Omega}}[(r-1)^2]=\delta_{\Omega}^2$.
Also $|E_{\nu^{\Omega}}[r]-1|=|E_{\nu^{\Omega}}[r-1]|\le \delta_{\Omega}$ by Cauchy--Schwarz,
so $E_{\nu^{\Omega}}[r]\ge 1-\delta_{\Omega}$.
This gives \eqref{eq:chi2_from_L2_on_event}.
Finally, $\delta_{\Omega}^2=E_{\nu}[(r-1)^2\mathbf 1_{\Omega}]/\nu(\Omega)\le \|r-1\|_{L^2(\nu)}^2/\nu(\Omega)$,
and \eqref{eq:chi2_from_L2_on_event_simple} follows by inserting this bound into \eqref{eq:chi2_from_L2_on_event}.
\end{proof}

\begin{lemma}[Main-chain covariance pinning under cut-local tilts]\label{lem:routeC_covariance_pinning}
Work on a planar reference tile of side length $L$ with boundary holonomy $\Phi_L$ and planar reference single-plaquette density $w_1^{\parallel}$.
Let $\mathcal F_{\mathrm{cut}}$ be a sub-$\sigma$-algebra (the ``cut'' information) such that:
\begin{enumerate}
\item (\emph{Cut size}) $\mathcal F_{\mathrm{cut}}$ is generated by at most $m_L$ plaquette variables with $m_L\le C_{\mathrm{cut}}L$.
\item (\emph{Tilt locality}) The transverse-energy tilt observable $X$ entering $\overline R:=e^{-tX}$ is $\mathcal F_{\mathrm{cut}}$-measurable.
\item (\emph{$L^2$ influence bound}) Writing $X^{(a)}$ for the version of $X$ in which the $a$th generating cut variable is resampled
(independently, holding all other cut variables fixed), assume
\begin{equation}\label{eq:routeC_influence_assumption}
\bigl\|X-X^{(a)}\bigr\|_{L^2(\parallel)}\ \le\ \frac{\Delta_X}{\beta}
\qquad\text{for each }a=1,\dots,m_L .
\end{equation}
\end{enumerate}
Then for any irrep $R\neq\mathbf 1$ and any $t\ge 0$,
\begin{equation}\label{eq:routeC_cov_bound_general}
\Bigl|\mathrm{Cov}_{\parallel}\bigl(\varphi_R(\Phi_L), e^{-tX}\bigr)\Bigr|
\ \le\
\sqrt{\mathrm{Var}_{\parallel}\bigl(e^{-tX}\bigr)}\cdot \bigl\|\mathbb{E}_{\parallel}[\varphi_R(\Phi_L)\mid\mathcal F_{\mathrm{cut}}]\bigr\|_{L^2(\parallel)}.
\end{equation}
Moreover,
\begin{equation}\label{eq:routeC_tilt_L2_bound}
\sqrt{\mathrm{Var}_{\parallel}\bigl(e^{-tX}\bigr)}\ \le\ t\,\frac{\Delta_X}{\beta}\,\sqrt{m_L},
\end{equation}
and, writing $k_L:=|T_L|-m_L$, one has the representation-mixing bound
\begin{equation}\label{eq:routeC_condexp_decay}
\Big\|\mathbb{E}_{\parallel}[\varphi_R(\Phi_L)\mid\mathcal F_{\mathrm{cut}}]\Big\|_{L^2(\parallel)}
\;\le\;
\frac{|\lambda_R(w_1^{\parallel})|^{k_L}}{d_R}\;
\sqrt{\big\|(w_1^{\parallel})^{*m_L}\big\|_{L^\infty(G)}}.
\end{equation}
Consequently,
\begin{equation}\label{eq:routeC_covariance_bound}
\Big|\mathrm{Cov}_{\parallel}\big(\varphi_R(\Phi_L), e^{-tX}\big)\Big|
\;\le\;
\frac{t\,\Delta_X}{\beta}\,\sqrt{m_L}\;
\frac{|\lambda_R(w_1^{\parallel})|^{k_L}}{d_R}\;
\sqrt{\big\|(w_1^{\parallel})^{*m_L}\big\|_{L^\infty(G)}}.
\end{equation}
In particular, if $m_L\le \tfrac12|T_L|$ (hence $k_L\ge \tfrac12|T_L|$) then
\[
\Bigl|\mathrm{Cov}_{\parallel}\bigl(\varphi_R(\Phi_L), e^{-tX}\bigr)\Bigr|
\ \le\ 
\frac{t\Delta_X}{\beta}\,\sqrt{m_L}\cdot 
\frac{|\lambda_R(w_1^{\parallel})|^{|T_L|/2}}{d_R}\;
\sqrt{\big\|(w_1^{\parallel})^{*m_L}\big\|_{L^\infty(G)}}.
\]

Moreover, since $w_1^{\parallel}$ is of positive type, each convolution power $(w_1^{\parallel})^{*m}$ is positive definite,
hence $\|(w_1^{\parallel})^{*m}\|_{L^\infty(G)}=(w_1^{\parallel})^{*m}(e)$.
Using the one-plaquette coefficient bound $\lambda_S(w_1^{\parallel})\le e^{-c_{\mathrm{plaq}}C_2(S)/\beta}$ we obtain
\begin{equation}\label{eq:routeC_wm_sup_heatkernel}
\big\|(w_1^{\parallel})^{*m}\big\|_{L^\infty(G)}
\;\le\;
\sum_{S} d_S^2\,e^{-c_{\mathrm{plaq}} C_2(S)\,m/\beta}
\;=\;
H_G\!\left(c_{\mathrm{plaq}}\,\frac{m}{\beta}\right).
\end{equation}

\end{lemma}

\begin{proof}
Since $X$ is $\mathcal F_{\mathrm{cut}}$--measurable, so is $\overline R=e^{-tX}$, hence
\[
\mathrm{Cov}_{\parallel}\big(\varphi_R(\Phi_L),\overline R\big)
=\mathrm{Cov}_{\parallel}\Big(\mathbb E_{\parallel}[\varphi_R(\Phi_L)\mid\mathcal F_{\mathrm{cut}}],\overline R\Big).
\]
Cauchy--Schwarz gives the general bound \eqref{eq:routeC_cov_bound_general}.

The influence assumption \eqref{eq:routeC_influence_assumption} together with the Efron--Stein inequality
(Lemma~\ref{lem:routeC_influence_variance}) yields
\begin{equation}\label{eq:routeC_tilt_L2_bound_proof}
\big\|\overline R-\mathbb E_{\parallel}[\overline R]\big\|_{L^2(\parallel)}
\;\le\;\frac{t\,\Delta_X}{\beta}\,\sqrt{m_L},
\end{equation}
which is \eqref{eq:routeC_tilt_L2_bound}.

It remains to bound $\|\mathbb E_{\parallel}[\varphi_R(\Phi_L)\mid\mathcal F_{\mathrm{cut}}]\|_{L^2(\parallel)}$.
Write $N:=|T_L|$ and index the plaquettes in a fixed surface order as $U_1,\dots,U_N$ so that
$\Phi_L=\prod_{\ell=1}^{N} U_\ell$.
Let $C\subseteq\{1,\dots,N\}$ be the cut indices, $|C|=m_L$, and $C^c$ its complement, $|C^c|=k_L$.

Let $\pi_R$ be a unitary realization of $R$.
Centrality of $w_1^{\parallel}$ implies that the matrix
\[
M_R:=\int_G \pi_R(g)\, w_1^{\parallel}(g)\,dg
\]
commutes with $\pi_R(h)$ for all $h\in G$, hence $M_R=\lambda_R(w_1^{\parallel})\,\mathrm{Id}$ by Schur's lemma.
Conditioning on $\mathcal F_{\mathrm{cut}}=\sigma(U_\ell:\ell\in C)$ and integrating the variables
$\{U_\ell:\ell\in C^c\}$ sequentially therefore gives the exact identity
\[
\mathbb E_{\parallel}[\pi_R(\Phi_L)\mid\mathcal F_{\mathrm{cut}}]
=\lambda_R(w_1^{\parallel})^{k_L}\,\pi_R(\Phi_{\mathrm{cut}}),
\qquad
\Phi_{\mathrm{cut}}:=\prod_{\ell\in C}^{\!\!\mathrm{(surface\ order)}} U_\ell .
\]
Taking normalized traces yields
$\mathbb E_{\parallel}[\varphi_R(\Phi_L)\mid\mathcal F_{\mathrm{cut}}]
=\lambda_R(w_1^{\parallel})^{k_L}\,\varphi_R(\Phi_{\mathrm{cut}})$.
Since $\Phi_{\mathrm{cut}}$ is the product of $m_L$ i.i.d.\ variables with density $w_1^{\parallel}$, it has density
$(w_1^{\parallel})^{*m_L}$ with respect to Haar, so
\[
\Big\|\mathbb E_{\parallel}[\varphi_R(\Phi_L)\mid\mathcal F_{\mathrm{cut}}]\Big\|_{L^2(\parallel)}^2
= \lambda_R(w_1^{\parallel})^{2k_L}
\int_G |\varphi_R(g)|^2\,(w_1^{\parallel})^{*m_L}(g)\,dg.
\]
Bounding the integral by $\|(w_1^{\parallel})^{*m_L}\|_{L^\infty(G)}\int_G |\varphi_R(g)|^2\,dg$
and using character orthogonality $\int_G |\chi_R(g)|^2\,dg = 1$ together with $\varphi_R=\chi_R/d_R$ (hence $\int_G |\varphi_R(g)|^2\,dg = 1/d_R^2$)
gives \eqref{eq:routeC_condexp_decay}.

Finally, insert \eqref{eq:routeC_tilt_L2_bound_proof} and \eqref{eq:routeC_condexp_decay} into
\eqref{eq:routeC_cov_bound_general} to obtain \eqref{eq:routeC_covariance_bound}.
\end{proof}

\paragraph{A concrete cut $\sigma$--algebra and the $1/\beta$ influence bound}\label{sec:routeC_cut_instantiation}

Lemma~\ref{lem:routeC_covariance_pinning} was stated for an abstract cut $\sigma$--algebra $\mathcal F_{\mathrm{cut}}$.
For Main chain we take $\mathcal F_{\mathrm{cut}}$ to be generated by a unit-width strip of plaquettes that separates the tile into two large sub-tiles.

Fix $L\ge 2$ and, for notational simplicity, assume $L$ is even.
Let $T_L$ be the $L\times L$ set of in-plane plaquettes in the planar reference tile.
Write $T_L=T_L^{-}\cup T_L^{+}$ for the left/right decomposition by the vertical midline and let
\[
T_{\mathrm{cut}}
:=
\{p\in T_L:\ p\text{ intersects the midline between }T_L^{-}\text{ and }T_L^{+}\}.
\]
Then $|T_{\mathrm{cut}}|\le 2L$.
Define the cut $\sigma$--algebra and the cut-local energy
\begin{equation}\label{eq:routeC_cut_sigma_def}
\mathcal F_{\mathrm{cut}}:=\sigma\bigl(\{U_p\}_{p\in T_{\mathrm{cut}}}\bigr),
\qquad
X_{\mathrm{cut}}:=\sum_{p\in T_{\mathrm{cut}}} A(U_p).
\end{equation}
(The choice of a vertical midline is inessential; any rectilinear cut of length $O(L)$ gives the same bounds up to constants.)

\begin{lemma}[Verification of Lemma~\ref{lem:routeC_covariance_pinning} hypotheses for $X_{\mathrm{cut}}$]\label{lem:routeC_cut_instantiation}
Under the planar reference product representation of the plaquette variables on $T_L$:
\begin{enumerate}
\item $\mathcal F_{\mathrm{cut}}$ is generated by $m_L:=|T_{\mathrm{cut}}|\le 2L$ plaquette variables, hence $m_L\le C_{\mathrm{cut}}L$ with $C_{\mathrm{cut}}=2$;
\item $X_{\mathrm{cut}}$ is $\mathcal F_{\mathrm{cut}}$-measurable;
\item the $L^2$ influence bound \eqref{eq:routeC_influence_assumption} holds with $\Delta_X:=2\sqrt{C_{\mathrm{plaq},2}}$, where $C_{\mathrm{plaq},2}$ is the constant in Lemma~\ref{lem:BM_plaq_defect_L2}.
\end{enumerate}
\end{lemma}
\begin{proof}
Items (1)--(2) are immediate from \eqref{eq:routeC_cut_sigma_def} and the cardinality estimate $|T_{\mathrm{cut}}|\le 2L$.

For (3), fix $p\in T_{\mathrm{cut}}$ and let $U_p'$ be an independent resample of $U_p$ (with the same marginal $w_1^{\parallel}$).
Since $X_{\mathrm{cut}}$ is a sum over $T_{\mathrm{cut}}$, resampling $U_p$ affects exactly one term:
\[
X_{\mathrm{cut}}-X_{\mathrm{cut}}^{(p)} = A(U_p)-A(U_p').
\]
Therefore
\[
\|X_{\mathrm{cut}}-X_{\mathrm{cut}}^{(p)}\|_{L^2(\parallel)}
\le \|A(U_p)\|_{L^2(\parallel)}+\|A(U_p')\|_{L^2(\parallel)}
=2\|A(U_p)\|_{L^2(\parallel)}
\le 2\frac{\sqrt{C_{\mathrm{plaq},2}}}{\beta},
\]
where the last inequality is Lemma~\ref{lem:BM_plaq_defect_L2}.
This is exactly \eqref{eq:routeC_influence_assumption} with $\Delta_X=2\sqrt{C_{\mathrm{plaq},2}}$.
\end{proof}

% : Canonical $X_{\mathrm{out}}$ closes \eqref{eq:routeC_Xout_influence} for the midline instantiation.
\begin{lemma}[Canonical choice of $X_{\mathrm{out}}$ and verification of \eqref{eq:routeC_Xout_influence}]\label{lem:routeC_Xout_influence_instantiation}
Fix $T_{\mathrm{cut}}\subset T_L$ as in Lemma~\ref{lem:routeC_cut_instantiation} and write the complementary set of in-plane plaquettes
\[
T_{\mathrm{out}}:=T_L\setminus T_{\mathrm{cut}}\qquad (|T_{\mathrm{out}}|=|T_L|-|T_{\mathrm{cut}}|=L^2+O(L)).
\]
Define the ``outside'' planar energy
\begin{equation}\label{eq:routeC_Xout_def}
X_{\mathrm{out}}:=\sum_{p\in T_{\mathrm{out}}} A(U_p).
\end{equation}
Then $X_{\mathrm{out}}$ depends only on the planar plaquette coordinates indexed by
\begin{equation}\label{eq:routeC_Iout_def}
I_{\mathrm{out}}:=T_{\mathrm{out}},\qquad |I_{\mathrm{out}}|=|T_{\mathrm{out}}|=L^2+O(L).
\end{equation}
Moreover, for every $a\in I_{\mathrm{out}}$ one has the $L^2$ influence bound
\begin{equation}\label{eq:routeC_Xout_influence_instantiated}
\bigl\|X_{\mathrm{out}}-X_{\mathrm{out}}^{(a)}\bigr\|_{L^2(\nu_L^{\parallel})}\le \frac{\Delta_{\mathrm{out}}}{\beta},
\qquad \Delta_{\mathrm{out}}:=2\sqrt{C_{\mathrm{plaq},2}},
\end{equation}
with $C_{\mathrm{plaq},2}$ from Lemma~\ref{lem:BM_plaq_defect_L2}. In particular, \eqref{eq:routeC_Xout_influence} holds for this choice of $X_{\mathrm{out}}$ (and the auxiliary $\mu_\omega$ is irrelevant here since $X_{\mathrm{out}}$ is planar).
\end{lemma}
\begin{proof}
By construction $X_{\mathrm{out}}$ is a sum of the in-plane plaquette actions over $T_{\mathrm{out}}$, hence depends only on the coordinates in $I_{\mathrm{out}}$.
Fix $a\in I_{\mathrm{out}}$ and let $p_a\in T_{\mathrm{out}}$ denote the corresponding plaquette.
Under Efron--Stein resampling at coordinate $a$ we have
\[
X_{\mathrm{out}}-X_{\mathrm{out}}^{(a)}=A(U_{p_a})-A(U'_{p_a}),
\]
where $U_{p_a}$ and $U'_{p_a}$ are i.i.d. with the one-plaquette Wilson law $w_1^{\parallel}$ (tree-gauge independence, Lemma~\ref{lem:planar_exact_convolution}).
Therefore, by the triangle inequality in $L^2$ and Lemma~\ref{lem:BM_plaq_defect_L2},
\[
\bigl\|X_{\mathrm{out}}-X_{\mathrm{out}}^{(a)}\bigr\|_2
\le \|A(U_{p_a})\|_2+\|A(U'_{p_a})\|_2
=2\|A(U_{p_a})\|_2
\le \frac{2\sqrt{C_{\mathrm{plaq},2}}}{\beta}.
\]
This is \eqref{eq:routeC_Xout_influence_instantiated}.
\end{proof}


\noindent
With this choice of $\mathcal F_{\mathrm{cut}}$ and $X_{\mathrm{cut}}$, the covariance bound \eqref{eq:routeC_covariance_bound} in Lemma~\ref{lem:routeC_covariance_pinning}
is now completely explicit.













%%%%%%%%%%%%%%%%%%%%%%%%%%%%%%%%%%%%%%%%%%%%%%%%%%%%%%%%%%%%%%%%%%%%%%%%%%%%%%%
\subsubsection{Explicit weak-window bound for the per-step transfer defect \texorpdfstring{$\|\Delta_n\|_2$}{||Delta\_n||\_2}}
\label{sec:explicit_Delta_n_bound}

In the finite-depth entrance/transfer criterion of Section~\ref{sec:finite_transfer} one needs
explicit, scale-indexed bounds on the $L^2$--size of the \emph{normalized} effective-tilt error.
The Main chain surrogate machinery already provides such bounds in terms of (i) outside/cut influences
and (ii) the shared-edge correction. We now package these ingredients into an explicit
$n$--dependent estimate (dyadic instantiation $L=2^n$) that is directly usable as a transfer budget
$\epsilon_n(\beta)$.

\medskip

\noindent\textbf{Notation.}
For a tile of side length $L$ (with the midline strip cut), set
\begin{equation}\label{eq:h52_DeltaL_def}
\Delta_L \;:=\; \widetilde\Psi_L^{(\alpha)}-1.
\end{equation}
where $\widetilde\Psi_L^{(\alpha)}$ is the normalized effective tilt from
Lemma~\ref{lem:routeC_surrogate_implies_chi2_flatness} (Main chain).

\begin{proposition}[Master Main chain weak-window bound for $\|\Delta_L\|_2$]
\label{prop:h52_explicit_DeltaL_bound}
Assume the Main chain setup of Remark~\ref{rem:routeC_midline_scaling_epsilonC}, with the midline strip cut
and associated parameters $t=\beta(1-\alpha)$, $m_L=|T_{\mathrm{cut}}|$ and constants
$c_\Psi>0$, $C_{\mathrm{sh}}>0$, $\Delta_{\mathrm{out}},\Delta_X>0$ as in
Remark~\ref{rem:routeC_midline_scaling_epsilonC} and
Lemma~\ref{lem:routeC_Xout_influence_instantiation}.
Then for every $L\ge 2$,
\begin{equation}
\label{eq:h52_master_bound_DeltaL}
\|\Delta_L\|_2
\;\le\;
\frac{t}{\beta\,c_\Psi}\Big(\Delta_{\mathrm{out}}\,L+\Delta_X\,\sqrt{L}\Big)
\;+\;
\frac{C_{\mathrm{sh}}}{c_\Psi}\,L\,z_{2\alpha}(\beta).
\end{equation}
\end{proposition}

\begin{proof}
By Lemma~\ref{lem:routeC_surrogate_implies_chi2_flatness},
\[
\|\Delta_L\|_2
=
\|\widetilde\Psi_L^{(\alpha)}-1\|_2
\le
\epsilon_L^{(C)}
+
\frac{t\,\Delta_X}{\beta\sqrt2\,c_\Psi}\sqrt{m_L}.
\]
In the midline strip geometry one has $m_L\le 2L$, hence
$\frac{1}{\sqrt2}\sqrt{m_L}\le \sqrt{L}$.
Insert the midline surrogate bound of Remark~\ref{rem:routeC_midline_scaling_epsilonC},
\[
\epsilon_L^{(C)}\le \frac{t\,\Delta_{\mathrm{out}}}{\beta\,c_\Psi}L+\frac{C_{\mathrm{sh}}}{c_\Psi}L\,z_{2\alpha}(\beta),
\]
and combine the two estimates to obtain~\eqref{eq:h52_master_bound_DeltaL}.
\end{proof}

\begin{corollary}[Dyadic instantiation $L=2^n$ under the pivot schedule $\alpha(L)$]
\label{cor:h52_Delta_n_bound_dyadic}
Fix $\kappa>0$ and define the pivot
\[
\alpha(L):=1-\frac{\kappa}{|P_\perp(\mathbb B_L)|}
\qquad\text{with}\qquad
|P_\perp(\mathbb B_L)|=8L(L+1)
\]
(as in Remark~\ref{rem:routeC_midline_scaling_epsilonC}), so that
$t=\beta(1-\alpha(L))=\beta\kappa/(8L(L+1))$.
For dyadic $L=2^n$ ($n\in\mathbb N$) write $\alpha_n:=\alpha(2^n)$ and
$\Delta_n:=\Delta_{2^n}$. Then, for every $n\ge 1$,
\begin{equation}
\label{eq:h52_dyadic_bound_Delta_n_exact}
\|\Delta_n\|_2
\;\le\;
\frac{\kappa\,\Delta_{\mathrm{out}}}{8\,c_\Psi}\frac{1}{2^n+1}
+
\frac{\kappa\,\Delta_X}{8\,c_\Psi}\frac{1}{2^{n/2}(2^n+1)}
+
\frac{C_{\mathrm{sh}}}{c_\Psi}\,2^n\,z_{2\alpha_n}(\beta).
\end{equation}
In particular, using $\frac{1}{2^n+1}\le 2^{-n}$ and $\frac{1}{2^{n/2}(2^n+1)}\le 2^{-3n/2}$,
\begin{equation}
\label{eq:h52_dyadic_bound_Delta_n_simplified}
\|\Delta_n\|_2
\;\le\;
\frac{\kappa\,\Delta_{\mathrm{out}}}{8\,c_\Psi}\,2^{-n}
+
\frac{\kappa\,\Delta_X}{8\,c_\Psi}\,2^{-3n/2}
+
\frac{C_{\mathrm{sh}}}{c_\Psi}\,2^n\,z_{2\alpha_n}(\beta).
\end{equation}
\end{corollary}

\begin{proof}
Insert $t/\beta=\kappa/(8L(L+1))$ into~\eqref{eq:h52_master_bound_DeltaL} and set $L=2^n$.
This yields~\eqref{eq:h52_dyadic_bound_Delta_n_exact}. The simplified bound
\eqref{eq:h52_dyadic_bound_Delta_n_simplified} follows from the elementary inequalities
$\frac{1}{2^n+1}\le 2^{-n}$ and $\frac{1}{2^{n/2}(2^n+1)}\le 2^{-3n/2}$.
\end{proof}

\begin{lemma}[An explicit weak-coupling bound for $z_{2\alpha}(\beta)$]
\label{lem:h52_z2a_explicit_bound}
Let $G$ be a compact connected Lie group with $\dim(G)=d_G$ and Haar measure $dg$.
Assume the one-plaquette Wilson action $A_\pi(g)$ has a unique minimum at $e$ and admits:
\begin{enumerate}
\item a coordinate chart $g=\exp X$ on a neighborhood $\mathcal U_0$ of $e$ with radius $r_0>0$,
\item constants $m_->0$ and $J_{\max}<\infty$ such that for all $\|X\|\le r_0$,
\[
A_\pi(\exp X)\ge \frac{m_-}{2}\|X\|^2,
\qquad
dg = J(X)\,dX,
\qquad
0\le J(X)\le J_{\max},
\]
\item a constant $a_0>0$ such that $A_\pi(g)\ge a_0$ for all $g\notin \mathcal U_0$.
\end{enumerate}
Then for every $\beta>0$ and $\alpha\in(0,1]$,
\begin{equation}
\label{eq:h52_z2a_explicit_bound}
z_{2\alpha}(\beta)
=
\int_G e^{-2\alpha\beta A_\pi(g)}\,dg
\;\le\;
J_{\max}\Big(\frac{\pi}{\alpha\beta\,m_-}\Big)^{d_G/2}
+
e^{-2\alpha\beta a_0}.
\end{equation}
\end{lemma}

\begin{proof}
Split the integral over $\mathcal U_0$ and its complement. On $\mathcal U_0$ write $g=\exp X$ and use
$A_\pi(\exp X)\ge \frac{m_-}{2}\|X\|^2$ and $J(X)\le J_{\max}$:
\[
\begin{aligned}
\int_{\mathcal U_0} e^{-2\alpha\beta A_\pi(g)}\,dg
&\le\int_{\|X\|\le r_0} J_{\max}\,e^{-\alpha\beta m_- \|X\|^2}\,dX\\
&\le J_{\max}\int_{\mathbb R^{d_G}} e^{-\alpha\beta m_- \|X\|^2}\,dX\\
&=J_{\max}\Big(\frac{\pi}{\alpha\beta\,m_-}\Big)^{d_G/2}.
\end{aligned}
\]
On $G\setminus\mathcal U_0$ use $A_\pi(g)\ge a_0$ and $dg$ has total mass $1$ to obtain
$\int_{G\setminus\mathcal U_0} e^{-2\alpha\beta A_\pi(g)}\,dg\le e^{-2\alpha\beta a_0}$.
Summing the two contributions yields~\eqref{eq:h52_z2a_explicit_bound}.
\end{proof}

%%%%%%%%%%%%%%%%%%%%%%%%%%%%%%%%%%%%%%%%%%%%%%%%%%%%%%%%%%%%%%%%%%%%%%%%%%%%%%%
\subsubsection{Closing the finite-depth transfer budget using explicit \texorpdfstring{$\|\Delta_n\|_2$}{||Delta\_n||\_2} bounds}
\label{sec:transfer_budget_closure}

The dyadic estimate~\eqref{eq:h52_dyadic_bound_Delta_n_exact} can be used directly as an admissible
transfer budget $\epsilon_n(\beta)$ in the finite-depth entrance criterion
(Section~\ref{sec:finite_transfer}), once a mixing depth $k(\beta)$ is chosen (e.g.\ via
Proposition~\ref{prop:h51_kbeta_opt}). The key point is that the \emph{sum} of these explicit defects up to the mixing
depth is small for large $\beta$, with fully transparent dependence on $\kappa$ and the
Wilson-action Laplace constants.

\begin{proposition}[An explicit summable transfer budget up to the mixing depth]
\label{prop:h53_transfer_budget_summable}
Fix $\delta\in(0,1)$ and let $k(\beta)$ be any mixing depth satisfying
$4^{k(\beta)}\ge m(\beta)$ with $m(\beta)$ as in Proposition~\ref{prop:h50_kbeta}.
Assume the pivot schedule of Corollary~\ref{cor:h52_Delta_n_bound_dyadic} and set
$\alpha_{\min}:=1-\kappa/48$ (so that $\alpha_n\ge \alpha_{\min}$ for all $n\ge 1$).
Define
\[
Z_\star(\beta)
:=
J_{\max}\Big(\frac{\pi}{\alpha_{\min}\beta\,m_-}\Big)^{d_G/2}
+
e^{-2\alpha_{\min}\beta a_0}.
\]
Then the transfer defects satisfy, for all $\beta>0$,
\begin{equation}
\label{eq:h53_transfer_sum_bound}
\sum_{n<k(\beta)} \|\Delta_n\|_2
\;\le\;
\frac{\kappa}{8c_\Psi}\Bigg(
2\,\Delta_{\mathrm{out}}
+
\frac{\Delta_X}{1-2^{-3/2}}
\Bigg)
+
\frac{2C_{\mathrm{sh}}}{c_\Psi}\,\sqrt{m(\beta)}\,Z_\star(\beta).
\end{equation}
In particular, since $Z_\star(\beta)=O(\beta^{-d_G/2})$ as $\beta\to\infty$, the right-hand side
is $O(\kappa)+O\!\big(\beta^{\frac12-\frac{d_G}{2}}\sqrt{\log(1/\delta)}\big)$ and hence tends to $0$
as $\beta\to\infty$ for every nonabelian $G$ (where $d_G\ge 3$).
\end{proposition}

\begin{proof}
Sum~\eqref{eq:h52_dyadic_bound_Delta_n_exact} over $n<k(\beta)$.
For the first two sums use the bounds
\[
\sum_{n\ge 1}\frac{1}{2^n+1}\le \sum_{n\ge 1}2^{-n}=1,
\qquad
\sum_{n\ge 1}\frac{1}{2^{n/2}(2^n+1)}\le \sum_{n\ge 1}2^{-3n/2}
=\frac{2^{-3/2}}{1-2^{-3/2}}
\le \frac{1}{1-2^{-3/2}}.
\]
For the shared-edge term, note that $\alpha_n\ge\alpha_{\min}$ for all $n\ge 1$, hence by
Lemma~\ref{lem:h52_z2a_explicit_bound} we have $z_{2\alpha_n}(\beta)\le Z_\star(\beta)$ uniformly in $n$.
Therefore,
\[
\sum_{n<k(\beta)} 2^n z_{2\alpha_n}(\beta)
\le
\Big(\sum_{n<k(\beta)}2^n\Big)\,Z_\star(\beta)
\le
2^{k(\beta)}\,Z_\star(\beta).
\]
Finally, $4^{k(\beta)}\ge m(\beta)$ implies $2^{k(\beta)}\le 2\sqrt{m(\beta)}$, giving
the last term in~\eqref{eq:h53_transfer_sum_bound}. The stated asymptotics follow by inspection of
$Z_\star(\beta)$ and the explicit form of $m(\beta)$ from Proposition~\ref{prop:h50_kbeta}.
\end{proof}

\begin{remark}[How to read Proposition~\ref{prop:h53_transfer_budget_summable}]
The first term on the right-hand side of~\eqref{eq:h53_transfer_sum_bound} is purely a
$\kappa$--budget (pivot slack) contribution.
The second term is the explicit weak-coupling decay of the shared-edge correction.
For every compact connected simple $G$ one has $d_G=\dim(G)\ge 3$, so this term decays at least as $O(\beta^{-1}\sqrt{\log(1/\delta)})$; the decay is faster in higher dimension.
Thus for any fixed KP entrance threshold $\delta$ one can choose $\kappa$ and then choose
$\beta$ large enough so that the full transfer-budget sum is $\le \delta$.
\end{remark}


%%%%%%%%%%%%%%%%%%%%%%%%%%%%%%%%%%%%%%%%%%%%%%%%%%%%%%%%%%%%%%%%%%%%%%%%%%%%%%%
\subsubsection{Explicit feasibility of the finite-depth \texorpdfstring{$\chi^2$}{chi2} transfer budget (closed-form \texorpdfstring{$\kappa(\delta)$}{kappa(delta)} and \texorpdfstring{$\beta_0(\delta)$}{beta\_0(delta)})}
\label{sec:transfer_budget_feasibility}

Proposition~\ref{prop:h53_transfer_budget_summable} shows that the cumulative early-step transfer slack
$\sum_{n<k(\beta)}\|\Delta_n\|_2$ is the sum of (i) a \emph{pure pivot-budget} term proportional to $\kappa$ and
(ii) a \emph{shared-edge} term that decays polynomially in $\beta$.
In this subsection we make the parameter selection completely explicit:
for a fixed target $\delta\in(0,1)$ we define deterministic functions $\kappa_{\delta}$ and $\beta_0(\delta)$
so that the transfer sum is $\le \delta$ for all $\beta\ge \beta_0(\delta)$.
No new probabilistic assumptions are introduced; only the constants already present in Proposition~\ref{prop:h53_transfer_budget_summable} are used.

\begin{lemma}[Elementary domination: exponential by a polynomial]
\label{lem:h54_exp_by_poly}
Let $p>0$. Then for all $x>0$,
\begin{equation}
\label{eq:h54_exp_by_poly}
 e^{-x} \le \left(\frac{p}{e}\right)^p x^{-p}.
\end{equation}
\end{lemma}

\begin{proof}
Consider $g(x):=x^p e^{-x}$ on $(0,\infty)$. Then
$g'(x)=x^{p-1}e^{-x}(p-x)$, so the unique maximizer is $x=p$, where $g(p)=p^p e^{-p}$.
Thus $x^p e^{-x}\le p^p e^{-p}$ for all $x>0$, and~\eqref{eq:h54_exp_by_poly} follows after dividing by $x^p$.
\end{proof}

\begin{corollary}[Uniform polynomial envelope for $Z_\star(\beta)$]
\label{cor:h54_Zstar_poly}
With $Z_\star(\beta)$ as in Proposition~\ref{prop:h53_transfer_budget_summable} and $p=d_G/2$, one has for all $\beta>0$,
\begin{equation}
\label{eq:h54_Zstar_poly_bound}
 Z_\star(\beta)\le C_Z(\kappa)\,\beta^{-d_G/2},
\end{equation}
where
\begin{equation}
\label{eq:h54_CZ_def}
C_Z(\kappa)
:=
J_{\max}\Big(\frac{\pi}{\alpha_{\min} m_-}\Big)^{d_G/2}
+
\Big(\frac{d_G}{2e}\Big)^{d_G/2}\,(2\alpha_{\min} a_0)^{-d_G/2},
\qquad \alpha_{\min}=1-\kappa/48.
\end{equation}
\end{corollary}

\begin{proof}
The first term in $Z_\star(\beta)$ already has the form const$\cdot \beta^{-d_G/2}$.
For the second term, apply Lemma~\ref{lem:h54_exp_by_poly} with $p=d_G/2$ and $x=2\alpha_{\min}\beta a_0$ to obtain
$e^{-2\alpha_{\min}\beta a_0}\le (\frac{d_G}{2e})^{d_G/2}(2\alpha_{\min}a_0)^{-d_G/2}\,\beta^{-d_G/2}$.
Summing the two contributions gives~\eqref{eq:h54_Zstar_poly_bound}.
\end{proof}

\begin{proposition}[Explicit $(\kappa,\beta)$ choice forcing $\sum_{n<k_{\mathrm{opt}}(\beta)}\|\Delta_n\|_2\le \delta$]
\label{prop:h54_transfer_feasible}
Fix a target $\delta\in(0,1)$.
Let $k_{\mathrm{opt}}(\beta)$ be the optimal mixing depth from Proposition~\ref{prop:h51_kbeta_opt}
for block factor $L=2$, and write $m_{\mathrm{opt}}(\beta)=\max\{1,\beta s_{\delta}\}$.
Define the deterministic constant
\begin{equation}
\label{eq:h54_Aout_def}
A_{\mathrm{out}}:=2\,\Delta_{\mathrm{out}}+\frac{\Delta_X}{1-2^{-3/2}}.
\end{equation}
Set
\begin{equation}
\label{eq:h54_kappa_delta_def}
\kappa_{\delta}:=\min\Big\{24,\ \frac{4 c_\Psi}{A_{\mathrm{out}}}\,\delta\Big\},
\qquad \alpha_{\min}:=1-\kappa_{\delta}/48\ge \frac12,
\end{equation}
and define $C_Z:=C_Z(\kappa_{\delta})$ as in~\eqref{eq:h54_CZ_def}.
Finally define
\begin{equation}
\label{eq:h54_beta0_def}
\beta_0(\delta)
:=
\max\Bigg\{\frac{1}{s_{\delta}},\ 1,\ \Big(\frac{4 C_{\mathrm{sh}}}{c_\Psi}\,\sqrt{s_{\delta}}\,C_Z\,\delta^{-1}\Big)^{\!\frac{2}{d_G-1}}\Bigg\}.
\end{equation}
Then for all $\beta\ge \beta_0(\delta)$, under the pivot schedule
$\alpha(L)=1-\kappa_{\delta}/|P_\perp(\mathbb B_L)|$, the dyadic defects satisfy
\begin{equation}
\label{eq:h54_transfer_sum_le_delta}
\sum_{n<k_{\mathrm{opt}}(\beta)}\|\Delta_n\|_2 \le \delta.
\end{equation}
\end{proposition}

\begin{proof}
Apply Proposition~\ref{prop:h53_transfer_budget_summable} with $k(\beta)=k_{\mathrm{opt}}(\beta)$ and $\kappa=\kappa_{\delta}$.
The first term on the right-hand side of~\eqref{eq:h53_transfer_sum_bound} equals
$\kappa_{\delta}A_{\mathrm{out}}/(8c_\Psi)$ and is $\le \delta/2$ by~\eqref{eq:h54_kappa_delta_def}.

For the shared-edge term, note that $\beta\ge 1/s_{\delta}$ implies $m_{\mathrm{opt}}(\beta)=\beta s_{\delta}$.
Corollary~\ref{cor:h54_Zstar_poly} gives $Z_\star(\beta)\le C_Z\,\beta^{-d_G/2}$.
Therefore the second term in~\eqref{eq:h53_transfer_sum_bound} is bounded by
\[
\frac{2C_{\mathrm{sh}}}{c_\Psi}\,\sqrt{m_{\mathrm{opt}}(\beta)}\,Z_\star(\beta)
\le
\frac{2C_{\mathrm{sh}}}{c_\Psi}\,\sqrt{\beta s_{\delta}}\,C_Z\,\beta^{-d_G/2}
=
\frac{2C_{\mathrm{sh}}}{c_\Psi}\,\sqrt{s_{\delta}}\,C_Z\,\beta^{-(d_G-1)/2}.
\]
With $\beta\ge \beta_0(\delta)$ from~\eqref{eq:h54_beta0_def}, this bound is $\le \delta/2$.
Summing the two contributions yields~\eqref{eq:h54_transfer_sum_le_delta}.
\end{proof}

\begin{corollary}[One-shot extension of the transfer budget]\label{cor:h54_transfer_sum_one_shot}
Under the hypotheses and notation of Proposition~\ref{prop:h54_transfer_feasible}, the same choice of $(\kappa_\delta,\beta_0(\delta))$ also yields
\[
\sum_{n\le k_{\mathrm{opt}}(\beta)} \|\Delta_n\|_2 \le \delta
\qquad (\beta\ge \beta_0(\delta)).
\]
Equivalently, the dyadic defect control is available at the one-shot scale $L=B(\beta)=2^{k_{\mathrm{opt}}(\beta)}$ itself.
\end{corollary}

\begin{proof}
Apply Proposition~\ref{prop:h53_transfer_budget_summable} with the slightly enlarged depth
\[
 k^{\sharp}(\beta):=k_{\mathrm{opt}}(\beta)+1.
\]
Because $k_{\mathrm{opt}}(\beta)$ is the smallest integer with $4^{k_{\mathrm{opt}}(\beta)}\ge m_{\mathrm{opt}}(\beta)$, one has
\[
2^{k^{\sharp}(\beta)}=2^{k_{\mathrm{opt}}(\beta)+1}\le 4\sqrt{m_{\mathrm{opt}}(\beta)}.
\]
Repeating the proof of Proposition~\ref{prop:h54_transfer_feasible} with $k^{\sharp}$ in place of $k_{\mathrm{opt}}$ therefore changes only the shared-edge contribution, which acquires at most an additional factor $2$.
The definition of $\beta_0(\delta)$ in \eqref{eq:h54_beta0_def} already contains this safety factor, so the same choice of $(\kappa_\delta,\beta_0(\delta))$ forces the enlarged sum to be at most $\delta$.
Since
\[
\sum_{n<k^{\sharp}(\beta)}\|\Delta_n\|_2
=
\sum_{n\le k_{\mathrm{opt}}(\beta)}\|\Delta_n\|_2,
\]
the claim follows.
\end{proof}

\begin{remark}[Scaling with the group dimension]
The exponent in~\eqref{eq:h54_beta0_def} is $2/(d_G-1)$, hence the entrance scale becomes \emph{easier} as $d_G$ grows.
For arbitrary compact connected simple $G$ the exponent in~\eqref{eq:h54_beta0_def} is always strictly less than $1$ because $d_G=\dim(G)\ge 3$. In the classical audited examples the explicit exponents can be read off from the corresponding dimensions recorded in the family appendices.
\end{remark}

%%%%%%%%%%%%%%%%%%%%%%%%%%%%%%%%%%%%%%%%%%%%%%%%%%%%%%%%%%%%%%%%%%%%%%%%%%%%%%%

\subsubsection{Splicing the explicit transfer-feasibility bound into the Plan-A \texorpdfstring{$\chi^2$}{chi2} gate}
\label{sec:planA_gate}

Proposition~\ref{prop:B9_planA_reduction} reduces the Plan-A variance transfer to the conditional $\chi^2$ gate
\eqref{eq:B9_planA_chi2_assump} with a constant $\varepsilon_\ast$ satisfying $\sqrt{\varepsilon_\ast}\le c_1/(2C_A)$.
We now connect this gate directly to the explicit weak-window bounds of Section~\ref{sec:transfer_budget_feasibility}, eliminating the remaining
``choose $\beta$ large enough'' ambiguity and making the $(\kappa,\beta)$ entrance choice audit-ready.

\begin{definition}[Plan-A $\chi^2$ gate parameters]\label{def:h55_planA_gate_params}
Fix the constants $c_1(G)$ from Lemma~\ref{lem:planar_one_plaq_var_lower} and $C_A(G)$ from Proposition~\ref{prop:B9_planA_reduction}.
Define
\begin{equation}\label{eq:h55_epsstar_def}
\varepsilon_\ast:=\Big(\frac{c_1}{2C_A}\Big)^2,
\qquad
\delta_{\mathrm{PA}}:=\frac{c_1}{4\sqrt2\,C_A}.
\end{equation}
\end{definition}

\begin{lemma}[$L^2$-flatness at level $\delta_{\mathrm{PA}}$ implies the Plan-A conditional $\chi^2$ gate]\label{lem:h55_L2_to_planA_chi2_gate}
Fix a planar tile of side length $L$ and write $\nu:=\nu_L^{\parallel}|_{\mathscr F_{\parallel}}$,
$\mu:=\nu_L^{(\alpha)}|_{\mathscr F_{\parallel}}$.
Let $\Omega\in\mathscr F_{\parallel}$ satisfy $\nu(\Omega)\ge \tfrac12$.
Set $\Delta_L:=\widetilde\Psi_L^{(\alpha)}-1$ as in \eqref{eq:h52_DeltaL_def}.
If $\|\Delta_L\|_{L^2(\nu)}\le \delta_{\mathrm{PA}}$, then
\begin{equation}\label{eq:h55_planA_gate_conclusion}
D_{\chi^2}\!\bigl(\mu^{\Omega}\,\|\,\nu^{\Omega}\bigr)
\le \varepsilon_\ast .
\end{equation}
\end{lemma}

\begin{proof}
By Lemma~\ref{lem:chi2_from_L2_on_event}, if $\|\Delta_L\|_{L^2(\nu)}\le \frac12\sqrt{\nu(\Omega)}$ then
\[
D_{\chi^2}\!\bigl(\mu^{\Omega}\,\|\,\nu^{\Omega}\bigr)
\le \frac{4}{\nu(\Omega)}\,\|\Delta_L\|_{L^2(\nu)}^2.
\]
Since $\nu(\Omega)\ge \tfrac12$, the hypothesis $\|\Delta_L\|_{L^2(\nu)}\le \delta_{\mathrm{PA}}$ implies
$\|\Delta_L\|_{L^2(\nu)}\le 1/(2\sqrt2)$, hence the lemma applies and yields
\[
D_{\chi^2}\!\bigl(\mu^{\Omega}\,\|\,\nu^{\Omega}\bigr)
\le 8\,\|\Delta_L\|_{L^2(\nu)}^2
\le 8\,\delta_{\mathrm{PA}}^2
=\varepsilon_\ast,
\]
using \eqref{eq:h55_epsstar_def} in the last step.
\end{proof}

\begin{proposition}[A uniform $(\kappa,\beta)$ choice meeting the Plan-A $\chi^2$ gate in the weak window]\label{prop:h55_planA_gate_closed}
Assume the weak-window setting and pivot schedule of Proposition~\ref{prop:h54_transfer_feasible}.
Fix any $\delta\in(0,\delta_{\mathrm{PA}}]$ and choose $\kappa=\kappa_\delta$ and $\beta_0(\delta)$ as in
\eqref{eq:h54_kappa_delta_def}--\eqref{eq:h54_beta0_def}.
Then for every $\beta\ge \beta_0(\delta)$ and every dyadic scale $L=2^n$ with $n\le k_{\mathrm{opt}}(\beta)$,
\[
\|\Delta_L\|_{L^2(\nu_L^{\parallel}|_{\mathscr F_{\parallel}})}\le \delta\le \delta_{\mathrm{PA}}.
\]
Consequently, for every such $L$ and every event $\Omega\in\mathscr F_{\parallel}$ with $\nu_L^{\parallel}(\Omega)\ge \tfrac12$
(in particular for the standard weak-field event $\Omega_{\varepsilon_0}$),
the Plan-A conditional $\chi^2$ requirement \eqref{eq:B9_planA_chi2_assump} holds with $\varepsilon_\ast$ from \eqref{eq:h55_epsstar_def}.
\end{proposition}

\begin{proof}
By Corollary~\ref{cor:h54_transfer_sum_one_shot},
\[
\sum_{m\le k_{\mathrm{opt}}(\beta)}\|\Delta_m\|_2\le \delta.
\]
Each summand is nonnegative, hence in particular $\|\Delta_n\|_2\le \delta$ for every $n\le k_{\mathrm{opt}}(\beta)$.
The conditional $\chi^2$ conclusion is Lemma~\ref{lem:h55_L2_to_planA_chi2_gate}.
\end{proof}

\begin{corollary}[Plan-A gate at the one-shot block size]\label{cor:h55_planA_gate_one_shot}
Under the hypotheses of Proposition~\ref{prop:h55_planA_gate_closed}, let
\[
B(\beta):=2^{k_{\mathrm{opt}}(\beta)}.
\]
Then the exact-to-reference comparison at the same block size $L=B(\beta)$ satisfies the Plan-A conditional $\chi^2$ gate \eqref{eq:B9_planA_chi2_assump} with defect $\varepsilon_\ast$, uniformly in finite volume, admissible boundary conditions, and all events $\Omega$ of reference probability at least $1/2$.
\end{corollary}

\begin{proof}
Take $n=k_{\mathrm{opt}}(\beta)$ in Proposition~\ref{prop:h55_planA_gate_closed}.
\end{proof}



%%%%%%%%%%%%%%%%%%%%%%%%%%%%%%%%%%%%%%%%%%%%%%%%%%%%%%%%%%%%%%%%%%%%%%%%%%%%%%%


%%%%%%%%%%%%%%%%%%%%%%%%%%%%%%%%%%%%%%%%%%%%%%%%%%%%%%%%%%%%%%%%%%%%%%%%%%%%%%%

\begin{proposition}[Crossover coefficient decay $\Rightarrow$ reference KP-smallness at a computable scale]\label{prop:TBmix_to_f}
Assume the coefficient decay \eqref{eq:TBmix_coeff_decay} for the thick-block reference family.
Define the nonnegative $L^2$ activity
\begin{equation}\label{eq:f_activity_def}
 f_L^{(\alpha)}:=\|w_L^{(\alpha)}-1\|_{L^2(G)}^2.
\end{equation}
Then, using the character orthogonality identity for central functions,
\begin{equation}\label{eq:f_rep_sum}
 f_L^{(\alpha)}
 =\sum_{R\in\mathrm{Irr}(G)\setminus\{\mathbf 1\}} d_R^{\,2}\,\bigl|\lambda_R(w_L^{(\alpha)})\bigr|^2.
\end{equation}
In particular, for $\beta\ge\beta_{\mathrm{mix}}$ and $\alpha(L)\in[1/2,1]$,
\begin{equation}\label{eq:f_ref_bound_crossover_G}
 f_L^{(\alpha(L))}
 \ \le\
 \sum_{R\neq\mathbf 1} d_R^{\,2}
 \exp\!\Big(-2|T_L|\,c_{\mathrm{A1}}(G,\rho)\,\frac{C_2(R)}{\alpha(L)\beta+\sqrt{C_2(R)}}\Big).
\end{equation}
Setting $s:=|T_L|/\beta=L^2/\beta$, the sum admits the genuine low-/high-representation split
\begin{align}\label{eq:f_ref_bound_split_G}
 f_L^{(\alpha(L))}
 &\le
 \sum_{\substack{R\neq\mathbf 1\\ C_2(R)\le \beta^2}} d_R^{\,2}\,\exp\!\bigl(-c_{\mathrm{A1}}(G,\rho)\,s\,C_2(R)\bigr)
 \notag\\
 &\qquad+
 \sum_{\substack{R\neq\mathbf 1\\ C_2(R)> \beta^2}} d_R^{\,2}\,\exp\!\bigl(-c_{\mathrm{A1}}(G,\rho)\,|T_L|\,\sqrt{C_2(R)}\bigr).
\end{align}
Define the two summability constants
\begin{equation}\label{eq:K0lo_def}
 K^{\mathrm{lo}}_{0,G}:=\sum_{R\neq\mathbf 1} d_R^{\,2}\,\exp\!\Big(-\tfrac{c_{\mathrm{A1}}(G,\rho)}{2}\,C_2(R)\Big)
 \ \in\ (0,\infty),
\end{equation}
\begin{equation}\label{eq:H0hi_def}
 H^{\mathrm{hi}}_{0,G}:=\sum_{R\neq\mathbf 1} d_R^{\,2}\,\exp\!\Big(-\tfrac{c_{\mathrm{A1}}(G,\rho)}{2}\,\sqrt{C_2(R)}\Big)
 \ \in\ (0,\infty),
\end{equation}
and let
\[
C_{\min}(G):=\min\{C_2(R):R\neq\mathbf 1\}
\]
be the first positive Laplace--Beltrami eigenvalue on $G$.
For every target $f_{\star}\in(0,1)$ set
\begin{equation}\label{eq:TBmix_Sstar_def_G}
 S_\star\ :=\ \max\!\Bigg\{
 \frac{1}{c_{\mathrm{A1}}(G,\rho)\,C_{\min}(G)}\,\log\!\Big(\frac{2K^{\mathrm{lo}}_{0,G}}{f_\star}\Big)+\frac12,\ 
 \frac{1}{c_{\mathrm{A1}}(G,\rho)\,\sqrt{C_{\min}(G)}}\,\log\!\Big(\frac{2H^{\mathrm{hi}}_{0,G}}{f_\star}\Big)+\frac12
 \Bigg\}.
\end{equation}
Then
\[
\frac{|T_L|}{\beta}=\frac{L^2}{\beta}\ge S_\star
\quad\Longrightarrow\quad
f_L^{(\alpha(L))}\le f_\star.
\]
\end{proposition}

\begin{proof}
Identity \eqref{eq:f_rep_sum} is standard Peter--Weyl orthogonality for central functions:
if $w=\sum_R d_R\,\lambda_R(w)\,\chi_R$ and $\int w\,dg=1$, then
$\|w-1\|_2^2=\sum_{R\neq\mathbf 1} d_R^2|\lambda_R(w)|^2$.

Bound \eqref{eq:f_ref_bound_crossover_G} follows immediately from \eqref{eq:f_rep_sum} and the coefficient decay
\eqref{eq:TBmix_coeff_decay}.
To obtain \eqref{eq:f_ref_bound_split_G}, split the irrep sum into the low-representation regime $C_2(R)\le\beta^2$ and the high-representation regime $C_2(R)>\beta^2$.
If $C_2(R)\le\beta^2$, then $\sqrt{C_2(R)}\le\beta$ and $\alpha(L)\beta\le\beta$, so
\[
\begin{aligned}
\alpha(L)\beta+\sqrt{C_2(R)}&\le 2\beta,\\
2|T_L|\,c_{\mathrm{A1}}(G,\rho)
\frac{C_2(R)}{\alpha(L)\beta+\sqrt{C_2(R)}}
&\ge c_{\mathrm{A1}}(G,\rho)\,\frac{|T_L|}{\beta}\,C_2(R)
=c_{\mathrm{A1}}(G,\rho)\,s\,C_2(R).
\end{aligned}
\]
If $C_2(R)>\beta^2$, then $\beta<\sqrt{C_2(R)}$ and therefore
\[
\alpha(L)\beta+\sqrt{C_2(R)}\le 2\sqrt{C_2(R)},
\qquad
2|T_L|\,c_{\mathrm{A1}}(G,\rho)\,\frac{C_2(R)}{\alpha(L)\beta+\sqrt{C_2(R)}}\ge c_{\mathrm{A1}}(G,\rho)\,|T_L|\,\sqrt{C_2(R)}.
\]
This yields \eqref{eq:f_ref_bound_split_G}.

For the low-representation tail and $s\ge 1/2$,
\[
 e^{-c_{\mathrm{A1}}(G,\rho)sC_2(R)}
 = e^{-\frac{c_{\mathrm{A1}}(G,\rho)}{2}C_2(R)}\,e^{-c_{\mathrm{A1}}(G,\rho)(s-\frac12)C_2(R)}
 \le e^{-\frac{c_{\mathrm{A1}}(G,\rho)}{2}C_2(R)}\,e^{-c_{\mathrm{A1}}(G,\rho)(s-\frac12)C_{\min}(G)}.
\]
Hence
\[
\sum_{\substack{R\neq\mathbf 1\\ C_2(R)\le\beta^2}} d_R^{\,2}\,e^{-c_{\mathrm{A1}}(G,\rho)sC_2(R)}
\le
K^{\mathrm{lo}}_{0,G}\,e^{-c_{\mathrm{A1}}(G,\rho)(s-\frac12)C_{\min}(G)}.
\]

For the high-representation tail, use $|T_L|=\beta s$, $\beta\ge1$, $s\ge1/2$, and $C_2(R)\ge C_{\min}(G)$ to write
\[
 e^{-c_{\mathrm{A1}}(G,\rho)|T_L|\sqrt{C_2(R)}}
 =
 e^{-c_{\mathrm{A1}}(G,\rho)\beta s\sqrt{C_2(R)}}
 \le
 e^{-\frac{c_{\mathrm{A1}}(G,\rho)}{2}\sqrt{C_2(R)}}\,
 e^{-c_{\mathrm{A1}}(G,\rho)(s-\frac12)\sqrt{C_{\min}(G)}}.
\]
Therefore
\[
\sum_{\substack{R\neq\mathbf 1\\ C_2(R)>\beta^2}} d_R^{\,2}\,e^{-c_{\mathrm{A1}}(G,\rho)|T_L|\sqrt{C_2(R)}}
\le
H^{\mathrm{hi}}_{0,G}\,e^{-c_{\mathrm{A1}}(G,\rho)(s-\frac12)\sqrt{C_{\min}(G)}}.
\]
The finiteness of $H^{\mathrm{hi}}_{0,G}$ follows from Weyl polynomial growth of $d_R$ together with the Casimir asymptotic $C_2(R)\asymp |\lambda_R|^2$, which makes the tail comparable to a polynomial times $e^{-c|\lambda_R|}$.

If $s\ge S_\star$, each displayed tail is at most $f_\star/2$ by \eqref{eq:TBmix_Sstar_def_G}, hence $f_L^{(\alpha(L))}\le f_\star$.
\end{proof}

\begin{corollary}[Reference KP-smallness is closed on the load-bearing Route~1 path]\label{cor:TBmix_reference_budget_closed}
For every compact connected simple $G$ and every target $f_\star\in(0,1)$, the computable two-tail threshold $S_\star$ from \eqref{eq:TBmix_Sstar_def_G} has the following property:
whenever $L$ is dyadic and $L^2/\beta\ge S_\star$, the exact thick-block reference activity satisfies
\[
 f_L^{(\alpha(L))}\le f_\star.
\]
This is the theorem-level QE1 reference-budget input used in the one-shot entrance proof.
\end{corollary}

\begin{proof}
Corollary~\ref{cor:TBmix_heatkernel_closed} supplies the coefficient decay \eqref{eq:TBmix_coeff_decay}; Proposition~\ref{prop:TBmix_to_f} then yields the stated bound with the explicit threshold \eqref{eq:TBmix_Sstar_def_G}.
\end{proof}

\begin{remark}[Why a thick reference is still needed]
Lemma~\ref{lem:leaf_factorization_thin_block} is exact because $\mathcal B_L$ is a \emph{thin} neighborhood:
the translated opposite edges of $P_\perp$ plaquettes are leaves, so the out-of-plane sector carries no
structural feedback into the in-plane holonomy marginal on this block.
In the \emph{true} $4$-dimensional lattice, those translated edges are not leaves: they participate in additional plaquettes
in shifted slices and in perpendicular cubes.
Therefore the thin-block computation is used only as a diagnostic for the slack bookkeeping in \eqref{eq:alpha_schedule_L}; it is not the reference used in the entrance proof.
The entrance construction in this appendix is implemented with a positivity-preserving coarse reference in which the translated edges are no longer leaves
(see Sections~\ref{sec:global} and~\ref{sec:majorant}), and is packaged in the one-shot KP entrance theorem (Theorem~\ref{thm:h60_one_shot_entrance}).
\end{remark}

\subsection{Infinite-volume exponential clustering at fixed large \texorpdfstring{$\beta$}{beta} via finite-depth KP entrance}\label{sec:iv_clustering}

This section records the next mathematically forced implication of the already completed modules:
\begin{enumerate}[leftmargin=*]
\item the downstream KP module (Dobrushin/KP contraction inside the admissible window);
\item the exact-step contraction and mixing bound (Proposition~12.1 and its corollaries in the current numbering);
\item the $\alpha$-pivot slack analysis and $\chi^2$-transfer module, including the weak-coupling $\chi^2$ decay
of Section~\ref{sec:chi2_decay_alpha}.
\end{enumerate}
We do \emph{not} re-derive those results; we only assemble them into an infinite-volume statement.

\subsubsection{Blocked specifications and influence coefficients}

Let $\mathbb Z_L^4:=L\mathbb Z^4$ be the block lattice. For a finite block region $\Lambda\subset \mathbb Z_L^4$,
write $\mu_{\Lambda}^{\mathrm{ex}}$ for the exact blocked coarse measure produced by the RG map at block size $L$
(with any fixed admissible boundary condition on $\partial\Lambda$), and $\mu_{\Lambda}^{(\alpha)}$ for the corresponding
$\alpha$-reference blocked measure. Denote by $\gamma^{\mathrm{ex}}=\{\gamma_i^{\mathrm{ex}}\}_{i\in\mathbb Z_L^4}$
(resp.\ $\gamma^{(\alpha)}$) the associated single-site conditional kernels (specifications).

For sites $i\neq j$, define the total-variation influence coefficient
\[
c_{ij}(\gamma):=\sup_{\eta,\eta':\ \eta_{\mathbb Z_L^4\setminus\{j\}}=\eta'_{\mathbb Z_L^4\setminus\{j\}}}
\bigl\|\gamma_i(\cdot\mid \eta)-\gamma_i(\cdot\mid \eta')\bigr\|_{\mathrm{TV}}.
\]
Write $C(\gamma)$ for the influence matrix $(c_{ij}(\gamma))$ and the Dobrushin coefficient
\[
\delta(\gamma):=\sup_{i}\sum_{j\neq i} c_{ij}(\gamma).
\]

\begin{proposition}[Finite-range KP/Dobrushin criterion in oscillation-matrix form]\label{prop:dobrushin_standard}
Let $\gamma$ be a single-site specification on a graph $(V,E)$ with graph distance $\dist_{\mathcal G}$ and influence matrix $C(\gamma)=(c_{ij})$.
For a bounded local observable $F$, define the single-site oscillation seminorm
\[
\operatorname{osc}_i(F):=
\sup_{\eta,\eta':\ \eta_{V\setminus\{i\}}=\eta'_{V\setminus\{i\}}}
|F(\eta)-F(\eta')|,
\qquad
\supp_{\mathrm{osc}}(F):=\{i\in V:\operatorname{osc}_i(F)\neq 0\}.
\]
Assume that
\[
c_{ij}(\gamma)=0\qquad\text{whenever }\dist_{\mathcal G}(i,j)>1,
\qquad
\delta(\gamma):=\sup_i\sum_{j\ne i} c_{ij}(\gamma)\le q<1.
\]
Then the Gibbs measure specified by $\gamma$ exists and is unique in infinite volume.
Let
\[
D:=\sum_{m\ge 0} C(\gamma)^m=(I-C(\gamma))^{-1}.
\]
For bounded local observables $F,G$ and the unique infinite-volume Gibbs state $\langle\cdot\rangle_\gamma$ one has the standard oscillation-matrix covariance bound
\begin{equation}\label{eq:dobrushin_oscillation_bound}
|\mathrm{Cov}_{\langle\cdot\rangle_\gamma}(F,G)|
\le
\sum_{i\in V}\sum_{j\in V}
\operatorname{osc}_i(F)\,D_{ij}\,\operatorname{osc}_j(G).
\end{equation}
Moreover, the finite-range hypothesis implies the explicit path-sum estimate
\begin{equation}\label{eq:dobrushin_D_decay}
D_{ij}
\le
\sum_{m\ge \dist_{\mathcal G}(i,j)} q^m
=
\frac{e^{-(-\log q)\,\dist_{\mathcal G}(i,j)}}{1-q}.
\end{equation}
Consequently, if $S:=\supp_{\mathrm{osc}}(F)$ and $T:=\supp_{\mathrm{osc}}(G)$, then
\begin{equation}\label{eq:dobrushin_support_weighted_bound}
|\mathrm{Cov}_{\langle\cdot\rangle_\gamma}(F,G)|
\le
\frac{e^{-(-\log q)\,\dist_{\mathcal G}(S,T)}}{1-q}
\Big(\sum_{i\in S}\operatorname{osc}_i(F)\Big)
\Big(\sum_{j\in T}\operatorname{osc}_j(G)\Big).
\end{equation}
In particular, since $\operatorname{osc}_i(F)\le 2\|F\|_\infty$ and $\operatorname{osc}_j(G)\le 2\|G\|_\infty$,
\begin{equation}\label{eq:dobrushin_support_cardinality_bound}
|\mathrm{Cov}_{\langle\cdot\rangle_\gamma}(F,G)|
\le
\frac{4|S||T|}{1-q}\,\|F\|_\infty\|G\|_\infty
e^{-(-\log q)\,\dist_{\mathcal G}(S,T)}.
\end{equation}
\end{proposition}

\begin{proof}
The uniqueness statement is the standard Dobrushin criterion.
The covariance estimate \eqref{eq:dobrushin_oscillation_bound} is the standard Dobrushin--Shlosman oscillation-matrix bound for the unique Gibbs state.
Because $c_{ij}(\gamma)=0$ unless $\dist_{\mathcal G}(i,j)\le 1$, every nonzero contribution to $(C(\gamma)^m)_{ij}$ comes from a path of length $m$ joining $i$ to $j$, so $m\ge \dist_{\mathcal G}(i,j)$ whenever $(C(\gamma)^m)_{ij}\neq 0$.
Using the row-sum bound $\sup_i\sum_j c_{ij}(\gamma)\le q$ therefore gives
\[
D_{ij}
=\sum_{m\ge0}(C(\gamma)^m)_{ij}
\le\sum_{m\ge \dist_{\mathcal G}(i,j)} q^m,
\]
which is exactly \eqref{eq:dobrushin_D_decay}.
Substituting \eqref{eq:dobrushin_D_decay} into \eqref{eq:dobrushin_oscillation_bound} yields \eqref{eq:dobrushin_support_weighted_bound}, and the cardinality estimate \eqref{eq:dobrushin_support_cardinality_bound} follows from the trivial bounds on single-site oscillations.
\end{proof}

\subsubsection{Stability of the influence matrix under \texorpdfstring{$\chi^2$}{chi2} transfer}

The completed $\chi^2$-transfer module gives, for each finite $\Lambda$ and each block site $i$,
a uniform bound on the divergence between the exact and reference \emph{single-site} kernels, after the finite
entrance depth has been reached. We record the elementary consequence needed here.

\begin{lemma}[$\chi^2$ control implies TV kernel stability]\label{lem:chi2_to_influence}
Let $\mu\ll\nu$ be probability measures on a common measurable space with
$D_{\chi^2}(\mu\|\nu)\le \varepsilon$. Then $\|\mu-\nu\|_{\mathrm{TV}}\le \tfrac12\sqrt{\varepsilon}$.
In particular, for two specifications $\gamma,\tilde\gamma$ with
\[
\sup_{i}\sup_{\eta}\ D_{\chi^2}\bigl(\gamma_i(\cdot\mid\eta)\,\|\,\tilde\gamma_i(\cdot\mid\eta)\bigr)\le \varepsilon,
\]
one has the influence stability bound
\[
\delta(\gamma)\le \delta(\tilde\gamma)+\sqrt{\varepsilon}.
\]
\end{lemma}

\begin{proof}
The TV bound is \eqref{eq:tv_vs_chi2_group} in the present notation.
For influence stability, fix $i\neq j$ and boundary conditions $\eta,\eta'$ differing only at $j$.
Insert and subtract $\tilde\gamma_i(\cdot\mid\eta)$ and $\tilde\gamma_i(\cdot\mid\eta')$ and use the triangle inequality:
\[
\|\gamma_i(\cdot\mid\eta)-\gamma_i(\cdot\mid\eta')\|_{\mathrm{TV}}
\le
\|\tilde\gamma_i(\cdot\mid\eta)-\tilde\gamma_i(\cdot\mid\eta')\|_{\mathrm{TV}}
+2\sup_{\xi}\|\gamma_i(\cdot\mid\xi)-\tilde\gamma_i(\cdot\mid\xi)\|_{\mathrm{TV}}.
\]
Take suprema and apply the TV--$\chi^2$ inequality.
\end{proof}

\subsubsection{Infinite-volume blocked exponential clustering at fixed \texorpdfstring{$\beta$}{beta}}

\begin{theorem}[Blocked infinite-volume uniqueness and clustering after finite entrance depth]\label{thm:blocked_iv_clustering}
Fix $d=4$ and compact $G$. Assume $\beta$ is in the weak-coupling regime of the entrance construction and
that the schedule $\alpha(L)=1-\kappa/|P_\perp(\mathcal B_L)|$ is used.
Suppose that at some finite entrance depth (already established in the completed KP module) the $\alpha$-reference
blocked specification satisfies the uniform contraction bound
\[
\delta(\gamma^{(\alpha)})\le q <1,
\]
with $q$ independent of $L$, the finite block volume $\Lambda$, and the boundary condition.
Then there exists $L_0(\beta,G,\kappa)$ such that for all $L\ge L_0$:
\begin{enumerate}[leftmargin=*]
\item the exact blocked specification satisfies $\delta(\gamma^{\mathrm{ex}})<1$ uniformly in $\Lambda$ and boundary conditions;
\item the infinite-volume exact blocked Gibbs measure exists and is unique;
\item block-local truncated correlations decay exponentially in block distance.
\end{enumerate}
\end{theorem}

\begin{proof}
By the weak-coupling $\chi^2$-decay lemma of Section~\ref{sec:chi2_decay_alpha} and the $\chi^2$-transfer module,
for $L$ large one has a uniform conditional divergence bound
\[
\sup_i\sup_\eta\ D_{\chi^2}\bigl(\gamma_i^{\mathrm{ex}}(\cdot\mid\eta)\,\|\,\gamma_i^{(\alpha)}(\cdot\mid\eta)\bigr)
\le \varepsilon_L,\qquad \varepsilon_L\to 0.
\]
Lemma~\ref{lem:chi2_to_influence} then gives $\delta(\gamma^{\mathrm{ex}})\le q+\sqrt{\varepsilon_L}$.
Choose $L_0$ so that $q+\sqrt{\varepsilon_L}<1$ for all $L\ge L_0$.
Apply Proposition~\ref{prop:dobrushin_standard} to $\gamma^{\mathrm{ex}}$.
\end{proof}

\subsubsection{Descent to microscopic observables}

To convert blocked clustering to microscopic clustering, it is enough to control the block-conditional
fluctuation $O-\mathbb E[O\mid\mathcal F_{\mathrm{blocked}}]$ for local gauge-invariant observables $O$.

\begin{lemma}[Uniform block-conditional Poincar\'e inequality at weak coupling]\label{lem:block_cond_poincare_proved}
Fix compact $G$ and $d=4$. There exist $\beta_1(G)<\infty$ and constants $c_{\mathrm{PI}}(G),C_{\mathrm{PI}}(G)<\infty$
such that for all $\beta\ge\beta_1(G)$, for every block $B_L$ and every admissible boundary condition on $\partial B_L$,
the conditional law of the microscopic edge variables in $B_L$ (in spanning-tree gauge) satisfies the Poincar\'e inequality
\[
\mathrm{Var}_{\mu_{B_L}^{\mathrm{cond}}}(f)\le \frac{C_{\mathrm{PI}}(G)}{\beta}
\sum_{e\in E_f(B_L)}\int \|\nabla_e f\|^2\,d\mu_{B_L}^{\mathrm{cond}},
\]
for all smooth $f$, where $E_f(B_L)$ are the free edges after tree gauge. In particular, the Poincar\'e constant is uniform
in $L$ and the boundary condition.
\end{lemma}

\begin{proof}
Work in spanning-tree gauge and exponential coordinates $U_e=\exp(iA_e)$ on a fixed normal neighborhood of the identity in $G$,
with $A_e\in\mathfrak g$ for $e\in E_f(B_L)$. For $\beta$ large, the conditional block density takes the standard weak-coupling form
\[
\frac{1}{Z}\exp\!\big(-\beta\,\Phi(A)\big)\,dA,\qquad
\Phi(A)=\tfrac12\langle A,\mathcal M A\rangle + R(A),
\]
where $\mathcal M$ is the gauge-fixed magnetic quadratic form and $R(A)=O(\|A\|^3)$ is smooth (uniformly in $L$ and boundary condition
on the chart). Tree gauge removes the gauge null space; hence there exists $c_0(G)>0$ such that $\mathcal M\succeq c_0(G) I$
uniformly in $L$ and the boundary condition (discrete Hodge theory with fixed gauge).

Choose $\beta_1$ so large that on the chart region $\{\|A\|\le r\}$ (with fixed small $r$) the Hessian obeys the uniform strong convexity bound
\[
\nabla^2\Phi(A)\succeq \tfrac{c_0(G)}{2}\,I.
\]
On this region, the Brascamp--Lieb inequality yields
\[
\mathrm{Var}(f)\le \frac{2}{\beta c_0(G)}\int \|\nabla f\|^2\,d\mu_{B_L}^{\mathrm{cond}}.
\]
Finally, the Wilson density supplies Gaussian tails at scale $\beta$ outside the chart; a standard localization/cutoff argument
extends the bound globally at the price of a multiplicative constant depending only on $G$ (absorbed into $C_{\mathrm{PI}}(G)$).
\end{proof}


\begin{lemma}[Block-conditional Poincar\'e at weak coupling]\label{lem:block_cond_poincare}
Fix $d=4$ and compact $G$. There exist $\beta_{\mathrm{PI}}(G)<\infty$ and $C_{\mathrm{PI}}(G)<\infty$ such that for all
$\beta\ge\beta_{\mathrm{PI}}$, all block sizes $L$, all $\alpha\in[\alpha_*,1]$, and all gauge-invariant functions $F$ supported
on a single block $B_L$, the conditional law of the edge variables in $B_L$ given the exterior configuration (in spanning-tree gauge)
obeys the Poincar\'e inequality
\[
\mathrm{Var}(F\mid \mathcal F_{\mathrm{ext}})
\;\le\; \frac{C_{\mathrm{PI}}(G)}{\beta}\,\sum_{e\in E_f(B_L)} \mathbb E\big[\|\nabla_e F\|^2\mid \mathcal F_{\mathrm{ext}}\big],
\]
with constants independent of $L$ and of the exterior configuration.
\end{lemma}

\begin{proof}
Work in spanning-tree gauge on $B_L$ so the conditional density is a smooth positive function on $G^{E_f(B_L)}$.
For $\beta$ large, restrict to the exponential chart $U_e=\exp(iA_e)$ on each free edge. The (conditional) negative log-density is
\[
H(A)=\frac{\beta}{2}\langle A,\mathcal M_{\alpha,\mathrm{cond}} A\rangle + R(A),
\]
where $\mathcal M_{\alpha,\mathrm{cond}}$ is the gauge-fixed quadratic form induced by the plaquette terms inside $B_L$ (with transverse weights $\alpha$),
and the remainder satisfies $\|\nabla^2 R(A)\|\le \tfrac{1}{4}\beta\,\lambda_{\min}(\mathcal M_{\alpha,\mathrm{cond}})$ on a high-probability region
$\|A\|\le c(\beta)$, while the complement has exponentially small conditional mass (localization via the plaquette tail bound).
Discrete Hodge theory in tree gauge gives a uniform coercivity bound
$\mathcal M_{\alpha,\mathrm{cond}}\succeq c_0(G) I$ independent of $L$ and exterior data.
Therefore on the localized region,
$\nabla^2 H\succeq \tfrac{\beta c_0(G)}{4} I$.
Brascamp--Lieb then yields a Poincar\'e inequality with constant $\le 4/(\beta c_0(G))$ for the localized measure.
The exterior region is handled by a standard two-region decomposition (Holley--Stroock gluing): the tail probability is
$\le e^{-c\beta}$ and only inflates the constant by a bounded factor.
Collecting constants gives the stated bound with $C_{\mathrm{PI}}(G)$ and $\beta_{\mathrm{PI}}(G)$.
\end{proof}

\begin{lemma}[Descent inequality]\label{lem:descent_inequality}
Assume Theorem~\ref{thm:blocked_iv_clustering} and Lemma~\ref{lem:block_cond_poincare_proved}.
Let $O$ be a bounded local gauge-invariant observable supported in a microscopic set of diameter $O(1)$.
Let $O^{(L)}:=\mathbb E[O\mid \mathcal F_{\mathrm{blocked}}]$ and $\delta O:=O-O^{(L)}$.
Then for all $x\in\mathbb Z^4$,
\[
|\langle O(x)O(0)\rangle_c|
\le
C_1\,e^{-m_1 |x|/L}
+
C_2\,\|\delta O\|_{L^2}^2,
\]
where $m_1>0$ depends only on the blocked Dobrushin constant and $C_1,C_2$ depend on $O,\beta,G$.
Moreover, Lemma~\ref{lem:block_cond_poincare_proved} yields $\|\delta O\|_{L^2}\le C(O,\beta)\,e^{-c L}$ (or at least
a polynomial decay in $L$), so optimizing in $L$ gives a microscopic exponential bound
$|\langle O(x)O(0)\rangle_c|\le C\,e^{-m |x|}$ at fixed $\beta$.
\end{lemma}

\begin{proof}
Write $O=O^{(L)}+\delta O$ and expand the truncated correlation.
The term $\langle O^{(L)}(x)O^{(L)}(0)\rangle_c$ is a truncated correlation of block-measurable observables and
decays exponentially in block distance by Theorem~\ref{thm:blocked_iv_clustering}.
The remaining terms are controlled by Cauchy--Schwarz and the $L^2$ bound on $\delta O$.
The final claim follows by applying Lemma~\ref{lem:block_cond_poincare} to the block-conditional variance term in the martingale decomposition.
\end{proof}

\begin{remark}[Status of the fixed-$\beta$ descent lemma]
Lemma~\ref{lem:descent_inequality} remains a qualitative fixed-$\beta$ diagnostic and is not part of the front-facing QE2$_{\mathrm{lat}}$ proof spine.
The actual QE2 closure below no longer seeks a separate fluctuation-remainder estimate at the one-shot scale.
Instead it uses exact terminal-block conditional independence at that scale, supplied by Corollary~\ref{cor:wilson_terminal_block_product} together with Proposition~\ref{prop:route1_blocked_sigma_identification}, to eliminate the fluctuation term entirely for separated local probes.
\end{remark}

% [v123 edit: removed nontechnical status remark]

\subsection{Continuum limit and OS reconstruction: interface statements and closure}

\subsubsection{Continuum scaling setup}
Fix a lattice spacing parameter $a>0$ and identify the rescaled lattice $a\,\mathbb Z^4$ with $\mathbb Z^4$ by the map
$x\mapsto x/a$. For each $\beta$ write $\mu_{\beta,\Lambda}$ for the finite-volume Wilson measure in a box
$\Lambda\subset\mathbb Z^4$ with a standard reflection-positive boundary condition (e.g.\ free; or periodic handled
via a standard doubling trick).
When the thermodynamic limit exists, write $\mu_\beta$ for the infinite-volume Gibbs state.

We distinguish three notions of gap:
\begin{itemize}[leftmargin=*]
\item $\Delta_{\mathrm{lat}}(\beta)$: the \emph{full fixed-lattice OS/transfer-matrix gap} in lattice units, i.e.
\[
\Delta_{\mathrm{lat}}(\beta):=\inf\bigl(\mathrm{Spec}(H_\beta)\setminus\{0\}\bigr),
\]
where $(\mathcal H_\beta,\Omega_\beta,H_\beta)$ is the OS reconstruction of the fixed-lattice Wilson Gibbs state on the bounded gauge-invariant positive-time cylinder algebra.
This is the quantity used on the primary QE2/QE3 spine.
\item $m_{\mathrm{lat}}(\beta)$: an \emph{auxiliary probe-family/channel decay rate} in lattice units.
It is retained only as a diagnostic shorthand for particular local probe families; the primary spine no longer identifies it with the full lattice spectral gap unless a separate density/cyclicity theorem is stated.
\item $m_{\mathrm{phys}}$: the physical mass in continuum units, obtained after choosing a scaling trajectory $a=a(\beta)\downarrow 0$
and a renormalization prescription for observables.
\end{itemize}

\subsubsection{OS axioms and reconstruction target}
Let $\Theta$ denote reflection in a time hyperplane and let $\mathcal A_+$ be the algebra generated by observables supported
in the positive-time half-space. Reflection positivity is the requirement
\[
\langle F\,\Theta F\rangle_{\mu_{\beta,\Lambda}}\ge 0\qquad\forall F\in\mathcal A_+.
\]
For the Wilson action this holds at finite volume. The continuum construction must preserve it in the infinite-volume and scaling limits so that
Osterwalder--Schrader reconstruction yields a Hilbert space, a Hamiltonian with spectrum bounded below, and a vacuum.
A \emph{mass gap} in the mass gap sense is a strictly positive spectral gap above the vacuum in this reconstructed Hamiltonian.

\subsubsection{Interface theorem chain to the continuum mass gap}
The mass-gap interfaces are separated as follows:
(i) one-shot KP entrance together with a tuned-sequence physical entrance bound is proved unconditionally earlier in Appendix~F;
(ii) the tuned-sequence full fixed-lattice OS gap is proved below by exact terminal-block covariance transfer at the one-shot scale together with the fixed-lattice spectral-gap theorem; and
(iii) QE3 now requires only a continuum-time OS reconstruction theorem stronger than the dyadic statement of Theorem~\ref{thm:continuum_tightness}.
The dense-core Route~1/Lane~B package is recorded later only as a secondary reformulation/consistency check and is not part of the primary front-facing spine.

\begin{theorem}[Provisional fixed-$\beta$ full-lattice-gap diagnostic (not used on the primary spine)]
Fix $\beta\ge\beta_\ast$ in the range of the one-shot entrance theorem (Theorem~\ref{thm:h60_one_shot_entrance}).
If the fixed-$\beta$ fluctuation-control step in Lemma~\ref{lem:descent_inequality} is upgraded to a genuine rate-preserving microscopic descent theorem for \emph{all} bounded gauge-invariant local observables, then the lattice OS Hamiltonian of the infinite-volume Gibbs state $\mu_\beta$ has a full vacuum gap $\Delta_{\mathrm{lat}}(\beta)>0$.
\end{theorem}

\begin{proof}
This theorem is retained only as a diagnostic side-branch.
Concretely, Theorem~\ref{thm:h60_one_shot_entrance} yields KP entrance at the one-shot block scale $B(\beta)$ and hence blocked exponential clustering.
A quantitative fixed-$\beta$ descent theorem for all bounded local observables would then feed Theorem~\ref{thm:F_decay_implies_gap} together with Proposition~\ref{prop:lattice_OS_cyclicity_local_algebra}, yielding a genuine full fixed-lattice OS gap.
Because the primary QE2$_{\mathrm{lat}}$ route below proves that gap directly along the tuned dyadic path, this diagnostic statement is not used on the paper-level spine.
\end{proof}

\begin{theorem}[Bounded positive-time continuum Wilson state and dyadic OS matrix elements along the tuned trajectory]
\label{thm:continuum_tightness}
Fix a weak-window target value $u_{\mathrm{ren}}\in(0,u_\ast]$ and let
\[
 a_n:=\ell_0 2^{-n},
 \qquad
 \beta_n:=\beta(a_n;u_{\mathrm{ren}}).
\]
Define the dense dyadic translation set
\[
\mathbb D_{\mathrm{dyad}}:=\{m\ell_0 2^{-k}:m,k\in\mathbb N_0\}.
\]
Then the following hold along this fixed tuned dyadic sequence:
\begin{enumerate}[label=(\roman*),leftmargin=*]
\item For every bounded cylinder observable $F\in\mathcal A_+$ the limit
\[
\omega_{\mathrm{bd}}^{(u_{\mathrm{ren}})}(F):=\lim_{n\to\infty}\omega_{a_n}(F)
\]
exists. The resulting functional $\omega_{\mathrm{bd}}^{(u_{\mathrm{ren}})}$ is a state on the bounded positive-time algebra $\mathcal A_+$ and is uniquely determined by the chosen target value $u_{\mathrm{ren}}$.
\item For every bounded $F,G\in\mathcal A_+$ and every dyadic shift $s\in\mathbb D_{\mathrm{dyad}}$, the OS matrix-element limit
\begin{equation}\label{eq:continuum_tightness_dyadic_OS_limit}
\mathcal K_{u_{\mathrm{ren}}}(F,G;s)
:=
\lim_{n\to\infty}\omega_{a_n}(\Theta F\,\tau_s G)
\end{equation}
exists (for all sufficiently large $n$ such that $s/a_n\in\mathbb N$).
In particular,
\begin{equation}\label{eq:continuum_tightness_RP}
\mathcal K_{u_{\mathrm{ren}}}(F,F;0)\ge 0,
\end{equation}
so $\omega_{\mathrm{bd}}^{(u_{\mathrm{ren}})}$ is reflection positive on $\mathcal A_+$.
\end{enumerate}
\end{theorem}

\begin{proof}
For part~(i), Corollary~\ref{cor:C3U_Aplus_tuned} gives joint convergence for every finite family of bounded cylinder observables from $\mathcal A_+$ along the tuned dyadic sequence. Applying that corollary to singleton families defines
\[
\omega_{\mathrm{bd}}^{(u_{\mathrm{ren}})}(F):=\lim_{n\to\infty}\omega_{a_n}(F)
\qquad (F\in\mathcal A_+).
\]
Linearity and unitality are inherited from the finite-volume states $\omega_{a_n}$ term by term, and positivity follows by applying the same convergence statement to $F^*F$. Hence the limit is a state on $\mathcal A_+$, uniquely determined by the chosen target value $u_{\mathrm{ren}}$ because the tuned sequence itself is fixed.

For part~(ii), Corollary~\ref{cor:C3U_reflected_tuned} gives joint convergence of every finite family of reflected/transformed bounded cylinder observables built from $\mathcal A_+$ using dyadic time translations. Applying that corollary to the single observable $\Theta F\,\tau_s G$ yields the limit \eqref{eq:continuum_tightness_dyadic_OS_limit}. When $s=0$ and $G=F$, Theorem~\ref{thm:RP_Wilson} gives
\[
\omega_{a_n}(\Theta F\,F)\ge 0
\qquad\text{for every }n,
\]
and taking the limit gives \eqref{eq:continuum_tightness_RP}. Thus the limiting bounded positive-time state is reflection positive on $\mathcal A_+$, with the actual OS matrix elements on the dense dyadic translation set recorded explicitly by \eqref{eq:continuum_tightness_dyadic_OS_limit}.
\end{proof}

\dependencyledger
  {Corollary~\ref{cor:C3U_Aplus_tuned}, Corollary~\ref{cor:C3U_reflected_tuned}, and Theorem~\ref{thm:RP_Wilson}}
  {none}
  {along the tuned dyadic trajectory the bounded Wilson state on $\mathcal A_+$ and its dyadic reflected OS matrix elements converge; the theorem records reflection positivity on $\mathcal A_+$ only and does not yet construct the all-time sharp-local semigroup}
  {Theorem~\ref{thm:continuum_time_OS_upgrade} and Theorem~\ref{thm:phys_gap_noncollapse}}

\begin{theorem}[Continuum-time OS upgrade of the bounded base state via Lane~A]\label{thm:continuum_time_OS_upgrade}
Fix a target value $u_{\mathrm{ren}}\in(0,u_\ast]$ and let $\omega^{(u_{\mathrm{ren}})}$ be the full sharp-local state furnished by Corollary~\ref{cor:LaneA_full_scope} over the same tuned dyadic bounded base state.
Let $(\mathcal H,\Omega,H)$ be its OS reconstruction.
Then:
\begin{enumerate}[label=(\roman*),leftmargin=*]
\item the restriction of $\omega^{(u_{\mathrm{ren}})}$ to $\mathcal A_+$ is exactly the bounded state of Theorem~\ref{thm:continuum_tightness}:
\[
\omega^{(u_{\mathrm{ren}})}\!\restriction_{\mathcal A_+}=\omega_{\mathrm{bd}}^{(u_{\mathrm{ren}})};
\]
\item for every bounded $F,G\in\mathcal A_+$ and every dyadic shift $s\in\mathbb D_{\mathrm{dyad}}$,
\[
\omega^{(u_{\mathrm{ren}})}(\Theta F\,\tau_s G)=\mathcal K_{u_{\mathrm{ren}}}(F,G;s);
\]
\item for every bounded $F,G\in\mathcal A_+$ and every $t\ge 0$,
\[
\omega^{(u_{\mathrm{ren}})}(\Theta F\,\tau_t G)=\langle [F],e^{-tH}[G]\rangle_{\mathcal H},
\]
and the map $t\mapsto \omega^{(u_{\mathrm{ren}})}(\Theta F\,\tau_t G)$ is continuous on $[0,\infty)$.
\end{enumerate}
Consequently the bounded positive-time state of Theorem~\ref{thm:continuum_tightness} has a canonical continuum-time OS semigroup/reconstruction obtained by transport through the Lane~A sharp-local state; no separate non-dyadic lattice-shift convergence theorem is required on the primary Route~1 spine.
\end{theorem}

\begin{proof}
The role split from Lane~A must be read literally.
First, part~(i) is the bounded-state restriction statement isolated in Corollary~\ref{cor:LaneA_bounded_state_restriction}: after the full sharp-local state has been assembled, that corollary identifies its restriction to $\mathcal A_+$ with the bounded state of Theorem~\ref{thm:continuum_tightness}.
Second, the dyadic identity in part~(ii) is the second clause of the same bounded-state compatibility theorem.
In particular, Corollary~\ref{cor:LaneA_bounded_state_restriction} is not the inductive-union theorem.

The finite-cap sharp-local OS structure used for part~(iii) comes instead from Theorem~\ref{thm:LaneA_closure}(b) on each finite engineering cap.
Corollary~\ref{cor:LaneA_full_scope} is the separate passage that transports those compatible capwise OS structures to the full sharp-local inductive union.
Hence the sharp-local state $\omega^{(u_{\mathrm{ren}})}$ satisfies the OS axioms on the relevant local positive-time algebra containing $\mathcal A_+$, so OS reconstruction yields the Hilbert space $(\mathcal H,\Omega,H)$ and the semigroup identity
\[
\omega^{(u_{\mathrm{ren}})}(\Theta F\,\tau_t G)=\langle [F],e^{-tH}[G]\rangle_{\mathcal H}
\qquad (F,G\in\mathcal A_+,\ t\ge 0).
\]
Strong continuity of $e^{-tH}$ gives the continuity statement in part~(iii).
The later passage from bounded classes to the full sharp-local Hilbert space is not used here; it is isolated separately in Corollary~\ref{cor:sharp_base_dense_in_full_H} and Lemma~\ref{lem:QE3_core_topology}.
The final sentence records that the bounded base state is the restriction of this OS-positive sharp-local state, so its positive-time semigroup is canonically inherited from the Lane~A reconstruction and no extra continuum-time lattice theorem is needed.
\end{proof}

\dependencyledger
  {Theorem~\ref{thm:continuum_tightness}, Corollary~\ref{cor:LaneA_bounded_state_restriction}, Theorem~\ref{thm:LaneA_closure}(b), and Corollary~\ref{cor:LaneA_full_scope}}
  {the standard OS reconstruction theorem for reflection-positive states on a positive-time algebra, applied only after the Lane~A seam package has produced the full sharp-local state}
  {the bounded base state is identified as the $\mathcal A_+$ restriction of the full sharp-local state, and its dyadic matrix elements extend to the canonical continuum-time semigroup for all $t\ge 0$; the theorem does not yet perform the density/graph-core or gap transfer}
  {Corollary~\ref{cor:sharp_base_dense_in_full_H}, Lemma~\ref{lem:QE3_core_topology}, and Theorem~\ref{thm:phys_gap_noncollapse}}

\begin{remark}[Phase-II seam note for Theorem~\ref{thm:continuum_time_OS_upgrade}]
Corollary~\ref{cor:LaneA_bounded_state_restriction} is the bounded-state compatibility theorem only.
The finite-cap sharp-local OS structure used above is Theorem~\ref{thm:LaneA_closure}(b), the passage to the full inductive union is Corollary~\ref{cor:LaneA_full_scope}, and the later density/graph-core handoff is isolated separately in Corollary~\ref{cor:sharp_base_dense_in_full_H} and Lemma~\ref{lem:QE3_core_topology}.
\end{remark}

\begin{remark}[Phase-IV dyadic-to-all-time route note]
Theorem~\ref{thm:continuum_tightness} is only the bounded dyadic OS-limit theorem on $\mathcal A_+$. Theorem~\ref{thm:continuum_time_OS_upgrade} is the only public-spine step that upgrades those dyadic matrix elements to all $t\ge 0$, and it does so only after Theorem~\ref{thm:LaneA_closure}(b) and Corollary~\ref{cor:LaneA_full_scope} have produced the sharp-local OS semigroup. No direct all-time lattice theorem is used or needed.
\end{remark}

\begin{corollary}[The bounded base algebra is dense in the full sharp-local vacuum sector]\label{cor:sharp_base_dense_in_full_H}
In the setting of Theorem~\ref{thm:continuum_time_OS_upgrade}, let $(\mathcal H,\Omega,H)$ be the OS reconstruction of the full sharp-local state.
Then the image of $\mathcal A_+$ is dense in $\mathcal H$, and
\[
\overline{\mathrm{span}}\{[F-\omega^{(u_{\mathrm{ren}})}(F)]:F\in\mathcal A_+\}=\Omega^\perp.
\]
In particular, the non-collapse argument may be carried out on bounded positive-time observables and then extended by density to the full sharp-local Hilbert space.
\end{corollary}

\begin{proof}
Theorem~\ref{thm:LaneA_bounded_base_cyclic} proves that, for the full sharp-local OS reconstruction, the image of the bounded positive-time algebra is dense and its centered classes span the vacuum-orthogonal sector. Theorem~\ref{thm:continuum_time_OS_upgrade} singles out exactly the same full sharp-local state as the one whose restriction to $\mathcal A_+$ is the tuned bounded base state. Rewriting the conclusions of Theorem~\ref{thm:LaneA_bounded_base_cyclic} in the notation of the present theorem gives the stated density and centered-span claims.
\end{proof}

\begin{lemma}[Hilbert/graph-core approximation package used in QE3]\label{lem:QE3_core_topology}
In the setting of Theorem~\ref{thm:continuum_time_OS_upgrade}, two approximation facts are used in the final gap transfer.
\begin{enumerate}[label=(\roman*),leftmargin=*]
\item The bounded positive-time algebra $\mathcal A_+$ is dense in $\mathcal H$, and its centered classes are dense in $\Omega^\perp$, by Corollary~\ref{cor:sharp_base_dense_in_full_H}.
\item Whenever a finite family of renormalized sharp-local fields belongs to a fixed engineering cap $\Delta_{\max}$, the corresponding field-polynomial vectors are realized as limits of vectors from the common invariant core $\mathcal D_{\mathrm{flow}}^{(\le \Delta_{\max})}$ of Lemma~\ref{lem:finite_cap_common_core}; the same approximants converge not only in the OS Hilbert norm but also in every finite graph norm relevant to that cap.
\end{enumerate}
For each $t\ge 0$ one has $\|e^{-tH}\|\le 1$, so any semigroup estimate established first on the bounded positive-time core extends by Hilbert-norm continuity to the capwise sharp-local limits from part~(ii).
Thus the QE3 passage from bounded observables to the full sharp-local Hilbert space uses only the explicit density/core package already isolated in Lane~A and no hidden domain argument.
\end{lemma}

\begin{proof}
Part~(i) is Corollary~\ref{cor:sharp_base_dense_in_full_H}, which already identifies the bounded positive-time algebra as dense in $\mathcal H$ and its centered classes as dense in $\Omega^\perp$.
For part~(ii), Lemma~\ref{lem:finite_cap_common_core} proves on each fixed engineering cap that the field-polynomial vectors lie in the closure of the common invariant core $\mathcal D_{\mathrm{flow}}^{(\le \Delta_{\max})}$ in every finite graph norm relevant to that cap. Corollary~\ref{cor:LaneA_full_scope} transports those finite-cap realizations compatibly through the inductive union, so the same approximants may be viewed inside the full sharp-local reconstruction. Finally, because $e^{-tH}$ is a contraction on the OS Hilbert space for each $t\ge 0$, any semigroup estimate first proved on the bounded positive-time core extends by Hilbert-norm continuity to the capwise sharp-local limit vectors. Thus the QE3 passage uses exactly the explicit density and graph-core package stated here and no hidden domain-completion step.
\end{proof}

\begin{remark}[Phase-IV finiteness note for Lemma~\ref{lem:QE3_core_topology}]
Every sharp-local vector later inserted into the QE3 transfer first lies in one fixed engineering cap, and only the finite list of graph norms attached to that cap is used. The lemma therefore does not invoke a graph-norm completion uniform over the full inductive union; the full-union step enters only through Corollary~\ref{cor:sharp_base_dense_in_full_H} and the boundedness of the semigroup.
\end{remark}

\dependencyledger
  {the tuned full fixed-lattice gap input, Theorem~\ref{thm:Wilson_patch_transport}, Theorem~\ref{thm:continuum_tightness}, Theorem~\ref{thm:continuum_time_OS_upgrade}, and Lemma~\ref{lem:QE3_core_topology}}
  {the spectral theorem only, used through the local semigroup and Laplace-transform lemmas isolated in this appendix}
  {the fixed-lattice lower bound transfers to the full sharp-local OS Hamiltonian, yielding $\ker(H)=\mathbb C\Omega$ and $\mathrm{Spec}(H)\subset\{0\}\cup[m_*,\infty)$; the theorem is the gap-transfer node itself, not merely a later reformulation}
  {Corollary~\ref{cor:route1_QE3_gap}, Theorem~\ref{thm:OS_axioms_package} via Corollary~\ref{cor:OS_cluster_vacuum}, and the main-theorem assembly}

\paragraph{Primary non-collapse route from the full fixed-lattice gap.}
The next theorem compresses the remaining Route~1 continuum transfer into one theorem-level statement.
Its inputs are: the tuned-sequence lower bound on the physical full fixed-lattice OS gap; the actual Wilson bounded-state and dyadic reflected limits recorded in Theorem~\ref{thm:continuum_tightness} (whose Wilson character is supplied upstream by the same-scale transport theorem~\ref{thm:Wilson_patch_transport}); the continuum-time OS upgrade of Theorem~\ref{thm:continuum_time_OS_upgrade}; and the explicit Hilbert/graph-core approximation package just recorded in Lemma~\ref{lem:QE3_core_topology}.
The dense-core Route~1/Lane~B package retained later is therefore a downstream reformulation rather than an additional closure interface.

\begin{remark}[Packet~3/8 transport certification ledger]\label{rem:QE3_dependency_map}\label{rem:route1_packet38_transport_certification}
On the frozen Route~1 spine, the fixed-lattice-to-continuum gap route is
\[
\begin{aligned}
\text{Corollary~\ref{cor:A1_compact_simple_spine}}
&\Longrightarrow \text{Proposition~\ref{prop:route1_QE1}}
\Longrightarrow \text{Proposition~\ref{prop:route1_QE2}}\\
&\Longrightarrow \text{Theorem~\ref{thm:Wilson_patch_transport}}
\Longrightarrow \text{Theorem~\ref{thm:continuum_tightness}}\\
&\Longrightarrow \text{Theorem~\ref{thm:continuum_time_OS_upgrade}}
\Longrightarrow \text{Theorem~\ref{thm:phys_gap_noncollapse}}
\Longrightarrow \text{Corollary~\ref{cor:route1_QE3_gap}}.
\end{aligned}
\]
The only live weak-window input on this chain is the repaired curvature-sector package certified by Theorem~\ref{thm:weak_window_sufficiency}; after that certificate, no edge-Hessian placeholder or auxiliary weak-window branch re-enters the gap proof.
The fixed-lattice gap entering the transport step is the tuned-sequence full OS gap package Proposition~\ref{prop:route1_QE2}, whose theorem-level closure is Theorem~\ref{thm:physical_gap_from_gluing}; Theorem~\ref{thm:Wilson_patch_transport} is the same-scale Wilson-to-patched return, and Theorem~\ref{thm:continuum_tightness} is the bounded Wilson/dyadic OS limit theorem.
The first Lane~A theorem-level compatibility input actually consumed by this transport route is Corollary~\ref{cor:LaneA_bounded_state_restriction}. In Theorem~\ref{thm:continuum_time_OS_upgrade} it supplies only the identification of the sharp-local state with the bounded base state on $\mathcal A_+$ and the dyadic OS matrix-element identity. The finite-cap sharp-local OS structure used there is Theorem~\ref{thm:LaneA_closure}(b) on each finite cap, while Corollary~\ref{cor:LaneA_full_scope} is only the later inductive-union passage. The later Hilbert/graph-core passage used by Theorem~\ref{thm:phys_gap_noncollapse} is exactly Corollary~\ref{cor:sharp_base_dense_in_full_H} together with Lemma~\ref{lem:QE3_core_topology}.
Lane~B is not part of the proof of the gap: Theorem~\ref{thm:route1_E2_internal} and Theorem~\ref{thm:R1_HS} are downstream reformulations of the already established QE3 conclusion.
\end{remark}

\begin{theorem}[Continuum gap transfer from the tuned full fixed-lattice Route~1 input]\label{thm:phys_gap_noncollapse}
Fix a target value $u_{\mathrm{ren}}\in(0,u_\ast]$ and let
\[
a_n:=\ell_0 2^{-n},\qquad \beta_n:=\beta(a_n;u_{\mathrm{ren}})
\]
be the tuned dyadic trajectory.
Let $\omega^{(u_{\mathrm{ren}})}$ be the full sharp-local state furnished by Corollary~\ref{cor:LaneA_full_scope} over the same bounded base state, and let $(\mathcal H,\Omega,H)$ be its OS reconstruction on the full sharp-local positive-time algebra.
Assume in addition that there exists $m_*>0$ such that
\[
\liminf_{n\to\infty}\frac{\Delta_{\mathrm{lat}}(\beta_n)}{a_n}\ge m_*>0.
\]
Then the canonical continuum-time OS semigroup on the full sharp-local Hilbert space has spectral gap at least $m_*$, equivalently
\[
\mathrm{Spec}(H)\subset\{0\}\cup[m_*,\infty),
\qquad
\ker(H)=\mathbb C\Omega.
\]
The only topology used in the passage from bounded positive-time observables to sharp-local vectors is the Hilbert-norm / finite-graph-norm approximation package of Lemma~\ref{lem:QE3_core_topology}.
\end{theorem}

\paragraph*{OS spectral-gap extraction from Euclidean decay (no hidden constants)}

\begin{lemma}[OS semigroup representation of time-separated two-point functions]
\label{lem:OS_semigroup_representation}
Let $\langle\cdot\rangle_{\mathrm{cont},\ell}$ be an OS-positive continuum Schwinger functional on a translation-invariant
algebra $\mathcal A$ with positive-time subalgebra $\mathcal A_+$ and reflection $\Theta$.
Let $(\mathcal H,\Omega,H)$ be the OS reconstruction: $\Omega$ is the vacuum, $H\ge 0$ the Hamiltonian,
and $\tau_t$ denotes Euclidean time translation.
Then for every $F\in\mathcal A_+$ and every $t\ge 0$,
\begin{equation}
 \langle F\,\Theta\,\tau_t F\rangle_{\mathrm{cont},\ell}
 \;=\;
 \langle [F],\,e^{-tH}\,[F]\rangle_{\mathcal H},
 \label{eq:OS_semigroup_identity}
\end{equation}
where $[F]$ is the OS equivalence class of $F$ in the reconstructed Hilbert space.
Moreover, $\langle F\rangle_{\mathrm{cont},\ell}=\langle \Omega,[F]\rangle_{\mathcal H}$, so
\begin{equation}
 [F_0]\perp\Omega\qquad\text{for }F_0:=F-\langle F\rangle_{\mathrm{cont},\ell}.
 \label{eq:OS_mean_zero_orth}
\end{equation}
\end{lemma}

\begin{proof}
This is the standard OS reconstruction identity.
The pre-Hilbert space is $\mathcal A_+/\mathcal N$ where $\mathcal N$ is the OS null ideal, with inner product
$(F,G):=\langle F\,\Theta G\rangle_{\mathrm{cont},\ell}$.
Time translations act by $T(t)[F]=[\tau_t F]$ and are well-defined and contractive by translation invariance and OS positivity.
The OS axioms imply $\{T(t)\}_{t\ge 0}$ is a strongly continuous symmetric contraction semigroup on $\mathcal H$;
hence $T(t)=e^{-tH}$ for a self-adjoint $H\ge 0$.
Then
\[\langle [F],e^{-tH}[F]\rangle=( [F],T(t)[F])=(F,\tau_t F)=\langle F\,\Theta\,\tau_t F\rangle,\]
giving \eqref{eq:OS_semigroup_identity}.
Finally $\Omega=[1]$, so $\langle F\rangle=(F,1)=\langle [F],\Omega\rangle$, and \eqref{eq:OS_mean_zero_orth} follows.
\end{proof}

\begin{lemma}[Laplace-transform gap criterion]
\label{lem:laplace_gap_criterion}
Let $H\ge 0$ be self-adjoint on a Hilbert space $\mathcal H$ and let $\psi\in\mathcal H$.
Assume there exist constants $m>0$ and $C<\infty$ such that for all $t\ge 0$,
\begin{equation}
 \langle \psi, e^{-tH}\psi\rangle\ \le\ C\,e^{-m t}.
 \label{eq:laplace_decay_assumption}
\end{equation}
Then the spectral measure $\mu_{\psi}$ of $H$ associated to $\psi$ satisfies
\begin{equation}
 \mu_{\psi}\big([0,m)\big)=0.
 \label{eq:spectral_support_gap}
\end{equation}
In particular, if $\psi\perp\ker(H)$ then $\psi$ has spectral support contained in $[m,\infty)$.
\end{lemma}

\begin{proof}
By the spectral theorem there exists a finite positive measure $\mu_{\psi}$ on $[0,\infty)$ such that
\begin{equation}
 \langle \psi, e^{-tH}\psi\rangle\;=\;\int_{0}^{\infty} e^{-t\lambda}\,d\mu_{\psi}(\lambda),
 \qquad t\ge 0.
 \label{eq:spectral_laplace}
\end{equation}
Fix $\varepsilon\in(0,m)$. Then for all $t\ge 0$,
\[
\int_{[0,m-\varepsilon]} e^{-t\lambda}\,d\mu_{\psi}(\lambda)
\ \ge\ e^{-t(m-\varepsilon)}\,\mu_{\psi}([0,m-\varepsilon]).
\]
Combine this lower bound with \eqref{eq:laplace_decay_assumption} and \eqref{eq:spectral_laplace} to obtain
\[
\mu_{\psi}([0,m-\varepsilon])
\ \le\ C\,e^{-\varepsilon t}\qquad\forall t\ge 0.
\]
Let $t\to\infty$ to get $\mu_{\psi}([0,m-\varepsilon])=0$.
Since $\varepsilon\downarrow 0$ and $[0,m)=\bigcup_{n\ge 1}[0,m-1/n]$, we obtain \eqref{eq:spectral_support_gap}.
If $\psi\perp\ker(H)$ then $\mu_{\psi}(\{0\})=0$, hence the spectral support is contained in $[m,\infty)$.
\end{proof}

\begin{lemma}[Centered positive-time classes are dense in the vacuum-orthogonal sector]\label{lem:OS_centered_dense}
Let $(\mathcal H,\Omega)$ be the OS reconstruction of a reflection-positive state on $\mathcal A_+$.
Then
\[
\overline{\mathrm{span}}\{[F-\omega(F)]:F\in\mathcal A_+\}=\Omega^\perp.
\]
\end{lemma}

\begin{proof}
By construction of the OS Hilbert space, the span of classes $[F]$ with $F\in\mathcal A_+$ is dense in $\mathcal H$.
For any such $F$ one has
\[
[F]=[F-\omega(F)]+\omega(F)\,\Omega,
\]
so the span of centered classes together with $\Omega$ is dense in $\mathcal H$.
Taking orthogonal projection to $\Omega^\perp$ gives the claim.
\end{proof}

\begin{proof}[Proof of Theorem~\ref{thm:phys_gap_noncollapse}]
By Theorem~\ref{thm:continuum_time_OS_upgrade}, the sharp-local state $\omega^{(u_{\mathrm{ren}})}$ restricts to the bounded base state on $\mathcal A_+$ and provides the canonical continuum-time OS semigroup for that base state.
Theorem~\ref{thm:continuum_tightness}, whose Wilson character is supplied upstream by the same-scale transport closure in Theorem~\ref{thm:Wilson_patch_transport}, records the actual Wilson one-point limits and dyadic reflected OS matrix elements needed below.
Fix $\varepsilon\in(0,m_*)$.
By the hypothesis on the tuned dyadic sequence, there exists $n_\varepsilon<\infty$ such that
\[
\frac{\Delta_{\mathrm{lat}}(\beta_n)}{a_n}\ge m_*-\varepsilon
\qquad (n\ge n_\varepsilon).
\]

\paragraph{Step 1 (fixed-lattice OS semigroup estimate on the dyadic time set).}
Fix a bounded gauge-invariant local observable $F\in\mathcal A_+$ and set
\[
F_0:=F-\omega^{(u_{\mathrm{ren}})}(F),
\qquad
F_{0,n}:=F-\omega_{a_n}(F).
\]
For each $n$, let $(\mathcal H_n,\Omega_n,H_n^{\mathrm{lat}})$ be the fixed-lattice OS reconstruction of the Wilson Gibbs state $\omega_{a_n}$ on the bounded gauge-invariant positive-time cylinder algebra, so that
\[
\Delta_{\mathrm{lat}}(\beta_n)=\inf\bigl(\mathrm{Spec}(H_n^{\mathrm{lat}})\setminus\{0\}\bigr).
\]
Let $s\in\mathbb D_{\mathrm{dyad}}$.
For every sufficiently large $n$ one has $s/a_n\in\mathbb N$.
Since $\Delta_{\mathrm{lat}}(\beta_n)$ is the full fixed-lattice OS gap in lattice units, the spectral theorem gives
\[
0\le
\omega_{a_n}(\Theta F_{0,n}\,\tau_s F_{0,n})
=\langle [F_{0,n}],e^{-(s/a_n)H_n^{\mathrm{lat}}}[F_{0,n}]\rangle
\le
\omega_{a_n}(\Theta F_{0,n}F_{0,n})\,e^{-\Delta_{\mathrm{lat}}(\beta_n)s/a_n}.
\]
For $n\ge n_\varepsilon$ this yields
\begin{equation}
\omega_{a_n}(\Theta F_{0,n}\,\tau_s F_{0,n})
\le
\omega_{a_n}(\Theta F_{0,n}F_{0,n})\,e^{-(m_*-\varepsilon)s}
\qquad (s\in\mathbb D_{\mathrm{dyad}}).
\label{eq:dyadic_lattice_decay_eps}
\end{equation}

\paragraph{Step 2 (pass to the continuum on dyadic times).}
For each fixed dyadic $s$, Theorem~\ref{thm:continuum_tightness} gives convergence of the one-point functions and of the dyadic OS matrix element.
Using Theorem~\ref{thm:continuum_time_OS_upgrade} to identify the limit with $\omega^{(u_{\mathrm{ren}})}$, we may send $n\to\infty$ in \eqref{eq:dyadic_lattice_decay_eps} and obtain
\begin{equation}
\omega^{(u_{\mathrm{ren}})}(\Theta F_0\,\tau_s F_0)
\le
\omega^{(u_{\mathrm{ren}})}(\Theta F_0F_0)\,e^{-(m_*-\varepsilon)s}
\qquad (s\in\mathbb D_{\mathrm{dyad}}).
\label{eq:dyadic_cont_decay_eps}
\end{equation}

\paragraph{Step 3 (extend from dyadic times to all $t\ge 0$).}
By Theorem~\ref{thm:continuum_time_OS_upgrade}, for every $t\ge 0$,
\[
\omega^{(u_{\mathrm{ren}})}(\Theta F_0\,\tau_t F_0)=\langle [F_0],e^{-tH}[F_0]\rangle_{\mathcal H},
\]
and the right-hand side is continuous in $t$.
Choose any dyadic sequence $s_k\to t$.
Passing to the limit in \eqref{eq:dyadic_cont_decay_eps} gives
\begin{equation}
\langle [F_0],e^{-tH}[F_0]\rangle_{\mathcal H}
\le
\|[F_0]\|_{\mathcal H}^2\,e^{-(m_*-\varepsilon)t}
\qquad (t\ge 0).
\label{eq:alltime_cont_decay_eps}
\end{equation}

\paragraph{Step 4 (spectral support on the bounded positive-time core).}
Apply Lemma~\ref{lem:laplace_gap_criterion} to \eqref{eq:alltime_cont_decay_eps} with $\psi=[F_0]$, $C=\|[F_0]\|^2$, and $m=m_*-\varepsilon$.
Since $[F_0]\perp\Omega\subset\ker(H)$, we conclude that $[F_0]$ has spectral support contained in $[m_*-\varepsilon,\infty)$.
Because $\varepsilon\downarrow 0$ is arbitrary, every centered class $[F_0]$ has spectral support contained in $[m_*,\infty)$.

\paragraph{Step 5 (closure on the full sharp-local Hilbert space).}
By Lemma~\ref{lem:QE3_core_topology}(i), centered bounded positive-time classes are dense in $\Omega^\perp$.
Hence
\[
P_{(0,m_*)}\big|_{\Omega^\perp}=0.
\]
Moreover, letting $t\to\infty$ in \eqref{eq:alltime_cont_decay_eps} gives
\[
\|P_{\{0\}}[F_0]\|^2=\lim_{t\to\infty}\langle [F_0],e^{-tH}[F_0]\rangle=0,
\]
so $P_{\{0\}}[F_0]=0$ for every centered class.
By density and boundedness of $P_{\{0\}}$, this implies $P_{\{0\}}\big|_{\Omega^\perp}=0$.
Thus
\[
\ker(H)=\mathbb C\Omega,
\qquad
\mathrm{Spec}(H)\subset\{0\}\cup[m_*,\infty).
\]
If one wants the same conclusion on vectors created by sharp-local field polynomials in a fixed engineering cap rather than only on bounded cylinder classes, Lemma~\ref{lem:QE3_core_topology}(ii) gives the required Hilbert/graph-norm approximants and the boundedness of $e^{-tH}$ passes the same semigroup estimate to those capwise limits.
This is the claimed continuum gap transfer.
\end{proof}

\begin{remark}[Continuum-scaling inputs]
The front-facing non-collapse route now uses three continuum-scaling ingredients along the chosen tuned dyadic trajectory:
(i) selecting and controlling the tuned sequence $a_n=\ell_0 2^{-n}$ and $\beta_n=\beta(a_n;u_{\mathrm{ren}})$;
(ii) dyadic bounded-state convergence together with the continuum-time OS transport furnished by Theorems~\ref{thm:continuum_tightness} and~\ref{thm:continuum_time_OS_upgrade};
(iii) non-collapse of the physical mass scale, i.e.\ a quantitative lower bound on $\Delta_{\mathrm{lat}}(\beta_n)/a_n$.
These ingredients are now all closed on the primary Route~1 spine, and Theorem~\ref{thm:phys_gap_noncollapse} converts them into the continuum spectral gap.
\end{remark}


\paragraph*{OS/Wightman packaging pointer}
The formal translation from the Euclidean OS theory to the Minkowski/Wightman formulation is no longer proved here inline.
Instead, Appendix~\ref{app:axioms_package} collects the existence dossier in one place: it packages the Euclidean axioms already established by Lane~A, the continuum-time OS reconstruction used in QE3, the positive-time clustering/vacuum-uniqueness consequences of the Route~1 gap, the standard Osterwalder--Schrader analytic continuation to a Poincar\'e-covariant Wightman/Haag--Kastler theory, and the field correspondence for the gauge-invariant local differential polynomials represented by the sharp-local algebra.
The Minkowski mass-gap corollary used for the final theorem is therefore cited from Appendix~\ref{app:axioms_package}.


\subsubsection{Step-scaling approach to the physical gap: a concrete plan}
The non-collapse condition in Theorem~\ref{thm:phys_gap_noncollapse} is best attacked by binding the RG entrance scale
(to the KP basin proved above) to a \emph{renormalized, scheme-defined physical length}.
A standard and analytically well-behaved choice is the gradient-flow coupling.

\paragraph{A finite-volume coupling and step scaling}
Fix a finite four-torus $\mathbb T_L^4$ of linear size $L$ (in lattice units) with periodic boundary conditions.
Let $B_\mu(t,x)$ be the Wilson/gradient flow field at flow time $t$ started from the lattice gauge field.
Define the flowed action density $E(t,x)$ and the finite-volume gradient-flow coupling (one standard normalization)
\begin{equation}\label{eq:gGF_def}
 g_{\mathrm{GF}}^2(L) := \mathcal N^{-1}\, t^2\,\langle E(t,0)\rangle\Big|_{t=c L^2},
\end{equation}
with fixed $c\in(0,1)$ and group-dependent constant $\mathcal N$ chosen so that $g_{\mathrm{GF}}^2(L)=g_0^2+O(g_0^4)$ in perturbation theory.
Define the step-scaling function
\begin{equation}\label{eq:stepscaling_def_plan}
 \Sigma(u) := g_{\mathrm{GF}}^2(2L)\ \text{given}\ g_{\mathrm{GF}}^2(L)=u.
\end{equation}

\paragraph{Entrance to the KP basin as a \emph{large}-coupling threshold}
The KP-basin entrance proved in this document is formulated as a deterministic \emph{smallness} condition
on a block-local activity norm \emph{after} coarse graining to scale $L$.
In physical terms this corresponds to the blocked single-plaquette law becoming sufficiently \emph{spread} on $G$
(close to Haar in the positive-type sector), which is an \emph{infrared} effect.
Accordingly, when the scale-dependent coupling is defined by a flowed energy density,
KP admissibility should be triggered once the coupling has grown past a fixed threshold.

A standard way to bind this to \eqref{eq:gGF_def} is therefore to prove the one-way implication
\begin{equation}\label{eq:KP_from_small_u}
 g_{\mathrm{GF}}^2(L)\ge u_0\ \Longrightarrow\ \text{KP admissibility holds for the blocked measure at scale $L$},
\end{equation}
for some universal threshold $u_0=u_0(d,G,\theta,b)$.
In the weak-to-intermediate coupling window, \eqref{eq:KP_from_small_u} is expected to follow from a character-expansion
bound that converts a lower bound on the local fluctuation scale (as measured by $g_{\mathrm{GF}}^2(L)$)
into an upper bound on the nontrivial Fourier--Peter--Weyl eigenvalues of the blocked plaquette density,
hence a KP-norm bound.

\paragraph{Scale setting and tuned dyadic normalization.}
Define the lattice spacing via a flow scale by choosing $t_0=t_0(\beta)$ through
\begin{equation}\label{eq:t0_def}
 t_0^2\,\langle E(t_0,0)\rangle = c_0
\end{equation}
with fixed $c_0$, and set
\begin{equation}\label{eq:a_from_t0}
 a(\beta):=\frac{\ell_0}{\sqrt{t_0(\beta)}}.
\end{equation}
For the final $\mathrm{QE1}\Rightarrow\mathrm{QE2}\Rightarrow\mathrm{QE3}$ chain we do \emph{not} use a separate first-threshold scale $L_*(\beta)$ or a fixed-$\beta$ coupling-to-KP implication in the main appendix.
The only weak-window ingredients retained here are the theorem-level step-scaling control proved below (Proposition~\ref{prop:stepscaling_control_proved}), the continuity statement of Lemma~\ref{lem:stepscaling_continuous}, and the explicit bare-to-flow comparison of Theorem~\ref{thm:h49_bare_flow_window}, which are used later in the tuned dyadic recursion of Proposition~\ref{prop:h59_uv_tuning_exists}.
The older fixed-$\beta$ coupling-to-KP threshold and $L_*(\beta)/\sqrt{t_0(\beta)}$ comparison package has been moved to the project supplement and is not cited on the load-bearing path.

\subsection{Step scaling and non-collapse: a weak-coupling assembled proof that closes the remaining wall}
\label{sec:stepscaling_proof}

This section records the exact weak-window input retained on the load-bearing path: a two-sided cubic remainder estimate for the finite-volume step-scaling map together with the local coupling-identification statement used later in the tuned dyadic recursion.
The older fixed-$\beta$ threshold package has been moved to the supplement and is no longer cited in the main Appendix~F proof spine.

\subsubsection{Finite-volume gradient-flow coupling and its step scaling}
Fix a physical flow scheme (e.g. Wilson flow) and define the finite-volume coupling
\begin{equation}\label{eq:F_gf_coupling_def}
 g^2_{\mathrm{GF}}(L) := \mathcal N^{-1} \; t^2\,\langle E(t)\rangle_{\Lambda(L)}\Big|_{t=c L^2},
\end{equation}
with $\Lambda(L)$ a periodic box of linear size $L$ (in lattice units) and fixed $c\in(0,1)$.
Let $\Sigma(u)$ denote the step-scaling function at scale factor $2$:
\begin{equation}\label{eq:stepscaling_def}
\Sigma(u):= g^2_{\mathrm{GF}}(2L)\ \text{ when }\ g^2_{\mathrm{GF}}(L)=u.
\end{equation}

\subsubsection{Removing the remaining analytic hypothesis: localization and curvature-dominated Brascamp--Lieb}
\label{subsec:convexity_to_stepscaling}
\label{subsec:curvature_domination}
This subsection records the live weak-window package in the exact form used later.
Historical discussion of the earlier edge-field formulation has been moved to the supplementary Route~1 archive.
The only ingredients retained on the load-bearing path are small-curvature localization and curvature
Brascamp--Lieb domination in the plaquette sector---exactly the sector seen by the flowed energy density.
The downstream sufficiency certificate is stated later as Theorem~\ref{thm:weak_window_sufficiency}.

\paragraph{Spanning-tree gauge and exponential chart.}
Fix a spanning tree $T$ and impose tree gauge. Parametrize free edges by exponential coordinates
$U_e=\exp(iA_e)$ in a normal chart of radius $\varepsilon_0>0$.
For a plaquette $p$ write $F_p(A)\in\mathfrak g$ for the logarithm of the plaquette holonomy in this chart:
\(U_p=\exp(iF_p(A))\).
Define the small-field event
\begin{equation}\label{eq:Omega_sf_def}
\Omega_{\varepsilon_0}:=\bigl\{A: \ \|F_p(A)\|\le \varepsilon_0\ \text{for all plaquettes }p\bigr\}.
\end{equation}

\paragraph{Plaquette potential.}
Let
\begin{equation}\label{eq:V_def}
V(X):=1-\frac{1}{\dim}\,\Re\,\mathrm{Tr}\,e^{iX},\qquad X\in\mathfrak g,
\end{equation}
where $\dim$ is the dimension of the fixed Wilson representation $\rho$ and the inner product on $\mathfrak g$ is fixed once and for all.
Since $G$ is compact and $V$ is smooth with $\nabla^2 V(0)\succ 0$, there exists $\varepsilon_0>0$ and constants
$0<m_-\le m_+<\infty$ such that for all $\|X\|\le \varepsilon_0$,
\begin{equation}\label{eq:V_hess_bounds}
m_- I\ \preceq\ \nabla^2 V(X)\ \preceq\ m_+ I.
\end{equation}

\begin{lemma}[Weak-coupling localization to the exponential chart]\label{lem:weak_localization}
Fix $\varepsilon_0$ as above. There exist constants $c_\mathrm{loc}(G,\varepsilon_0)>0$ and $\beta_\mathrm{loc}$ such that
for all finite volumes $\Lambda$ and all $\beta\ge\beta_\mathrm{loc}$,
\begin{equation}\label{eq:localization_bound}
\mu_{\beta,\Lambda}(\Omega_{\varepsilon_0}^c)\ \le\ |P(\Lambda)|\,e^{-c_\mathrm{loc}\,\beta}.
\end{equation}
\end{lemma}

\begin{proof}
If $\|F_p(A)\|>\varepsilon_0$ for some plaquette $p$, then $U_p$ lies outside a fixed neighborhood of the identity,
and hence the plaquette energy satisfies $V(F_p(A))\ge v_0(G,\varepsilon_0)>0$.
Since the Wilson density is proportional to $\exp\{-\beta\sum_{q}V(F_q(A))\}$, conditioning on the complement of any single plaquette event
and integrating out the remaining variables yields a bound of the form
$\mu_{\beta,\Lambda}(\|F_p\|>\varepsilon_0)\le e^{-\beta v_0}$.
Summing over $p\in P(\Lambda)$ gives \eqref{eq:localization_bound}.
\end{proof}

\paragraph{Curvature Brascamp--Lieb domination.}
On $\Omega_{\varepsilon_0}$, the action is a sum of strictly convex plaquette potentials in the curvature variables $F_p(A)$.
The needed uniform estimate is obtained directly in the curvature sector, because the map $A\mapsto dA$ isolates the plaquette modes seen by the flowed energy density.

Let $d$ be the cochain differential mapping (tree-gauge) 1-forms $(A_e)$ to 2-forms on plaquettes (linearized curvature);
write $d^*$ for its adjoint in the standard counting inner products on cochains.

\begin{lemma}[Projection identity for the linearized curvature covariance]\label{lem:projection_identity}
Let $\Delta_1:=d^*d$ on 1-forms. Then the operator
\begin{equation}\label{eq:proj_exact}
\Pi_{\mathrm{ex}}:=d\,\Delta_1^{-1}d^*
\end{equation}
is the orthogonal projection from 2-forms onto the exact subspace $\mathrm{im}(d)$.
In particular $\|\Pi_{\mathrm{ex}}\|_{2\to 2}=1$.
\end{lemma}

\begin{proof}
On a finite complex, $\Delta_1$ is invertible on tree-gauge 1-forms. The Moore--Penrose identity implies
$d\Delta_1^{-1}d^*$ is the orthogonal projection onto $\mathrm{im}(d)$.
\end{proof}

\begin{proposition}[Curvature Brascamp--Lieb domination]\label{prop:curvature_BL}
Let $\nu_{\beta,\Lambda}$ be the gauge-fixed Wilson measure in tree gauge on a finite volume $\Lambda$.
On the event $\Omega_{\varepsilon_0}$, for every finitely supported 2-form $J=(J_p)_{p\in P(\Lambda)}$ and every $s\in\mathbb R$,
\begin{equation}\label{eq:curvature_mgf}
\log\,\mathbb E_{\nu_{\beta,\Lambda}}\Big[\exp\big(s\langle J, dA\rangle\big)\ \big|\ \Omega_{\varepsilon_0}\Big]
\ \le\ \frac{s^2}{2\beta m_-}\,\langle J,\Pi_{\mathrm{ex}}J\rangle
\ \le\ \frac{s^2}{2\beta m_-}\,\|J\|_2^2.
\end{equation}
Consequently, all centered linear curvature functionals are sub-Gaussian with variance proxy $(\beta m_-)^{-1}$,
uniformly in $\Lambda$.
\end{proposition}

\begin{proof}
On $\Omega_{\varepsilon_0}$ the plaquette potential Hessian bounds \eqref{eq:V_hess_bounds} imply that the full gauge-fixed action
has (conditional) Hessian bounded below by $\beta m_- d^*d$ in the sense of quadratic forms on the image of $d^*$.
Applying the Brascamp--Lieb inequality to the linear functional $A\mapsto\langle d^*J, A\rangle$ yields
\[\mathrm{Var}(\langle d^*J, A\rangle\mid\Omega_{\varepsilon_0})\le \frac{1}{\beta m_-}\,\langle d^*J,\Delta_1^{-1} d^*J\rangle
=\frac{1}{\beta m_-}\,\langle J,\Pi_{\mathrm{ex}}J\rangle.\]
The sub-Gaussian mgf bound \eqref{eq:curvature_mgf} follows by the usual Herbst argument for convex measures.
\end{proof}

\paragraph{From conditional to unconditional estimates.}
By Lemma~\ref{lem:weak_localization}, the complement of $\Omega_{\varepsilon_0}$ is exponentially small in $\beta$.
Thus any inequality proved conditionally on $\Omega_{\varepsilon_0}$ for bounded curvature observables may be upgraded to an
unconditional inequality by paying an additive error $O(|P(\Lambda)|e^{-c\beta})$, which is negligible in the step-scaling and
non-collapse regime where $\beta\to\infty$.


\paragraph*{Uniform weak-window step-scaling remainder control}
\label{subsec:hardened44_stepscaling_remainder}

The entrance/non-collapse package uses one quantitative UV input: a two-scale step-scaling estimate for the chosen finite-volume
coupling $g^2_{\mathrm{GF}}(L)$ at scale factor $2$.
In earlier roadmap form this estimate was recorded via an informal ``Laplace/Wick expansion'' description.
Here we make the dependence on \emph{curvature} Brascamp--Lieb domination explicit.
For the mass-gap/non-collapse chain we only need that the quadratic coefficient in the weak-window step-scaling expansion is strictly
positive; its closed form is immaterial.

\medskip
\noindent
\textbf{Fixed constants.}
Fix once and for all: the chart radius $\varepsilon_0$ and convexity constants $m_\pm$ from \eqref{eq:V_hess_bounds},
the flow parameter $c\in(0,1)$ and normalization $\mathcal N$ from \eqref{eq:gf_coupling_def},
and the quadratic step-scaling coefficient $b_0>0$ from Lemma~\ref{lem:b0_positive} (positivity only; closed form not used below).
Set the curvature-variance proxy
\begin{equation}\label{eq:h44_sigma_def}
\sigma_\beta^2:=\frac{1}{\beta m_-}.
\end{equation}

\begin{lemma}[Weak-window step scaling with a uniform $O(u^3)$ remainder]
\label{lem:flow_energy_expansion}
Assume Lemma~\ref{lem:weak_localization} and Proposition~\ref{prop:curvature_BL}.
There exist $\beta_1\ge\beta_{\mathrm{loc}}$, a weak-window threshold $u_{\mathrm{AF}}>0$, and a constant $C_{\Sigma}<\infty$,
depending only on $(d,G,c,\varepsilon_0)$ and the chosen discretization of $E(t)$, such that the following holds.

\smallskip
\noindent
For every $\beta\ge\beta_1$ and every $L$ with
\[
u:=g^2_{\mathrm{GF}}(L)\ \le\ u_{\mathrm{AF}},
\]
the step-scaling function $\Sigma$ from \eqref{eq:stepscaling_def} obeys the two-sided estimate
\begin{equation}\label{eq:h44_stepscaling_two_sided}
u + b_0\,u^2\log 2 - C_{\Sigma}\,u^3
\ \le\ \Sigma(u)\ \le\
u + b_0\,u^2\log 2 + C_{\Sigma}\,u^3.
\end{equation}
Equivalently,
\[
\Sigma(u)=u+b_0u^2\log 2 + R_{\Sigma}(u),
\qquad |R_{\Sigma}(u)|\le C_{\Sigma}u^3.
\]
All constants are uniform in $L$ and in admissible boundary conditions (after tree-gauge reduction).
\end{lemma}

\begin{proof}
\textbf{Step 1: curvature cumulant control from Brascamp--Lieb.}
On the small-curvature event $\Omega_{\varepsilon_0}$, Proposition~\ref{prop:curvature_BL} gives the sub-Gaussian mgf bound
\eqref{eq:curvature_mgf} with variance proxy $\sigma_\beta^2$ from \eqref{eq:h44_sigma_def}.
By differentiation of the mgf at the origin and the standard Leonov--Shiryaev cumulant formula,
for every $k\ge 2$ and every finitely supported $2$--forms $J_1,\dots,J_k$,
\begin{equation}\label{eq:h44_cumulant_bound}
\bigl|\kappa_k(\langle J_1,dA\rangle,\dots,\langle J_k,dA\rangle\mid \Omega_{\varepsilon_0})\bigr|
\ \le\ (k-1)!\,(\sigma_\beta^2)^{k/2}\,\prod_{i=1}^k\|J_i\|_2,
\end{equation}
where $\kappa_k$ denotes the joint cumulant.

\smallskip
\textbf{Step 2: analytic curvature expansion of the flow observable.}
Fix $L$ and set $t=cL^2$.
In tree gauge, the Wilson flow map $A\mapsto V_t(A)$ and the clover energy density $E(t)$ are smooth gauge-covariant functionals of the initial
links.
Restricted to $\Omega_{\varepsilon_0}$ these functionals are real-analytic in the finitely many curvature variables
$\{F_p(A)\}_{p\in P(\Lambda(L))}$ with an analyticity radius depending only on $(G,\varepsilon_0)$.
Consequently $t^2E(t)$ admits an absolutely convergent expansion on $\Omega_{\varepsilon_0}$ into curvature polynomials:
\begin{equation}\label{eq:h44_E_poly}
t^2E(t)\;=\;\sum_{n\ge 1}\mathcal P_{2n}(dA),
\end{equation}
where each $\mathcal P_{2n}$ is a homogeneous polynomial of degree $2n$ in the curvature field $dA$.
Moreover, because the ratio $t/L^2=c$ is fixed, the coefficients of $\mathcal P_{2n}$ are bounded uniformly in $L$.

\smallskip
\textbf{Step 3: connected (cumulant) expansion and the $O(u^3)$ remainder.}
Take expectations in \eqref{eq:h44_E_poly} under the gauge-fixed Wilson measure.
Each polynomial expectation can be written as a finite sum of joint cumulants of linear curvature functionals.
The cumulant bound \eqref{eq:h44_cumulant_bound}, together with $\|\Pi_{\mathrm{ex}}\|_{2\to 2}=1$ (Lemma~\ref{lem:projection_identity}),
implies that for $u$ in a sufficiently small window $u\le u_{\mathrm{AF}}$ the series \eqref{eq:h44_E_poly} may be termwise
integrated and is dominated by a convergent majorant of the form $\sum_{n\ge 1}C_E^{n}u^{n}$ with $C_E=C_E(d,G,c,\varepsilon_0)$.
In particular, truncating \eqref{eq:h44_E_poly} at $n=1,2$ produces an error bounded by $C_{\Sigma}u^3$
uniformly in $L$.

\smallskip
\textbf{Step 4: extraction of the quadratic coefficient.}
The degree-$4$ term $\mathcal P_4$ in \eqref{eq:h44_E_poly} yields the quadratic contribution to step scaling.
The coefficient extraction is isolated below in Theorem~\ref{thm:b0_coefficient_extraction}: after linearizing the flowed field strength,
the connected degree-$4$ term is written as a positive Lie factor times the explicit squared-kernel shell integral
\eqref{eq:b0_shell_integral}.
In particular the coefficient multiplying $u^2$ at dyadic scale factor $2$ is $b_0\log 2$ with $b_0>0$.

\smallskip
\textbf{Step 5: removal of the small-curvature conditioning.}
Lemma~\ref{lem:weak_localization} gives $\mu_{\beta,\Lambda(L)}(\Omega_{\varepsilon_0}^c)\le |P(\Lambda(L))|e^{-c_{\mathrm{loc}}\beta}$.
On $\Omega_{\varepsilon_0}^c$ the Wilson density is suppressed by at least $\exp(-\beta v_0)$ for some $v_0=v_0(G,\varepsilon_0)>0$,
while $t^2E(t)$ grows at most polynomially in the number of plaquettes for fixed discretization.
Because $u\le u_{\mathrm{AF}}$ forces $\beta$ into the weak window, one may increase $\beta_1$ (equivalently shrink $u_{\mathrm{AF}}$)
so that the contribution of $\Omega_{\varepsilon_0}^c$ is absorbed into the $C_{\Sigma}u^3$ remainder bound.
\end{proof}
\subsubsection{Step scaling bound (closed in the exact form used later)}
\begin{proposition}[Weak-coupling step scaling with a two-sided cubic remainder]\label{prop:stepscaling_control_proved}
Assume Lemma~\ref{lem:flow_energy_expansion}. Then there exists $u_{\mathrm{AF}}>0$ and constants $b_0>0$, $C_\Sigma<\infty$ such that for all $u\in(0,u_{\mathrm{AF}}]$,
\begin{equation}\label{eq:stepscaling_bound_two_sided}
 u + b_0 u^2 \log 2 - C_\Sigma u^3
 \ \le\ \Sigma(u)\ \le\
 u + b_0 u^2 \log 2 + C_\Sigma u^3.
\end{equation}
After shrinking $u_{\mathrm{AF}}$ so that $C_\Sigma u_{\mathrm{AF}}\le \tfrac12 b_0\log 2$, one has the forward-increment bounds
\begin{equation}\label{eq:stepscaling_increment_window}
 \frac12 b_0 u^2\log 2
 \ \le\ \Sigma(u)-u\ \le\
 \frac32 b_0 u^2\log 2
 \qquad (0<u\le u_{\mathrm{AF}}).
\end{equation}
\end{proposition}

\begin{proof}
The two-sided estimate is exactly \eqref{eq:h44_stepscaling_two_sided} from Lemma~\ref{lem:flow_energy_expansion}.
Once $u_{\mathrm{AF}}$ is chosen so that $C_\Sigma u_{\mathrm{AF}}\le \tfrac12 b_0\log 2$, the cubic remainder is bounded by
$C_\Sigma u^3\le \tfrac12 b_0 u^2\log 2$ on $(0,u_{\mathrm{AF}}]$, which yields \eqref{eq:stepscaling_increment_window}.
\end{proof}

\begin{lemma}[Weak-window continuity of the step-scaling map]\label{lem:stepscaling_continuous}
On the weak window, the step-scaling map $\Sigma$ is continuous.
\end{lemma}

\begin{proof}
By Proposition~\ref{prop:small_polymer_neighborhood}, the exact dyadic blocking map is analytic in the extracted coordinates $(u,K)$ on the small polymer neighborhood.
Restricting that analytic map to the invariant graph $K=h(u)$ from Theorem~\ref{thm:rg_stable_manifold} gives an analytic scalar update
\[
F_h(u):=u'(u,h(u)).
\]
By the coordinate choice of Proposition~\ref{ass:coord_match}, this scalar update is exactly the weak-window step-scaling map $\Sigma$.
Hence $\Sigma$ is analytic, in particular continuous, on the weak window.
\end{proof}

\subsubsection{Archival fixed-$\beta$ threshold package (moved to supplement)}
The older fixed-$\beta$ coupling-to-KP threshold and the associated comparison of the first threshold scale $L_*(\beta)$ with $\sqrt{t_0(\beta)}$
are no longer part of the main Appendix~F proof spine.
That historical material has been moved to the supplementary Route~1 archive.
The load-bearing Route~1 chain now uses only the one-shot entrance theorem at the explicit scale $B(\beta)$ together with the tuned dyadic physical-scale corollary.


\subsection{Bounded macroscopic observables and upgrade to renormalized local fields}
\label{sec:what_remains_continuum}

This section implements the passage from bounded macroscopic observables to renormalized local fields in a way compatible with reflection positivity.
We first take a continuum limit on a separating algebra of bounded macroscopic observables $\mathcal O_{\ell}$ at fixed coarse scale $\ell$,
and then upgrade to renormalized point fields (operator-valued distributions) by a controlled small-flow-time/shrinking procedure based on the
gauge-covariant gradient flow.
The uniqueness/universality mechanism used below is contractivity of the exact RG in the KP basin together with the weak-window step-scaling control established earlier.

\subsubsection{A continuum limit on a bounded observable algebra}
\label{subsec:bounded_algebra_tightness}

The Wilson measure is reflection positive at finite lattice spacing, and bounded gauge-invariant observables automatically satisfy
uniform moment bounds. This suggests a clean two-stage construction:

\begin{enumerate}[leftmargin=*]
\item First construct the continuum limit on an algebra of \emph{bounded macroscopic observables} (no renormalization needed);
\item Then upgrade to renormalized local fields by a controlled shrinking/small-flow-time limit.
\end{enumerate}

Fix a physical coarse scale $\ell>0$. For lattice spacing $a$ and a physical point $x\in\mathbb R^4$, let $Q_{a,\ell}(x)$ denote
the lattice cube of side $\ell/a$ centered at the nearest lattice point to $x$.
Define the block-averaged plaquette energy density
\begin{equation}\label{eq:block_energy_def}
E_{a,\ell}(x)
:= \frac{1}{|\mathcal P(Q_{a,\ell}(x))|}\sum_{p\subset Q_{a,\ell}(x)}\Bigl(1-\frac{1}{\dim}\Re\,\mathrm{Tr}\, U_p\Bigr)\in[0,2].
\end{equation}
Let $\mathcal O_{\ell}$ be the algebra generated by finite products of the fields $E_{a,\ell}(x)$ at finitely many points.

\begin{lemma}[Polynomial moments determine weak limits on compact cubes]
\label{lem:poly_moments_determine_weak}
Let $K\subset\mathbb R^m$ be compact and let $(X_n)_{n\ge 1}$ be $K$--valued random vectors.
Assume that for every polynomial $P$ on $\mathbb R^m$ the limit $\lim_{n\to\infty}\mathbb E[P(X_n)]$ exists.
Then the laws of $X_n$ converge weakly to a unique probability measure on $K$.
\end{lemma}

\begin{proof}
Write $\mathcal P(K)$ for the algebra of restrictions to $K$ of polynomials on $\mathbb R^m$.
By the Stone--Weierstrass theorem, $\mathcal P(K)$ is uniformly dense in $C(K)$.
Define a linear functional $L$ on $\mathcal P(K)$ by $L(P):=\lim_{n\to\infty}\mathbb E[P(X_n)]$.
If $P\ge 0$ on $K$, then $\mathbb E[P(X_n)]\ge 0$ for all $n$, hence $L(P)\ge 0$.
Thus $L$ is positive and bounded on $\mathcal P(K)$, and extends uniquely by density to a positive bounded linear functional
(still denoted $L$) on $C(K)$.

By the Riesz representation theorem there exists a unique probability measure $\mu$ on $K$ such that
$L(f)=\int_K f\,d\mu$ for all $f\in C(K)$.
Finally, for any $f\in C(K)$ and $\varepsilon>0$, choose $P\in\mathcal P(K)$ with $\|f-P\|_{\infty}\le\varepsilon$.
Then for all $n$,
\[
\big|\mathbb E[f(X_n)]-L(f)\big|
\le \big|\mathbb E[(f-P)(X_n)]\big|+\big|\mathbb E[P(X_n)]-L(P)\big|+\big|L(P-f)\big|
\le 2\varepsilon+\big|\mathbb E[P(X_n)]-L(P)\big|.
\]
Letting $n\to\infty$ and then $\varepsilon\downarrow 0$ gives $\mathbb E[f(X_n)]\to \int f\,d\mu$.
This is exactly weak convergence of the laws of $X_n$ to $\mu$.
\end{proof}

\begin{lemma}[Reflection positivity is closed under pointwise convergence of moments]
\label{lem:rp_closed_under_limits}
Let $\langle\cdot\rangle_n$ be expectation functionals on a $*$--algebra $\mathcal A$ with a positive-time subalgebra $\mathcal A_+$
and a reflection $\Theta$.
Assume each $\langle\cdot\rangle_n$ is reflection positive: $\langle F\,\Theta F\rangle_n\ge 0$ for all $F\in\mathcal A_+$.
Assume moreover that for every finite family $F_1,\dots,F_k\in\mathcal A_+$, the limits
$\lim_{n\to\infty}\langle F_i\,\Theta F_j\rangle_n$ exist for all $i,j$.
Then the pointwise limit functional $\langle\cdot\rangle_{\infty}$ is reflection positive on $\mathcal A_+$.
\end{lemma}

\begin{proof}
Fix $F_1,\dots,F_k\in\mathcal A_+$ and define the $k\times k$ matrices
$M^{(n)}_{ij}:=\langle F_i\,\Theta F_j\rangle_n$.
Reflection positivity of $\langle\cdot\rangle_n$ implies $M^{(n)}\succeq 0$ for all $n$.
By assumption, $M^{(n)}\to M^{(\infty)}$ entrywise.
The cone of positive semidefinite matrices is closed under entrywise limits, hence $M^{(\infty)}\succeq 0$.
For any $c\in\mathbb C^k$ and $F:=\sum_i c_i F_i$ we have
\[
\langle F\,\Theta F\rangle_{\infty}=c^* M^{(\infty)} c\ge 0,
\]
which is reflection positivity on the span of $\{F_i\}$.
Since $k$ and $\{F_i\}$ were arbitrary, $\langle\cdot\rangle_{\infty}$ is reflection positive on $\mathcal A_+$.
\end{proof}

\begin{theorem}[Continuum limit on $\mathcal O_{\ell}$ without subsequences]
\label{thm:bounded_algebra_limit}
Fix $\ell>0$ and a scaling trajectory $a(\beta)\downarrow 0$ as $\beta\to\infty$ in the fixed scheme \eqref{eq:t0_def}.
Assume Theorem~\ref{thm:universality_on_Oell}.
Then for every $m\in\mathbb N$ and every $x_1,\dots,x_m\in\mathbb R^4$, the joint law of
\[
\bigl(E_{a(\beta),\ell}(x_1),\dots,E_{a(\beta),\ell}(x_m)\bigr)\in[0,2]^m
\]
converges weakly as $\beta\to\infty$ \emph{without} extracting subsequences.
The limiting Schwinger functional on $\mathcal O_{\ell}$ satisfies reflection positivity and inherits the lattice symmetries.
\end{theorem}

\begin{proof}
Fix $m$ and $x_1,\dots,x_m$.
The corresponding random vector is valued in the compact cube $K=[0,2]^m$.
By Theorem~\ref{thm:universality_on_Oell}, expectations of \emph{all} polynomial observables in these variables converge as $\beta\to\infty$.
Therefore Lemma~\ref{lem:poly_moments_determine_weak} yields weak convergence of the joint laws.

For reflection positivity, fix $F_1,\dots,F_k$ in the finitely generated positive-time algebra built from $\mathcal O_{\ell}$.
At each $\beta$, the Wilson measure satisfies reflection positivity, hence the matrices
$M^{(\beta)}_{ij}:=\langle F_i\,\Theta F_j\rangle_{\mu_{\beta}}$ are positive semidefinite.
Since each entry is an expectation of a bounded observable in $\mathcal O_{\ell}$, it has a limit as $\beta\to\infty$ by
Theorem~\ref{thm:universality_on_Oell}.
Lemma~\ref{lem:rp_closed_under_limits} now implies reflection positivity of the limit functional.
Lattice translation and hypercubic rotation symmetries are exact at each $\beta$ (in infinite volume or on the torus),
hence are inherited by the limit by the same pointwise-limit argument.
\end{proof}

\subsubsection{Upgrading to renormalized point fields via gradient flow}
\label{subsec:point_fields_flow}

The bounded macroscopic algebra $\mathcal O_{\ell}$ provides an OS-positive continuum limit on a separating bounded class and suffices to transfer a positive \emph{mass gap} once non-collapse is known.
To state the main theorem on a renormalized local field algebra (operator-valued distributions), we now upgrade from macroscopic observables to renormalized point fields using the gauge-covariant gradient flow.
The resulting dyadic Cauchy bounds are recorded in Proposition~\ref{prop:flow_pointfield_cauchy} and completed in Theorem~\ref{thm:field_algebra_truncation}.

\medskip
\noindent\textbf{Flow equation.}
We use the Wilson (gradient) flow on lattice links: given an initial configuration $U=(U_e)$, define flowed links $V_t(e)$ as the 
solution of
\begin{equation}\label{eq:wilson_flow}
 \partial_t V_t(e)= -\,\nabla_e S_W(V_t)\;V_t(e),\qquad V_0(e)=U_e,
\end{equation}
where $S_W$ is the Wilson action and $\nabla_e$ is the left-invariant gradient on the link variable.
The flow is gauge-covariant and smooth for all $t\ge 0$.

\medskip
\noindent\textbf{Energy density.}
Let $G_{t,\mu\nu}(x)$ be any standard local discretization of the flowed field strength (e.g. the clover definition) built from the 
flowed plaquettes. Define the flowed energy density
\begin{equation}\label{eq:flow_energy_def}
 E(t,x):=\frac{1}{2}\,\mathrm{Tr}\,G_{t,\mu\nu}(x)G_{t,\mu\nu}(x)\ \ge 0.
\end{equation}
For each fixed $t>0$, $E(t,x)$ is a bounded, gauge-invariant local functional of $V_t$ and hence a bounded functional of the initial 
links $U$.

\begin{definition}[Smearing at fixed physical scale]\label{def:flow_smearing}
Let $a=a(\beta)$ be the lattice spacing along the chosen scaling trajectory. For a test function $\phi\in C_c^{\infty}(\mathbb R^4)$
set
\begin{equation}\label{eq:smeared_energy_def}
 E_a(t,\phi)
 := a^4\sum_{x\in a\mathbb Z^4} \phi(x)\,E_a\bigl(t/a^2, x/a\bigr),
\end{equation}
where $E_a(\tau,\cdot)$ denotes the flowed energy density at lattice spacing $a$ and lattice flow-time $\tau$.
We also write the centered variable $\widetilde E_a(t,\phi):=E_a(t,\phi)-\langle E_a(t,\phi)\rangle$.
\end{definition}

\begin{remark}
For each fixed physical $t>0$, the smearing radius is $\sqrt t$. In lattice units the flow-time is $\tau=t/a^2\to\infty$ as 
$a\downarrow 0$, so $E_a(t,\phi)$ is automatically a bounded \emph{macroscopic} observable at scale comparable to $\sqrt t$.
Consequently, the tightness/OS-positivity argument of Theorem~\ref{thm:bounded_algebra_limit} applies verbatim to the family 
$\{E_a(t,\phi):\phi\in C_c^{\infty}\}$ for each fixed $t>0$.
\end{remark}

\paragraph*{Remaining UV ingredient: a uniform small-flow-time renormalization bound}

The nontrivial content is the $t\downarrow 0$ limit.
As in standard SFTE/OPE analyses, after subtracting the vacuum piece one expects the leading local operator to be the renormalized 
energy density $\mathrm{Tr}\,F_{\mu\nu}^2$ (dimension $4$), with higher-dimension operators suppressed by powers of $t$.
The correct Hilbert-space formulation is: the renormalized field exists as an \emph{operator-valued distribution}, meaning that for 
fixed test functions $\phi$ the smeared variables converge in $L^2$ after a multiplicative renormalization.

\begin{proposition}[Dyadic Cauchy estimate for a renormalized $\mathrm{Tr}\,F^2$ distribution]\label{prop:flow_pointfield_cauchy}
Assume the weak-coupling curvature control (Lemma~\ref{lem:weak_localization} and Proposition~\ref{prop:curvature_BL}) and the
step-scaling control of Proposition~\ref{prop:stepscaling_control_proved}.
Then there exist:
\begin{enumerate}[leftmargin=*]
\item a positive function $Z_E(t)$ on $(0,t_*)$ with $Z_E(t)=1+O(g^2(1/\sqrt t))$ as $t\downarrow 0$;
\item constants $\eta>0$ and, for each $\phi\in C_c^{\infty}(\mathbb R^4)$, a constant $C_\phi<\infty$;
\end{enumerate}
with the property that for dyadic times $t_j:=2^{-2j}t_0$ (with $t_0\le t_*$), the centered smeared variables
\begin{equation}\label{eq:ren_series_terms}
 X_j(\phi):= Z_E(t_j)\,\widetilde E(t_j,\phi)
\end{equation}
obey the summable increment bound
\begin{equation}\label{eq:dyadic_increment_bound}
 \|X_{j+1}(\phi)-X_j(\phi)\|_{L^2}\ \le\ C_\phi\,(1+j)^{-1-\eta}.
\end{equation}
Consequently $X_j(\phi)$ converges in $L^2$ as $j\to\infty$ to a limit, denoted
\begin{equation}\label{eq:F2ren_by_flow}
 \mathcal E_{\mathrm{ren}}(\phi)
 :=\lim_{t\downarrow 0} Z_E(t)\,\bigl(E(t,\phi)-\langle E(t,\phi)\rangle\bigr),
\end{equation}
which defines an OS field (operator-valued distribution) on the unique continuum limit on $\mathcal O_{\ell}$.
\end{proposition}

\paragraph*{One-shell estimate: explicit dyadic control}

The increment bound \eqref{eq:dyadic_increment_bound} has two inputs:
(i) a \emph{one-shell} $L^2$ estimate for the flowed-energy difference at scales $t_j$ and $t_{j+1}$,
and (ii) a weak-coupling bound showing that the running coupling along dyadic scales decays like $1/j$,
so that the shell contributions are summable.

\begin{lemma}[Two-sided weak-window step scaling]\label{lem:stepscaling_two_sided}
Assume Lemma~\ref{lem:flow_energy_expansion}. Then there exist $u_{\mathrm{AF}}>0$ and constants $b_0>0,C_\Sigma>0$ such that for all
$u\in(0,u_{\mathrm{AF}}]$,
\begin{equation}\label{eq:stepscaling_two_sided}
 u + b_0 u^2 \log 2 - C_\Sigma u^3 \ \le\ \Sigma(u)\ \le\ u + b_0 u^2 \log 2 + C_\Sigma u^3 .
\end{equation}
\end{lemma}

\begin{proof}
Lemma~\ref{lem:flow_energy_expansion} gives the one-loop expansion for the defining observable of the gradient-flow coupling with a
uniform remainder bound of order $O(u^3)$.
Substituting into the definition of $\Sigma$ yields
$\Sigma(u)=u+b_0u^2\log 2+R(u)$ with $|R(u)|\le C_\Sigma u^3$ for $u$ in the weak window, which is exactly
\eqref{eq:stepscaling_two_sided}.
\end{proof}

\begin{lemma}[Dyadic decay and square-summability of the weak-window coupling]\label{lem:dyadic_coupling_decay}
Fix $t_0>0$ and set $t_j:=2^{-2j}t_0$ and $L_j:=c^{-1/2}\sqrt{t_j}$ (so that $t_j=cL_j^2$ in the flow-coupling definition).
Let $u_j:=g_{\mathrm{GF}}^2(L_j)$, so that $u_{j-1}=\Sigma(u_j)$.
Assume \eqref{eq:stepscaling_two_sided} holds on $(0,u_{\mathrm{AF}}]$ and that $u_0\le u_{\mathrm{AF}}$.
Then there exists $c_*>0$ such that for all $j\ge 0$,
\begin{equation}\label{eq:uj_decay}
 u_j \ \le\ \frac{c_*}{1+j},
\qquad\text{and hence}\qquad
 \sum_{j\ge 0} u_j^2 < \infty .
\end{equation}
\end{lemma}

\begin{proof}
From $u_{j-1}=\Sigma(u_j)$ and \eqref{eq:stepscaling_two_sided} we have, for $u_j$ small,
\[
u_{j-1}-u_j \ \ge\ b_0 u_j^2\log 2 - C_\Sigma u_j^3.
\]
Choose $u_{\mathrm{AF}}$ smaller if needed so that $C_\Sigma u_{\mathrm{AF}}\le \tfrac12 b_0\log 2$.
Then for all $j$ with $u_j\le u_{\mathrm{AF}}$,
\[
u_{j-1}-u_j \ \ge\ \frac{b_0\log 2}{2}\,u_j^2.
\]
Divide by $u_j u_{j-1}$ and use $u_{j-1}\ge u_j$ to obtain
\[
\frac{1}{u_j}-\frac{1}{u_{j-1}}
=\frac{u_{j-1}-u_j}{u_j u_{j-1}}
\ \ge\ \frac{b_0\log 2}{2}.
\]
Summing from $1$ to $j$ yields $u_j^{-1}\ge u_0^{-1}+\tfrac{b_0\log 2}{2}j$, i.e. \eqref{eq:uj_decay} with
$c_*:=2/(b_0\log 2)$ up to adjusting for $u_0$.
Square-summability follows immediately from $u_j\lesssim (1+j)^{-1}$.
\end{proof}

\begin{lemma}[One-shell $L^2$ increment bound for the renormalized flowed energy]\label{lem:one_shell_increment}
Assume the curvature control (Lemma~\ref{lem:weak_localization} and Proposition~\ref{prop:curvature_BL}).
Let $t_{j+1}=t_j/4$ and $u_j=g_{\mathrm{GF}}^2(L_j)$ as in Lemma~\ref{lem:dyadic_coupling_decay}.
Then for each test function $\phi\in C_c^\infty(\mathbb R^4)$ there exist constants $C_\phi<\infty$ and $j_\phi$ such that for all
$j\ge j_\phi$ one can choose $Z_E(t_j)$ (with $Z_E(t)=1+O(u_j)$) so that
\begin{equation}\label{eq:one_shell_bound}
 \|X_{j+1}(\phi)-X_j(\phi)\|_{L^2}\ \le\ C_\phi\,u_j^2 \ +\ C_\phi\,e^{-c_{\mathrm{loc}}\beta}.
\end{equation}
\end{lemma}

\begin{proof}
Work in spanning-tree gauge and on the small-curvature event $\Omega_{\varepsilon_0}$.
On $\Omega_{\varepsilon_0}$ the Wilson action is uniformly strictly convex in the curvature sector (Section~\ref{subsec:curvature_domination}),
and Proposition~\ref{prop:curvature_BL} implies sub-Gaussian control for all linear curvature functionals with variance proxy
$\sigma^2\asymp u_j$ at physical scale $\sqrt{t_j}$.
The flowed field at times $t_j$ and $t_{j+1}$ depends only on the initial field in a neighborhood of radius $O(\sqrt{t_j})$,
and the map $U\mapsto E(t,\phi)$ is smooth on $\Omega_{\varepsilon_0}$ with derivatives bounded by constants depending on finitely many
heat-kernel norms.

Represent the difference as a single ``shell'' contribution using the semigroup property and a finite-range decomposition of the heat kernel
on 2-forms:
\begin{equation}\label{eq:finite_range_decomp}
 e^{-t\Delta_2}=\sum_{k\ge 0} K_{t,k},\qquad
 \mathrm{supp}\,K_{t,k}\subset\{(x,y):\dist(x,y)\le c_{\mathrm{fr}}2^k\sqrt t\},\qquad
 \|K_{t,k}\|_{\ell^1\to\ell^1}\le C_{\mathrm{fr}}2^{-Mk},
\end{equation}
for any fixed $M$ (standard on $\mathbb Z^4$).
Set $S_j:=e^{-t_{j+1}\Delta_2}-e^{-t_j\Delta_2}$ and decompose $S_j=\sum_{k\ge 0}S_{j,k}$ with $S_{j,k}:=K_{t_{j+1},k}-K_{t_j,k}$.

On $\Omega_{\varepsilon_0}$ one expands $E(t,\phi)$ in a local Taylor series in the curvature variables,
with coefficients given by the kernels $K_{t,k}$ and their discrete derivatives.
The $t$--increment $\widetilde E(t_{j+1},\phi)-\widetilde E(t_j,\phi)$ is therefore a sum of finitely-supported curvature polynomials
each containing at least one factor of the shell operator $S_j$.
By Brascamp--Lieb (Proposition~\ref{prop:curvature_BL}) and standard polynomial moment bounds for sub-Gaussian fields
(Hanson--Wright for quadratic terms, plus hypercontractive bounds for higher-degree remainders),
each such local contribution has $L^2$ norm bounded by $C_\phi\,u_j^{r/2}$ where $r\ge 2$ is its curvature degree.

The renormalization factor $Z_E(t)$ is chosen to cancel the \emph{order-$u_j$} part of the shell drift (the only logarithmically
accumulating contribution), leaving a remainder of order $u_j^2$.
Concretely, choose $Z_E(t_j)$ so that the shell expectation of the leading (dimension-$4$) curvature polynomial is annihilated,
equivalently: fix $Z_E$ by imposing a scale-independent normalization of one nondegenerate two-point function
$\langle X_j(\phi_*)X_j(\phi_*)\rangle$ for a fixed reference test function $\phi_*$.
The curvature cumulant bounds on $\Omega_{\varepsilon_0}$ imply that the adjustment $Z_E(t_{j+1})-Z_E(t_j)$ is $O(u_j)$ and that after this
choice the remaining shell contribution to $X_{j+1}-X_j$ is $O(u_j^2)$ in $L^2$.

The complement $\Omega_{\varepsilon_0}^c$ contributes at most $C_\phi\,\mu(\Omega_{\varepsilon_0}^c)^{1/2}$ to the $L^2$ norm,
which is bounded by $C_\phi e^{-c_{\mathrm{loc}}\beta}$ by Lemma~\ref{lem:weak_localization}.
Combining these bounds yields \eqref{eq:one_shell_bound}.
\end{proof}

\begin{proof}[Proof of Proposition~\ref{prop:flow_pointfield_cauchy}]
By Lemma~\ref{lem:one_shell_increment} and Lemma~\ref{lem:dyadic_coupling_decay},
\[
\|X_{j+1}(\phi)-X_j(\phi)\|_{L^2}\ \le\ C_\phi\,u_j^2 + C_\phi e^{-c_{\mathrm{loc}}\beta}
\ \le\ \widetilde C_\phi\,(1+j)^{-2}
\]
for all $j$ large (absorbing the exponentially small tail into $\widetilde C_\phi$).
This has the form \eqref{eq:dyadic_increment_bound} with $\eta=1$.
Summability implies $(X_j(\phi))_{j\ge 0}$ is Cauchy in $L^2$, hence converges to an $L^2$ limit.
\end{proof}

\begin{remark}[What is \emph{actually} proved vs. what is bookkeeping]
The analytic input is entirely local and uses only: (i) small-curvature localization; (ii) curvature Brascamp--Lieb domination; and 
(iii) a finite-range decomposition for the heat kernel (standard on $\mathbb Z^4$).
The renormalization factor $Z_E(t)$ is fixed by requiring that $\mathcal E_{\mathrm{ren}}$ have finite nonzero two-point function 
when paired with smooth test functions (equivalently: cancel the logarithmic divergence in the second moment).
The precise identification of $\mathcal E_{\mathrm{ren}}$ with $(\mathrm{Tr}\,F^2)_{\mathrm{ren}}$ is then a standard matching statement 
(using perturbation theory at small running coupling).
\end{remark}


\paragraph*{Field-algebra completion: a finite operator basis and mixing}

The single-field construction above extends to any \emph{finite} family of gauge-invariant local fields of a fixed engineering
dimension once one allows a \emph{matrix}-valued renormalization to handle operator mixing.
This is the standard constructive-QFT upgrade from one distinguished composite field to a finite operator basis.

Fix an integer $\Delta\ge 4$ and let $\{\mathcal O^{(\Delta)}_1,\dots,\mathcal O^{(\Delta)}_m\}$ be bounded gauge-invariant local lattice
observables supported in a fixed unit cube (e.g.\ a Symanzik basis of local gauge-invariant operators of dimension $\Delta$ modulo total
derivatives and lattice Bianchi relations).
Let $\mathbf{\mathcal O}(t,\phi):=(\mathcal O^{(\Delta)}_1(t,\phi),\dots,\mathcal O^{(\Delta)}_m(t,\phi))$ denote their flowed, smeared versions
defined exactly as in Definition~\ref{def:flow_smearing}, and let $\widetilde{\mathbf{\mathcal O}}(t,\phi)$ be the centered vector.

\begin{proposition}[Matrix dyadic Cauchy estimate for a renormalized operator multiplet]\label{prop:flow_operator_multiplet_cauchy}
Assume the weak-coupling curvature control (Lemma~\ref{lem:weak_localization} and Proposition~\ref{prop:curvature_BL})
and the step-scaling control of Proposition~\ref{prop:stepscaling_control_proved}.
Then there exist:
\begin{enumerate}[leftmargin=*]
\item an invertible matrix function $Z^{(\Delta)}(t)\in \mathrm{GL}_m(\mathbb R)$ on $(0,t_*)$ with $Z^{(\Delta)}(t)=I+O(g^2(1/\sqrt t))$ as $t\downarrow 0$;
\item constants $\eta>0$ and, for each $\phi\in C_c^{\infty}(\mathbb R^4)$, a constant $C_\phi<\infty$;
\end{enumerate}
such that for dyadic times $t_j:=2^{-2j}t_0$ (with $t_0\le t_*$), the $\mathbb R^m$--valued variables
\[
X^{(\Delta)}_j(\phi):=Z^{(\Delta)}(t_j)\,\widetilde{\mathbf{\mathcal O}}(t_j,\phi)
\]
obey the summable increment bound
\[
\|X^{(\Delta)}_{j+1}(\phi)-X^{(\Delta)}_j(\phi)\|_{L^2(\mathbb R^m)}\ \le\ C_\phi\,(1+j)^{-1-\eta}.
\]
Consequently $X^{(\Delta)}_j(\phi)$ converges in $L^2$ as $j\to\infty$, defining an $m$--component OS field multiplet
$\mathcal O^{(\Delta)}_{\mathrm{ren}}(\phi)$ on the unique continuum limit on $\mathcal O_{\ell}$.
\end{proposition}

\begin{remark}[How mixing enters]
The proof is identical to the scalar case (Proposition~\ref{prop:flow_pointfield_cauchy}).
One applies the one-shell $L^2$ estimate componentwise (curvature domination controls all local curvature polynomials uniformly),
and chooses the renormalization matrix $Z^{(\Delta)}(t_j)$ to cancel the \emph{linear} (order-$u_j$) shell drift in the multiplet,
equivalently by holding a fixed nondegenerate Gram matrix of two-point functions at a reference set of test functions constant across $j$.
After this matrix counterterm is imposed the residual shell increment is $O(u_j^2)$ in operator norm, hence summable.
\end{remark}

\subsubsection{Field-algebra completion across dimensions (inductive closure)}
\label{subsec:field_algebra_across_dimensions}

The remaining ingredient is to extend the finite-basis construction
(Proposition~\ref{prop:flow_operator_multiplet_cauchy}) to a \emph{generating} family of local gauge-invariant fields across
engineering dimensions, and to upgrade universality/uniqueness from bounded macroscopic observables to this unbounded local field algebra.

\medskip
\noindent
\textbf{Standard strategy (in standard terminology).}
Fix a graded local operator basis \emph{dimension by dimension}. At each fixed dimension $\Delta$, allow matrix-valued renormalization to
handle operator mixing inside that dimension and with lower dimensions (triangular structure by power counting).
Then prove a dyadic Cauchy estimate \emph{uniformly for the whole finite truncation} $\Delta\le \Delta_{\max}$.
Compatibility in $\Delta_{\max}$ yields an inductive limit defining the renormalized local field algebra.

\paragraph*{A graded operator basis}

\begin{definition}[Local gauge-invariant operator spaces]
\label{def:local_operator_spaces}
Fix an integer $\Delta\ge 0$.
Let $\mathscr V_{\Delta}$ be the real vector space spanned by gauge-invariant local lattice observables supported in a unit cube
that admit a small-curvature expansion with leading engineering dimension $\Delta$ (built from $F_{\mu\nu}$ and covariant derivatives)
and that are quotiented by lattice symmetries, total derivatives and lattice Bianchi identities.

Decompose $\mathscr V_{\Delta}$ into its exact-dimension symmetry blocks
\[
\mathscr V_{\Delta}=\bigoplus_{(\epsilon,\lambda)} \mathscr V_{\Delta,\epsilon,\lambda},
\]
where $\epsilon$ denotes CP parity and $\lambda$ the Euclidean tensor type.
Choose once and for all the canonical blockwise basis described in Appendix~\ref{app:normalization_crosscheck}; in particular, each block is written in a fixed complete-contraction basis and the blocks are ordered by increasing engineering dimension and then by a fixed symmetry-type order.
For $\Delta_{\max}\ge 0$ define the truncated operator multiplet
\[
\mathbf{\mathcal O}^{(\le \Delta_{\max})}
:=\bigl(\mathcal O^{(0)}_1,\dots,\mathcal O^{(0)}_{m_0};\ \mathcal O^{(1)}_1,\dots;\ \dots;\ \mathcal O^{(\Delta_{\max})}_1,\dots,\mathcal O^{(\Delta_{\max})}_{m_{\Delta_{\max}}}\bigr),
\]
an $M(\Delta_{\max})$--vector with $M(\Delta_{\max})=\sum_{\Delta\le \Delta_{\max}} m_{\Delta}$.
\end{definition}

\begin{remark}[Why a finite basis exists]
For each fixed $\Delta$ there are only finitely many gauge-invariant monomials of dimension $\Delta$ modulo symmetries and identities.
On the lattice, one may take a Symanzik-type basis built from small loops and covariant difference operators.
Appendix~\ref{app:normalization_crosscheck} records the canonical blockwise choice used later in Lane~A.
The project uses only that each exact-dimension symmetry block is finite dimensional and that every canonical basis element becomes a bounded local observable after positive flow time and smearing.
\end{remark}

\paragraph*{Triangular mixing and coherent renormalization}

The gradient-flow smearing at flow time $t>0$ makes each $\mathcal O^{(\Delta)}_i(t,\phi)$ bounded.
As $t\downarrow 0$ the family develops UV divergences.
Power counting implies that at fixed $\Delta$ the flowed observable admits a small-flow-time expansion into local fields of dimensions
$\le \Delta$; equivalently, in the natural ordering by increasing engineering dimension, the renormalization matrix in the graded basis
is \emph{block lower triangular}.

\begin{definition}[Graded triangular renormalization matrices]
\label{def:graded_Z}
Fix $\Delta_{\max}\ge 0$.
A renormalization matrix $Z^{(\le \Delta_{\max})}(t)\in\mathrm{GL}_{M(\Delta_{\max})}(\mathbb R)$ is called \emph{graded triangular}
if, when written in blocks by engineering dimension,
\[
Z^{(\le \Delta_{\max})}(t)=\bigl(Z_{\Delta\Delta'}(t)\bigr)_{0\le \Delta,\Delta'\le \Delta_{\max}},
\]
it satisfies $Z_{\Delta\Delta'}(t)=0$ whenever $\Delta'>\Delta$.
\end{definition}

\begin{remark}[Normalization to fix $Z$]
For each $\Delta_{\max}$ the graded triangular matrix is fixed by imposing a finite set of scale-independent normalization conditions
on a nondegenerate Gram matrix of two-point functions for a reference set of test functions.
This is the same mechanism used in Proposition~\ref{prop:flow_operator_multiplet_cauchy}, but now enforced simultaneously
across all dimensions up to $\Delta_{\max}$.
\end{remark}

\paragraph*{Uniform dyadic Cauchy estimates for finite truncations}

\begin{theorem}[Dyadic construction of the renormalized local field algebra up to $\Delta_{\max}$]
\label{thm:field_algebra_truncation}
Assume:
\begin{enumerate}[leftmargin=*]
\item weak-coupling curvature control (Lemma~\ref{lem:weak_localization} and Proposition~\ref{prop:curvature_BL});
\item the theorem-level weak-window package already proved (Proposition~\ref{prop:stepscaling_control_proved}, Lemma~\ref{lem:stepscaling_continuous}, Theorem~\ref{thm:h49_bare_flow_window}, and Proposition~\ref{prop:h59_uv_tuning_exists}).
\end{enumerate}
Fix $\Delta_{\max}\ge 0$ and use the canonical blockwise basis of Definition~\ref{def:local_operator_spaces} and Appendix~\ref{app:normalization_crosscheck}.
Then there exist:
\begin{enumerate}[leftmargin=*]
\item a graded triangular matrix function $Z^{(\le \Delta_{\max})}(t)$ on $(0,t_*)$;
\item for each test function $\phi\in C_c^{\infty}(\mathbb R^4)$ constants $C_{\phi,\Delta_{\max}}<\infty$ and $\eta>0$;
\end{enumerate}
such that for dyadic times $t_j=2^{-2j}t_0$ (with $t_0\le t_*$) the renormalized multiplet
\[
X^{(\le \Delta_{\max})}_j(\phi)
\ :=\ Z^{(\le \Delta_{\max})}(t_j)\,\widetilde{\mathbf{\mathcal O}}^{(\le \Delta_{\max})}(t_j,\phi)
\]
obeys the summable increment bound
\[
\|X^{(\le \Delta_{\max})}_{j+1}(\phi)-X^{(\le \Delta_{\max})}_j(\phi)\|_{L^2(\mathbb R^{M(\Delta_{\max})})}
\ \le\ C_{\phi,\Delta_{\max}}\,(1+j)^{-1-\eta}.
\]
Consequently $X^{(\le \Delta_{\max})}_j(\phi)$ converges in $L^2$ as $j\to\infty$ for each $\phi$, defining an
OS field multiplet $\mathbf{\mathcal O}^{(\le \Delta_{\max})}_{\mathrm{ren}}(\phi)$.

Moreover, if $\mathcal A^{(\le \Delta_{\max})}_{+,\mathrm{flow}}$ denotes the bounded positive-time algebra generated by the
canonical flowed lifts in the cap and
\[
\mathcal D^{(\le \Delta_{\max})}_{\mathrm{flow}}
:=\mathrm{span}\{[F]:F\in \mathcal A^{(\le \Delta_{\max})}_{+,\mathrm{flow}}\},
\]
then each smeared component $\mathcal O^{(\le \Delta_{\max})}_{\mathrm{ren},i}(\phi)$ is obtained on
$\mathcal D^{(\le \Delta_{\max})}_{\mathrm{flow}}$ as the graph-norm limit of bounded dyadic approximants, and the
polynomial $*$--algebra generated by the components is well defined on this explicit common dense invariant core
(see Lemma~\ref{lem:finite_cap_common_core}).
\end{theorem}

\begin{proof}
The proof is the finite-cap version of Proposition~\ref{prop:flow_operator_multiplet_cauchy}, with the earlier basis-by-basis placeholder replaced by the canonical block transport of Corollary~\ref{cor:graded_triangular_transport_all_blocks}.

\smallskip
\noindent\emph{(i) Whole-cap canonical one-shell transport.}
Corollary~\ref{cor:graded_triangular_transport_all_blocks} provides a block-lower-triangular one-shell update
\[
\mathbf Y^{(\le \Delta_{\max})}_{j+1}(\phi)
=
\bigl(I+u_jA^{(\le \Delta_{\max})}_j\bigr)\mathbf Y^{(\le \Delta_{\max})}_j(\phi)
+\mathbf Q^{(\le \Delta_{\max})}_j(\phi),
\]
with $\|\mathbf Q^{(\le \Delta_{\max})}_j(\phi)\|_{L^2}\le C_{\phi,\Delta_{\max}}u_j^2$.
This is already basis-free at the block level: Theorem~\ref{thm:LaneA_associated_graded_identity} identifies the exact-dimension quotient intrinsically before any basis is chosen, and the lower-block terms are then forced by power counting in the grading order.

\smallskip
\noindent\emph{(ii) Graded triangular renormalization.}
Impose the same scale-independent normalization moments used in Proposition~\ref{prop:flow_operator_multiplet_cauchy}, but now recursively in the grading order of the canonical blocks.
Because the one-shell transport is block lower triangular, each diagonal block is normalized after the lower blocks have already been fixed, producing a graded triangular matrix $Z^{(\le\Delta_{\max})}(t_j)$.
The order-$u_j$ shell drift is removed block-by-block, leaving only the $O(u_j^2)$ remainder from the canonical one-shell theorem.
Compatibility under cap enlargement is ensured by Lemma~\ref{lem:LaneA_block_cap_compat} and later recorded abstractly in Lemma~\ref{lem:coherent_Z}.

\smallskip
\noindent\emph{(iii) Summability.}
The step-scaling decay (Lemma~\ref{lem:dyadic_coupling_decay}) gives $\sum_j u_j^2<\infty$ in the AF window.
Hence the renormalized shell increments are summable in $L^2$, yielding the stated dyadic Cauchy property.

\smallskip
\noindent\emph{(iv) Algebra/domain statement.}
The explicit common-core argument is isolated in Lemma~\ref{lem:finite_cap_common_core}.
In brief, the dyadic approximants $X^{(\le \Delta_{\max})}_j(\phi)$ are bounded flowed observables and therefore preserve the dense span of bounded flowed OS-classes.
Theorem~\ref{thm:ins_to_OS} converts the finite-cap insertion estimates to OS-norm bounds against any fixed bounded flowed vector,
and the summable increment bound from part~(iii) then gives Cauchy convergence in the graph norms relevant to any fixed finite field polynomial.
This yields a common dense invariant core for all finite-cap field polynomials.
\end{proof}

\begin{lemma}[Explicit finite-cap common invariant core]\label{lem:finite_cap_common_core}
Fix $\Delta_{\max}\ge 0$ and let $\mathcal A^{(\le \Delta_{\max})}_{+,\mathrm{flow}}$ and
$\mathcal D^{(\le \Delta_{\max})}_{\mathrm{flow}}$ be as in Theorem~\ref{thm:field_algebra_truncation}.
Then:
\begin{enumerate}[label=(\roman*),leftmargin=2.3em]
\item $\mathcal D^{(\le \Delta_{\max})}_{\mathrm{flow}}$ is dense in the finite-cap OS Hilbert space and is invariant under every bounded dyadic approximant $X^{(\le \Delta_{\max})}_j(\phi)$.
\item For each canonical field component $\mathcal O^{(\le \Delta_{\max})}_{\mathrm{ren},i}(\phi)$ and each $F\in \mathcal A^{(\le \Delta_{\max})}_{+,\mathrm{flow}}$, the classes
\[
\bigl[ X^{(\le \Delta_{\max})}_{j,i}(\phi)F \bigr]
\]
form a Cauchy sequence in the OS Hilbert norm as $j\to\infty$; moreover the same estimate remains valid after adjoining any fixed finite family of additional field insertions, so the convergence is Cauchy in every finite graph norm relevant to the polynomial field algebra.
\item Consequently each $\mathcal O^{(\le \Delta_{\max})}_{\mathrm{ren},i}(\phi)$ is closable, $\mathcal D^{(\le \Delta_{\max})}_{\mathrm{flow}}$ is a common invariant core for the closures, and every finite polynomial in the smeared fields is well defined on that core.
\end{enumerate}
\end{lemma}

\begin{proof}
Density is by construction: before the $j\to\infty$ limit is taken, the finite-cap OS Hilbert space is generated from bounded flowed positive-time observables.
Each $X^{(\le \Delta_{\max})}_j(\phi)$ is a linear combination of bounded flowed observables at positive flow time, so multiplying by an element of $\mathcal A^{(\le \Delta_{\max})}_{+,\mathrm{flow}}$ stays inside the same algebra; this proves invariance.

Now fix $F\in \mathcal A^{(\le \Delta_{\max})}_{+,\mathrm{flow}}$.
Theorem~\ref{thm:ins_to_OS} upgrades the finite-cap insertion estimates to OS norms against the fixed bounded flowed insertion family determined by $F$.
Applying that upgrade to the dyadic increment bound of Theorem~\ref{thm:field_algebra_truncation} gives, for each component index $i$,
\[
\bigl\|\bigl[(X^{(\le \Delta_{\max})}_{j+1,i}(\phi)-X^{(\le \Delta_{\max})}_{j,i}(\phi))F\bigr]\bigr\|
\le C_{F,\phi,\Delta_{\max}}(1+j)^{-1-\eta}.
\]
The same estimate holds after adjoining any fixed finite family of additional dyadic approximants on the left, because those extra factors are again bounded flowed observables and are therefore covered by the same OS-seminorm control.
Summability in $j$ yields the asserted Cauchy property in the Hilbert norm and in every finite graph norm used by a fixed field polynomial.
Hence the graph-norm limits defining the renormalized fields exist on $\mathcal D^{(\le \Delta_{\max})}_{\mathrm{flow}}$, preserve that core, and are closable there.
\end{proof}

\begin{corollary}[Canonical finite-truncation SFTE and intrinsic remainder control]\label{cor:finite_truncation_SFTE}
Fix $\Delta_{\max}\ge 0$ and the canonical blockwise basis of Definition~\ref{def:local_operator_spaces}.
Then the canonical flowed lifts and the renormalized local fields produced by Theorem~\ref{thm:field_algebra_truncation} satisfy, on that finite truncation, the small-flow-time expansion of Lane~A with a block-lower-triangular coefficient matrix
\[
\mathbf Y_t^{(\le \Delta_{\max})}(f)
=
C^{(\le \Delta_{\max})}(t;\mu)\,\mathbf O_{\mathrm{ren}}^{(\le \Delta_{\max})}(f;\mu)
+\mathbf R_t^{(\le \Delta_{\max})}(f;\mu).
\]
Let $\mathcal H^{(\le \Delta_{\max})}_{\mathrm{cap}}$ be the finite-cap OS Hilbert space underlying Theorem~\ref{thm:field_algebra_truncation}, let $\Omega_{\mathrm{cap}}:=[1]$ be its vacuum vector, let
$\mathcal D^{(\le \Delta_{\max})}_{\mathrm{flow}}$ be the common invariant core of Lemma~\ref{lem:finite_cap_common_core}, and write
\[
\mathcal B^{(\le \Delta_{\max})}_{+,\mathrm{cap}}
:=
\ast\!\bigl(\mathcal A^{(\le \Delta_{\max})}_{+,\mathrm{flow}},\mathcal A_+\bigr)
\]
for the bounded positive-time capwise algebra generated by the capwise flowed lifts and the bounded positive-time Wilson observables.
Then there exists $\eta_{\Delta_{\max}}>0$ such that for every component index $i$, every fixed finite insertion family $\{\mathcal X_\ell\}$ from the same truncation, and every
$\Psi,\Phi\in\mathcal D^{(\le \Delta_{\max})}_{\mathrm{flow}}$,
\[
\Big|
\big\langle
\Psi,
\big(\mathcal R^{(\le \Delta_{\max})}_i(t;f;\mu)-r^{(\le \Delta_{\max})}_i(t;f;\mu)\big)
\prod_{\ell=1}^k \mathcal X_\ell\,\Phi
\big\rangle_{\mathrm{cap}}
\Big|
\le C_{\Delta_{\max},f,\Psi,\Phi,\mathcal X,\mu}\,t^{\eta_{\Delta_{\max}}}
\]
for all sufficiently small $t$, where
\[
r^{(\le \Delta_{\max})}_i(t;f;\mu)
:=
\big\langle \Omega_{\mathrm{cap}},\mathcal R^{(\le \Delta_{\max})}_i(t;f;\mu)\,\Omega_{\mathrm{cap}}\big\rangle_{\mathrm{cap}}.
\]
In particular, taking $\Psi=\Phi=\Omega_{\mathrm{cap}}$ yields the capwise vacuum remainder bound, and the remainder/inverse compatibility condition
\eqref{eq:LaneA_inverse_remainder_compat} holds on every finite truncation against any polylogarithmic inverse bound.
Equivalently, for every fixed capwise vacuum channel, the centered inverse-flow approximants are Cauchy as $t\downarrow 0$.

Moreover, for every fixed positive-time capwise word
\[
\begin{aligned}
W_+
={}&
B_0\,\widehat O_{j_1}^{(\le \Delta_{\max})}(f_1;\mu)\,B_1\cdots
\widehat O_{j_m}^{(\le \Delta_{\max})}(f_m;\mu)\,B_m,\\
& B_r\in\mathcal B^{(\le \Delta_{\max})}_{+,\mathrm{cap}},
\qquad
\supp(f_r)\subset\{x_0>0\},
\end{aligned}
\]
let
\[
\begin{aligned}
W_+(\mathbf t)
:={}&
B_0\,\mathcal O_{j_1}^{\mathrm{app},(\le \Delta_{\max})}(t_1;f_1;\mu)\,B_1\cdots\\
&\mathcal O_{j_m}^{\mathrm{app},(\le \Delta_{\max})}(t_m;f_m;\mu)\,B_m.
\end{aligned}
\]
Then the compatible finite-observable-set limits packaged one finite family at a time by Corollary~\ref{cor:finite_cap_mixed_family_packaging}, and later abbreviated locally in Theorem~\ref{thm:finite_cap_flowed_sector_bridge} by
$\omega_{\mathrm{cap,bnd}}^{(u_{\mathrm{ren}})}$, satisfy
\[
\omega_{\mathrm{cap,bnd}}^{(u_{\mathrm{ren}})}\big(W_+(\mathbf t)-W_+(\mathbf s)\big)
\xrightarrow[\mathbf t,\mathbf s\downarrow 0]{}0.
\]
More precisely, Lemma~\ref{lem:finite_cap_mixed_channel_reduction} expands each literal telescoping summand of
$W_+(\mathbf t)-W_+(\mathbf s)$ into finitely many bounded-word remainder channels with coefficient envelopes governed by
Theorem~\ref{thm:finite_truncation_inverse_control}; the resulting coefficient envelope is the polylogarithmic one displayed in Lemma~\ref{lem:finite_cap_mixed_channel_reduction}, and Corollary~\ref{cor:finite_truncation_SFTE} supplies the finite-cap remainder / inverse compatibility used there.
Hence the mixed-channel limit exists for every such $W_+$ and is independent of the order or relative rates with which the coordinates of $\mathbf t$ tend to $0$.
\end{corollary}

\begin{proof}
The coefficient matrix is the finite-cap version of the extracted coefficient matrix from Theorem~\ref{thm:ins_decomp}, written in the canonical block basis and then normalized by
Theorem~\ref{thm:field_algebra_truncation}; lower-block mixing is exactly the graded triangular structure already built into both results.
The extracted-irrelevant remainder obeys the power bound \eqref{eq:rem_tpower} in insertion seminorm, and Theorem~\ref{thm:ins_to_OS} upgrades that control to OS matrix elements against any fixed bounded
flowed insertion family.
Lemma~\ref{lem:finite_cap_common_core} identifies $\mathcal D^{(\le \Delta_{\max})}_{\mathrm{flow}}$ as the graph-norm closure of such bounded flowed vectors, so the same estimate extends to every pair
$\Psi,\Phi\in\mathcal D^{(\le \Delta_{\max})}_{\mathrm{flow}}$ after adjoining any fixed finite capwise insertion family on the right.
Because the truncation is finite, one may choose a single exponent $\eta_{\Delta_{\max}}$ and absorb the finitely many vectors and insertions into the constant.
Any polylogarithmic inverse factor is dominated by $t^{-\varepsilon}$ for arbitrarily small $\varepsilon>0$, so the $t^{\eta_{\Delta_{\max}}}$ remainder is compatible with the inverse bounds used later in Lane~A.
Finally, if $t,t'\le t_0$, subtracting the SFTE identities at scales $t$ and $t'$ and using cancellation of the common renormalized local term shows that the difference of the centered inverse-flow approximants is a linear combination of the centered remainders with coefficients controlled by
Theorem~\ref{thm:finite_truncation_inverse_control}; the displayed bound therefore gives the announced Cauchy property in every fixed capwise vacuum channel.

For the additional positive-time mixed-channel clause, fix $W_+$ and two scale vectors $\mathbf t,\mathbf s$, and write
\[
W_+(\mathbf t)-W_+(\mathbf s)=\sum_{r=1}^m T_r(\mathbf t,\mathbf s)
\]
for the multilinear telescoping decomposition.
The literal telescoping summands $T_r(\mathbf t,\mathbf s)$ are not fixed-word channels: every non-difference local slot is itself an inverse-flow approximant at one of the scales from $\mathbf t$ or $\mathbf s$.
Lemma~\ref{lem:finite_cap_mixed_channel_reduction} isolates the required reduction.
Its displayed bound \eqref{eq:finite_cap_literal_summand_bound} is the literal telescoping-summand estimate: \eqref{eq:finite_cap_literal_summand_coeff_env} supplies the coefficient envelope coming from Theorem~\ref{thm:finite_truncation_inverse_control}, and Proposition~\ref{prop:finite_cap_bounded_factor_insertion}, obtained from Corollary~\ref{cor:ins_to_OS_bounded_word} and hence from Theorem~\ref{thm:ins_to_OS}, supplies the factor $t_r^{\eta_{\Delta_{\max}}}+s_r^{\eta_{\Delta_{\max}}}$ after the $r$th slot is converted into a centered remainder difference.
Equation~\eqref{eq:finite_cap_literal_summand_o1} is the corresponding slotwise $o(1)$ handoff furnished by the finite-cap remainder / inverse compatibility clause of Corollary~\ref{cor:finite_truncation_SFTE} applied to that polylogarithmic envelope.
Hence each
$\omega_{\mathrm{cap,bnd}}^{(u_{\mathrm{ren}})}\bigl(T_r(\mathbf t,\mathbf s)\bigr)$ tends to $0$ as $\mathbf t,\mathbf s\downarrow 0$.
Summing over the finitely many slots $r$ gives the announced mixed-channel Cauchy property and therefore the order/rate-independent limit, with the single-slot and telescoping inputs supplied by Proposition~\ref{prop:finite_cap_bounded_factor_insertion} and Lemma~\ref{lem:finite_cap_mixed_channel_reduction}.
\end{proof}

\begin{propFspecial}[State-free positive-time bounded-factor insertion estimate]\label{prop:finite_cap_bounded_factor_insertion}
Work in the setting and notation of Corollary~\ref{cor:finite_truncation_SFTE}.
Fix a component index $i$, a test function $f$ with $\supp(f)\subset\{x_0>0\}$, bounded factors
\[
B_0,\dots,B_q\in \mathcal B^{(\le \Delta_{\max})}_{+,\mathrm{cap}},
\]
and a compact positive-time neighborhood $K_{B,f}\subset\{x_0>0\}$ containing $\supp(f)$ together with the supports of every bounded cylinder / flowed factor appearing in the fixed word $B_0,\dots,B_q$.
Then for every $\Psi,\Phi\in \mathcal D^{(\le \Delta_{\max})}_{\mathrm{flow}}$ one has
\[
\begin{aligned}
&\Big|
\big\langle
\Psi,\,
B_0\bigl(\mathcal R_i^{(\le \Delta_{\max})}(t;f;\mu)-r_i^{(\le \Delta_{\max})}(t;f;\mu)\bigr)
B_1\cdots B_q\,\Phi
\big\rangle_{\mathrm{cap}}
\Big|\\
&\qquad\le
C_{\Psi,\Phi,K_{B,f},i,\Delta_{\max},\mu}
\Bigl(\prod_{r=0}^q \|B_r\|_{\infty}\Bigr)
 t^{\eta_{\Delta_{\max}}}.
\end{aligned}
\]
for all sufficiently small $t$.
After subtracting the centered finite-cap SFTE identities at scales $t,s$, the same fixed bounded word satisfies
\[
\begin{aligned}
&\Big|
\big\langle
\Psi,\,
B_0\bigl(
\mathcal O_i^{\mathrm{app},(\le \Delta_{\max})}(t;f;\mu)
-\mathcal O_i^{\mathrm{app},(\le \Delta_{\max})}(s;f;\mu)
\bigr)
B_1\cdots B_q\,\Phi
\big\rangle_{\mathrm{cap}}
\Big|\\
&\qquad\le
C'_{\Psi,\Phi,K_{B,f},i,\Delta_{\max},\mu}
\Bigl(\prod_{r=0}^q \|B_r\|_{\infty}\Bigr)
\bigl(t^{\eta_{\Delta_{\max}}}+s^{\eta_{\Delta_{\max}}}\bigr).
\end{aligned}
\]
for all sufficiently small $t,s$.
In particular, taking $\Psi=\Phi=\Omega_{\mathrm{cap}}$ gives the vacuum-channel estimate later specialized in the proof of Theorem~\ref{thm:finite_cap_flowed_sector_bridge}(i).
\end{propFspecial}

\begin{proof}
Set
\[
F_t:=\mathcal R_i^{(\le \Delta_{\max})}(t;f;\mu)-r_i^{(\le \Delta_{\max})}(t;f;\mu)
\]
and let $J_{i,t}$ denote the matching-scale insertion activity whose summed observable is $F_t$.
By the capwise remainder bound \eqref{eq:rem_tpower} at the matching scale,
\begin{equation}\label{eq:finite_cap_word_remainder_ins}
\|J_{i,t}\|_{\mathrm{ins};\mu;k,K_f,\varepsilon}
\le
C_{i,\Delta_{\max},\mu}\,t^{\eta_{\Delta_{\max}}},
\end{equation}
where $K_f\subset K_{B,f}$ is the compact neighborhood attached to the test function support.
Fix the bounded positive-time word
\[
\mathbf W:=(B_0,\dots,B_q)
\]
and let $J_{i,t}^{\mathbf W}$ be the modified insertion activity from Corollary~\ref{cor:ins_to_OS_bounded_word}.
Its summed observable is
\[
F_{J_{i,t}^{\mathbf W}}=B_0F_tB_1\cdots B_q.
\]
Applying Corollary~\ref{cor:ins_to_OS_bounded_word} at the finite-observable-set level to the fixed bounded word $\mathbf W$ --- and only to that positive-time word --- then passing through the compatible capwise limit gives
\[
\begin{aligned}
&\big\langle [F_{J_{i,t}^{\mathbf W}}],[F_{J_{i,t}^{\mathbf W}}]\big\rangle_{\mathrm{cap}}^{1/2}\\
&\qquad\le
C_{\mathrm{OS}\leftarrow\mathrm{ins}}
C_{\mathrm{supp}}(K_{B,f})
C_{\mathrm{word}}(K_{B,f};\mathbf W^{\sharp})
\Bigl(\prod_{r=0}^q \|B_r\|_{\infty}\Bigr)
\|J_{i,t}\|_{\mathrm{ins};\mu;k,K_f,\varepsilon}.
\end{aligned}
\]
Using \eqref{eq:finite_cap_word_remainder_ins},
\begin{equation}\label{eq:finite_cap_word_remainder_os}
\big\langle [F_{J_{i,t}^{\mathbf W}}],[F_{J_{i,t}^{\mathbf W}}]\big\rangle_{\mathrm{cap}}^{1/2}
\le
C_{K_{B,f},i,\Delta_{\max},\mu}
\Bigl(\prod_{r=0}^q \|B_r\|_{\infty}\Bigr)
t^{\eta_{\Delta_{\max}}}
\end{equation}
for all sufficiently small $t$.

It now suffices to treat core vectors of the form $\Psi=[H]$ and $\Phi=[G]$ with
$H,G\in\mathcal A^{(\le \Delta_{\max})}_{+,\mathrm{flow}}$, since these span
$\mathcal D^{(\le \Delta_{\max})}_{\mathrm{flow}}$.
Set
\[
A_t^{\mathbf W}:=F_{J_{i,t}^{\mathbf W}}=B_0F_tB_1\cdots B_q.
\]
By the finite-cap Cauchy--Schwarz inequality,
\[
\Big|
\big\langle [H],[A_t^{\mathbf W}G]\big\rangle_{\mathrm{cap}}
\Big|
\le
\big\langle [H],[H]\big\rangle_{\mathrm{cap}}^{1/2}
\big\langle [A_t^{\mathbf W}G],[A_t^{\mathbf W}G]\big\rangle_{\mathrm{cap}}^{1/2}.
\]
Because $G$ is a bounded positive-time observable, the reflected product carries an extra factor $|G|^2\le \|G\|_\infty^2$ pointwise, so
\[
\big\langle [A_t^{\mathbf W}G],[A_t^{\mathbf W}G]\big\rangle_{\mathrm{cap}}^{1/2}
\le
\|G\|_\infty\,
\big\langle [A_t^{\mathbf W}],[A_t^{\mathbf W}]\big\rangle_{\mathrm{cap}}^{1/2}.
\]
Combining this with \eqref{eq:finite_cap_word_remainder_os} yields
\[
\Big|
\big\langle [H],B_0F_tB_1\cdots B_q[G]\big\rangle_{\mathrm{cap}}
\Big|
\le
C_{H,G,K_{B,f},i,\Delta_{\max},\mu}
\Bigl(\prod_{r=0}^q \|B_r\|_{\infty}\Bigr)
 t^{\eta_{\Delta_{\max}}}
\]
for all sufficiently small $t$.
This is the only place where the dependence on $H$ and $G$ enters; Corollary~\ref{cor:ins_to_OS_bounded_word} has been applied solely to the fixed bounded positive-time word $\mathbf W$.
Linearity now extends the same estimate to arbitrary
$\Psi,\Phi\in\mathcal D^{(\le \Delta_{\max})}_{\mathrm{flow}}$, proving the first displayed bound.

For the second estimate, subtract the centered finite-cap SFTE identities at scales $t$ and $s$.
The common centered local symbol $\widehat O_i^{(\le \Delta_{\max})}(f;\mu)$ cancels, and the resulting difference of inverse-flow approximants is a finite linear combination of centered remainders with coefficients controlled by Theorem~\ref{thm:finite_truncation_inverse_control}.
Applying the first displayed estimate term-by-term to that linear combination and summing over the finitely many capwise components yields the announced bound.
\end{proof}

\begin{lemFspecial}[Actual mixed-channel telescoping reduction on a finite truncation]\label{lem:finite_cap_mixed_channel_reduction}
Work in the setting of Corollary~\ref{cor:finite_truncation_SFTE} and fix a positive-time capwise word
\[
\begin{aligned}
W_+
={}&
B_0\,\widehat O_{j_1}^{(\le \Delta_{\max})}(f_1;\mu)\,B_1\cdots
\widehat O_{j_m}^{(\le \Delta_{\max})}(f_m;\mu)\,B_m,\\
& B_r\in\mathcal B^{(\le \Delta_{\max})}_{+,\mathrm{cap}},
\qquad
\supp(f_r)\subset\{x_0>0\}.
\end{aligned}
\]
For $1\le r\le m$, let $T_r(\mathbf t,\mathbf s)$ denote the $r$th multilinear telescoping summand of $W_+(\mathbf t)-W_+(\mathbf s)$, namely
\[
\begin{aligned}
T_r(\mathbf t,\mathbf s)
:={}&
B_0\,\mathcal O_{j_1}^{\mathrm{app},(\le \Delta_{\max})}(t_1;f_1;\mu)\,B_1\cdots
\mathcal O_{j_{r-1}}^{\mathrm{app},(\le \Delta_{\max})}(t_{r-1};f_{r-1};\mu)\,B_{r-1}\\
&\cdot
\bigl(
\mathcal O_{j_r}^{\mathrm{app},(\le \Delta_{\max})}(t_r;f_r;\mu)
-
\mathcal O_{j_r}^{\mathrm{app},(\le \Delta_{\max})}(s_r;f_r;\mu)
\bigr)
B_r\,\mathcal O_{j_{r+1}}^{\mathrm{app},(\le \Delta_{\max})}(s_{r+1};f_{r+1};\mu)\\
&\cdots
\mathcal O_{j_m}^{\mathrm{app},(\le \Delta_{\max})}(s_m;f_m;\mu)\,B_m.
\end{aligned}
\]
Then there exist a compact positive-time neighborhood $K_{W_+}$, finite index sets $\mathcal A_r$ independent of $\mathbf t,\mathbf s$, and uniformly bounded positive-time word families
\[
\mathbf B_{r,\alpha}(\mathbf t,\mathbf s)
:=
\bigl(B_{r,\alpha,0}(\mathbf t,\mathbf s),\dots,B_{r,\alpha,q_{r,\alpha}}(\mathbf t,\mathbf s)\bigr),
\qquad
B_{r,\alpha,\ell}(\mathbf t,\mathbf s)\in\mathcal B^{(\le \Delta_{\max})}_{+,\mathrm{cap}},
\]
supported in $K_{W_+}$, such that
\[
\begin{aligned}
T_r(\mathbf t,\mathbf s)
={}&
\sum_{\alpha\in\mathcal A_r}
 c_{r,\alpha}(\mathbf t,\mathbf s)\,
B_{r,\alpha,0}(\mathbf t,\mathbf s)
\bigl(
\mathcal O_{i_{r,\alpha}}^{\mathrm{app},(\le \Delta_{\max})}(t_r;f_r;\mu)
-
\mathcal O_{i_{r,\alpha}}^{\mathrm{app},(\le \Delta_{\max})}(s_r;f_r;\mu)
\bigr)\\
&\qquad\qquad\cdot
B_{r,\alpha,1}(\mathbf t,\mathbf s)\cdots B_{r,\alpha,q_{r,\alpha}}(\mathbf t,\mathbf s),
\end{aligned}
\]
with coefficient bound
\[
\begin{aligned}
|c_{r,\alpha}(\mathbf t,\mathbf s)|
&\le
C_{r,\alpha,W_+}\,\Lambda^{\mathrm{mix}}_{r}(\mathbf t,\mathbf s),\\
\Lambda^{\mathrm{mix}}_{r}(\mathbf t,\mathbf s)
&:=
\prod_{\ell<r}\big\|\big(C^{(\le \Delta_{\max})}(t_\ell,\mu)\big)^{-1}\big\|_{\mathrm{op}}\\
&\qquad\cdot
\prod_{\ell>r}\big\|\big(C^{(\le \Delta_{\max})}(s_\ell,\mu)\big)^{-1}\big\|_{\mathrm{op}}.
\end{aligned}
\]
Hence Theorem~\ref{thm:finite_truncation_inverse_control} gives the explicit coefficient envelope
\begin{equation}\label{eq:finite_cap_literal_summand_coeff_env}
\Lambda^{\mathrm{mix}}_{r}(\mathbf t,\mathbf s)
\le
C^{\mathrm{inv}}_{r,W_+,\Delta_{\max},\mu}
\prod_{\ell<r}\Bigl(\frac{u(\mu)}{u(\mu_{t_\ell})}\Bigr)^{p_{\Delta_{\max}}}
\prod_{\ell>r}\Bigl(\frac{u(\mu)}{u(\mu_{s_\ell})}\Bigr)^{p_{\Delta_{\max}}}.
\end{equation}
After subtracting the centered finite-cap SFTE identities in the $r$th slot, every term in the last display is a finite linear combination of uniformly bounded positive-time word families multiplying a single centered remainder channel at scales $t_r$ and $s_r$.
Consequently, for the capwise vacuum channel there exists a constant $C^{\mathrm{tel}}_{r,W_+,\Delta_{\max},\mu}<\infty$ such that
\begin{equation}\label{eq:finite_cap_literal_summand_bound}
\Big|
\big\langle \Omega_{\mathrm{cap}},T_r(\mathbf t,\mathbf s)\,\Omega_{\mathrm{cap}}\big\rangle_{\mathrm{cap}}
\Big|
\le
C^{\mathrm{tel}}_{r,W_+,\Delta_{\max},\mu}\,
\Lambda^{\mathrm{mix}}_{r}(\mathbf t,\mathbf s)\,
\bigl(t_r^{\eta_{\Delta_{\max}}}+s_r^{\eta_{\Delta_{\max}}}\bigr),
\end{equation}
where the coefficient factor is the explicit polylogarithmic envelope carried by the non-difference slots.
Combining \eqref{eq:finite_cap_literal_summand_coeff_env} with the finite-cap remainder / inverse compatibility clause of Corollary~\ref{cor:finite_truncation_SFTE} gives the displayed handoff
\begin{equation}\label{eq:finite_cap_literal_summand_o1}
\Lambda^{\mathrm{mix}}_{r}(\mathbf t,\mathbf s)\,
\bigl(t_r^{\eta_{\Delta_{\max}}}+s_r^{\eta_{\Delta_{\max}}}\bigr)
\xrightarrow[\mathbf t,\mathbf s\downarrow 0]{}0,
\end{equation}
so the right-hand side of \eqref{eq:finite_cap_literal_summand_bound} is an $o(1)$ quantity.
\end{lemFspecial}

\begin{proof}
For $\ell\ne r$, write the non-difference slot at its relevant scale $u_\ell$ ($u_\ell=t_\ell$ for $\ell<r$ and $u_\ell=s_\ell$ for $\ell>r$) using the inverse-flow definition \eqref{eq:inverse_flow_approx}:
\[
\mathcal O_{j_\ell}^{\mathrm{app},(\le \Delta_{\max})}(u_\ell;f_\ell;\mu)
=
\sum_{a=1}^{N(\Delta_{\max})}
\bigl(C^{(\le \Delta_{\max})}(u_\ell,\mu)^{-1}\bigr)_{j_\ell a}\,
Y_a^{(\le \Delta_{\max})}(u_\ell;f_\ell),
\]
where
\[
Y_a^{(\le \Delta_{\max})}(u_\ell;f_\ell)
:=
\mathcal O_{a,u_\ell}^{(\le \Delta_{\max})}(f_\ell)
-
\omega_{\mathrm{flow}}^{(u_{\mathrm{ren}})}\!\bigl(\mathcal O_{a,u_\ell}^{(\le \Delta_{\max})}(f_\ell)\bigr)\mathbf 1
\in\mathcal B^{(\le \Delta_{\max})}_{+,\mathrm{cap}}.
\]
Because the truncation is finite, expanding the product over all $\ell\ne r$ produces only finitely many index choices $\alpha$; this yields the displayed decomposition of $T_r(\mathbf t,\mathbf s)$.
The scalar coefficient attached to a given choice is a product of matrix entries of the inverse coefficient matrices from the non-difference slots, so it is bounded by the corresponding product of operator norms, which gives the displayed coefficient bound.

Each bounded factor $Y_a^{(\le \Delta_{\max})}(u_\ell;f_\ell)$ lies in the positive-time capwise algebra and, on the fixed truncation, belongs to a uniformly bounded family supported in a common compact positive-time neighborhood depending only on the original word $W_+$.
Multiplying these factors with the fixed bounded separators $B_0,\dots,B_m$ gives the uniformly bounded word families $\mathbf B_{r,\alpha}(\mathbf t,\mathbf s)$ supported in some common neighborhood $K_{W_+}$.

Finally, subtract the centered finite-cap SFTE identities at scales $t_r$ and $s_r$ in the $r$th slot.
The common centered local symbol cancels, so the remaining difference slot is a finite linear combination of centered remainder channels.
For each $\alpha\in\mathcal A_r$, write the resulting single-slot difference channel as
\[
\begin{aligned}
\Delta_{r,\alpha}(\mathbf t,\mathbf s)
:={}&
B_{r,\alpha,0}(\mathbf t,\mathbf s)
\bigl(
\mathcal O_{i_{r,\alpha}}^{\mathrm{app},(\le \Delta_{\max})}(t_r;f_r;\mu)
-
\mathcal O_{i_{r,\alpha}}^{\mathrm{app},(\le \Delta_{\max})}(s_r;f_r;\mu)
\bigr)\\
&\cdot
B_{r,\alpha,1}(\mathbf t,\mathbf s)\cdots B_{r,\alpha,q_{r,\alpha}}(\mathbf t,\mathbf s).
\end{aligned}
\]
Because $\mathcal A_r$ is finite and the word families $\mathbf B_{r,\alpha}(\mathbf t,\mathbf s)$ are uniformly bounded and supported in the common neighborhood $K_{W_+}$, Proposition~\ref{prop:finite_cap_bounded_factor_insertion}, specialized to $\Psi=\Phi=\Omega_{\mathrm{cap}}$, yields constants $C^{\mathrm{rem}}_{r,\alpha,W_+,\Delta_{\max},\mu}$ independent of $\mathbf t,\mathbf s$ such that
\[
\Big|
\big\langle \Omega_{\mathrm{cap}},\Delta_{r,\alpha}(\mathbf t,\mathbf s)\,\Omega_{\mathrm{cap}}\big\rangle_{\mathrm{cap}}
\Big|
\le
C^{\mathrm{rem}}_{r,\alpha,W_+,\Delta_{\max},\mu}
\bigl(t_r^{\eta_{\Delta_{\max}}}+s_r^{\eta_{\Delta_{\max}}}\bigr).
\]
Therefore
\[
\begin{aligned}
\Big|
\big\langle \Omega_{\mathrm{cap}},T_r(\mathbf t,\mathbf s)\,\Omega_{\mathrm{cap}}\big\rangle_{\mathrm{cap}}
\Big|
&\le
\sum_{\alpha\in\mathcal A_r}
|c_{r,\alpha}(\mathbf t,\mathbf s)|\,
C^{\mathrm{rem}}_{r,\alpha,W_+,\Delta_{\max},\mu}
\bigl(t_r^{\eta_{\Delta_{\max}}}+s_r^{\eta_{\Delta_{\max}}}\bigr)\\
&\le
C^{\mathrm{tel}}_{r,W_+,\Delta_{\max},\mu}\,
\Lambda^{\mathrm{mix}}_{r}(\mathbf t,\mathbf s)\,
\bigl(t_r^{\eta_{\Delta_{\max}}}+s_r^{\eta_{\Delta_{\max}}}\bigr),
\end{aligned}
\]
which is exactly \eqref{eq:finite_cap_literal_summand_bound}.
Together with \eqref{eq:finite_cap_literal_summand_coeff_env}, the finite-cap remainder / inverse compatibility clause already recorded in Corollary~\ref{cor:finite_truncation_SFTE} yields \eqref{eq:finite_cap_literal_summand_o1}.
Thus \eqref{eq:finite_cap_literal_summand_bound} is the displayed archive-grade estimate for the literal scale-dependent telescoping summand, and \eqref{eq:finite_cap_literal_summand_o1} is the explicit final $o(1)$ handoff.
\end{proof}

\paragraph*{Inductive limit and the full local field algebra}

\begin{definition}[Renormalized local field algebra]
\label{def:local_field_algebra}
Let $\mathfrak A_{\mathrm{loc}}^{(\le \Delta_{\max})}$ be the $*$--algebra generated by
\(\{\mathcal O^{(\le \Delta_{\max})}_{\mathrm{ren},i}(\phi):\phi\in C_c^{\infty}(\mathbb R^4)\}\)
as provided by Theorem~\ref{thm:field_algebra_truncation}.
Define the renormalized local field algebra as the inductive limit
\[
\mathfrak A_{\mathrm{loc}}:=\varinjlim_{\Delta_{\max}\to\infty}\mathfrak A_{\mathrm{loc}}^{(\le \Delta_{\max})},
\]
with embeddings given by enlarging the operator basis.
\end{definition}

\begin{theorem}[Temperedness verification for the inductive-union sharp-local Schwinger family]\label{thm:LaneA_temperedness_verification}
Fix $u_{\mathrm{ren}}\in(0,u_\ast]$ and let $\omega^{(u_{\mathrm{ren}})}$ be the full sharp-local state of Corollary~\ref{cor:LaneA_full_scope}.
For any finite family of renormalized fields
\[
\mathcal O^{\mathrm{ren}}_{[P_1]}(f_1;\mu),\dots,\mathcal O^{\mathrm{ren}}_{[P_m]}(f_m;\mu)\in \mathfrak A_{\mathrm{loc}},
\]
there exist a finite engineering cap $\Delta_{\max}$ containing that family, an integer $N=N(\Delta_{\max},m)$, and a constant $C<\infty$ such that
\begin{equation}\label{eq:LaneA_temperedness_bound}
\Big|
\omega^{(u_{\mathrm{ren}})}\!\Big(
\mathcal O^{\mathrm{ren}}_{[P_1]}(f_1;\mu)\cdots \mathcal O^{\mathrm{ren}}_{[P_m]}(f_m;\mu)
\Big)
\Big|
\le C\prod_{r=1}^m p_N(f_r),
\end{equation}
where
\[
p_N(f):=\sum_{|\alpha|\le N}\sup_{x\in\mathbb R^4}(1+|x|)^N|\partial^\alpha f(x)|
\]
is a Schwartz seminorm.
Hence each finite-family Schwinger functional extends uniquely to a tempered distribution on $(\mathbb R^4)^m$.
Equivalently, the full inductive-union Schwinger family of Lane~A satisfies OS$_0$.
\end{theorem}

\begin{proof}
Fix the field family and choose $\Delta_{\max}$ so that all of its members lie in $\mathfrak A_{\mathrm{loc}}^{(\le \Delta_{\max})}$.
For each field, Theorem~\ref{thm:field_algebra_truncation} provides dyadic bounded approximants
$X^{(\le \Delta_{\max})}_{j,r}(f_r)$ converging to the corresponding renormalized field.
At each fixed dyadic scale these are bounded flowed observables, so the Lane~C moment/cumulant package
(Theorem~\ref{thm:poly_cumulant}, Proposition~\ref{ass:precompact}, and Theorem~\ref{thm:C3U}) together with the heat-kernel bounds of Appendices~B--D implies that, for some finite $N$, the base-scale correlators satisfy
\[
\Big|\omega\Big(\prod_{r=1}^m X^{(\le \Delta_{\max})}_{0,r}(f_r)\Big)\Big|
\le C_0\prod_{r=1}^m p_N(f_r).
\]
Because the cap and the word length $m$ are finite, one may take the same seminorm order $N$ for the whole family.

Next, the multilinear telescoping identity, Lemma~\ref{lem:finite_cap_common_core}, and the summable increment bound of Theorem~\ref{thm:field_algebra_truncation} show that the dyadic $m$-point forms are Cauchy in the seminorm determined by $\prod_r p_N(f_r)$:
\[
\Big|
\omega\Big(\prod_{r=1}^m X^{(\le \Delta_{\max})}_{J+1,r}(f_r)\Big)
-
\omega\Big(\prod_{r=1}^m X^{(\le \Delta_{\max})}_{J,r}(f_r)\Big)
\Big|
\le C_1(1+J)^{-1-\eta}\prod_{r=1}^m p_N(f_r).
\]
Summing over $J$ gives a finite bound for the limiting finite-cap Schwinger function,
\[
\Big|S^{(\le \Delta_{\max})}_{[P_1],\dots,[P_m]}(f_1,\dots,f_m)\Big|
\le \Big(C_0+C_1\sum_{J\ge 0}(1+J)^{-1-\eta}\Big)\prod_{r=1}^m p_N(f_r).
\]
Corollary~\ref{cor:finite_truncation_SFTE} identifies the same limit with the Lane~A inverse-flow representatives and shows that the SFTE remainders obey the same finite-cap bound, so \eqref{eq:LaneA_temperedness_bound} is the finite-cap OS$_0$ estimate.
Finally, Corollary~\ref{cor:LaneA_full_scope} passes that estimate to the inductive union because every finite field family belongs to some finite cap and the capwise Schwinger functions agree on overlaps.
\end{proof}

\begin{remark}[About OPE]
The inductive-limit construction defines an OS-positive local field algebra of smeared gauge-invariant composite fields.
It does \emph{not} by itself give an operator product expansion at coincident points.
For mass gap purposes, OS reconstruction and a spectral gap follow from reflection positivity plus exponential clustering.
An OPE (if desired) is an additional regularity structure that can be formulated and proved within this framework using multiscale
short-distance estimates; in this paper it is treated as an optional refinement rather than a prerequisite.
\end{remark}

\paragraph*{Universality upgrade from $\mathcal O_{\ell}$ to $\mathfrak A_{\mathrm{loc}}$}

\begin{theorem}[Universality for the renormalized local field algebra (finite truncations)]
\label{thm:universality_local_field_algebra}
Work within the local action class of Definition~\ref{def:action_universality_class} and along the flow renormalization scheme
defined by \eqref{eq:t0_def}.
Assume the uniform entrance/non-collapse hypothesis of Theorem~\ref{thm:uniqueness_on_Oell} for this action class.
Then for each fixed $\Delta_{\max}$:
\begin{enumerate}[leftmargin=*]
\item the renormalized fields $\mathbf{\mathcal O}^{(\le \Delta_{\max})}_{\mathrm{ren}}$ constructed in
Theorem~\ref{thm:field_algebra_truncation} have continuum limits that exist without subsequence and are independent of the
discretization coefficients $\mathbf c\in\mathcal C$;
\item the joint laws of finitely many smeared fields in $\mathfrak A_{\mathrm{loc}}^{(\le \Delta_{\max})}$ are unique and universal
within the action class, up to the choice of the normalization conditions that fix the renormalization matrix.
\end{enumerate}
\end{theorem}

\begin{proof}
Fix $\Delta_{\max}$.
At each positive flow time $t>0$ the flowed observables are bounded local functionals of the link variables.
Therefore, by Theorem~\ref{thm:universality_on_Oell} and the Duhamel stability estimates
(Lemma~\ref{lem:duhamel_covariance} and Proposition~\ref{prop:gf_stability_discretization}), their macroscopic laws are universal.

The renormalized fields are obtained as $L^2$ limits of linear combinations of these bounded flowed observables at dyadic times.
The one-shell increment bounds are uniform in the action coefficients $\mathbf c\in\mathcal C$ because the small-curvature coercivity
constants are uniform on $\mathcal C$ (Definition~\ref{def:action_universality_class}) and the curvature Brascamp--Lieb constant is
controlled uniformly at large $\beta$.
Hence the dyadic Cauchy construction is stable under varying $\mathbf c$.

Finally, the normalization conditions fixing $Z^{(\le\Delta_{\max})}$ are expressed in terms of a finite set of two-point functions
at a fixed physical scale; these are universal by the preceding bounded-observable universality step.
Thus the limiting renormalized fields and all their joint moments (for any fixed polynomial) are independent of $\mathbf c$.
\end{proof}

\subsubsection{Uniqueness and universality (no subsequence)}
\label{subsec:uniqueness_universality}

The compactness argument underlying Theorem~\ref{thm:bounded_algebra_limit} yields existence of subsequential limits on $\mathcal O_{\ell}$. We now prove uniqueness (no subsequence) and universality within the action class by combining contractivity of the exact RG in the KP basin with the weak-window step-scaling control of the UV coupling.

In the present paper, the natural uniqueness mechanism is \emph{contractivity} of the exact RG in the KP basin combined with the
step-scaling control of the UV coupling.

\begin{proposition}[Uniqueness on $\mathcal O_{\ell}$ reduces to a single stability estimate]\label{prop:uniqueness_reduction}
Fix $\ell>0$ and let $\beta,\beta'$ be large.
Assume that there exists a deterministic scale $L_\mathrm{enter}(\beta)$ with bounded physical size $a(\beta)L_\mathrm{enter}(\beta)\le C\ell$
such that the exact blocked measures at scale $L_\mathrm{enter}(\beta)$ lie in the KP admissible window.
Then for any bounded observable $O\in\mathcal O_{\ell}$,
\[\big|\langle O\rangle_{\mu_{\beta}}-\langle O\rangle_{\mu_{\beta'}}\big|\ \le\ C(O)\,\rho_*^{\,k(\ell)}+\mathrm{UV\,Err}(\beta,\beta'),\]
where $\rho_*\in(0,1)$ is the KP contraction rate and $\mathrm{UV\,Err}(\beta,\beta')\to 0$ as $\beta,\beta'\to\infty$.
In particular, if $\mathrm{UV\,Err}(\beta,\beta')$ can be bounded by a function of the flow coupling difference
$|g^2_{\mathrm{GF}}(\ell/a(\beta)) - g^2_{\mathrm{GF}}(\ell/a(\beta'))|$, uniqueness follows from the step-scaling control.
\end{proposition}

\begin{theorem}[Uniqueness on $\mathcal O_{\ell}$ along the flow renormalization scheme]\label{thm:uniqueness_on_Oell}
Fix $\ell>0$ and define the scaling trajectory $a(\beta)=\ell_0/\sqrt{t_0(\beta)}$ by \eqref{eq:t0_def}.
Assume:
\begin{enumerate}[leftmargin=*]
\item (Non-collapse / one-shot entrance) there exists $\beta_\mathrm{ent}$ and constants $C_{\mathrm{ent}},\rho_*\in(0,1)$ such that for all $\beta\ge\beta_\mathrm{ent}$ the one-shot entrance block size $B(\beta)$ from \eqref{eq:h60_Bbeta_def} satisfies
\[
 a(\beta)\,B(\beta)\le C_{\mathrm{ent}}\ell
\]
and the exact blocked measure at that same scale lies in the KP basin with contraction rate $\rho_*$ (uniform in volume and boundary conditions);
\item (UV matching at the one-shot entrance scale) for the renormalization condition \eqref{eq:t0_def}, the blocked relevant parameter at scale $B(\beta)$ is fixed (up to an $o_{\beta\to\infty}(1)$ error) by the flow coupling value at the reference scale $\ell_0$.
\end{enumerate}
Then the laws of the macroscopic observables in $\mathcal O_{\ell}$ converge as $\beta\to\infty$ without subsequence, and the limiting OS-positive continuum theory on $\mathcal O_{\ell}$ is unique (in this renormalization scheme).
\end{theorem}

\begin{proof}
Let $O\in\mathcal O_{\ell}$.
Choose $k(\ell)$ RG steps so that the support of $O$ becomes contained in $O(1)$ blocks at the one-shot entrance scale $B(\beta)$.
By uniform KP contraction inside the basin, the map taking the entrance-scale blocked specification to $\langle O\rangle$ is Lipschitz with factor $\rho_*^{k(\ell)}$.
Thus for $\beta,\beta'$ large,
\[
|\langle O\rangle_{\mu_\beta}-\langle O\rangle_{\mu_{\beta'}}|
\le C(O)\,\rho_*^{k(\ell)}
+ \big|\langle O\rangle_{\mathrm{ent},\beta}-\langle O\rangle_{\mathrm{ent},\beta'}\big|,
\]
where $\langle\cdot\rangle_{\mathrm{ent},\beta}$ denotes expectation under the entrance-scale blocked measure.
By the renormalization condition \eqref{eq:t0_def}, the flow coupling at the reference physical scale $\ell_0$ is fixed.
Step scaling transports this to the entrance scale, so the entrance-scale relevant parameter is the same for all sufficiently large $\beta$ (up to $o(1)$), and the remaining directions are irrelevant and uniformly small by the entrance assumption.
Hence the entrance-scale expectations converge and the second term tends to $0$.
Letting first $\beta,\beta'\to\infty$ and then taking $k(\ell)$ large (or equivalently refining the entrance-to-observable descent) yields Cauchy convergence.
Joint laws follow by repeating the preceding contraction estimate for products of bounded macroscopic observables.
\end{proof}

\begin{remark}
In words: once two trajectories enter the same KP basin at a bounded physical scale, their long-distance expectations become
insensitive to further UV details by contraction. Uniqueness is therefore reduced to matching the renormalized coupling at that entrance scale.
This is the standard universality mechanism in constructive RG.
\end{remark}


\paragraph*{Quantitative UV matching error and discretization stability}

Proposition~\ref{prop:uniqueness_reduction} isolates the only non-contracting input needed for uniqueness: control of the
\emph{UV matching error} at the entrance scale.
Here we record a standard way to bound this error and to obtain universality across a conventional class of local lattice actions.

\begin{definition}[A standard universality class of local gauge actions]\label{def:action_universality_class}
Fix a finite set of oriented Wilson loops $\mathfrak L$ (plaquettes, rectangles, chairs, \dots) contained in a unit cube.
For each loop $\gamma\in\mathfrak L$ let $U_\gamma$ denote its holonomy and let $V_\gamma:G\to\mathbb R$ be a smooth class function
with $V_\gamma(\mathbf 1)=0$ and small-curvature expansion
\begin{equation}\label{eq:loop_potential_expansion}
 V_\gamma(e^{iX})=\frac{1}{2}\kappa_\gamma\,\|X\|^2 + O(\|X\|^4)\qquad (\|X\|\le \varepsilon_0),
\end{equation}
for some $\kappa_\gamma>0$.
A lattice action in this class has the form
\begin{equation}\label{eq:general_local_action}
 S_{\beta,\mathbf c}(U)=\beta\sum_{x}\sum_{\gamma\in\mathfrak L} c_\gamma\,V_\gamma\!\bigl(U_\gamma(x)\bigr),
\end{equation}
with coefficients $\mathbf c=(c_\gamma)_{\gamma\in\mathfrak L}$ in a fixed compact set $\mathcal C\subset[0,\infty)^{\mathfrak L}$.
We say the action is \emph{continuum-normalized} if $\sum_{\gamma} c_\gamma\kappa_\gamma=1$, so that the quadratic part agrees with the Wilson
$F^2$ normalization.
\end{definition}

\begin{remark}[Sparse coefficient families are allowed]
\label{rem:action_class_sparse}
The proofs below use only compactness of $\mathcal C$, smoothness of the local potentials, and the continuum-normalized quadratic coercivity
$\sum_{\gamma}c_\gamma\kappa_\gamma=1$.
We therefore allow some coefficients to vanish.
Appendix~\ref{app:global_form_resolution} uses this flexibility to place several faithful Wilson plaquette regularizations of the same compact simple group into a single finite action family by labeling the plaquette slots separately and turning on one label at a time.
\end{remark}

\begin{lemma}[Duhamel covariance formula]\label{lem:duhamel_covariance}
Let $\langle\cdot\rangle_{\mathbf c}$ denote expectation in finite volume (with any fixed boundary condition) for the action
$S_{\beta,\mathbf c}$.
For any bounded observable $F$ and any loop coefficient $c_\gamma$,
\begin{equation}\label{eq:duhamel_formula}
 \frac{\partial}{\partial c_\gamma}\,\langle F\rangle_{\mathbf c}
 \;=\; -\,\beta\,\mathrm{Cov}_{\mathbf c}\!\Bigl(F,\ \sum_{x} V_\gamma\!\bigl(U_\gamma(x)\bigr)\Bigr).
\end{equation}
\end{lemma}

\begin{proof}
Differentiate under the integral sign:
$\partial_{c_\gamma}\langle F\rangle_{\mathbf c}
= -\langle F\,\partial_{c_\gamma}S_{\beta,\mathbf c}\rangle_{\mathbf c}
+\langle F\rangle_{\mathbf c}\langle \partial_{c_\gamma}S_{\beta,\mathbf c}\rangle_{\mathbf c}$,
with $\partial_{c_\gamma}S_{\beta,\mathbf c}=\beta\sum_x V_\gamma(U_\gamma(x))$.
\end{proof}

\begin{proposition}[Stability of the gradient-flow coupling under discretization changes]\label{prop:gf_stability_discretization}
Fix $c\in(0,1)$ in the flow-coupling definition and consider two continuum-normalized actions $S_{\beta,\mathbf c}$ and $S_{\beta,\mathbf c'}$
in the class of Definition~\ref{def:action_universality_class}.
Let $g^2_{\mathrm{GF},\mathbf c}(L)$ and $g^2_{\mathrm{GF},\mathbf c'}(L)$ denote the corresponding finite-volume gradient-flow couplings
(using the same Wilson flow \eqref{eq:wilson_flow} and the same observable definition \eqref{eq:gf_coupling_def}).
Assume the weak-coupling curvature control (Lemma~\ref{lem:weak_localization} and Proposition~\ref{prop:curvature_BL}) holds uniformly for
$\mathbf c,\mathbf c'\in\mathcal C$.
Then for all sufficiently large $\beta$ and all $L$ in the weak window,
\begin{equation}\label{eq:gf_stability_bound}
 \big|g^2_{\mathrm{GF},\mathbf c}(L)-g^2_{\mathrm{GF},\mathbf c'}(L)\big|
 \ \le\ C_{\mathcal C}\,\|\mathbf c-\mathbf c'\|\,\Big(\frac{a}{L}\Big)^{2}\,\big(1+g^2_{\mathrm{GF}}(L)\big),
\end{equation}
where $a=a(\beta)$ is the lattice spacing and $C_{\mathcal C}$ depends only on $\mathcal C$, $d$, and $G$.
In particular, at fixed physical scale $\ell$ (so $L=\ell/a$), the difference is $O(a^2)$ as $a\downarrow 0$.
\end{proposition}

\begin{proof}
On the small-curvature event, each loop term admits the Taylor expansion \eqref{eq:loop_potential_expansion}.
Continuum-normalization ensures that the \emph{quadratic} parts agree, so the first difference between the two actions begins at quartic order in
the curvature.
By curvature Brascamp--Lieb domination (Proposition~\ref{prop:curvature_BL}), all local curvature moments in the weak window are uniformly bounded
and scale like powers of the running coupling at scale $L$.
The flowed energy density at time $t=cL^2$ is a local curvature polynomial smeared at scale $L$; its expectation therefore admits a controlled
Symanzik-type expansion in $a/L$ whose leading discretization correction is $O((a/L)^2)$.
Applying the Duhamel formula (Lemma~\ref{lem:duhamel_covariance}) and summing the exponentially decaying covariance contributions over $x$
(using the fixed-$\beta$ exponential clustering module once the entrance scale is reached) yields \eqref{eq:gf_stability_bound}.
\end{proof}

\begin{corollary}[Explicit UV matching error bound]\label{cor:UV_matching_error}
Fix a physical scale $\ell>0$ and tune $\beta$ for each action in the class of Definition~\ref{def:action_universality_class} by the same flow
renormalization condition \eqref{eq:t0_def}, so that $a(\beta)=\ell_0/\sqrt{t_0(\beta)}$.
Then for each fixed $\ell$ the UV matching error in Proposition~\ref{prop:uniqueness_reduction} satisfies
\[
\mathrm{UV\,Err}(\beta,\beta')\ \le\ C(O)\,\Big(\frac{a(\beta)}{\ell}\Big)^{2}\ +\ C(O)\,\Big(\frac{a(\beta')}{\ell}\Big)^{2}
\]
for all sufficiently large $\beta,\beta'$.
Hence $\mathrm{UV\,Err}(\beta,\beta')\to 0$ as $\beta,\beta'\to\infty$.
\end{corollary}

\begin{theorem}[Universality and uniqueness on $\mathcal O_{\ell}$ within the action class]\label{thm:universality_on_Oell}
Fix $\ell>0$ and consider any two continuum-normalized actions in the class of Definition~\ref{def:action_universality_class},
each tuned along the same flow renormalization scheme \eqref{eq:t0_def}.
Assume the non-collapse/entrance results of the paper (Sections~\ref{sec:stepscaling_proof} and \ref{sec:what_remains_continuum})
hold uniformly for $\mathbf c\in\mathcal C$.
Then:
\begin{enumerate}[leftmargin=*]
\item the continuum limits of all observables in $\mathcal O_{\ell}$ exist \emph{without subsequences};
\item the limiting OS-positive continuum theory on $\mathcal O_{\ell}$ is independent of the chosen action coefficients $\mathbf c\in\mathcal C$.
\end{enumerate}
\end{theorem}

\begin{proof}
By Corollary~\ref{cor:UV_matching_error}, the entrance-scale relevant parameter (equivalently the entrance-scale flow coupling) is the same for all
large $\beta$ and all $\mathbf c\in\mathcal C$ up to an $o(1)$ error.
Once in the KP basin, contraction implies that expectations of bounded macroscopic observables depend Lipschitz-continuously on the entrance-scale
relevant coordinate with contraction factor $\rho_*^{k(\ell)}$.
Letting $\beta,\beta'\to\infty$ and then taking $k(\ell)$ large yields Cauchy convergence and independence of $\mathbf c$.
\end{proof}


\subsubsection{Nontriviality: a quantitative diagnostic}
\label{subsec:nontriviality}

For completeness we record a convenient diagnostic for nontriviality of the continuum limit: the step-scaling function.

\begin{theorem}[Explicit coefficient extraction for the connected degree-$4$ step-scaling term]\label{thm:b0_coefficient_extraction}
Assume Lemma~\ref{lem:weak_localization} and Proposition~\ref{prop:curvature_BL}.
Fix the gradient-flow coupling scheme \eqref{eq:gf_coupling_def} and the scale factor $2$.
For a finite periodic box $\Lambda(L)$ with flow time $t=cL^2$, let
\[
\mathcal K_{s,L}(x,p):\mathfrak g\to \mathfrak g
\]
denote the linearized Wilson-flow/clover kernel sending an exact plaquette-curvature insertion at $p\in P(\Lambda(L))$
to the corresponding linearized flowed field-strength insertion contributing to $E(s,x)$.
Let $\Gamma_L$ be the exact conditional covariance operator of the curvature field $dA$ on exact plaquette $2$-forms under the conditional measure on $\Omega_{\varepsilon_0}$, and let $\Gamma_L^{1/2}$ be its positive square root.
Define the covariance-dressed kernel
\[
\widetilde{\mathcal K}_{s,L}(x,p)
:= \sum_{q\in P(\Lambda(L))}\mathcal K_{s,L}(x,q)\,\Gamma_L^{1/2}(q,p),
\]
and the finite-volume squared-kernel shell integral
\begin{equation}\label{eq:b0_shell_integral}
\mathcal I_{L,c}
:=
\sum_{x\in \Lambda(L)}\sum_{p\in P(\Lambda(L))}
\left\|\int_{cL^2}^{4cL^2}\widetilde{\mathcal K}_{s,L}(x,p)\,\frac{ds}{s}\right\|_{\mathfrak g}^{2}.
\end{equation}
Let $\mathcal C^{(4)}_{L,c}$ denote the connected degree-$4$ contribution appearing at order $u^2$ in the weak-window Taylor/cumulant expansion of $\Sigma(u)-u$.
Then
\begin{equation}\label{eq:b0_shell_coefficient}
\mathcal C^{(4)}_{L,c}
=
\frac{\mathfrak c_{\mathrm{Lie}}(G)}{\mathcal N^2}\,\mathcal I_{L,c},
\qquad
\mathfrak c_{\mathrm{Lie}}(G):=\sum_{a,b,c}(f^{abc})^2=(\dim \mathfrak g)\,C_A.
\end{equation}
In particular $\mathcal I_{L,c}<\infty$, $\mathcal I_{L,c}>0$, and $\mathcal C^{(4)}_{L,c}>0$ for every compact connected simple $G$.
\end{theorem}

\begin{proof}
Linearizing the Wilson flow \eqref{eq:wilson_flow} about the trivial configuration gives the heat evolution on exact plaquette $2$-forms.
Combined with the clover discretization in \eqref{eq:flow_energy_def}, this yields the kernel $\mathcal K_{s,L}$.
Appendices~\ref{app:lattice_heat_kernel} and~\ref{app:lattice_parametrix} give the required Gaussian and finite-difference bounds for that linearized kernel, uniformly at fixed ratio $s/L^2$.
Consequently the integral in \eqref{eq:b0_shell_integral} converges absolutely, and $\mathcal I_{L,c}<\infty$.

Write the linearized flowed field strength as
\[
G^{\mathrm{lin}}_s(x)=\sum_{p\in P(\Lambda(L))}\mathcal K_{s,L}(x,p)\,F_p(A),
\]
so the quadratic part of the flowed energy density is
\[
E^{(2)}(s,x)=\frac12\,\langle G^{\mathrm{lin}}_s(x),G^{\mathrm{lin}}_s(x)\rangle.
\]
The connected degree-$4$ contribution to the coefficient of $u^2$ in $\Sigma(u)-u$ is obtained by pairing two copies of this quadratic term.
Every term involving a cubic-or-higher Taylor coefficient of the flow map, or a genuine curvature cumulant of order $\ge 3$, is $O(u^3)$ by Proposition~\ref{prop:curvature_BL} together with the standard polynomial cumulant bounds for sub-Gaussian curvature fields.
Thus the order-$u^2$ coefficient is exhausted by the covariance pairing of the linearized quadratic piece.

Passing that covariance through the positive square root $\Gamma_L^{1/2}$ gives exactly the squared-kernel formula \eqref{eq:b0_shell_integral}; this is simply the Hilbert--Schmidt norm of the dyadic shell operator built from the covariance-dressed linearized kernel.
This is the explicit finite-volume squared-kernel integral whose existence was only sketched previously.

For the Lie factor, choose an orthonormal basis $\{T^a\}$ of $\mathfrak g$ with
$[T^a,T^b]=\sum_c f^{abc}T^c$.
In the connected degree-$4$ pairing, the color contraction is
\[
\sum_{a,b,c}(f^{abc})^2=(\dim \mathfrak g)\,C_A,
\]
by the standard identity $\sum_{b,c}f^{abc}f^{dbc}=C_A\delta^{ad}$.
Hence the entire group dependence is a positive multiple of $C_A$, proving \eqref{eq:b0_shell_coefficient}.

Finally, $\mathcal I_{L,c}>0$ because the linearized flowed field strength is not identically zero.
Choose any plaquette $p_0$ and any site $x_0$ incident to $p_0$.
The clover field-strength discretization has a nonzero linear coefficient in that plaquette channel at the identity, and the lattice heat kernel is strictly positive on the diagonal for every positive time.
Therefore $s\mapsto \widetilde{\mathcal K}_{s,L}(x_0,p_0)$ is not identically zero on $[cL^2,4cL^2]$, so the square integral in \eqref{eq:b0_shell_integral} is strictly positive.
This yields $\mathcal C^{(4)}_{L,c}>0$.
\end{proof}

\begin{lemma}[Quadratic step-scaling coefficient is strictly positive for simple $G$]\label{lem:b0_positive}
Let $G$ be compact connected simple with adjoint Casimir $C_A>0$.
For the gradient-flow finite-volume coupling in the fixed scheme of \eqref{eq:gf_coupling_def}, the weak-window step-scaling expansion takes the form
\begin{equation}\label{eq:b0_value}
 \Sigma(u)=u + b_0\,u^2\log 2 + O(u^3),\qquad b_0>0.
\end{equation}
Moreover $b_0$ is proportional to $C_A$ with a strictly positive proportionality constant depending only on the chosen coupling scheme.
\end{lemma}

\begin{proof}
By Theorem~\ref{thm:b0_coefficient_extraction}, the connected degree-$4$ contribution to the coefficient of $u^2$ in $\Sigma(u)-u$ is
\[
\mathcal C^{(4)}_{L,c}=\frac{\mathfrak c_{\mathrm{Lie}}(G)}{\mathcal N^2}\,\mathcal I_{L,c}>0.
\]
Since the weak-window expansion has the form $\Sigma(u)-u=b_0u^2\log 2+O(u^3)$ and $\log 2>0$, it follows that
\[
 b_0=\frac{\mathcal C^{(4)}_{L,c}}{\log 2}>0.
\]
Absorbing the harmless factor $(\dim\mathfrak g)/\mathcal N^2$ into the scheme constant gives the stated proportionality
$b_0=c_{\mathrm{flow}}\,C_A$ with $c_{\mathrm{flow}}>0$.
\end{proof}

\begin{remark}[Optional closed form]
The standard perturbative evaluation of the same connected degree-$4$ contribution gives the familiar closed form
$b_0=\tfrac{11}{3}\,\tfrac{C_A}{16\pi^2}$.
For the present submission the load-bearing point is the explicit positive formula \eqref{eq:b0_shell_coefficient} together with the squared-kernel integral \eqref{eq:b0_shell_integral};
the closed form is recorded only as a normalization cross-check and is not used downstream.
\end{remark}


\begin{proposition}[Nontriviality from nonzero one-loop coefficient]\label{prop:nontriviality_from_stepscaling}
Assume the step-scaling expansion of Proposition~\ref{prop:stepscaling_control_proved} with $b_0>0$.
Then any continuum limit for which $g^2_{\mathrm{GF}}(L)$ has a nondegenerate scaling limit cannot be a free (Gaussian) theory.
\end{proposition}

\begin{proof}
For a Gaussian fixed point, any finite-volume running coupling defined by a quadratic energy observable is exactly scale invariant at leading order,
so its step scaling satisfies $\Sigma(u)=u$.
The strict inequality $\Sigma(u)=u+b_0u^2\log 2+O(u^3)$ with $b_0>0$ forces a nonzero beta function at small $u$ and hence an interacting
renormalization group flow.
\end{proof}




\begin{proposition}[Continuum step scaling inherits the one-loop coefficient]\label{prop:continuum_stepscaling_b0}
Assume Theorem~\ref{thm:universality_on_Oell} and the weak-window step-scaling control of Proposition~\ref{prop:stepscaling_control_proved}.
Define the continuum gradient-flow coupling $g^2_{\mathrm{GF,cont}}(L)$ by taking the continuum limit of the flowed-energy definition
\eqref{eq:gf_coupling_def} on $\mathcal O_{\ell}$ (which exists and is unique by Theorem~\ref{thm:universality_on_Oell} for each fixed physical $L$).
Let $\Sigma_{\mathrm{cont}}$ be the corresponding continuum step-scaling function.
Then, as $u\downarrow 0$,
\[
\Sigma_{\mathrm{cont}}(u)=u+b_0u^2\log 2 + O(u^3),
\]
with the same universal coefficient $b_0$ as in Lemma~\ref{lem:b0_positive}.
In particular, the continuum beta function is nonzero near the origin and the continuum limit cannot be Gaussian.
\end{proposition}

\begin{proof}
By Theorem~\ref{thm:universality_on_Oell}, for each fixed physical $L$ the lattice couplings $g^2_{\mathrm{GF}}(L)$ converge (without subsequence)
to a unique continuum value $g^2_{\mathrm{GF,cont}}(L)$.
The finite-volume lattice step scaling $\Sigma_a$ converges pointwise to $\Sigma_{\mathrm{cont}}$ in the weak window because both arguments
$g^2_{\mathrm{GF}}(L)$ and $g^2_{\mathrm{GF}}(2L)$ have unique limits as $a\downarrow 0$.
Proposition~\ref{prop:stepscaling_control_proved} gives the weak-window expansion of $\Sigma_a$ with a uniform $O(u^3)$ remainder.
Passing to the $a\downarrow 0$ limit yields the same expansion for $\Sigma_{\mathrm{cont}}$, and the coefficient $b_0$ is universal by
Lemma~\ref{lem:b0_positive}.
\end{proof}


\subsection{Universality beyond finite truncations and scheme independence}
\label{sec:universality_beyond_truncations}

Theorem~\ref{thm:universality_local_field_algebra} yields universality for each fixed truncation
$\mathfrak A_{\mathrm{loc}}^{(\le \Delta_{\max})}$.
To obtain a fully local continuum field algebra suitable for OS reconstruction and scheme comparisons, one performs two standard upgrades:

\begin{enumerate}[leftmargin=*]
\item \emph{Beyond finite truncations:} promote universality from each fixed $\Delta_{\max}$ to the inductive limit
 $\mathfrak A_{\mathrm{loc}}=\varinjlim_{\Delta_{\max}}\mathfrak A_{\mathrm{loc}}^{(\le\Delta_{\max})}$, with quantitative control of how the
 constants depend on $\Delta_{\max}$;
\item \emph{Scheme independence:} compare different admissible renormalization schemes (different flow discretizations / normalizations) and show they
 define canonically isomorphic continuum algebras, related by finite renormalizations and coupling reparametrizations.
\end{enumerate}

\subsubsection{Growth control in \texorpdfstring{$\Delta_{\max}$}{Delta\_max} and coherence of truncations}

\begin{lemma}[Exponential bound on basis sizes]
\label{lem:basis_growth}
In fixed dimension $d=4$, for any choice of local gauge-invariant operator spaces $\mathscr V_{\Delta}$ supported in a unit cube as in
Definition~\ref{def:local_operator_spaces}, one may choose bases with
\begin{equation}\label{eq:basis_growth}
 m_{\Delta}\ \le\ C_{\mathrm{basis}}\,\Lambda_{\mathrm{basis}}^{\Delta}
\end{equation}
for constants $C_{\mathrm{basis}}<\infty$ and $\Lambda_{\mathrm{basis}}<\infty$ depending only on $(d,G)$ and the chosen support size.
\end{lemma}

\begin{proof}
Any gauge-invariant local lattice observable supported in a fixed cube can be written as a linear combination of products of traces of Wilson loops
contained in that cube (together with finitely many covariant difference insertions to represent derivatives).
The number of lattice paths of length $n$ in a fixed cube is $O((2d-1)^n)$, hence the number of loop monomials of total length $\le c\Delta$
is at most exponential in $\Delta$.
Quotienting by lattice symmetries and identities reduces the count.
Absorbing all support-dependent prefactors into $C_{\mathrm{basis}}$ yields \eqref{eq:basis_growth}.
\end{proof}

\begin{lemma}[One-shell increment bounds with explicit dimension growth]
\label{lem:one_shell_dimension_growth}
Assume the weak-coupling curvature control (Lemma~\ref{lem:weak_localization} and Proposition~\ref{prop:curvature_BL}) and the finite-range shell
decomposition used in Lemma~\ref{lem:one_shell_increment}.
Then there exist constants $\Lambda_{\mathrm{shell}}<\infty$ and, for each test function $\phi$, a constant $C_{\phi}<\infty$ such that for every
basis element $\mathcal O_i^{(\Delta)}$ and dyadic shell $j$ in the AF window,
\begin{equation}\label{eq:one_shell_dim_growth}
 \bigl\|\,\Delta_j\mathcal O_i^{(\Delta)}(\phi)\,\bigr\|_{L^2}
 \ \le\ C_{\phi}\,\Lambda_{\mathrm{shell}}^{\Delta}\,u_j^2,
\end{equation}
where $\Delta_j$ denotes the renormalized (order-$u_j^2$) shell increment from scale $t_j$ to $t_{j+1}$.
\end{lemma}

\begin{proof}
On the small-curvature event $\Omega_{\varepsilon_0}$ each flowed local operator is a finite sum of local curvature polynomials with uniformly bounded
coefficients depending on heat-kernel norms on a fixed cube.
For an operator of engineering dimension $\Delta$, the total degree of the curvature polynomial is $O(\Delta)$.
Curvature Brascamp--Lieb domination gives sub-Gaussian control for linear curvature functionals; by standard polynomial moment bounds for finite-degree
polynomials of sub-Gaussian fields, the $L^2$ norm grows at most exponentially in the degree, hence at most exponentially in $\Delta$.
The renormalized shell increment is quadratic in the shell fluctuations after the normalization conditions cancel the order-$u_j$ drift, hence gains a
factor $u_j^2$.
The complement $\Omega_{\varepsilon_0}^c$ contributes an exponentially small tail, absorbed into $C_{\phi}$.
\end{proof}

\begin{definition}[Dimension-weighted seminorms]
\label{def:dimension_weighted_norm}
Fix $\lambda>1$ and $\Delta_{\max}\ge 0$.
For an $M(\Delta_{\max})$--vector $v$ written in blocks $v^{(\Delta)}\in\mathbb R^{m_{\Delta}}$, set
\begin{equation}\label{eq:weighted_norm}
 \|v\|_{\lambda}:=\max_{0\le \Delta\le \Delta_{\max}}\ \lambda^{-\Delta}\,\|v^{(\Delta)}\|_{\ell^2}.
\end{equation}
This defines an equivalent norm on each truncation and induces an LF-topology on the inductive limit.
\end{definition}

\begin{lemma}[Coherence of triangular renormalization across truncations]
\label{lem:coherent_Z}
Fix a family of normalization conditions that are imposed dimension-by-dimension and, at dimension $\Delta$, involve only two-point functions of
operators of dimension $\le \Delta$.
Let $Z^{(\le\Delta_{\max})}(t)$ be the triangular renormalization matrices produced by Theorem~\ref{thm:field_algebra_truncation}.
Then for any $\Delta_0\le \Delta_{\max}$, the restriction of $Z^{(\le\Delta_{\max})}(t)$ to the sub-block of dimensions $\le\Delta_0$ coincides with
$Z^{(\le\Delta_0)}(t)$.
Consequently, the renormalized fields of dimension $\le\Delta_0$ are independent of the truncation level $\Delta_{\max}\ge \Delta_0$.
\end{lemma}

\begin{proof}
Triangularity in the engineering-dimension filtration implies that counterterms at dimension $\Delta$ only mix with dimensions $\le\Delta$.
Because the normalization conditions at dimension $\Delta\le\Delta_0$ reference only fields of dimensions $\le\Delta$, the recursion fixing the
counterterms up to $\Delta_0$ closes in the subsystem $\Delta\le\Delta_0$.
Uniqueness of the resulting linear system implies that enlarging the basis at higher dimensions cannot modify the already-fixed lower-dimensional blocks.
\end{proof}

\begin{theorem}[Universality for the full inductive-limit algebra $\mathfrak A_{\mathrm{loc}}$]
\label{thm:universality_full_inductive_limit}
Work within the local action class of Definition~\ref{def:action_universality_class} and along the flow renormalization scheme
defined by \eqref{eq:t0_def}.
Assume the uniform entrance/non-collapse hypothesis of Theorem~\ref{thm:uniqueness_on_Oell} for this action class.
Fix coherent normalization conditions as in Lemma~\ref{lem:coherent_Z}.
Then:
\begin{enumerate}[leftmargin=*]
\item the continuum limits of all finite collections of smeared fields in the inductive-limit algebra $\mathfrak A_{\mathrm{loc}}$ exist without
subsequence and are universal within the action class;
\item there exists $\lambda_*>1$ depending only on $(d,G)$ such that for any $\lambda>\lambda_*$ and any test function $\phi$, the dyadic increments
in Theorem~\ref{thm:field_algebra_truncation} are summable \emph{uniformly in $\Delta_{\max}$} in the weighted norm \eqref{eq:weighted_norm}.
\end{enumerate}
\end{theorem}

\begin{proof}
Fix finitely many elements of $\mathfrak A_{\mathrm{loc}}$.
Each lies in some truncation $\mathfrak A_{\mathrm{loc}}^{(\le\Delta_{\max})}$, so universality without subsequence follows from
Theorem~\ref{thm:universality_local_field_algebra}.
Coherence (Lemma~\ref{lem:coherent_Z}) ensures these joint laws do not depend on the choice of a larger truncation.

For the uniform-in-$\Delta_{\max}$ summability, apply Lemma~\ref{lem:one_shell_dimension_growth} in each dimension block and use the weighted norm.
Choose $\lambda>\Lambda_{\mathrm{shell}}$.
Then
\[\|\Delta_j\mathbf{\mathcal O}^{(\le\Delta_{\max})}(\phi)\|_{\lambda}\ \le\ C_{\phi}\,u_j^2\ \sup_{\Delta\le\Delta_{\max}}\Bigl(\frac{\Lambda_{\mathrm{shell}}}{\lambda}\Bigr)^{\Delta}\ \le\ C_{\phi}\,u_j^2.\]
Summability follows from $\sum_j u_j^2<\infty$ (Lemma~\ref{lem:dyadic_coupling_decay}).
\end{proof}

\subsubsection{Scheme change: finite renormalization and coupling reparametrization}

\begin{definition}[Admissible flow-based renormalization schemes]
\label{def:admissible_scheme}
An \emph{renormalization scheme} $\mathsf S$ consists of:
\begin{enumerate}[leftmargin=*]
\item a choice of lattice flow and energy density discretization entering \eqref{eq:wilson_flow}--\eqref{eq:flow_energy_def};
\item a scale-setting condition of the form $t_0^2\langle E(t_0)\rangle=c_0$ defining the trajectory $a(\beta)=\ell_0/\sqrt{t_0(\beta)}$;
\item a coherent family of normalization conditions (Lemma~\ref{lem:coherent_Z}) defining triangular renormalization matrices
$Z^{(\le\Delta)}_{\mathsf S}(t)$ for all $\Delta$.
\end{enumerate}
We restrict to schemes for which the weak-coupling curvature control holds uniformly (this is stable under local changes of the flow
discretization and energy density).
\end{definition}

\begin{lemma}[Finite renormalization matrix between two admissible schemes]
\label{lem:finite_renorm_matrix}
Let $\mathsf S$ and $\mathsf S'$ be two admissible schemes.
For each fixed $\Delta_{\max}$, define
\begin{equation}\label{eq:finite_renorm_def}
 R^{(\le\Delta_{\max})}
 :=\lim_{j\to\infty} Z^{(\le\Delta_{\max})}_{\mathsf S'}(t_j)\,\bigl(Z^{(\le\Delta_{\max})}_{\mathsf S}(t_j)\bigr)^{-1},
\end{equation}
where $t_j=2^{-2j}t_0$.
Then the limit exists in operator norm (and in the weighted norms of Definition~\ref{def:dimension_weighted_norm} for any $\lambda>\lambda_*$),
is invertible, and is triangular in the dimension filtration.
Moreover, $R^{(\le\Delta_{\max})}$ is coherent in $\Delta_{\max}$.
\end{lemma}

\begin{proof}
The ratio at dyadic scale $j$ is a product of triangular matrices with increments that are $O(u_j)$.
After imposing the respective normalization conditions in the two schemes, the difference in renormalized one-shell drifts is $O(u_j^2)$, so the
increment of the ratio is $O(u_j^2)$ in weighted operator norm.
Summability of $\sum_j u_j^2$ implies the ratio is Cauchy and converges to an invertible limit.
Triangularity and coherence follow from the same properties of each $Z^{(\le\Delta_{\max})}_{\mathsf S}$.
\end{proof}

\begin{lemma}[Coupling reparametrization between two admissible schemes]
\label{lem:coupling_reparam}
Let $\mathsf S$ and $\mathsf S'$ be two admissible schemes with flow couplings $u=g^2_{\mathrm{GF},\mathsf S}(L)$ and
$u'=g^2_{\mathrm{GF},\mathsf S'}(L)$ defined at the same physical scale $L$.
In the weak window there exists an analytic map $\varphi$ with
\begin{equation}\label{eq:coupling_reparam}
 u' = \varphi(u)=u+c_1 u^2+O(u^3),
\end{equation}
such that the step-scaling functions satisfy the conjugacy relation
\[\Sigma_{\mathsf S'}=\varphi\circ \Sigma_{\mathsf S}\circ \varphi^{-1}\]
on a neighborhood of $0$.
The coefficient $b_0$ of the $u^2\log 2$ term is scheme independent.
\end{lemma}

\begin{proof}
Both couplings are defined by flowed energy observables at the same physical scale.
In the weak window, each has an analytic expansion in the bare coupling (equivalently in $u$) with the same universal one-loop coefficient.
Therefore $u'$ is an analytic function of $u$ with the stated form.
The conjugacy relation follows by expressing $u'(2L)$ as a function of $u(L)$ in the two schemes and eliminating the intermediate variable.
Universality of $b_0$ is the standard statement that the one-loop beta function coefficient is scheme independent.
\end{proof}

\begin{theorem}[Scheme-change isomorphism of the continuum local field algebra]
\label{thm:scheme_change_isomorphism}
Let $\mathsf S$ and $\mathsf S'$ be two admissible schemes.
Assume the hypotheses of Theorem~\ref{thm:universality_full_inductive_limit}.
Then there exists a unique $*$--algebra isomorphism
\[\Phi_{\mathsf S\to\mathsf S'}:\ \mathfrak A_{\mathrm{loc},\mathsf S}\ \to\ \mathfrak A_{\mathrm{loc},\mathsf S'}\]
such that, for each finite truncation $\Delta_{\max}$ and each smeared generator,
\[
\Phi_{\mathsf S\to\mathsf S'}\bigl(\mathbf{\mathcal O}^{(\le\Delta_{\max})}_{\mathrm{ren},\mathsf S}(\phi)\bigr)
 = R^{(\le\Delta_{\max})}\,\mathbf{\mathcal O}^{(\le\Delta_{\max})}_{\mathrm{ren},\mathsf S'}(\phi),
\]
with $R^{(\le\Delta_{\max})}$ as in Lemma~\ref{lem:finite_renorm_matrix}.
Moreover, the corresponding running couplings are related by the analytic reparametrization $u'=\varphi(u)$ of
Lemma~\ref{lem:coupling_reparam}.
\end{theorem}

\begin{proof}
For each fixed $\Delta_{\max}$, the finite renormalization matrix $R^{(\le\Delta_{\max})}$ gives a linear change of generators of the truncated
algebra.
Coherence in $\Delta_{\max}$ implies these maps are compatible with the inductive system, hence define an isomorphism on
$\mathfrak A_{\mathrm{loc}}$.
The coupling reparametrization identifies the two flows in the weak window and preserves the universal coefficient $b_0$.
\end{proof}

\begin{corollary}[Scheme independence of the physical mass gap]
\label{cor:scheme_independent_gap}
Under the hypotheses of Theorem~\ref{thm:scheme_change_isomorphism}, the OS-reconstructed Hamiltonians in the two schemes are unitarily equivalent and
their spectral gaps coincide.
In particular, positivity of the mass gap is scheme independent within the admissible class.
\end{corollary}

\begin{proof}
The algebra isomorphism intertwines all Schwinger functions.
By the Osterwalder--Schrader reconstruction theorem, the reconstructed Hilbert space and Hamiltonian are unique up to unitary equivalence once the
Schwinger functions are fixed on a separating local algebra.
Therefore the spectra, and in particular the spectral gap, coincide.
\end{proof}


\subsection{Banach-space RG formulation and a stable-manifold theorem for exact blocking}
\label{sec:banach_rg_stable_manifold}

This section formalizes a standard constructive-RG viewpoint: represent the (gauge-fixed) effective action at a given scale by
\emph{one marginal/relevant coordinate} (the plaquette coupling) plus an \emph{irrelevant polymer remainder} living in a Banach space.
We then define the exact blocking map as an analytic self-map of this Banach space and prove an invariant-graph (stable foliation)
result by a graph-transform argument.

\medskip
\noindent\textbf{Purpose.}
The earlier modules supply (i) an entrance mechanism into a small neighborhood in a polymer norm (via deterministic majorantes and
$\chi^2$ transfer), and (ii) quantitative control of the coupling update (step scaling).
The only additional abstract ingredient needed to eliminate ``subsequence'' escape loopholes is a \emph{dynamical-systems} statement:
\emph{in a small neighborhood, the RG map contracts all irrelevant directions uniformly, hence irrelevant perturbations are slaved to the
running coupling.}

\subsubsection{Polymer activities and Banach norms}
\label{subsec:polymer_banach}

Fix an integer block factor $L\ge 2$.
Work on a finite $4$-torus $\Lambda$ (periodic boundary conditions) for definiteness; the estimates are uniform in $|\Lambda|$.
Let $\mathcal B_L$ be the partition of $\Lambda$ into disjoint $L$-blocks.
A \emph{polymer} $X$ is a finite connected union of $L$-blocks; write $|X|$ for the number of $L$-blocks in $X$.
Let $E(X)$ be the set of oriented fine edges contained in $X$.

We work in spanning-tree gauge on each $L$-block so that gauge orbits are coordinatized by the remaining ``free'' edges.
All functions below are gauge invariant; the gauge fixing is used only to define derivatives and norms.

\begin{definition}[Polymer activity]
A polymer activity is an assignment
\[K:\{X\ \text{connected polymer}\}\to \{K(X,\cdot):G^{E(X)}\to\mathbb R\}\]
with $K(\emptyset)=0$ and each $K(X,\cdot)$ smooth and gauge invariant.
\end{definition}

For $r\in\mathbb N$ define the local $C^r$ norm using left-invariant derivatives on $G$ in each edge variable:
\[
\|f\|_{C^r(X)}:=\max_{0\le m\le r}\ \max_{e_1,\dots,e_m\in E(X)}\ \sup_{U\in G^{E(X)}}\|\nabla_{e_1}\cdots\nabla_{e_m}f(U)\|.
\]
Fix a polymer-decay parameter $\kappa>0$ and define the Banach norm
\begin{equation}\label{eq:polymer_norm}
\|K\|_{r,\kappa}:=\sup_{B\in\mathcal B_L}\ \sum_{\substack{X\ni B\\ X\ \text{connected}}} e^{\kappa |X|}\ \|K(X)\|_{C^r(X)}.
\end{equation}
Let $\mathcal K_{r,\kappa}$ be the corresponding Banach space.

\begin{lemma}[Algebra property of the polymer norm]\label{lem:polymer_algebra}
There exists $C_{\times}=C_{\times}(r,\kappa,d)$ such that for any two polymer activities $K_1,K_2$ one has
\[
\|K_1\cdot K_2\|_{r,\kappa}\le C_{\times}\,\|K_1\|_{r,\kappa}\,\|K_2\|_{r,\kappa},
\]
where $(K_1\cdot K_2)(X)$ denotes the standard polymer product obtained by summing over decompositions $X=Y\cup Z$ with $Y,Z$ connected and
setting $(K_1\cdot K_2)(X)=\sum_{Y\cup Z=X}K_1(Y)K_2(Z)$.
\end{lemma}

\begin{proof}
This is standard.
For fixed base block $B$ and polymer $X\ni B$, the number of decompositions $X=Y\cup Z$ with $Y,Z$ connected is at most exponential in $|X|$;
absorb this into $e^{\kappa|X|}$ by reducing $\kappa$ (or by an explicit combinatorial factor).
Derivatives satisfy a Leibniz rule, giving a factor depending on $r$.
Summing over $X\ni B$ yields the stated bound.
\end{proof}

\subsubsection{Coordinates \texorpdfstring{$(u,K)$}{(u,K)} for effective actions and extraction}
\label{subsec:coordinates_extraction}

Let $A(U_p)$ be the standard plaquette energy density (normalized as in the earlier $\alpha$--pivot section).
For a given scale, represent a gauge-fixed effective measure on edges by
\begin{equation}\label{eq:measure_uK}
 d\mu_{u,K}(U)
 = \frac{1}{Z(u,K)}\exp\Bigl(-u^{-1}\,V(U)\Bigr)
 \exp\Bigl(\sum_{X\subset\Lambda\ \mathrm{conn.}} K(X,U_{E(X)})\Bigr)\,d\mu_0(U),
\end{equation}
where $d\mu_0(U)$ is a fixed gauge-invariant reference measure at that scale (e.g. the gauge-fixed heat-kernel/Gaussian reference from weak-coupling localization), and $V(U):=\sum_{p\in P(\Lambda)}A(U_p)$ is the extracted dimension-4 local functional.

\medskip
\noindent\textbf{Extraction convention.}
Write $\mathcal P$ for the projection onto the local subspace spanned by the constant functional and the plaquette functional
$\sum_{p}A(U_p)$.
We impose
\begin{equation}\label{eq:extraction_convention}
 \mathcal P K = 0,
\end{equation}
so that $u$ is the unique plaquette coupling coordinate.
(This is the standard ``local part extraction'' step in constructive RG.)

\subsubsection{Exact blocking map on \texorpdfstring{$(u,K)$}{(u,K)}}
\label{subsec:exact_blocking_map}

Let $\mathcal R_L$ denote the exact blocking transformation by factor $L$ followed by rescaling back to unit lattice spacing.
Operationally:
(i) choose a gauge-covariant block map $B_L$ from fine edges to coarse edges;
(ii) define the coarse marginal by integrating out fine degrees of freedom conditional on the coarse configuration;
(iii) re-expand the resulting coarse measure in the form \eqref{eq:measure_uK} and apply extraction \eqref{eq:extraction_convention}.
This yields
\[
\boxed{\ \mathcal R_L:(u,K)\mapsto(u',K').\ }
\]


\subsubsection{Explicit blocking map and one-step re-expansion}
\label{subsec:hardened46_blocking_map}

This subsection makes the RG step in Section~\ref{subsec:exact_blocking_map} completely explicit.
It fixes a canonical gauge-covariant block map $B_L$ and a smoothing kernel $T_{L,\alpha}$ so that the ``exact blocked measure''
is a well-defined probability measure on coarse links, and it records the concrete re-expansion procedure that produces
the extracted coordinates $(u',K')$.

\paragraph*{Block lattice and canonical path-ordered block map}

Let $\Lambda\subset\mathbb Z^d$ be a finite periodic box of side length divisible by $L\ge2$.
Let $\Lambda_L$ denote the coarse lattice with sites $x^L\in(\mathbb Z/L\mathbb Z)^d$ and edges $E(\Lambda_L)$.
For each coarse oriented edge $e^L=(x^L\to x^L+\hat\mu)$, define its canonical fine path
\[
\mathrm{path}_L(e^L):=\bigl\{(Lx^L+j\hat\mu\to Lx^L+(j+1)\hat\mu)\ :\ j=0,1,\dots,L-1\bigr\}\subset E(\Lambda).
\]
Define the canonical block map $B_L:\mathcal U_\Lambda\to\mathcal U_{\Lambda_L}$ by
\begin{equation}\label{eq:h46_block_map}
(B_LU)_{e^L}\ :=\ \prod_{j=0}^{L-1} U_{(Lx^L+j\hat\mu)\to(Lx^L+(j+1)\hat\mu)},
\end{equation}
where the product is ordered along the path from the tail to the head of $e^L$.

\begin{lemma}[Gauge covariance of the canonical block map]\label{lem:h46_block_map_covariant}
Let $g:\Lambda\to G$ be a fine-lattice gauge transformation and define the induced coarse transformation $g^L:\Lambda_L\to G$
by $g^L(x^L):=g(Lx^L)$. Then
\[
B_L(g\cdot U)\ =\ g^L\cdot (B_LU).
\]
\end{lemma}

\begin{proof}
For a fine edge $(x\to y)$, $(g\cdot U)_{x\to y}=g(x)U_{x\to y}g(y)^{-1}$.
Along the path $\mathrm{path}_L(e^L)$, the intermediate $g(\cdot)^{-1}g(\cdot)$ factors telescope, leaving
$g(Lx^L)$ on the left and $g(Lx^L+L\hat\mu)^{-1}=g(L(x^L+\hat\mu))^{-1}$ on the right.
\end{proof}

\paragraph*{Smoothing kernel and exact blocked measure}

To avoid distributional group delta constraints, we define the blocked measure using a positive-type smoothing kernel.
Fix $\alpha>0$ and let $K_\alpha:G\to(0,\infty)$ be the heat kernel at time $\alpha$ (normalized so $\int_G K_\alpha(g)\,dg=1$).
Define the gauge-invariant Markov kernel on coarse links by
\begin{equation}\label{eq:h46_markov_kernel}
T_{L,\alpha}(U^L\mid U)\ :=\ \prod_{e^L\in E(\Lambda_L)} K_\alpha\!\left(U^L_{e^L}\,\bigl((B_LU)_{e^L}\bigr)^{-1}\right).
\end{equation}
Normalization of $K_\alpha$ implies $\int T_{L,\alpha}(U^L\mid U)\,dU^L=1$ for each fixed $U$.

Given a fine-lattice gauge measure $\mu_\Lambda$ on $\mathcal U_\Lambda$, define the exact blocked measure on coarse links by
\begin{equation}\label{eq:h46_exact_blocked_measure}
\mu^{\mathrm{ex}}_{\Lambda_L}(dU^L)\ :=\ \int_{\mathcal U_\Lambda} \mu_\Lambda(dU)\,T_{L,\alpha}(U^L\mid U).
\end{equation}

\begin{lemma}[Gauge invariance of the blocked measure]\label{lem:h46_blocked_gauge_invariant}
If $\mu_\Lambda$ is gauge invariant, then $\mu^{\mathrm{ex}}_{\Lambda_L}$ is gauge invariant on $\Lambda_L$.
\end{lemma}

\begin{proof}
Gauge invariance of $\mu_\Lambda$ and Lemma~\ref{lem:h46_block_map_covariant} reduce the claim to conjugation invariance of the heat kernel:
$K_\alpha(hgh^{-1})=K_\alpha(g)$, which holds because $K_\alpha$ is a class function.
\end{proof}

\paragraph*{Exact RG map as re-expansion in the $(u,K)$ coordinates}

Let the fine measure have the $(u,K)$ form \eqref{eq:measure_uK} at unit lattice spacing.
Define the blocked (coarse) density by \eqref{eq:h46_exact_blocked_measure}, then rescale $\Lambda_L$ back to unit spacing
(by the canonical identification of $\Lambda_L$ with a unit lattice). Denote the resulting unit-spacing measure by
$\mathcal B_{L,\alpha}\mu$.

The \emph{exact RG map} $\mathcal R_L$ is the composition:
\[
(u,K)\ \xrightarrow{\ \text{form } \mu\ }\ \mu
\ \xrightarrow{\ \mathcal B_{L,\alpha}\ }\ \mathcal B_{L,\alpha}\mu
\ \xrightarrow{\ \text{extract local part}\ }\ (u',K'),
\]
where the last arrow is the fixed extraction convention \eqref{eq:extraction_convention} (local plaquette sector absorbed into $u'V$).

\paragraph*{Target one-step polymer decomposition}

The definitions above are unconditional; the nontrivial content is the \emph{quantitative} statement that the blocked measure
$\mathcal B_{L,\alpha}\mu$ can be re-expanded with a polymer remainder in the Banach norm uniformly in volume.
We record the needed statement as a named proposition.

\begin{proposition}[One-step polymer re-expansion of the exact blocked measure]\label{prop:h46_one_step_reexpansion}
Fix $(d,G,r,\kappa)$ and choose $L\ge2$.
There exist $u_0>0$, $\varepsilon_0>0$ and constants $C_u,C_K<\infty$ (depending only on $(d,G,L,r,\kappa,\alpha)$) such that:
whenever $|u|\le u_0$ and $\|K\|_{r,\kappa}\le\varepsilon_0$, the blocked-and-rescaled measure $\mathcal B_{L,\alpha}\mu$
admits a representation \eqref{eq:measure_uK} with extracted coordinates $(u',K')$ satisfying
\begin{align}
|u' - u - b_0 u^2\log L| &\le C_u\,u^3 + C_u\,u\,\|K\|_{r,\kappa},\label{eq:h46_u_update}\\
\|K'\|_{r,\kappa} &\le C_K\,(L^{-2}\|K\|_{r,\kappa} + u^2).\label{eq:h46_K_update}
\end{align}
Moreover, $(u',K')$ depends analytically on $(u,K)$ on the neighborhood $[-u_0,u_0]\times B_{\varepsilon_0}(0)$.
\end{proposition}


\subsubsection{One-step re-expansion proof (small/large field decomposition + polymer cluster expansion)}
\label{subsec:hardened47_one_step_proof}

This subsection completes the proof of Proposition~\ref{prop:h46_one_step_reexpansion} by supplying two missing ingredients:
\begin{enumerate}[leftmargin=*]
\item a \emph{uniform} small-field/large-field decomposition for the conditional block integration defining the exact blocked measure
\eqref{eq:h46_exact_blocked_measure}, and
\item a polymer/cluster expansion for the resulting effective interaction, yielding an explicit KP majorant bound on the blocked remainder.
\end{enumerate}
The argument is standard in constructive RG, but we record it here with explicit constants and clear parameter-dependence.
No new probabilistic hypotheses are introduced: the only analytic inputs are the already-proved weak-localization tail bound
(Lemma~\ref{lem:weak_localization}) and the curvature Brascamp--Lieb domination on the small-field chart
(Proposition~\ref{prop:curvature_BL}), together with the compactness/boundedness of the heat kernel $K_\alpha$.

\paragraph*{(1) Small-field chart, block gauge-fixing, and conditional measures}

Fix the chart radius $\varepsilon_0$ from Lemma~\ref{lem:weak_localization} so that on $\Omega_{\varepsilon_0}$ each plaquette holonomy admits a
unique logarithm $\xi_p=\log U_p\in\mathfrak g$ with $\|\xi_p\|\le \varepsilon_0$.
Fix once and for all a spanning tree $\mathsf T_B$ inside each $L$--block $B\in\mathcal B_L$ and impose tree gauge on $\mathsf T_B$ so that the
remaining ``free'' edges coordinatize the gauge orbit.

For a fixed coarse configuration $U^L\in G^{E(\Lambda_L)}$, define the fine conditional weight
\begin{equation}\label{eq:h47_conditional_weight}
W_{U^L}(U)\ :=\ \exp\bigl(-uV(U)\bigr)\,\prod_{X}(1+K(X,U_X))\;T_{L,\alpha}(U^L\mid U),
\end{equation}
and the corresponding conditional probability measure
\begin{equation}\label{eq:h47_conditional_measure}
\mu_{u,K}^{U^L}(dU)\ :=\ \frac{1}{Z(U^L)}\,W_{U^L}(U)\,dU,\qquad Z(U^L):=\int W_{U^L}(U)\,dU.
\end{equation}

\begin{lemma}[Uniform boundedness of the conditional density]\label{lem:h47_conditional_bounded}
Let $M_\alpha:=\|K_\alpha\|_{\infty}<\infty$.
Then for every $U^L$ and every $(u,K)$,
\[
0\le W_{U^L}(U)\ \le\ M_\alpha^{|E(\Lambda_L)|}\,\exp(-uV(U))\,\prod_X(1+|K(X,U_X)|).
\]
In particular, on the small polymer neighborhood $\|K\|_{r,\kappa}\le\varepsilon_0\le 1/4$ one has
$\prod_X(1+|K(X,U_X)|)\le \exp(C_K^{(0)}\,\varepsilon_0\,|\Lambda|)$ for an explicit $C_K^{(0)}(d,L,\kappa)$.
\end{lemma}

\begin{proof}
The heat kernel is bounded by $M_\alpha$ and the kernel \eqref{eq:h46_markov_kernel} is a product over coarse edges.
The second claim follows by the crude estimate
$\log(1+|K(X,U_X)|)\le 2|K(X,U_X)|$ for $|K|\le 1/2$, together with the polymer norm definition and the bound
$\sum_{X\ni B}\|K(X)\|_\infty\le \|K\|_{0,\kappa}\le \varepsilon_0$ for each block $B$.
Summing over the $O(|\Lambda|/L^4)$ blocks gives the stated volume factor.
\end{proof}

\paragraph*{(2) Large-field suppression at the block-integration level}

Let $\Omega_{\varepsilon_0}$ denote the global small-curvature event (all plaquettes in chart).
Write $\Omega_{\varepsilon_0}^c$ for its complement.

\begin{lemma}[Large-field tail for conditional block integration]\label{lem:h47_large_field_tail}
There exist constants $c_{\mathrm{LF}}>0$ and $\beta_{\mathrm{LF}}<\infty$ (depending only on $(d,G,\varepsilon_0)$) such that for all
$\beta\ge\beta_{\mathrm{LF}}$, for all $U^L$ and all $(u,K)$ in the small polymer neighborhood,
\begin{equation}\label{eq:h47_large_field_tail}
\mu_{u,K}^{U^L}\bigl(\Omega_{\varepsilon_0}^c\bigr)\ \le\ C_{\mathrm{LF}}\,|\Lambda|\,e^{-c_{\mathrm{LF}}\beta},
\end{equation}
with $C_{\mathrm{LF}}$ depending only on $(d,G,L,\alpha)$.
Consequently, the contribution of $\Omega_{\varepsilon_0}^c$ to the blocked measure can be absorbed into the polymer remainder as an
error term of order $e^{-c_{\mathrm{LF}}\beta}$.
\end{lemma}

\begin{proof}
On $\Omega_{\varepsilon_0}^c$, at least one plaquette holonomy lies outside the normal neighborhood.
The Wilson plaquette energy satisfies a uniform positive lower bound away from identity, hence the action cost is at least
$c_{\mathrm{LF}}\beta$ per ``bad'' plaquette (after fixing normalizations once).
Lemma~\ref{lem:weak_localization} gives the corresponding unconditional bound
$\mu_{u,0}(\Omega_{\varepsilon_0}^c)\le |P(\Lambda)|e^{-c\beta}$.
By Lemma~\ref{lem:h47_conditional_bounded}, inserting the bounded kernel and the small $K$ perturbation changes probabilities only by a bounded
Radon--Nikodym factor on $\Omega_{\varepsilon_0}^c$; absorb this into $C_{\mathrm{LF}}$.
\end{proof}

\paragraph*{(3) Polymerization of the small-field conditional integral}

Work henceforth on $\Omega_{\varepsilon_0}$.
In block tree gauge, write the free edge variables in exponential coordinates $U_e=\exp(A_e)$ with $A_e\in\mathfrak g$ and
$\|A_e\|\le C(\varepsilon_0)$.
On $\Omega_{\varepsilon_0}$, the plaquette logarithms satisfy
$\xi_p = (dA)_p + \mathrm{Rem}_p(A)$ with $\|\mathrm{Rem}_p(A)\|\le C_{\mathrm{rem}}\|A\|^2$.

Let $\mathcal N$ be the Gaussian (quadratic) reference measure on the free edge variables induced by the Hessian lower bound
$m_-$ from \eqref{eq:V_hess_bounds}, so that on $\Omega_{\varepsilon_0}$ the conditional density is a perturbation of a strictly log-concave
Gaussian field on the block graph.

\begin{lemma}[Connected polymer activity bound for one RG step]\label{lem:h47_polymer_activity_bound}
Fix the block factor $L\ge 2$.
There exist constants $C_{\mathrm{pol}}$ and $\kappa_{\mathrm{pol}}>0$ (depending only on $(d,G,L,\varepsilon_0,r,\kappa,\alpha)$) such that the
small-field contribution to the blocked density admits a polymer product representation
\begin{equation}\label{eq:h47_polymer_product}
\frac{d\mathcal B_{L,\alpha}\mu_{u,K}}{d\mu^{\mathrm{loc}}_{u'}}(U^L)\Bigg|_{\Omega_{\varepsilon_0}}
\;=\;
\prod_{\substack{X\subset \mathcal B_L\\ X\ \mathrm{connected},\ |X|\ge 2}}
\bigl(1+R'_X(U^L_{E(X)})\bigr),
\end{equation}
where $\mu^{\mathrm{loc}}_{u'}$ is the extracted local reference measure at the coarse scale (one-plaquette sector absorbed into $u'V$),
and the polymer remainders satisfy the uniform activity bound
\begin{equation}\label{eq:h47_activity_bound}
\sup_{B\in\mathcal B_L}\ \sum_{X\ni B} e^{\kappa_{\mathrm{pol}}|X|}\ \|R'_X\|_{C^r(X)}
\ \le\ C_{\mathrm{pol}}\bigl(u^2 + |u|\,\|K\|_{r,\kappa}+\|K\|_{r,\kappa}^2\bigr).
\end{equation}
The constants are uniform in the volume and admissible boundary conditions.
\end{lemma}

\begin{proof}
We give a complete cluster/forest-expansion argument with explicit constants.

Throughout this proof we are in the extracted polymer neighborhood
$|u|\le u_0$ and $\|K\|_{r,\kappa}\le\varepsilon_0$ (with $u_0,\varepsilon_0>0$ chosen sufficiently small in the weak window),
and we work at fixed finite volume with admissible boundary conditions.
In particular, all subset/polymer regroupings below are purely algebraic (finite sums).
The polymer product representation is defined by the finite-volume linked-cluster (Ursell) formulas;
the smallness of $(u,\|K\|_{r,\kappa})$ is used only to obtain the quantitative activity bound
\eqref{eq:h47_activity_bound} with uniform constants.


For notational convenience set
\[
\mathfrak f := u^2 + |u|\,\|K\|_{r,\kappa}+\|K\|_{r,\kappa}^2.
\]

\emph{Step 1: local interaction decomposition on the small-field set.}
On $\Omega_{\varepsilon_0}$, after extracting the one-plaquette sector into the coarse local reference measure
$\mu^{\mathrm{loc}}_{u'}$, the remaining blocked density ratio can be written in the form
\begin{equation}\label{eq:h47_exp_sumW}
\frac{d\mathcal B_{L,\alpha}\mu_{u,K}}{d\mu^{\mathrm{loc}}_{u'}}(U^L)\Bigg|_{\Omega_{\varepsilon_0}}
\;=\;\exp\Big(\sum_{b\in\mathscr I} W_b(U^L)\Big),
\end{equation}
where $\mathscr I$ is a finite-range collection of \emph{interaction units} $b$ at the $L$-blocked scale
(boundary plaquettes and the finite-range couplings coming from $T_{L,\alpha}$), and each $W_b$ depends only on
$U^L$ restricted to at most $q=q(d,L)$ coarse blocks.
After the one-plaquette extraction, every $W_b$ is at least quadratic in the centered curvature coordinates on
$\Omega_{\varepsilon_0}$. Consequently, by the uniform curvature moment bounds on $\Omega_{\varepsilon_0}$
(Proposition~\ref{prop:curvature_BL}) and the product/regularity axioms of the local seminorms, there exist
constants $C_{W,0}$ and $C_{W,r}$ (depending only on $(d,G,L,\varepsilon_0,r,\kappa,\alpha)$) such that, uniformly in
$b$ and in the volume/boundary conditions,
\begin{equation}\label{eq:h47_W_bounds}
\sup_{\Omega_{\varepsilon_0}} |W_b|\ \le\ C_{W,0}\mathfrak f,
\qquad
\|W_b\|_{C^r(b)}\ \le\ C_{W,r}\mathfrak f.
\end{equation}
(Here $\|\cdot\|_{C^r(b)}$ denotes the $C^r$ seminorm in the block variables on the support of $b$.
The point is that, on $\Omega_{\varepsilon_0}$, the configuration space is compact in the chart and the interaction
units are uniformly local.)

\emph{Step 2: polymer expansion by regrouping monomials.}
For each interaction unit $b\in\mathscr I$ define the (bounded) local factor
\[
V_b := e^{W_b}-1.
\]
Then
\[
\exp\Big(\sum_{b\in\mathscr I} W_b\Big)
=\prod_{b\in\mathscr I} e^{W_b}
=\prod_{b\in\mathscr I} (1+V_b)
=\sum_{\Gamma\subset\mathscr I}\ \prod_{b\in\Gamma} V_b.
\]
Using $|e^z-1|\le e^{|z|}|z|$ and the standard Fa\`a di Bruno bound for derivatives of compositions, and the
uniform bounds \eqref{eq:h47_W_bounds} on $\Omega_{\varepsilon_0}$, there exist explicit constants
$C_{V,0}$ and $C_{V,r}$ (depending only on $(d,G,L,\varepsilon_0,r,\kappa,\alpha)$) such that, uniformly in $b$,
\begin{equation}\label{eq:h47_V_bounds}
\sup_{\Omega_{\varepsilon_0}} |V_b|\ \le\ C_{V,0}\mathfrak f,
\qquad
\|V_b\|_{C^r(b)}\ \le\ C_{V,r}\mathfrak f.
\end{equation}
Each monomial is supported on the union of coarse blocks touched by $\Gamma$; let $X(\Gamma)$ be this block support.
Regroup terms by $X$ and define the \emph{raw polymer weights}
\begin{equation}\label{eq:h47_raw_polymer}
A_X(U^L):=\sum_{\Gamma:\,X(\Gamma)=X}\ \prod_{b\in\Gamma} V_b(U^L),
\end{equation}
so that the density ratio is $1+\sum_X A_X$.
Because each interaction unit crosses a block boundary (after the one-plaquette extraction), we have $A_X\equiv 0$ for $|X|=1$.
Moreover, $A_X$ depends only on the restriction of $U^L$ to $E(X)$.

If $X$ is connected with $|X|=n$, then any $\Gamma$ with $X(\Gamma)=X$ must contain at least $n-1$ interaction units
(since the interaction units induce an edge set on the block-adjacency graph, and a connected $n$-vertex graph has at least $n-1$ edges).
Consequently, using \eqref{eq:h47_V_bounds} and the product/regularity axioms of the $C^r$ seminorms, there exists a constant
$C_A=C_A(d,G,L,\varepsilon_0,r,\kappa,\alpha)$ such that
\begin{equation}\label{eq:h47_A_bound}
\|A_X\|_{C^r(X)}\ \le\ (C_A\,\mathfrak f)^{|X|-1}
\qquad\text{for all connected }X\text{ with }|X|\ge 2.
\end{equation}


\emph{Step 3: connected extraction via the BKAR/forest formula.}
Apply the standard linked-cluster (Ursell) expansion to the polymer family $\{A_X\}$ to write
\begin{equation}\label{eq:h47_linked_cluster}
\log\Big(1+\sum_X A_X\Big)
=\sum_{\substack{X\ \mathrm{connected}\\ |X|\ge 2}} \Phi_X,
\end{equation}
where each $\Phi_X$ is supported on $X$ and is given by a finite-volume sum over forests (equivalently, spanning trees) via the
Brydges--Kennedy--Abdesselam--Rivasseau forest formula for Ursell coefficients.

A key combinatorial point is that for any finite family of polymers $X_1,\dots,X_m$ whose overlap graph is connected and whose union is $X$,
one has
\begin{equation}\label{eq:h47_union_ineq}
\sum_{j=1}^m (|X_j|-1)\ \ge\ |X|-1,
\end{equation}
since $|X|\le \sum_{j=1}^m |X_j|-(m-1)$ for a connected overlap graph.
Consequently, using \eqref{eq:h47_A_bound} and the product/regularity axioms of the $C^r$ seminorms, every cluster term contributing to $\Phi_X$
carries at least $|X|-1$ factors of $\mathfrak f$.

The forest formula bounds the Ursell coefficients by a sum over trees and derivative book-keeping at order $r$; this contributes only an explicit
combinatorial factor (absorbed below).
Therefore there exists a constant $C_\Phi=C_\Phi(d,G,L,\varepsilon_0,r,\kappa,\alpha)$ such that, for all connected $X$ with $|X|\ge 2$,
\begin{equation}\label{eq:h47_Phi_bound}
\|\Phi_X\|_{C^r(X)}\ \le\ \bigl(C_\Phi\,\mathfrak f\bigr)^{|X|-1}.
\end{equation}
(For instance, one may take $C_\Phi := e\,C_A\,C_{\mathrm{count}}^{\mathrm{tree}}(r)$, enlarging constants if desired.)

\emph{Step 4: product representation and activity bounds.}
Exponentiating \eqref{eq:h47_linked_cluster} gives
\[
\exp\Big(\sum_{\substack{X\ \mathrm{connected}\\ |X|\ge 2}} \Phi_X\Big)
=\prod_{\substack{X\ \mathrm{connected}\\ |X|\ge 2}} \exp(\Phi_X),
\]
so \eqref{eq:h47_polymer_product} holds with
\begin{equation}\label{eq:h47_Rprime_def}
1+R'_X := \exp(\Phi_X),
\qquad\text{i.e.}\qquad
R'_X := \exp(\Phi_X)-1.
\end{equation}
By shrinking the polymer neighborhood (i.e. choosing $u_0$ and $\varepsilon_0$ sufficiently small),
we may assume $C_\Phi\,\mathfrak f\le \tfrac12$. Since $|X|\ge 2$, this implies $\|\Phi_X\|_{C^0(X)}\le 1$ for all connected $X$.
Using $|e^z-1|\le e^{|z|}|z|$ and the standard Fa\`a di Bruno/Leibniz bounds for derivatives of compositions, there is a constant $C_R=C_R(r)\ge 1$
such that
\begin{equation}\label{eq:h47_Rprime_bound}
\|R'_X\|_{C^r(X)}
\ \le\ C_R\,e^{\|\Phi_X\|_{C^0(X)}}\,\|\Phi_X\|_{C^r(X)}
\ \le\ \bigl(C_{\Phi}'\,\mathfrak f\bigr)^{|X|-1},
\end{equation}
where one may take $C_{\Phi}':= (C_R e)\,C_\Phi$ (since $C_R e\ge 1$ and $|X|-1\ge 1$).

\emph{Step 5: polymer counting and the weighted sum bound.}
Let $N_n$ be the number of connected block-polymers $X$ with $|X|=n$ containing a fixed block $B$.
A standard lattice-animal bound gives $N_n\le (C_{\mathrm{count}}^{\mathrm{blk}})^n$ with $C_{\mathrm{count}}^{\mathrm{blk}}$
depending only on $d$.

Fix $\kappa_{\mathrm{pol}}>0$ (for instance $\kappa_{\mathrm{pol}}=1$).
Then for each block $B$, using \eqref{eq:h47_Rprime_bound},
\[
\sum_{X\ni B} e^{\kappa_{\mathrm{pol}}|X|}\,\|R'_X\|_{C^r(X)}
\le \sum_{n\ge 2} (C_{\mathrm{count}}^{\mathrm{blk}})^n\,e^{\kappa_{\mathrm{pol}}n}\,\bigl(C_{\Phi}'\,\mathfrak f\bigr)^{n-1}.
\]
Write $m=n-2$ and factor out the $n=2$ contribution. This is
\[
\le (C_{\mathrm{count}}^{\mathrm{blk}}e^{\kappa_{\mathrm{pol}}})^2\,(C_{\Phi}'\,\mathfrak f)\,
\sum_{m\ge 0}\bigl(C_{\mathrm{count}}^{\mathrm{blk}}e^{\kappa_{\mathrm{pol}}}C_{\Phi}'\,\mathfrak f\bigr)^m.
\]
Choose the weak-window constants $u_0$ and $\varepsilon_0$ so that
\begin{equation}\label{eq:h47_series_small}
C_{\mathrm{count}}^{\mathrm{blk}}e^{\kappa_{\mathrm{pol}}}C_{\Phi}'\,\mathfrak f\ \le\ \tfrac12
\qquad\text{throughout the polymer neighborhood.}
\end{equation}
Then the geometric series is bounded by $2$, and we obtain
\[
\sup_{B\in\mathcal B_L}\ \sum_{X\ni B} e^{\kappa_{\mathrm{pol}}|X|}\,\|R'_X\|_{C^r(X)}
\ \le\ 2\,(C_{\mathrm{count}}^{\mathrm{blk}}e^{\kappa_{\mathrm{pol}}})^2\,C_{\Phi}'\,\mathfrak f.
\]
This yields \eqref{eq:h47_activity_bound} upon defining $C_{\mathrm{pol}}:=2\,(C_{\mathrm{count}}^{\mathrm{blk}}e^{\kappa_{\mathrm{pol}}})^2\,C_{\Phi}'$.
Uniformity in the volume and admissible boundary conditions follows because all bounds were local on $\Omega_{\varepsilon_0}$ and the
counting bound is uniform.
\end{proof}


\paragraph*{(4) Completion of Proposition~\ref{prop:h46_one_step_reexpansion}}

\begin{proof}[Proof of Proposition~\ref{prop:h46_one_step_reexpansion}]
Fix $(d,G,L,r,\kappa,\alpha)$.

\emph{Step 1: split into small-field and large-field parts.}
Write the blocked density as
\[
\mu^{\mathrm{ex}}_{\Lambda_L}(dU^L)
=\int_{\Omega_{\varepsilon_0}} \mu_{u,K}(dU)\,T_{L,\alpha}(U^L\mid U)
\;+\;
\int_{\Omega_{\varepsilon_0}^c} \mu_{u,K}(dU)\,T_{L,\alpha}(U^L\mid U).
\]
By Lemma~\ref{lem:h47_large_field_tail}, the second term contributes an additive error to the blocked density bounded by
$O(|\Lambda|e^{-c_{\mathrm{LF}}\beta})$ uniformly in $U^L$.
Choose $\beta$ (equivalently, choose $u_0$ small in the weak window) so that $|\Lambda|e^{-c_{\mathrm{LF}}\beta}\le u^3$
on the neighborhood of interest; absorb this into the $C_u u^3$ term in \eqref{eq:h46_u_update} and into the $u^2$ source term
in \eqref{eq:h46_K_update} after enlarging constants.

\emph{Step 2: small-field polymer product representation.}
On $\Omega_{\varepsilon_0}$, Lemma~\ref{lem:h47_polymer_activity_bound} gives the product form \eqref{eq:h47_polymer_product}
relative to an extracted local reference $\mu^{\mathrm{loc}}_{u'}$ and an activity bound \eqref{eq:h47_activity_bound}.
Choose $u_0,\varepsilon_0$ small so that the RHS of \eqref{eq:h47_activity_bound} is $\le 1/4$.
Then the blocked remainder majorant $\qkp_{\mathrm{rem}}$ defined by \eqref{eq:qrem_def} satisfies
$\qkp_{\mathrm{rem}}\le C_{\mathrm{pol}}(u^2+|u|\|K\|+\|K\|^2)$.
Lemma~\ref{lem:qrem_to_banach} converts this into a Banach-norm polymer activity $K'$ with
\[
\|K'\|_{r,\kappa}\ \le\ C_r\,\qkp_{\mathrm{rem}}
\ \le\ C_r C_{\mathrm{pol}}\bigl(u^2+|u|\,\|K\|_{r,\kappa}+\|K\|_{r,\kappa}^2\bigr).
\]
Shrink $\varepsilon_0$ further so that $\|K\|^2\le \|K\|$ on the neighborhood (e.g. $\varepsilon_0\le 1$) and absorb the mixed and quadratic
terms into a constant multiple of $u^2+ \|K\|$.
Finally, the engineering-dimension rescaling bound (Lemma~\ref{lem:rescaling_dim_bound}) yields the explicit $L^{-2}$ factor on the inherited
irrelevant part of $K$, giving \eqref{eq:h46_K_update} after enlarging $C_K$.

\emph{Step 3: coupling update.}
By construction, $u'$ is the extracted coefficient of the one-plaquette sector in the blocked local reference.
At $K=0$, the weak-window step-scaling control (Proposition~\ref{prop:stepscaling_control_proved}) identifies the one-loop coefficient
$b_0$ and yields
$u'(u,0)=u+b_0u^2\log L+O(u^3)$ uniformly.
The linear dependence of $u'$ on $K$ is controlled by the same polymer activity bound \eqref{eq:h47_activity_bound} (a single insertion of
$K$ changes the extracted one-plaquette coefficient by at most $C|u|\|K\|$), yielding \eqref{eq:h46_u_update}.

\emph{Step 4: analyticity.}
On $\Omega_{\varepsilon_0}$ the integrand is real-analytic in the exponential coordinates and the polymer sums converge absolutely under
the smallness condition; on $\Omega_{\varepsilon_0}^c$ the contribution is exponentially small and analytic by dominated convergence.
Thus $(u',K')$ depends analytically on $(u,K)$ on the neighborhood.

This proves the proposition.
\end{proof}


\subsubsection{From KP-small polymer remainders to the Banach polymer norm}
\label{subsec:kp_to_banach_polymer}

The stable-manifold theorem below is stated in a polymer Banach norm \eqref{eq:polymer_norm}.
Earlier in the paper, measure-level stability inside the KP basin is formulated in terms of the remainder majorant
$\qkp_{\mathrm{rem}}$ defined from the sup-activity bounds $\rho(P)$ in
Definition~\ref{def:poly_remainder}.
This subsection supplies the bridge:
\begin{quote}
\emph{KP-smallness of the remainder (in the sense $\qkp_{\mathrm{rem}}<1$) implies smallness of the corresponding polymer activity in the
Banach norm $\|\cdot\|_{r,\kappa}$, after passing from the product form to an exponential polymer action.}
\end{quote}

\medskip
\noindent
\textbf{Standing simplification.}
As in Section~\ref{sec:measure_level}, we assume the one-plaquette sector has been absorbed into the local reference measure
$\mu^{\mathrm{loc}}$, i.e. $\rho(P)=0$ for $|P|=1$.

\begin{lemma}[From $\qkp_{\mathrm{rem}}$ to a Banach-norm polymer activity]
\label{lem:qrem_to_banach}
Assume $\qkp_{\mathrm{rem}}<1/4$ in \eqref{eq:qrem_def} and $\rho(P)=0$ for $|P|=1$.
Define the logarithmic polymer activity
\begin{equation}\label{eq:K_from_R}
K(P,U_P):=\log\bigl(1+R_P(U_P)\bigr),\qquad |P|\ge 2,
\end{equation}
and $K(P,\cdot)=0$ for $|P|\le 1$.
Then for each fixed derivative order $r\in\mathbb N$ there exist constants
\(C_r=C_r(G,r)\) and a polymer-decay parameter $\kappa=\kappa(\theta,r,G)>0$ such that $K\in\mathcal K_{r,\kappa}$ and
\begin{equation}\label{eq:Knorm_from_qrem}
\boxed{\ \|K\|_{r,\kappa}\ \le\ C_r\,\qkp_{\mathrm{rem}}.\ }
\end{equation}
The constants can be chosen uniformly in volume and boundary conditions.
\end{lemma}

\begin{proof}
Since $\qkp_{\mathrm{rem}}<1/4$, one has $\sup_P\rho(P)\le \qkp_{\mathrm{rem}}<1/4$ and therefore $\|R_P\|_{\infty}\le 1/4$.
For $|z|\le 1/2$, $|\log(1+z)|\le 2|z|$, hence
\begin{equation}\label{eq:log_bound}
\|K(P)\|_{\infty}\le 2\rho(P).
\end{equation}

To control derivatives, note that each $R_P$ is a smooth (indeed real-analytic) gauge-invariant function on the compact manifold
$G^{E(P)}$, obtained by finitely many Haar integrations of products of smooth single-plaquette weights.
Repeated Leibniz rules on products over the $O(|P|)$ local factors yield a uniform bound
\begin{equation}\label{eq:derivative_growth}
\|R_P\|_{C^r(P)}\le C_r^{|P|}\,\rho(P),
\end{equation}
where $C_r$ depends only on $(G,r)$.
By the chain rule for $\log(1+R_P)$ and the bound $\|R_P\|_{\infty}\le 1/4$, the same estimate holds for $K(P)$ up to adjusting $C_r$.

Choose $\kappa>0$ so that $e^{\kappa}C_r\le e^{\theta}$.
Then, for any base plaquette $p_0$,
\[
\sum_{P\ni p_0} e^{\kappa|P|}\,\|K(P)\|_{C^r(P)}
\le
\sum_{P\ni p_0} e^{\theta|P|}\,\rho(P)
\le\qkp_{\mathrm{rem}},
\]
which gives \eqref{eq:Knorm_from_qrem}.
\end{proof}

\begin{proposition}[Entrance into the small polymer neighborhood (Banach norm)]
\label{prop:small_polymer_neighborhood}
Fix $(d,G,b,\theta)$ and choose $r\ge 2$.
Assume the effective single-plaquette weight is positive type with activity $f$ and lies in the KP admissibility window
(Section~\ref{sec:kp_majorant}).
Assume further that the one-plaquette sector has been absorbed into the local reference $\mu^{\mathrm{loc}}$ (so $\rho(P)=0$ for $|P|=1$).
Then for all sufficiently small $f$ the exact blocked measure admits a Banach-space representation \eqref{eq:measure_uK} with
\begin{equation}\label{eq:small_polymer_bounds}
|u|\le u_0,\qquad \|K\|_{r,\kappa}\le \varepsilon_0,
\end{equation}
where one may take $u_0=O(f)$ and $\varepsilon_0=O(f^2)$.
In particular, for $f\le f_{\mathrm{adm}}$ the quadratic remainder majorant \eqref{eq:qrem_ge2} gives $\qkp_{\mathrm{rem}}=O(f^2)$ and
Lemma~\ref{lem:qrem_to_banach} yields an explicit $\varepsilon_0$.

Moreover, on the neighborhood \eqref{eq:small_polymer_bounds}, the exact blocking map $\mathcal R_L$ is analytic as a map
$[-u_0,u_0]\times B_{\varepsilon_0}(0)\subset\mathbb R\times\mathcal K_{r,\kappa}\to\mathbb R\times\mathcal K_{r,\kappa}$.
\end{proposition}

\begin{proof}
Inside the KP admissibility window, $f$ is small.
After absorbing the one-plaquette sector, the remainder majorant is quadratic in $f$ as in \eqref{eq:qrem_ge2}; hence
$\qkp_{\mathrm{rem}}=O(f^2)$.
Lemma~\ref{lem:qrem_to_banach} converts $\qkp_{\mathrm{rem}}$ into a Banach-norm bound $\|K\|_{r,\kappa}\le C_r\qkp_{\mathrm{rem}}=O(f^2)$.
The extracted coupling coordinate $u$ is the coefficient of the absorbed one-plaquette sector; in the activity parameterization it satisfies
$|u|=O(f)$.
This gives \eqref{eq:small_polymer_bounds}.

Analyticity of $\mathcal R_L$ is standard: in the polymer representation the blocked weights are obtained by finite Haar integrations of
analytic functions and the polymer sums converge absolutely when $\|K\|_{r,\kappa}$ is sufficiently small.
\end{proof}

\begin{corollary}[Entrance alignment to the invariant graph]
\label{cor:entrance_alignment_invariant_graph}
Let $h$ be the invariant graph from Theorem~\ref{thm:rg_stable_manifold}.
Under the hypotheses of Proposition~\ref{prop:small_polymer_neighborhood},
\[
\|K-h(u)\|_{r,\kappa}\le \|K\|_{r,\kappa}+\|h(u)\|_{r,\kappa}\le \varepsilon_0+C_h u^2,
\]
so in particular $\|K-h(u)\|_{r,\kappa}=O(u^2)$ when $u=O(f)$ and $\varepsilon_0=O(f^2)$.
\end{corollary}


\begin{lemma}[Local invertibility of the coupling update]\label{lem:coupling_update_invertible}
Let $u'(u,K)$ denote the first coordinate of $\mathcal R_L(u,K)$.
In the weak-coupling window where the step-scaling expansion holds (cf. Proposition~\ref{prop:stepscaling_control_proved}), one has on
$|u|\le u_0$, $\|K\|_{r,\kappa}\le\varepsilon_0$,
\[
\bigl|\partial_u u'(u,K)-1\bigr|\ +\ \|\partial_K u'(u,K)\|\ \le\ C\bigl(|u|+\|K\|_{r,\kappa}\bigr)
\]
for a constant $C$ depending only on $(d,G,L)$ and the chosen norms.
Consequently, for $u_0,\varepsilon_0$ sufficiently small, for every Lipschitz graph $h$ with $\sup\|h(u)\|\le\varepsilon_0$ the restricted map
$F_h(u):=u'(u,h(u))$ is a $C^1$ diffeomorphism of $[-u_0,u_0]$ onto its image and has a uniformly Lipschitz inverse.
\end{lemma}

\begin{proof}
The step-scaling control gives a representation
$u'(u,0)=u+\beta_L(u)$ with $\beta_L(u)=O(u^2\log L)$ and $\partial_u u'(u,0)=1+O(u)$.
The dependence on $K$ enters through cumulant corrections and is analytic with first derivative bounded by $O(|u|+\|K\|)$ in the small
polymer neighborhood (Proposition~\ref{prop:small_polymer_neighborhood}); this is the standard fact that inserting one polymer activity changes
only local conditional expectations.
Choosing $u_0,\varepsilon_0$ so that $C(u_0+\varepsilon_0)\le\tfrac12$ yields strict monotonicity and the inverse function theorem.
\end{proof}
\subsubsection{Irrelevant contraction: the \texorpdfstring{$L^{-2}$}{L-2} factor from engineering dimension}
\label{subsec:irrelevant_contraction}

The stable-manifold theorem requires a uniform contraction on the \emph{irrelevant} sector.
The contraction factor $L^{-2}$ in $d=4$ is the standard dimensional gap once the dimension-$\le 4$ local operators have been extracted.

\begin{definition}[Dimension filtration on local gauge-invariant actions]
A local gauge-invariant operator has \emph{engineering dimension} $\Delta$ if in the small-curvature chart it is a polynomial in curvature
$F$ and covariant derivatives $D$ with the usual assignment $\dim(F)=2$ and $\dim(D)=1$.
Let $\mathscr V_{\le 4}$ be the span of $\{\mathbf 1,\sum_p A(U_p)\}$.
Define the irrelevant subspace $\mathscr V_{\ge 6}$ to be the (closure of the) span of local gauge-invariant operators with
$\Delta\ge 6$.
\end{definition}

\begin{lemma}[Linear rescaling bound on irrelevant local terms]\label{lem:rescaling_dim_bound}
Let $\mathcal S_L$ be the rescaling map associated to blocking by factor $L$ (after extraction of $\mathscr V_{\le 4}$).
On the irrelevant sector one has
\[
\|\mathcal S_L\|_{\mathscr V_{\ge 6}}\le C_{\mathrm{dim}}\,L^{-2},
\]
where $C_{\mathrm{dim}}$ depends only on the choice of local basis and the fixed derivative order $r$ in \eqref{eq:polymer_norm}.
\end{lemma}

\begin{proof}
Write an irrelevant local functional as a finite sum $\sum_i c_i\,\mathcal O_i$ with each $\mathcal O_i$ of dimension $\Delta_i\ge 6$.
Under the change of lattice spacing $a\mapsto La$, the coefficient rescales as
$c_i\mapsto L^{4-\Delta_i}c_i$ by standard power counting in $d=4$.
Since $\Delta_i\ge 6$, $|L^{4-\Delta_i}|\le L^{-2}$.
The $C^r$ norm introduces only a fixed multiplicative constant because left-invariant derivatives commute with group translation and the
charts are fixed at each scale; absorb this into $C_{\mathrm{dim}}$.
\end{proof}

\begin{proposition}[Irrelevant linearization of $\mathcal R_L$ contracts by $O(L^{-2})$]\label{prop:irrelevant_linear_contraction}
Assume the small neighborhood of Proposition~\ref{prop:small_polymer_neighborhood}.
Let $D_K\mathcal R_L(u,0)$ denote the Fr\'echet derivative in the polymer direction at $K=0$.
Then there exists $C_{\mathrm{irr}}=C_{\mathrm{irr}}(r,\kappa,d,G)$ such that for all $|u|\le u_0$,
\[
\|D_K\mathcal R_L(u,0)\|_{\mathcal K_{r,\kappa}\to\mathcal K_{r,\kappa}}\le C_{\mathrm{irr}}\,L^{-2}.
\]
\end{proposition}

\begin{proof}
At $K=0$ the measure \eqref{eq:measure_uK} is the reference measure $\mu_0$ tilted by $uV$.
The derivative $D_K\mathcal R_L(u,0)$ corresponds to inserting a single polymer activity into the block integration and propagating it through
blocking and extraction.
Extraction removes the $\mathscr V_{\le 4}$ components, so only irrelevant local contributions survive in the linearization.
These contributions rescale according to Lemma~\ref{lem:rescaling_dim_bound}, yielding the factor $L^{-2}$.
The polymer norm \eqref{eq:polymer_norm} is stable under the finite-range dependence induced by a single block integration; the associated
combinatorial constants are absorbed into $C_{\mathrm{irr}}$.
\end{proof}

\subsubsection{Quadratic remainder bound in the polymer norm}
\label{subsec:quadratic_remainder}

\begin{proposition}[Nonlinear remainder is quadratic in the polymer norm]\label{prop:quadratic_remainder_polymer}
Under Proposition~\ref{prop:small_polymer_neighborhood}, there exists $C_{\Xi}$ such that for all $|u|\le u_0$ and
$\|K\|_{r,\kappa}\le \varepsilon_0$ one can write
\[
K' = D_K\mathcal R_L(u,0)\,K\ +\ \Xi(u,K),
\]
with
\[
\|\Xi(u,K)\|_{r,\kappa}
\le C_{\Xi}\bigl(u^2+|u|\,\|K\|_{r,\kappa}+\|K\|_{r,\kappa}^2\bigr).
\]
Moreover $\Xi$ is locally Lipschitz with the corresponding quadratic Lipschitz constant.
\end{proposition}

\begin{proof}
Expand the blocked partition function and conditional expectations in cumulants with respect to the $K=0$ reference measure.
The linear term is exactly $D_K\mathcal R_L(u,0)K$.
All higher cumulants contain at least two insertions of $K$ (or one insertion of $K$ and one insertion from the $u$-dependence), hence are
quadratic in the small parameters.
Bounds in the norm \eqref{eq:polymer_norm} follow from the standard polymer algebra property (Lemma~\ref{lem:polymer_algebra}) and
finite-range dependence of block integrations.
The $u^2$ term reflects the fact that after extraction the $u$-flow generates irrelevant operators at order $u^2$.
\end{proof}

\subsubsection{Stable-manifold theorem for \texorpdfstring{$\mathcal R_L$}{R\_L} (invariant graph)}
\label{subsec:stable_manifold_theorem}

We now state and prove a Banach-space invariant-graph theorem tailored to the RG map $\mathcal R_L$.
The coupling coordinate $u$ is marginally relevant (asymptotic freedom), so the result is a stable foliation only in the irrelevant sector.

\begin{theorem}[Invariant renormalized graph and stable foliation]\label{thm:rg_stable_manifold}
Assume Proposition~\ref{prop:small_polymer_neighborhood}.
Choose $L$ so large that $C_{\mathrm{irr}}L^{-2}\le \tfrac14$, where $C_{\mathrm{irr}}$ is from
Proposition~\ref{prop:irrelevant_linear_contraction}.
Then, for sufficiently small $u_0,\varepsilon_0$, there exists a unique analytic map
\[
 h:[-u_0,u_0]\to\mathcal K_{r,\kappa},\qquad h(0)=0,
\]
with $\|h(u)\|_{r,\kappa}\le C_h u^2$ such that the graph
\[
\mathcal M:=\{(u,h(u)):\ |u|\le u_0\}
\]
is invariant under $\mathcal R_L$.
Moreover, if $(u_n,K_n)=\mathcal R_L^n(u_0,K_0)$ remains in the neighborhood, then
\[
\|K_n-h(u_n)\|_{r,\kappa}\le 2^{-n}\,\|K_0-h(u_0)\|_{r,\kappa}.
\]
\end{theorem}

\begin{proof}
This is a standard graph-transform argument in the Banach space $\mathbb R\times\mathcal K_{r,\kappa}$.
Define the space $\mathcal H$ of Lipschitz maps $h$ with $h(0)=0$ and
$\sup_{|u|\le u_0}\|h(u)\|/u^2\le M$ and $\mathrm{Lip}(h)\le Mu_0$.
Given $h\in\mathcal H$, write the $u$-update along the graph as $u'=F_h(u)$ and denote its local inverse by $G_h$.
Define the graph transform
\[
(\mathcal T h)(u') := D_K\mathcal R_L(G_h(u'),0)\,h(G_h(u')) + \Xi(G_h(u'),h(G_h(u'))),
\]
where $\Xi$ is from Proposition~\ref{prop:quadratic_remainder_polymer}.

By Proposition~\ref{prop:irrelevant_linear_contraction}, the linear part has operator norm at most $\tfrac14$.
By Proposition~\ref{prop:quadratic_remainder_polymer}, the nonlinear remainder is $O(u^2)+O(u\|h\|)+O(\|h\|^2)$.
Choosing $u_0$ and $\varepsilon_0$ so that the quadratic Lipschitz constants are $\le \tfrac14$ makes $\mathcal T$ a contraction on
$(\mathcal H,\|\cdot\|_*)$.
Banach's fixed point theorem gives a unique fixed point $h$.
The contraction estimate for $K_n-h(u_n)$ follows by subtracting the invariance equation from the RG recursion and using the same linear/quadratic
bounds.
Analyticity follows from analyticity of $\mathcal R_L$ in the neighborhood.
\end{proof}

\begin{remark}[Interpretation]
The invariant graph $K=h(u)$ is the standard \emph{renormalized trajectory} near the weak-coupling (Gaussian) fixed point.
Irrelevant perturbations are exponentially damped under blocking at rate $\approx L^{-2}$ per step once the local extraction is imposed.
This is the mechanism behind universality and uniqueness of continuum limits once the running coupling is fixed by a renormalization condition.
\end{remark}


\subsubsection{Archival finite-depth weak-window branch (moved to supplement)}
\label{subsec:hardened48_weak_window_iteration}

The earlier finite-depth iterated Banach-space control up to a KP threshold is no longer part of the final one-shot Route~1 proof.
That historical material has been moved to the supplementary Route~1 archive.
The only retained coupling-identification statement used later in the tuned dyadic normalization is the explicit bare/flow comparison theorem recorded in the next subsection.

\subsubsection{Coupling identification between the bare coupling and the gradient-flow coupling}
\label{subsec:hardened49_coupling_identification}

The load-bearing tuned dyadic construction uses only an explicit comparison between the bare coupling $\beta^{-1}$
and the gradient-flow coupling $g^2_{\mathrm{GF}}$ at the same physical scale.
No inverse-function theorem on the invariant graph and no finite-depth iterated RG control are used in this subsection.

\medskip
\noindent\textbf{Bare-to-flow comparison.}
Fix the physical reference scale $\ell_0$ and evaluate the flowed observable directly in the bare Wilson measure at lattice spacing
$a=\ell_0 2^{-n}$ and box size $L(a)=\ell_0/a$.
The next theorem is the sole weak-window identification used later in the tuned dyadic normalization.
\begin{lemma}[Uniform coefficient majorant for the fixed-scale flowed energy]\label{lem:h49_flow_coeff_majorant}
Fix the physical reference scale $\ell_0$, a dyadic spacing $a=\ell_0 2^{-n}$, set $L=L(a)=\ell_0/a$, and let $t=cL^2$.
Under Lemma~\ref{lem:weak_localization} and Proposition~\ref{prop:curvature_BL}, there exist constants $A_{\mathrm{flow}},R_{\mathrm{flow}}<\infty$ and $\beta_{\mathrm{flow}}\ge\beta_1$, depending only on $(d,G,c,\varepsilon_0,\ell_0)$ and the chosen discretization of $E(t)$, such that on the small-curvature event $\Omega_{\varepsilon_0}$ one has the absolutely convergent expansion
\begin{equation}\label{eq:h49_E_poly}
 t^2E(t)=\mathcal Q_2(dA)+\sum_{m\ge 2}\mathcal Q_{2m}(dA),
\end{equation}
where each $\mathcal Q_{2m}$ is homogeneous of degree $2m$ in the curvature field and, uniformly in the dyadic box size $L$,
\begin{equation}\label{eq:h49_degree_majorant}
\big|\E_{\beta,\Lambda(L)}\big[\mathcal Q_{2m}(dA)\,\mathbf 1_{\Omega_{\varepsilon_0}}\big]\big|
\le
A_{\mathrm{flow}}\Big(\frac{R_{\mathrm{flow}}}{\beta}\Big)^m,
\qquad m\ge 1,
\end{equation}
for every $\beta\ge\beta_{\mathrm{flow}}$.
In particular, after enlarging $\beta_{\mathrm{flow}}$ so that $R_{\mathrm{flow}}/\beta_{\mathrm{flow}}\le \tfrac12$,
\begin{equation}\label{eq:h49_tail_majorant}
\Big|\E_{\beta,\Lambda(L)}\Big[\sum_{m\ge 2}\mathcal Q_{2m}(dA)\,\mathbf 1_{\Omega_{\varepsilon_0}}\Big]\Big|
\le
2A_{\mathrm{flow}}R_{\mathrm{flow}}^2\,\beta^{-2}.
\end{equation}
\end{lemma}

\begin{proof}
Step~2 of Lemma~\ref{lem:flow_energy_expansion} gives the analytic expansion \eqref{eq:h49_E_poly} with an $L$-independent analyticity radius because the ratio $t/L^2=c$ is fixed.
Thus there exist constants $A_0,R_0<\infty$, independent of $L$, such that the coefficient tensor of the degree-$2m$ piece is bounded by $A_0R_0^m$.
Every monomial expectation appearing in $\E[\mathcal Q_{2m}(dA)\mathbf 1_{\Omega_{\varepsilon_0}}]$ is a finite linear combination of joint cumulants of $2m$ linear curvature functionals.
By \eqref{eq:h44_cumulant_bound}, each such cumulant is bounded by $(2m-1)!(\sigma_\beta^2)^m$ times the product of the corresponding $L^2$ norms, with $\sigma_\beta^2=(\beta m_-)^{-1}$.
Absorbing the finitely many Wick/contraction patterns and the uniform coefficient bounds into new constants gives \eqref{eq:h49_degree_majorant} for suitable $A_{\mathrm{flow}},R_{\mathrm{flow}}$.
Summing the geometric series for $m\ge 2$ yields \eqref{eq:h49_tail_majorant} once $R_{\mathrm{flow}}/\beta\le\tfrac12$.
\end{proof}

\begin{lemma}[Bad-curvature complement for the fixed-scale flow observable]\label{lem:h49_bad_event_absorb}
Under Lemma~\ref{lem:weak_localization}, there exist constants $C_{\mathrm{bad}},c_{\mathrm{bad}}>0$, depending only on $(d,G,c,\varepsilon_0,\ell_0)$ and the chosen discretization of $E(t)$, such that for every dyadic $a=\ell_0 2^{-n}$ and every $\beta\ge\beta_{\mathrm{flow}}$,
\begin{equation}\label{eq:h49_bad_event_bound}
\big|\E_{\beta,\Lambda(L(a))}[t^2E(t)\,\mathbf 1_{\Omega_{\varepsilon_0}^c}]\big|
\le
C_{\mathrm{bad}}e^{-c_{\mathrm{bad}}\beta}.
\end{equation}
\end{lemma}

\begin{proof}
Lemma~\ref{lem:weak_localization} gives an exponentially small bound on $\Omega_{\varepsilon_0}^c$ uniformly in finite volume.
For the fixed discretization at flow time $t=cL(a)^2$, the observable $t^2E(t)$ is a bounded-degree local polynomial in the flowed links, hence on each finite box its absolute value is dominated by a polynomial of uniformly bounded degree in the local curvature variables.
Combining this polynomial bound with the exponential suppression of $\Omega_{\varepsilon_0}^c$ yields \eqref{eq:h49_bad_event_bound}; after enlarging $\beta_{\mathrm{flow}}$ if necessary, the right-hand side is also $O(\beta^{-2})$.
\end{proof}


\begin{proposition}[Quantitative decomposition of the fixed-scale flow coupling]\label{prop:h49_quantitative_decomposition}
Fix the physical reference scale $\ell_0$, a dyadic spacing $a=\ell_0 2^{-n}$, and $t=cL(a)^2$.
Under Lemma~\ref{lem:weak_localization}, Proposition~\ref{prop:curvature_BL}, Lemma~\ref{lem:h49_flow_coeff_majorant}, and Lemma~\ref{lem:h49_bad_event_absorb}, there exist constants $C_{\mathrm{good,GF}},C_{\mathrm{bad,GF}},c_{\mathrm{bad,GF}}>0$, depending only on $(d,G,c,\varepsilon_0,\ell_0)$ and the chosen discretization of $E(t)$, such that
\begin{equation}\label{eq:h49_quantitative_split}
\E_{\beta,\Lambda(L(a))}[t^2E(t)]
=
\beta^{-1}+R_{\mathrm{good}}(a,\beta)+R_{\mathrm{bad}}(a,\beta),
\end{equation}
with
\begin{equation}\label{eq:h49_quantitative_remainders}
|R_{\mathrm{good}}(a,\beta)|\le C_{\mathrm{good,GF}}\beta^{-2},
\qquad
|R_{\mathrm{bad}}(a,\beta)|\le C_{\mathrm{bad,GF}}e^{-c_{\mathrm{bad,GF}}\beta},
\end{equation}
uniformly in the dyadic box size $L(a)$.
\end{proposition}

\begin{proof}
By Lemma~\ref{lem:h49_flow_coeff_majorant}, on $\Omega_{\varepsilon_0}$ one has the absolutely convergent expansion
\[
 t^2E(t)=\mathcal Q_2(dA)+\sum_{m\ge 2}\mathcal Q_{2m}(dA),
\]
with the explicit tail bound \eqref{eq:h49_tail_majorant}.
The normalization in \eqref{eq:gGF_def} was chosen so that the full expectation of the quadratic term is exactly $\beta^{-1}$.
Therefore
\[
\E[t^2E(t)]-\beta^{-1}
=
\E\Big[\sum_{m\ge 2}\mathcal Q_{2m}(dA)\,\mathbf 1_{\Omega_{\varepsilon_0}}\Big]
+\E\Big[(t^2E(t)-\mathcal Q_2(dA))\,\mathbf 1_{\Omega_{\varepsilon_0}^c}\Big].
\]
The first term is bounded by \eqref{eq:h49_tail_majorant}; this gives $|R_{\mathrm{good}}(a,\beta)|\le 2A_{\mathrm{flow}}R_{\mathrm{flow}}^2\beta^{-2}$.
For the second term, both $t^2E(t)$ and the quadratic piece $\mathcal Q_2(dA)$ are bounded-degree local curvature polynomials at the fixed ratio $t/L(a)^2=c$; enlarging the constants in Lemma~\ref{lem:h49_bad_event_absorb} therefore yields
\[
\Big|\E\Big[(t^2E(t)-\mathcal Q_2(dA))\,\mathbf 1_{\Omega_{\varepsilon_0}^c}\Big]\Big|
\le C_{\mathrm{bad,GF}}e^{-c_{\mathrm{bad,GF}}\beta}.
\]
This is exactly \eqref{eq:h49_quantitative_split}--\eqref{eq:h49_quantitative_remainders}.
\end{proof}

\begin{theorem}[Explicit weak-window comparison of the bare and flow couplings]\label{thm:h49_bare_flow_window}
Fix the physical reference scale $\ell_0$ and, for each dyadic spacing $a=\ell_0 2^{-n}$, write
\[
L(a):=\ell_0/a.
\]
Under Lemma~\ref{lem:weak_localization}, Proposition~\ref{prop:curvature_BL}, Lemma~\ref{lem:h49_flow_coeff_majorant}, and Lemma~\ref{lem:h49_bad_event_absorb}, there exist $\beta_{\mathrm{GF}}\ge\beta_{\mathrm{flow}}$ and a constant $C_{\mathrm{GF}}<\infty$, depending only on $(d,G,c,\varepsilon_0,\ell_0)$ and the chosen discretization of $E(t)$, such that for every dyadic $a=\ell_0 2^{-n}$ and every $\beta\ge\beta_{\mathrm{GF}}$,
\begin{equation}\label{eq:h49_gf_beta_expansion}
 g_{\mathrm{GF}}^2(\ell_0;a,\beta)
 =
 \beta^{-1}+R_{\mathrm{GF}}(a,\beta),
 \qquad
 |R_{\mathrm{GF}}(a,\beta)|\le C_{\mathrm{GF}}\beta^{-2}.
\end{equation}
After enlarging $\beta_{\mathrm{GF}}$ if necessary so that $C_{\mathrm{GF}}\beta_{\mathrm{GF}}^{-1}\le\tfrac12$, one has the uniform comparison
\begin{equation}\label{eq:h49_gf_beta_comparison}
 \frac{1}{2\beta}
 \le
 g_{\mathrm{GF}}^2(\ell_0;a,\beta)
 \le
 \frac{3}{2\beta}
 \qquad (\beta\ge\beta_{\mathrm{GF}}).
\end{equation}
Consequently, for every target value $u\in\bigl(0,(2\beta_{\mathrm{GF}})^{-1}\bigr]$ and every dyadic $a=\ell_0 2^{-n}$, there exists at least one bare coupling
\begin{equation}\label{eq:h49_beta_selection_interval}
 \beta(a;u)\in\Bigl[\frac{1}{2u},\frac{3}{2u}\Bigr]
\end{equation}
such that
\begin{equation}\label{eq:h49_beta_selection_exact}
 g_{\mathrm{GF}}^2(\ell_0;a,\beta(a;u))=u.
\end{equation}
In particular, on the load-bearing tuned dyadic path one has $\beta(a;u)=\Theta(1/u)$ uniformly in $a$.

This theorem uses only the explicit bare-coupling expansion, the quantitative tail bounds \eqref{eq:h49_tail_majorant} and \eqref{eq:h49_bad_event_bound}, and continuity in $\beta$ at fixed finite $L$; no inverse-function theorem on the invariant graph is used.
\end{theorem}

\begin{proof}
Fix the physical scale $\ell_0$, write $L=L(a)=\ell_0/a$, and set $t=cL^2$.
Proposition~\ref{prop:h49_quantitative_decomposition} gives the uniform expansion
\[
\E_{\beta,\Lambda(L)}[t^2E(t)]
=
\beta^{-1}+R_{\mathrm{good}}(a,\beta)+R_{\mathrm{bad}}(a,\beta),
\]
with the explicit bounds
\[
|R_{\mathrm{good}}(a,\beta)|\le C_{\mathrm{good,GF}}\beta^{-2},
\qquad
|R_{\mathrm{bad}}(a,\beta)|\le C_{\mathrm{bad,GF}}e^{-c_{\mathrm{bad,GF}}\beta},
\]
uniformly in the dyadic spacing $a$.
Choose
\[
C_{\mathrm{GF}}:=C_{\mathrm{good,GF}}+1
\]
and enlarge $\beta_{\mathrm{GF}}$ so that
\[
C_{\mathrm{bad,GF}}e^{-c_{\mathrm{bad,GF}}\beta}\le \beta^{-2}
\qquad (\beta\ge \beta_{\mathrm{GF}}).
\]
Then \eqref{eq:h49_gf_beta_expansion} follows immediately.
After enlarging $\beta_{\mathrm{GF}}$ once more so that $C_{\mathrm{GF}}\beta_{\mathrm{GF}}^{-1}\le\tfrac12$, the two-sided comparison \eqref{eq:h49_gf_beta_comparison} follows by
\[
\beta^{-1}-C_{\mathrm{GF}}\beta^{-2}
\le g_{\mathrm{GF}}^2(\ell_0;a,\beta)
\le \beta^{-1}+C_{\mathrm{GF}}\beta^{-2}.
\]

For the selection statement, fix $u\le (2\beta_{\mathrm{GF}})^{-1}$ and set
\[
\beta_-:=\frac{1}{2u},\qquad \beta_+:=\frac{3}{2u}.
\]
Then $\beta_\pm\ge\beta_{\mathrm{GF}}$, and \eqref{eq:h49_gf_beta_comparison} gives
\[
 g_{\mathrm{GF}}^2(\ell_0;a,\beta_-)
 \ge
 \frac{1}{2\beta_-}=u,
 \qquad
 g_{\mathrm{GF}}^2(\ell_0;a,\beta_+)
 \le
 \frac{3}{2\beta_+}=u.
\]
For fixed finite $L$, the map $\beta\mapsto g_{\mathrm{GF}}^2(\ell_0;a,\beta)$ is continuous because it is the expectation of a bounded continuous observable against a finite-volume Gibbs density that depends continuously on $\beta$.
The intermediate-value theorem therefore supplies at least one $\beta(a;u)\in[\beta_-,\beta_+]$ satisfying \eqref{eq:h49_beta_selection_exact}, and the interval \eqref{eq:h49_beta_selection_interval} yields $\beta(a;u)=\Theta(1/u)$ uniformly in $a$.
\end{proof}

\begin{remark}[Role in the mass gap path]
Theorem~\ref{thm:h49_bare_flow_window} is the load-bearing weak-window identification used later in Proposition~\ref{prop:h59_uv_tuning_exists} and Corollary~\ref{cor:route1_bounded_physical_entrance}.
The older invariant-graph matching / bi-Lipschitz coordinate comparison and the iterate-level comparison between $(u_n,K_n)$ and the gradient-flow coupling have both been moved to the project supplement.
\end{remark}

\begin{definition}[Repaired weak-window package]\label{def:weak_window_package}
Let $\mathsf{WW}_{\mathrm{curv}}$ denote the theorem-level weak-window package consisting of
Lemma~\ref{lem:weak_localization}, Proposition~\ref{prop:curvature_BL},
Lemma~\ref{lem:flow_energy_expansion}, Proposition~\ref{prop:stepscaling_control_proved},
Lemma~\ref{lem:stepscaling_continuous}, and Theorem~\ref{thm:h49_bare_flow_window}.
These are the only weak-window statements retained on the live proof spine.
\end{definition}

\begin{theorem}[Weak-window sufficiency certificate]\label{thm:weak_window_sufficiency}
All live downstream weak-window uses in the manuscript factor through the repaired curvature-sector package
$\mathsf{WW}_{\mathrm{curv}}$ of Definition~\ref{def:weak_window_package} and through nothing stronger.
More precisely:
\begin{enumerate}[leftmargin=*]
\item the only weak-window export to Section~\ref{subsec:step_scaling} is the sign input $b_0>0$ from Theorem~\ref{thm:b0_coefficient_extraction} and Lemma~\ref{lem:b0_positive}; the dyadic beta-remainder proof in Section~\ref{sec:rg} is otherwise internal to the chart-truncated Lane~C map;
\item Proposition~\ref{prop:h59_uv_tuning_exists} uses only Proposition~\ref{prop:stepscaling_control_proved}, Lemma~\ref{lem:stepscaling_continuous}, and Theorem~\ref{thm:h49_bare_flow_window};
\item Corollary~\ref{cor:route1_bounded_physical_entrance} uses Proposition~\ref{prop:h59_uv_tuning_exists} together with Proposition~\ref{prop:h51_kbeta_opt};
\item Theorem~\ref{thm:h60_one_shot_entrance} and Theorem~\ref{thm:gluing_KP_entrance} do not use any weak-window convexity statement;
\item consequently the load-bearing Route~1 package Proposition~\ref{prop:route1_QE1}, Proposition~\ref{prop:route1_QE2}, Corollary~\ref{cor:route1_QE3_gap}, and Theorem~\ref{thm:route1_laneB_package} contains no hidden dependence on any edge-Hessian convexity hypothesis.
\end{enumerate}
Appendix~\ref{app:referee_audit} records the same dependency information in table form.
\end{theorem}

\begin{proof}
For the Section~\ref{subsec:step_scaling} export, inspect first Theorem~\ref{thm:b0_coefficient_extraction} and then Lemma~\ref{lem:b0_positive}.
The theorem uses only Lemma~\ref{lem:weak_localization}, Proposition~\ref{prop:curvature_BL}, and the lattice heat-kernel / parametrix inputs already isolated in Appendices~\ref{app:lattice_heat_kernel} and~\ref{app:lattice_parametrix}; Lemma~\ref{lem:b0_positive} is then an immediate positivity corollary.
The beta-remainder file in Section~\ref{sec:rg} states explicitly that the remainder estimate is internal to the chart-truncated Lane~C map,
with the same-scale Wilson-versus-patched return handled separately by Theorem~\ref{thm:Wilson_patch_transport}.
Item~(ii) is the explicit hypothesis list in Proposition~\ref{prop:h59_uv_tuning_exists}.
Item~(iii) is exactly the proof of Corollary~\ref{cor:route1_bounded_physical_entrance}.
For item~(iv), inspect the proofs of Theorem~\ref{thm:gluing_KP_entrance} and Theorem~\ref{thm:h60_one_shot_entrance}: they use the one-shot crossover/transfer package and never cite Lemma~\ref{lem:weak_localization}, Proposition~\ref{prop:curvature_BL}, or any edge-Hessian statement.
Finally, Proposition~\ref{prop:route1_QE1} cites Theorem~\ref{thm:h60_one_shot_entrance} together with Corollary~\ref{cor:route1_bounded_physical_entrance};
Proposition~\ref{prop:route1_QE2}, Corollary~\ref{cor:route1_QE3_gap}, and Theorem~\ref{thm:route1_laneB_package} then consume Proposition~\ref{prop:route1_QE1} and later Route~1 closure theorems.
This closes the dependency chain and leaves no remaining place where an edge-Hessian convexity hypothesis could re-enter.
\end{proof}

\subsection{Same-scale transfer lemma at the one-shot entrance scale}
\label{sec:global_gluing}

\begin{remark}[What is and is not load-bearing here]\label{rem:gluing_scope}
Only Theorem~\ref{thm:gluing_KP_entrance} is used in the final one-shot Route~1 proof.
The legacy entrance-scale comparison notes have been moved to the supplementary Route~1 archive and are not invoked in Theorem~\ref{thm:h60_one_shot_entrance}, Corollary~\ref{cor:route1_bounded_physical_entrance}, Proposition~\ref{prop:route1_blocked_representatives}, or Theorem~\ref{thm:route1_E2_internal}.
A referee auditing only the final load-bearing chain may read Theorem~\ref{thm:gluing_KP_entrance} and then skip directly to Section~\ref{sec:hardened60_one_shot}.
\end{remark}

\paragraph{Archival note.}
The older $L_{\mathrm{ent}}(\beta)$ comparison statements are no longer carried in the main Appendix~F.
The one-shot proof uses only the prescribed block size $L=B(\beta)$ appearing in Theorem~\ref{thm:h60_one_shot_entrance}.

This section retains one generic same-scale transfer lemma. The retained lemma says exactly: \emph{if at some prescribed block size $L$ the reference blocked specification is KP/Dobrushin-contracting and the Plan-A conditional $\chi^2$ gate holds for the exact-to-reference comparison, then the exact blocked specification is KP/Dobrushin-contracting at that same block size.}

\subsubsection{General transfer/gluing theorem at a prescribed block size}

This is a scale-local transfer theorem.
In the final load-bearing Route~1 chain it is applied at the one-shot entrance scale $L=B(\beta)$, so it does not define a second entrance mechanism.

\begin{theorem}[Transfer/gluing to the KP basin at a prescribed block size]\label{thm:gluing_KP_entrance}
Fix $d=4$ and compact $G$. Fix a block size $L\ge 2$.

Assume the $\alpha$--reference blocked specification at block size $L$ satisfies a uniform Dobrushin contraction bound
\[
\delta(\gamma^{(\alpha)})\le q_*<1
\]
uniformly in finite block volume and admissible boundary conditions.
Assume moreover that the exact-to-reference comparison at the same block size satisfies the Plan-A conditional $\chi^2$ gate
with deterministic defect $\varepsilon_L\ge 0$ (as produced in the $\chi^2$--transfer module), so that the single-site kernels obey
the influence diagnostic of Lemma~\ref{lem:chi2_to_influence} with error $\sqrt{\varepsilon_L}$.

Then the exact blocked specification satisfies
\[
\delta(\gamma^{\mathrm{ex}})\le q_*+\sqrt{\varepsilon_L}.
\]
In particular, if $q_*+\sqrt{\varepsilon_L}<1$, the infinite-volume exact blocked Gibbs state at block size $L$ exists and is unique, and
block-local truncated correlations decay exponentially in block distance with constants depending only on $(q_*,G)$.
\end{theorem}

\begin{proof}
By Lemma~\ref{lem:chi2_to_influence}, the assumed conditional $\chi^2$ gate implies the influence comparison
\[
\delta(\gamma^{\mathrm{ex}})
\le
\delta(\gamma^{(\alpha)})+\sqrt{\varepsilon_L}
\le q_*+\sqrt{\varepsilon_L}.
\]
If $q_*+\sqrt{\varepsilon_L}<1$, Proposition~\ref{prop:dobrushin_standard} applies directly to the exact blocked specification and yields the asserted infinite-volume existence, uniqueness, and exponential block-distance decay.
\end{proof}

\begin{corollary}[Infinite-volume uniqueness and boundary-condition uniformity]\label{cor:IV_uniqueness_uniformity}
If $\delta(\gamma^{\mathrm{ex}})\le q<1$ uniformly in finite volume and admissible boundary conditions, then the infinite-volume Gibbs/DLR state exists and is unique and block-local truncated correlations decay exponentially with constants depending only on $(q,G)$.
\end{corollary}

\begin{proof}
Immediate from Proposition~\ref{prop:dobrushin_standard}.
\end{proof}

\subsubsection{QE2$_{\mathrm{lat}}$ closure module: exact terminal-block transfer at the one-shot scale}

\medskip
\noindent
\textbf{Status.}
The phase-1 reduction isolated a missing fluctuation-remainder estimate in the law of total covariance.
At the exact one-shot scale that remainder can in fact be removed completely.
The key point is that the blocked sigma-algebra is the terminal skeleton sigma-algebra, and Corollary~\ref{cor:wilson_terminal_block_product} gives exact conditional independence of fine observables supported in disjoint entrance-block families once that sigma-algebra is frozen.
Combined with entrance-scale Dobrushin decay for the exact blocked specification, this yields QE2$_{\mathrm{lat}}$ directly.

\medskip
\noindent
The repaired implication now has the form
\[
\begin{aligned}
&\text{(one-shot block-scale Dobrushin contraction)}\\
&\qquad +\ \text{(exact terminal-block conditional independence)}\\
&\qquad \Longrightarrow\ \text{QE2$_{\mathrm{lat}}$}.
\end{aligned}
\]
The block-decay part is already available from the exact one-shot entrance theorem and Proposition~\ref{prop:dobrushin_standard}.
The second input comes from the exact Wilson recursion / terminal-block product structure and is already available elsewhere in the manuscript.

\begin{lemma}[Explicit block-distance decay from finite-range Dobrushin]\label{lem:dobrushin_to_block_decay}
Let $\gamma$ be a single-site specification on a graph with graph distance $\mathrm{dist}$ and assume the finite-range Dobrushin hypotheses of Proposition~\ref{prop:dobrushin_standard} with coefficient bound $\delta(\gamma)\le q<1$.
For a bounded local observable $F$, write
\[
\operatorname{osc}_i(F):=
\sup_{\eta,\eta':\ \eta_{V\setminus\{i\}}=\eta'_{V\setminus\{i\}}}
|F(\eta)-F(\eta')|,
\qquad
\supp_{\mathrm{osc}}(F):=\{i:\operatorname{osc}_i(F)\neq0\}.
\]
Then every pair of bounded local functions $F,G$ satisfies the support-weighted covariance estimate
\begin{equation}\label{eq:dobrushin_block_support_weighted}
\begin{aligned}
|\mathrm{Cov}(F,G)|
\le{}&\frac{e^{-m_{\mathrm{blk}}\,\mathrm{dist}(\supp_{\mathrm{osc}}(F),\supp_{\mathrm{osc}}(G))}}{1-q}
\Big(\sum_{i\in \supp_{\mathrm{osc}}(F)}\operatorname{osc}_i(F)\Big)\\
&\times\Big(\sum_{j\in \supp_{\mathrm{osc}}(G)}\operatorname{osc}_j(G)\Big),
\qquad m_{\mathrm{blk}}:=-\log q.
\end{aligned}
\end{equation}
In particular,
\begin{equation}\label{eq:dobrushin_block_cardinality}
|\mathrm{Cov}(F,G)|
\le
\frac{4\,|\supp_{\mathrm{osc}}(F)|\,|\supp_{\mathrm{osc}}(G)|}{1-q}
\|F\|_\infty\|G\|_\infty
 e^{-m_{\mathrm{blk}}\,\mathrm{dist}(\supp_{\mathrm{osc}}(F),\supp_{\mathrm{osc}}(G))}.
\end{equation}
\end{lemma}

\begin{proof}
This is exactly the support-weighted estimate \eqref{eq:dobrushin_support_weighted_bound} from Proposition~\ref{prop:dobrushin_standard}, specialized to the present graph and with $m_{\mathrm{blk}}:=-\log q$.
The cardinality estimate \eqref{eq:dobrushin_block_cardinality} follows from $\operatorname{osc}_i(F)\le 2\|F\|_\infty$ and $\operatorname{osc}_j(G)\le 2\|G\|_\infty$.
\end{proof}

\begin{remark}[Why the fluctuation remainder disappears at the one-shot scale]
The relevant blocked sigma-algebra at scale $B$ is the terminal skeleton sigma-algebra by Proposition~\ref{prop:route1_blocked_sigma_identification}.
For observables supported in disjoint families of entrance blocks, Corollary~\ref{cor:wilson_terminal_block_product} therefore gives exact conditional independence once that sigma-algebra is frozen.
Equivalently, the fluctuation covariance term in the law of total covariance vanishes identically for separated local probes at the exact one-shot scale.
\end{remark}

\begin{lemma}[Entrance-block geometry at a fixed one-shot lattice]\label{lem:route1_probe_geometry}
Fix $n$ and let $B_n$ be the one-shot block side on the tuned lattice $(a_n,\beta_n)$.
For a bounded local cylinder observable $H$, define
\[
\alpha_n(H):=\Bigl\lceil \frac{\operatorname{diam}_\infty(\supp H)}{B_n}\Bigr\rceil,
\qquad
N_n(H):=(\alpha_n(H)+2)^4.
\]
Then for every translate $H(x)$:
\begin{enumerate}[leftmargin=*]
\item the entrance-block hull satisfies
\[
|\mathsf B_n(H(x))|\le N_n(H);
\]
\item if $H_1,H_2$ are bounded local cylinder observables and
\[
r:=\dist_\infty(\supp H_1(x),\supp H_2(y)),
\]
then
\begin{equation}\label{eq:route1_probe_geometry_distance}
\dist_{\mathrm{blk}}\bigl(\mathsf B_n(H_1(x)),\mathsf B_n(H_2(y))\bigr)
\ge \frac{r}{2B_n}-\Gamma_n(H_1,H_2),
\end{equation}
where
\[
\Gamma_n(H_1,H_2):=\alpha_n(H_1)+\alpha_n(H_2)+4.
\]
\end{enumerate}
In particular, if
\[
r\ge 2B_n\bigl(\Gamma_n(H_1,H_2)+2\bigr),
\]
the two entrance-block hulls are disjoint.
\end{lemma}

\begin{proof}
Fix $H$.
In each coordinate direction, the projection of $\supp H(x)$ meets at most $\alpha_n(H)+2$ consecutive $B_n$-blocks, so the full entrance-block hull contains at most $(\alpha_n(H)+2)^4=N_n(H)$ blocks.
This proves item~(1).

For item~(2), let $d_{\infty}^{\mathrm{idx}}$ be the $\ell^\infty$-distance between the corresponding entrance-block index hulls.
The support of $H_i$ can enlarge the index hull by at most $\alpha_n(H_i)+1$ block units in each coordinate, so
\[
d_{\infty}^{\mathrm{idx}}
\ge \frac{r}{B_n}-\alpha_n(H_1)-\alpha_n(H_2)-2.
\]
One step in the interaction graph of Definition~\ref{def:route1_block_graph} changes the block index by at most $2$ in $\ell^\infty$, hence
\[
\dist_{\mathrm{blk}}\bigl(\mathsf B_n(H_1(x)),\mathsf B_n(H_2(y))\bigr)
\ge \frac12\bigl(d_{\infty}^{\mathrm{idx}}-2\bigr)
\ge \frac{r}{2B_n}-\Gamma_n(H_1,H_2),
\]
which is exactly \eqref{eq:route1_probe_geometry_distance}.
If $r\ge 2B_n(\Gamma_n(H_1,H_2)+2)$, the right-hand side is at least $2$, so the two hulls are disjoint.
\end{proof}

\begin{lemma}[Fixed-lattice block-to-microscopic covariance transfer with common rate]\label{lem:block_decay_to_micro}
Fix a target value $u_{\mathrm{ren}}\in(0,u_\ast]$ and let $a_n:=\ell_0 2^{-n}$, $\beta_n:=\beta(a_n;u_{\mathrm{ren}})$, and $B_n:=B(\beta_n)$.
Then there exists $n_{\mathrm{mic}}(u_{\mathrm{ren}})<\infty$ such that for every $n\ge n_{\mathrm{mic}}(u_{\mathrm{ren}})$, every pair of bounded gauge-invariant local cylinder observables $H_1,H_2$, and all translates $x,y$ with microscopic support separation
\[
r:=\dist_\infty(\supp H_1(x),\supp H_2(y)),
\]
one has
\begin{equation}\label{eq:block_to_micro_exact_transfer}
\big|\mathrm{Cov}_{\omega_{a_n}}\big(H_1(x),H_2(y)\big)\big|
\le C_n(H_1,H_2)\,
\exp\!\Big(-\frac{m_{\mathrm{blk}}}{2B_n}\,r\Big),
\qquad m_{\mathrm{blk}}:=-\log q_\ast,
\end{equation}
for some finite constant $C_n(H_1,H_2)$ depending on $n,H_1,H_2$ but not on the translates.
In particular, for each fixed $n\ge n_{\mathrm{mic}}(u_{\mathrm{ren}})$ the rate $m_{\mathrm{blk}}/(2B_n)$ applies simultaneously to \emph{every} centered bounded gauge-invariant local observable on the fixed-lattice positive-time local algebra.
\end{lemma}

\begin{proof}
Set
\[
n_{\mathrm{mic}}(u_{\mathrm{ren}}):=
\max\bigl\{N_{\mathrm{ent}}(u_{\mathrm{ren}}),\,n_{\mathrm{repr}}(u_{\mathrm{ren}})\bigr\},
\]
where $N_{\mathrm{ent}}(u_{\mathrm{ren}})$ is the entry index from \Cref{prop:route1_QE1} and $n_{\mathrm{repr}}(u_{\mathrm{ren}})$ comes from Proposition~\ref{prop:route1_blocked_representatives}.
Fix $n\ge n_{\mathrm{mic}}(u_{\mathrm{ren}})$ and write
\[
\alpha_i:=\alpha_n(H_i),\qquad
N_i:=N_n(H_i),\qquad
\Gamma:=\Gamma_n(H_1,H_2).
\]
Let $\Phi_{H_i(x),n}$ be the exact blocked representatives furnished by Proposition~\ref{prop:route1_blocked_representatives}.
By item~(1) of that proposition and Lemma~\ref{lem:route1_probe_geometry},
\[
|\supp_{\mathrm{osc}}(\Phi_{H_i(x),n})|
\le N_i,
\qquad
\|\Phi_{H_i(x),n}\|_\infty\le \|H_i\|_\infty.
\]

If
\[
r\ge 2B_n(\Gamma+2),
\]
the entrance-block hulls are disjoint by Lemma~\ref{lem:route1_probe_geometry}; hence Proposition~\ref{prop:route1_blocked_representatives} gives the exact covariance identity
\[
\mathrm{Cov}_{\omega_{a_n}}\big(H_1(x),H_2(y)\big)
=
\mathrm{Cov}_{\omega_{a_n}^{\langle B_n\rangle}}\big(\Phi_{H_1(x),n},\Phi_{H_2(y),n}\big).
\]
Apply Proposition~\ref{prop:route1_block_covariance} together with \eqref{eq:route1_probe_geometry_distance} to obtain
\begin{align*}
\big|\mathrm{Cov}_{\omega_{a_n}}\big(H_1(x),H_2(y)\big)\big|
&\le
\frac{4N_1N_2}{1-q_\ast}\,
\|H_1\|_\infty\|H_2\|_\infty
e^{-m_{\mathrm{blk}}\,\dist_{\mathrm{blk}}(\mathsf B_n(H_1(x)),\mathsf B_n(H_2(y)))}\\
&\le
\frac{4N_1N_2}{1-q_\ast}\,
\|H_1\|_\infty\|H_2\|_\infty
e^{m_{\mathrm{blk}}\Gamma}\,
e^{-m_{\mathrm{blk}}r/(2B_n)}.
\end{align*}

If instead $r<2B_n(\Gamma+2)$, use the trivial bound
\[
\big|\mathrm{Cov}_{\omega_{a_n}}\big(H_1(x),H_2(y)\big)\big|
\le 4\|H_1\|_\infty\|H_2\|_\infty
\le 4e^{m_{\mathrm{blk}}(\Gamma+2)}\|H_1\|_\infty\|H_2\|_\infty e^{-m_{\mathrm{blk}}r/(2B_n)}.
\]
Combining the two regimes yields \eqref{eq:block_to_micro_exact_transfer} after setting
\[
C_n(H_1,H_2):=
4\|H_1\|_\infty\|H_2\|_\infty
\max\!\Bigl\{
e^{m_{\mathrm{blk}}(\Gamma+2)},
\frac{N_1N_2}{1-q_\ast}\,e^{m_{\mathrm{blk}}\Gamma}
\Bigr\}.
\]
Because $n$ is now fixed, this applies to every bounded local pair $H_1,H_2$ on that lattice; only the prefactor changes with the observables.
\end{proof}

\begin{corollary}[Fixed-lattice common-rate decay criterion]\label{cor:fixed_n_decay_criterion}
Fix $n\ge n_{\mathrm{mic}}(u_{\mathrm{ren}})$.
Then for every centered bounded gauge-invariant local cylinder observable $H$ on the fixed-lattice positive-time algebra there exists $C_{n,H}<\infty$ such that
\[
0\le \omega_{a_n}\bigl(\Theta H\,\tau_t H\bigr)\le C_{n,H}\,e^{-m_{\mathrm{blk}}t/(2B_n)}
\qquad (t\in\mathbb N).
\]
Hence the hypotheses of Theorem~\ref{thm:F_decay_implies_gap} hold at the fixed lattice spacing $a_n$ with common decay rate $m_{\mathrm{blk}}/(2B_n)$.
\end{corollary}

\begin{proof}
Let $H$ be an arbitrary centered bounded gauge-invariant local cylinder observable in the fixed-lattice positive-time algebra.
Because $H$ has finite temporal diameter, there exists $c_H<\infty$ such that the microscopic support separation between $\Theta H$ and $\tau_t H$ is at least $t-c_H$ for every integer $t\ge 0$.
Apply Lemma~\ref{lem:block_decay_to_micro} to the pair $(\Theta H,H)$ and absorb the finite shift $c_H$ into the prefactor:
\[
\big|\omega_{a_n}(\Theta H\,\tau_t H)\big|
\le C_n(\Theta H,H)e^{m_{\mathrm{blk}}c_H/(2B_n)}e^{-m_{\mathrm{blk}}t/(2B_n)}.
\]
Reflection positivity gives the nonnegativity of the left-hand side, so the displayed bound is exactly the required fixed-lattice decay estimate.
Because $H$ was arbitrary, Theorem~\ref{thm:F_decay_implies_gap} applies with $m=m_{\mathrm{blk}}/(2B_n)$.
\end{proof}

\begin{corollary}[Uniform full fixed-lattice gap from bounded one-shot entrance]\label{cor:uniform_phys_gap_from_bounded_entrance}
Fix a target value $u_{\mathrm{ren}}\in(0,u_\ast]$ and let $a_n=\ell_0 2^{-n}$, $\beta_n=\beta(a_n;u_{\mathrm{ren}})$, and $B_n=B(\beta_n)$.
Assume the standing one-shot entrance hypotheses, so that the exact blocked specification has Dobrushin constant $\le q_\ast<1$ at scale $B_n$ and $a_n B_n\le C_{\mathrm{ent}}\ell_0$ for all large $n$.
Then for all sufficiently large $n$, the fixed-lattice OS gap satisfies
\[
\Delta_{\mathrm{lat}}(\beta_n)\ge \frac{m_{\mathrm{blk}}}{2B_n},
\qquad
\frac{\Delta_{\mathrm{lat}}(\beta_n)}{a_n}\ge \frac{m_{\mathrm{blk}}}{2C_{\mathrm{ent}}\ell_0},
\]
with $m_{\mathrm{blk}}:=-\log q_\ast$.
\end{corollary}

\begin{proof}
Fix $n\ge n_{\mathrm{mic}}(u_{\mathrm{ren}})$.
By Corollary~\ref{cor:fixed_n_decay_criterion}, every centered bounded gauge-invariant local observable on the fixed-lattice positive-time algebra satisfies the hypothesis of Theorem~\ref{thm:F_decay_implies_gap} with the same decay rate $m_{\mathrm{blk}}/(2B_n)$.
Proposition~\ref{prop:lattice_OS_cyclicity_local_algebra} therefore yields
\[
\Delta_{\mathrm{lat}}(\beta_n)\ge \frac{m_{\mathrm{blk}}}{2B_n}.
\]
Divide by $a_n$ and use $a_nB_n\le C_{\mathrm{ent}}\ell_0$.
\end{proof}

\begin{theorem}[Full fixed-lattice OS gap from one-shot entrance]\label{thm:physical_gap_from_gluing}
Within the standing choices of Subsection~\ref{sec:main_theorem_package}, there exist $n_\ast<\infty$ and
\[
\mathfrak m_\ast:=\frac{m_{\mathrm{blk}}}{2C_{\mathrm{ent}}\ell_0}>0,
\qquad m_{\mathrm{blk}}:=-\log q_\ast,
\]
such that for all $n\ge n_\ast$,
\[
\frac{\Delta_{\mathrm{lat}}(\beta_n)}{a_n}\ge \mathfrak m_\ast.
\]
Equivalently, at each fixed tuned lattice spacing $a_n$ with $n\ge n_\ast$, the full fixed-lattice OS Hamiltonian has vacuum gap at least $m_{\mathrm{blk}}/(2B_n)$.
In particular QE2$_{\mathrm{lat}}$ holds unconditionally along the tuned dyadic trajectory as a theorem about the full fixed-lattice OS gap.
\end{theorem}

\begin{proof}
Take $n_\ast$ to be an entry index for Corollary~\ref{cor:uniform_phys_gap_from_bounded_entrance}.
\end{proof}


This section upgrades the finite-depth KP entrance argument into a post-entrance monotone safety statement: once the effective plaquette density (and the full interaction) lies in the KP admissibility basin at some scale, subsequent blocking steps preserve (and strictly improve) that basin membership. This removes the need for a classical global stable-manifold theorem for the $(u,K)$-dynamics beyond entrance; the only remaining global work is to reach entrance along the tuned ultraviolet trajectory and to construct the continuum limit along that trajectory.

\subsubsection{KP basin definition in the present paper}
Fix a block factor $L=2$ and the canonical polymer norm $\|\cdot\|_{\mathrm{KP}}$ used in the strong-coupling lattice module. Define the KP basin $\mathcal B_{\mathrm{KP}}(\varepsilon_{\mathrm{KP}})$ to be the set of effective interactions whose polymer activity satisfies
\[
\|K\|_{\mathrm{KP}} \le \varepsilon_{\mathrm{KP}},
\]
with $\varepsilon_{\mathrm{KP}}$ chosen so that the KP cluster expansion converges and yields exponential clustering (as in the channel-gap module).

\subsubsection{One-step forward invariance estimate}
\begin{proposition}[Forward invariance of the KP basin]
\label{prop:h56_forward_invariance}
There exist constants $0<q_{\mathrm{KP}}<1$ and $C_{\mathrm{KP}}<\infty$ (depending only on $G$ and the fixed blocking kernel parameters) such that the exact blocking map $\mathcal B_{2,\alpha}$ satisfies the one-step estimate
\[
\|K'\|_{\mathrm{KP}} \le q_{\mathrm{KP}}\,\|K\|_{\mathrm{KP}} + C_{\mathrm{KP}}\,a_{\mathrm{plaq}}(\beta),
\]
where the plaquette activity control parameter is
\[
a_{\mathrm{plaq}}(\beta):=\sum_{R\neq\mathbf 1} d_R^2 \exp\!\left(-\frac{C_2(R)}{\beta}\right).
\]
In particular, for all $\beta\le\beta_{\mathrm{KP}}$ such that
\[
C_{\mathrm{KP}}\,a_{\mathrm{plaq}}(\beta)\le (1-q_{\mathrm{KP}})\,\varepsilon_{\mathrm{KP}},
\]
one has forward invariance of the KP basin:
\[
K\in\mathcal B_{\mathrm{KP}}(\varepsilon_{\mathrm{KP}})\ \Longrightarrow\ K'\in\mathcal B_{\mathrm{KP}}(\varepsilon_{\mathrm{KP}}).
\]
\end{proposition}

\begin{proof}
Inside the KP basin, the conditional block integral (Section~\ref{subsec:hardened47_one_step_proof}) admits an absolutely convergent polymer expansion. The resulting re-expanded interaction can be written as a convergent product over connected polymers, and the KP norm of the re-expanded remainder satisfies a Lipschitz estimate in the incoming KP norm, yielding the contraction term $q_{\mathrm{KP}}\|K\|_{\mathrm{KP}}$ with $q_{\mathrm{KP}}<1$. The additive forcing term is controlled by the bare plaquette activity $a_{\mathrm{plaq}}(\beta)$ because the block kernel is class-function heat-kernel smoothing and convolution monotonicity improves $a_{\mathrm{plaq}}$ under further blocking. Uniformity in volume and admissible boundary conditions follows because the KP norm is local and the block kernel factorizes over coarse edges.
\end{proof}

\subsubsection{Consequence: post-entrance monotone safety}
\begin{corollary}[Post-entrance monotone safety]
\label{cor:h56_post_entrance_monotone}
Assume $\beta\le\beta_{\mathrm{KP}}$ and $K_{n_*}\in\mathcal B_{\mathrm{KP}}(\varepsilon_{\mathrm{KP}})$ at some entrance step $n_*$. Then for all $j\ge 0$,
\[
K_{n_*+j}\in\mathcal B_{\mathrm{KP}}(\varepsilon_{\mathrm{KP}}),
\]
and moreover
\[
\|K_{n_*+j}\|_{\mathrm{KP}}\le q_{\mathrm{KP}}^j\|K_{n_*}\|_{\mathrm{KP}} + \frac{C_{\mathrm{KP}}}{1-q_{\mathrm{KP}}}\,a_{\mathrm{plaq}}(\beta).
\]
\end{corollary}

\begin{proof}
Iterate Proposition~\ref{prop:h56_forward_invariance} and sum the resulting geometric series.
\end{proof}

\subsubsection{Summary: post-entrance monotonicity and one-shot entrance}
With Corollary~\ref{cor:h56_post_entrance_monotone}, post-entrance RG control is monotone: once the effective interaction
lies in the KP admissibility basin at some scale, all further coarse graining preserves (and improves) that admissibility.
Theorem~\ref{thm:h60_one_shot_entrance} complements this by supplying a \emph{one-shot} pre-entrance entrance theorem at the explicit scale $B(\beta)$
without invoking an iterated RG dynamical system.
Together with the tuned ultraviolet recursion (Proposition~\ref{prop:h59_uv_tuning_exists}) and the sharp local limit\/OS reconstruction proved in the main manuscript,
the QE1 contraction proof uses no infinite-step RG trajectory, while the tuned-sequence physical normalization still uses the separate weak-window coupling-identification input recorded above.

%%%%%%%%%%%%%%%%%%%%%%%%%%%%%%%%%%%%%%%%%%%%%%%%%%%%%%%%%%%%%%%%%%%%%%%%%%%%%%%
% Tuned UV→Entrance trajectory (explicit construction)
%%%%%%%%%%%%%%%%%%%%%%%%%%%%%%%%%%%%%%%%%%%%%%%%%%%%%%%%%%%%%%%%%%%%%%%%%%%%%%%

\subsection{Forward UV$\to$entrance tuning note (archival; moved to supplement)}
\label{sec:hardened57}

The earlier forward-recursion presentation of the UV$\to$entrance tuning argument is no longer part of the load-bearing Appendix~F spine.
That historical material has been moved to the supplementary Route~1 archive.
The only load-bearing UV input retained in the main appendix is the discrete one-dimensional tuning result of Subsection~\ref{sec:hardened59_uv_tuning}, which is written directly in the exact two-sided step-scaling form used later.

\subsubsection{OS-positive base and sharp local limit}
\label{subsec:h57_local_algebra_limit}

\paragraph{Purpose of this subsection.}
mass-gap uses a fixed positive flow time only as a \emph{renormalization prescription} to define tuned ultraviolet couplings and a family of bounded probes.
Osterwalder--Schrader (OS) reflection positivity and OS reconstruction are \emph{not} taken from the flowed probe algebra at fixed flow time:
Yang--Mills gradient flow is parabolic and mixes degrees of freedom across a Euclidean time reflection plane.
Accordingly, this appendix does not invoke (and does not require) reflection positivity for the fixed-$t$ flowed algebra.

\paragraph{OS-positive base and sharp local limit.}
OS reflection positivity is taken from the unflowed Wilson positive-time algebra (standard Wilson reflection positivity) and passed to the continuum sharp local limit by stability of positivity under pointwise limits.
The passage from bounded probes to the renormalized local field algebra and the verification of the OS axioms are carried out in the main manuscript (Lane~A closure theorem);
we do not duplicate that construction here.
In particular, the continuum Hamiltonian whose spectral gap is claimed is the OS Hamiltonian reconstructed from the sharp local state on the \emph{unflowed} local positive-time algebra.

\paragraph{Compatibility with the gap constant.}
The target \emph{physical} mass gap along the tuned dyadic trajectory (Corollary~\ref{cor:h60_physical_gap_noncollapse}) is a statement about the spectrum of that reconstructed OS Hamiltonian once the QE2/QE3 interfaces are closed.
It is therefore independent of which fixed-$t$ bounded probe family is used to define the tuning scheme.

\subsection{UV tuning trajectory as a one--dimensional recursion}
\label{sec:hardened59_uv_tuning}

\subsubsection{Statement}
Fix the step--scaling map $\Sigma:[0,u_{\mathrm{AF}}]\to\mathbb R$ from Subsubsection~\ref{subsec:hardened44_stepscaling_remainder}.
On the weak window we use exactly the two ingredients proved above:
\begin{enumerate}[leftmargin=*]
\item the two-sided cubic remainder bound \eqref{eq:stepscaling_bound_two_sided} and its forward-increment consequence \eqref{eq:stepscaling_increment_window};
\item continuity of $\Sigma$ on the weak window (Lemma~\ref{lem:stepscaling_continuous}).
\end{enumerate}
No derivative bound on the remainder is used anywhere in the tuning argument.

\begin{proposition}[Existence of a tuned UV sequence from the exact step-scaling input]\label{prop:h59_uv_tuning_exists}
Assume Proposition~\ref{prop:stepscaling_control_proved}, Lemma~\ref{lem:stepscaling_continuous}, and the explicit bare-to-flow comparison of Theorem~\ref{thm:h49_bare_flow_window}.
Then there exists $u_\ast\in(0,u_{\mathrm{AF}}]$ such that the following holds.

For every target value $u_0\in(0,u_\ast]$ there exists a sequence $\{u_{-n}\}_{n\ge0}\subset(0,u_\ast]$ with $u_{0}=u_0$ and
\begin{equation}
\label{eq:h59_backward_recursion}
\Sigma(u_{-(n+1)})=u_{-n},\qquad n\ge0,
\end{equation}
which is strictly decreasing and converges to $0$.
Moreover, with
\[
 c_-:=\frac12 b_0\log 2,
 \qquad
 c_+:=\frac32 b_0\log 2,
\]
one has the explicit reciprocal bounds
\begin{equation}\label{eq:h59_u_explicit_bounds}
\frac{1}{u_0^{-1}+c_+ n}
\le
u_{-n}
\le
\frac{1}{u_0^{-1}+\frac{c_-}{2}n}
\qquad (n\ge 0).
\end{equation}
Consequently, if
\begin{equation}\label{eq:h59_entry_index}
N_0(u_0):=\Big\lceil \frac{1}{c_+u_0}\Big\rceil,
\end{equation}
then for every $n\ge N_0(u_0)$,
\begin{equation}\label{eq:h59_un_over_n}
\frac{1}{2c_+ n}\ \le\ u_{-n}\ \le\ \frac{2}{c_- n}.
\end{equation}
Writing $a_n:=\ell_0 2^{-n}$, one may choose a bare sequence $\beta(a_n;u_0)$ satisfying
\begin{equation}\label{eq:h59_beta_exact_match}
 g_{\mathrm{GF}}^2(\ell_0;a_n,\beta(a_n;u_0))=u_{-n}
\end{equation}
and the explicit comparison bounds
\begin{equation}\label{eq:h59_beta_vs_u_bounds}
 \frac{1}{2u_{-n}}
 \le
 \beta(a_n;u_0)
 \le
 \frac{3}{2u_{-n}}
 \qquad (n\ge 0).
\end{equation}
Hence for every $n\ge N_0(u_0)$,
\begin{equation}\label{eq:h59_beta_vs_n_bounds}
 \frac{c_-}{4}\,n
 \le
 \beta(a_n;u_0)
 \le
 3c_+\,n.
\end{equation}
Consequently, along the tuned dyadic trajectory attached to $u_0$,
\[
 u(a_n)=\Theta\!\big(1/n\big)=\Theta\!\big(1/\log(1/a_n)\big),
 \qquad
 \beta(a_n;u_0)=\Theta(n)=\Theta\!\big(\log(1/a_n)\big),
\]
with comparison constants depending only on $b_0$, $C_\Sigma$, and $C_{\mathrm{GF}}$, and with entry index depending on $u_0$ only through \eqref{eq:h59_entry_index}.
\end{proposition}

\subsubsection{Proof}
Choose $u_\ast\le u_{\mathrm{AF}}$ so small that \eqref{eq:stepscaling_increment_window} holds on $(0,u_\ast]$ and, in addition,
\[
 u_\ast\le (2\beta_{\mathrm{GF}})^{-1},
\]
where $\beta_{\mathrm{GF}}$ is the threshold from Theorem~\ref{thm:h49_bare_flow_window}.
Set
\[
 c_-:=\frac12 b_0\log 2,
 \qquad
 c_+:=\frac32 b_0\log 2.
\]
Then for every $u\in(0,u_\ast]$,
\begin{equation}
\label{eq:h59_forward_increment_bounds}
 c_-u^2\le \Sigma(u)-u\le c_+u^2.
\end{equation}

\paragraph{Step 1: existence of a backward preimage inside the weak window.}
Fix $y\in(0,u_\ast]$.
By Lemma~\ref{lem:stepscaling_continuous}, $\Sigma$ is continuous on $[0,u_\ast]$.
Since $\Sigma(0)=0<y$ and \eqref{eq:h59_forward_increment_bounds} gives $\Sigma(y)>y$, the intermediate-value theorem yields at least one $x\in(0,y)$ with $\Sigma(x)=y$.
Starting from the prescribed target $u_0$, choose $u_{-(n+1)}\in(0,u_{-n})$ recursively so that \eqref{eq:h59_backward_recursion} holds.
This produces a strictly decreasing sequence in $(0,u_\ast]$, hence $u_{-n}\downarrow u_\infty$ for some $u_\infty\ge 0$.
Passing to the limit in \eqref{eq:h59_backward_recursion} and using \eqref{eq:h59_forward_increment_bounds}, one gets $\Sigma(u_\infty)=u_\infty$ and therefore $u_\infty=0$.

\paragraph{Step 2: reciprocal increments and the correct eventual-in-$n$ bounds.}
Let $v_n:=1/u_{-n}$.
From $u_{-n}=\Sigma(u_{-(n+1)})$ and \eqref{eq:h59_forward_increment_bounds} we obtain
\[
 c_-u_{-(n+1)}^2\le u_{-n}-u_{-(n+1)}\le c_+u_{-(n+1)}^2.
\]
Divide by $u_{-n}u_{-(n+1)}$ to get
\[
 \frac{c_-}{1+c_+u_\ast}
 \le
 v_{n+1}-v_n
 \le
 c_+.
\]
After shrinking $u_\ast$ once more so that $c_+u_\ast\le1$, this becomes
\begin{equation}
\label{eq:h59_reciprocal_increment}
 \frac{c_-}{2}
 \le
 v_{n+1}-v_n
 \le
 c_+.
\end{equation}
Summing \eqref{eq:h59_reciprocal_increment} from $0$ to $n-1$ gives
\[
 u_0^{-1}+\frac{c_-}{2}n\le v_n\le u_0^{-1}+c_+n.
\]
Inverting yields the exact bounds \eqref{eq:h59_u_explicit_bounds}.
The upper bound in \eqref{eq:h59_un_over_n} follows immediately from the right-hand inequality in \eqref{eq:h59_u_explicit_bounds}.
For the lower bound, if $n\ge N_0(u_0)$ then $u_0^{-1}\le c_+n$, hence
\[
 u_{-n}\ge \frac{1}{u_0^{-1}+c_+n}\ge \frac{1}{2c_+n}.
\]
This proves \eqref{eq:h59_un_over_n} and makes the dependence on the target value explicit only through the entry index \eqref{eq:h59_entry_index}.

\paragraph{Step 3: selection of the bare tuning sequence.}
Because the backward trajectory stays inside $(0,u_\ast]$ and $u_\ast\le (2\beta_{\mathrm{GF}})^{-1}$, Theorem~\ref{thm:h49_bare_flow_window} applies to every value $u_{-n}$.
Hence for each $n\ge 0$ there exists a bare coupling $\beta(a_n;u_0)$ satisfying \eqref{eq:h59_beta_exact_match}, and the interval bound \eqref{eq:h49_beta_selection_interval} gives exactly \eqref{eq:h59_beta_vs_u_bounds}.
Combining \eqref{eq:h59_beta_vs_u_bounds} with \eqref{eq:h59_un_over_n} yields the explicit eventual-in-$n$ comparison
\[
 \frac{c_-}{4}\,n
 \le
 \beta(a_n;u_0)
 \le
 3c_+\,n
 \qquad (n\ge N_0(u_0)),
\]
which is \eqref{eq:h59_beta_vs_n_bounds}.
Since $n=(\log(\ell_0/a_n))/\log 2$, the final asymptotics follow.
\qed

\subsubsection{Consequence for the tuning sequence \texorpdfstring{$\beta(a)$}{beta(a)}}
\label{subsec:h59_beta_of_a}
Fix a target value $u_{\mathrm{ren}}\in(0,u_\ast]$.
Let $a_n=\ell_0 2^{-n}$, define $u(a_n):=u_{-n}$ from \eqref{eq:h59_backward_recursion}, and choose $\beta(a_n;u_{\mathrm{ren}})$ by Theorem~\ref{thm:h49_bare_flow_window} so that
\[
 g_{\mathrm{GF}}^2(\ell_0;a_n,\beta(a_n;u_{\mathrm{ren}})) = u(a_n).
\]
Then \eqref{eq:h59_un_over_n} and \eqref{eq:h59_beta_vs_n_bounds} imply
\[
 u(a_n)=\Theta(1/n)=\Theta(1/\log(1/a_n)),
 \qquad
 \beta(a_n;u_{\mathrm{ren}})=\Theta(n)=\Theta(\log(1/a_n)),
\]
with universal comparison constants and with entry index $N_0(u_{\mathrm{ren}})$.



% ==============================
% : One-shot entrance and transfer (eliminate global RG dynamical control)
% ==============================
\subsection{One-shot KP entrance from majorant mixing (no infinite-step RG control)}
\label{sec:hardened60_one_shot}

The remaining ``global RG'' barrier can be eliminated entirely by reframing the pre-entrance argument as a \emph{single}
coarse-graining to the entrance scale, rather than an indefinitely iterated Banach-space dynamical system.
The key point is that all post-entrance control is already monotone, while pre-entrance control
is finite-depth and uniform . Thus it suffices to define the blocked theory at the entrance
scale by a one-shot coarse-graining kernel and verify that it lies in the KP/Dobrushin basin.

\subsubsection{One-shot entrance scale}

Fix block factor $L=2$.
Let $k_{\mathrm{opt}}(\beta)$ be the explicit majorant mixing depth from Proposition~\ref{prop:h51_kbeta_opt}
(defined by $2^{2k_{\mathrm{opt}}(\beta)}\ge \max\{1,\beta s_{\delta}\}$).
Define the associated entrance block size
\begin{equation}
\label{eq:h60_Bbeta_def}
B(\beta):=2^{k_{\mathrm{opt}}(\beta)}.
\end{equation}
By definition,
\begin{equation}
\label{eq:h60_Bbeta_time}
\frac{B(\beta)^2}{\beta}\ \ge\ s_{\delta}.
\end{equation}

\subsubsection{Exact one-shot blocked variables}

Let $\Lambda\subset\mathbb Z^4$ be a finite region with admissible boundary conditions and let $B=2^k$ be dyadic.
For the \emph{exact} blocked theory used on the load-bearing Route~1 chain we do not insert any additional smoothing randomness.
Instead we take the terminal skeleton field produced by $k$ repetitions of the half-open dyadic blocking rule of Section~\ref{sec:patching} and group those surviving links by their owner entrance blocks.
The auxiliary smoothing kernel $\mathcal T_{B,\alpha}$ belongs only to the \emph{$\alpha$-reference} specification already constructed in the transfer module; it does not define the exact blocked sigma-algebra used in Proposition~\ref{prop:route1_blocked_representatives}.

\begin{definition}[Exact one-shot blocked variables and blocked measure at scale $B$]
\label{def:h60_one_shot_blocked_measure}
Let $\Sigma_B$ be the terminal skeleton obtained after $k=\log_2 B$ exact dyadic Wilson blockings, and let $\mathcal F_B$ be the sigma-algebra generated by the links on $\Sigma_B$.
For each half-open entrance block $\mathcal Q$ of side $B$, let $U^{\langle B\rangle}_{\mathcal Q}$ be the collection of terminal skeleton links whose anchor (equivalently: half-open owner block) lies in $\mathcal Q$.
The exact blocked field is
\[
U^{\langle B\rangle}:=\bigl(U^{\langle B\rangle}_{\mathcal Q}\bigr)_{\mathcal Q},
\]
and the exact scale-$B$ blocked measure is the push-forward
\begin{equation}
\label{eq:h60_one_shot_block_def}
\mu_{\beta,\Lambda}^{\langle B\rangle}
:=
(U\mapsto U^{\langle B\rangle})_\#\,\mu_{\beta,\Lambda}.
\end{equation}
We write $\gamma^{\langle B\rangle}_{\beta,\Lambda}$ for the corresponding exact blocked specification on the entrance-block lattice.
\end{definition}

\begin{proposition}[Exact blocked variables coincide with the terminal skeleton sigma-algebra]
\label{prop:route1_blocked_sigma_identification}
Let $B=2^k$ and set $j_0:=0$, $m:=k$ in Theorem~\ref{thm:wilson_exact_recursion}.
Then the sigma-algebra generated by the exact blocked field $U^{\langle B\rangle}$ of Definition~\ref{def:h60_one_shot_blocked_measure} is exactly the terminal skeleton sigma-algebra:
\[
\sigma(U^{\langle B\rangle})=\mathcal F_B=\mathcal F_{j_0+m}.
\]
Consequently, for every bounded fine observable $H$,
\begin{equation}\label{eq:route1_sigma_conditional_expectation}
\E_{\omega_a}[H\mid U^{\langle B\rangle}]
=
\E_{\omega_a}[H\mid \mathcal F_B].
\end{equation}
Under this identification, Corollary~\ref{cor:wilson_terminal_block_product} applies directly to observables whose fine supports lie in pairwise disjoint families of entrance blocks.
\end{proposition}

\begin{proof}
The terminal skeleton links are partitioned uniquely by the half-open owner entrance blocks.
Thus the family $\{U^{\langle B\rangle}_{\mathcal Q}\}_{\mathcal Q}$ is simply a relabeling of the full terminal skeleton field, so it generates the same sigma-algebra $\mathcal F_B$.
Equation~\eqref{eq:route1_sigma_conditional_expectation} is therefore just the identity of conditional expectations onto the same sigma-algebra.
The final sentence is Corollary~\ref{cor:wilson_terminal_block_product} rewritten in the owner-block coordinates of Definition~\ref{def:h60_one_shot_blocked_measure}.
\end{proof}

\begin{remark}[Reference versus exact blocked objects]
The $\alpha$-reference blocked specification entering Proposition~\ref{prop:TBmix_to_f} and Theorem~\ref{thm:gluing_KP_entrance} continues to be the smoothed/reference object constructed in the transfer module.
Definition~\ref{def:h60_one_shot_blocked_measure} records only the \emph{exact} blocked law on the load-bearing Route~1 path.
Both objects are taken at the same prescribed block size $B$, so no second entrance mechanism is introduced.
\end{remark}

\begin{definition}[Entrance-block interaction graph]\label{def:route1_block_graph}
For a dyadic entrance scale $B=2^k$, write the half-open entrance blocks as
\[
\mathcal Q(\mathbf m):=B\mathbf m+[0,B)^4,
\qquad \mathbf m\in\mathbb Z^4.
\]
Two entrance blocks are declared adjacent, written $\mathcal Q(\mathbf m)\sim_{\mathrm{blk}}\mathcal Q(\mathbf n)$, if
\[
\|\mathbf m-\mathbf n\|_{\infty}\le 2.
\]
Let $\dist_{\mathrm{blk}}$ denote the graph distance on this interaction graph.
Equivalently, one step in $\dist_{\mathrm{blk}}$ moves across at most two entrance blocks in each coordinate direction.
\end{definition}

\begin{proposition}[Local Gibbs form and finite-range property of the exact one-shot blocked specification]
\label{prop:route1_block_finite_range}
Let $B=2^k$ and consider the exact blocked measure $\mu_{\beta,\Lambda}^{\langle B\rangle}$ of Definition~\ref{def:h60_one_shot_blocked_measure}.
Then, with respect to the product Haar measure on the terminal skeleton links grouped by owner entrance blocks, the blocked law has the local Gibbs form
\begin{equation}\label{eq:route1_block_local_gibbs}
 d\mu_{\beta,\Lambda}^{\langle B\rangle}(U^{\langle B\rangle})
 =
 \frac{1}{\mathcal Z_{\beta,\Lambda}^{\langle B\rangle}}
 \prod_{\mathcal Q} W_{\mathcal Q}\bigl(U^{\langle B\rangle}|_{\operatorname{St}_1(\mathcal Q)}\bigr)
 \bigotimes_{\mathcal P} dH_{\mathcal P}\bigl(U^{\langle B\rangle}_{\mathcal P}\bigr),
\end{equation}
where
\[
\operatorname{St}_1(\mathcal Q):=\{\mathcal P:\|\mathbf p-\mathbf q\|_\infty\le 1\}
\]
for $\mathcal Q=\mathcal Q(\mathbf q)$, each $dH_{\mathcal P}$ is the product Haar measure on the terminal skeleton links owned by $\mathcal P$, and each local weight $W_{\mathcal Q}$ depends only on the blocked variables carried by blocks in $\operatorname{St}_1(\mathcal Q)$.
Consequently the exact blocked specification is finite range on the interaction graph of Definition~\ref{def:route1_block_graph}: for every entrance block $\mathcal Q$, the single-site conditional kernel
\[
\gamma^{\langle B\rangle}_{\mathcal Q}(\cdot\mid\eta)
\]
depends on the exterior configuration only through the values of $\eta$ on the graph-neighborhood
\[
\operatorname{Nb}_{\mathrm{blk}}(\mathcal Q):=\{\mathcal R:\dist_{\mathrm{blk}}(\mathcal Q,\mathcal R)\le 1\},
\]
and hence
\begin{equation}\label{eq:route1_block_finite_range}
 c_{\mathcal Q\mathcal R}\bigl(\gamma^{\langle B\rangle}_{\beta,\Lambda}\bigr)=0
 \qquad\text{whenever }\dist_{\mathrm{blk}}(\mathcal Q,\mathcal R)>1.
\end{equation}
\end{proposition}

\begin{proof}
By Corollary~\ref{cor:wilson_terminal_block_product}, once the terminal skeleton field is frozen, the fine interior-link law factorizes exactly over the terminal half-open blocks.
Integrating the fine interior links block by block therefore yields the push-forward density of the terminal skeleton field in the form \eqref{eq:route1_block_local_gibbs}, where $W_{\mathcal Q}$ is the terminal-block partition function with the skeleton boundary data held fixed.
Because a fine plaquette anchored in $\mathcal Q$ can touch only terminal skeleton links whose owner block has anchor at max-distance at most $1$ from that of $\mathcal Q$, each factor $W_{\mathcal Q}$ depends only on the blocked variables carried by blocks in $\operatorname{St}_1(\mathcal Q)$.

Now fix $\mathcal Q_0=\mathcal Q(\mathbf q_0)$.
In the single-site conditional density of $U^{\langle B\rangle}_{\mathcal Q_0}$ given the other blocked variables, only the factors $W_{\mathcal P}$ with $\|\mathbf p-\mathbf q_0\|_\infty\le 1$ can contain $U^{\langle B\rangle}_{\mathcal Q_0}$.
Each such factor depends only on blocks $\mathcal R=\mathcal Q(\mathbf r)$ with $\|\mathbf r-\mathbf p\|_\infty\le 1$, hence on blocks satisfying $\|\mathbf r-\mathbf q_0\|_\infty\le 2$.
By Definition~\ref{def:route1_block_graph}, these are exactly the graph-neighbors of $\mathcal Q_0$.
Therefore the single-site kernel for $\mathcal Q_0$ depends on the exterior configuration only through $\operatorname{Nb}_{\mathrm{blk}}(\mathcal Q_0)$, which is the finite-range statement \eqref{eq:route1_block_finite_range}.
\end{proof}

\subsubsection{One-shot entrance theorem}

We now show that for $\beta$ large enough the one-shot blocked measure at scale $B(\beta)$ lies in the KP basin.

\begin{theorem}[One-shot entrance into the KP basin]
\label{thm:h60_one_shot_entrance}
Fix the Plan-A / KP entrance threshold $\delta\in(0,\delta_{\mathrm{PA}}]$ small enough that
\[
s_{\delta}\ge S_\star,
\]
where $S_\star$ is the explicit KP-activity budget threshold from Proposition~\ref{prop:TBmix_to_f} for the final choice of $f_\star$.
(Existence of such a choice follows from Lemma~\ref{lem:h51_chi2_monotone}, since $s_\delta\to\infty$ as $\delta\downarrow0$.)
Choose $\kappa_{\delta}$ and $\beta_0(\delta)$ as in Proposition~\ref{prop:h54_transfer_feasible}.
Then for every $\beta\ge \beta_0(\delta)$, with $B(\beta)$ as in~\eqref{eq:h60_Bbeta_def}, the one-shot blocked
specification induced by $\mu_{\beta,\Lambda}^{\langle B(\beta)\rangle}$ lies in the KP/Dobrushin basin:
there exists $q_\ast\in(0,1)$ (independent of $\beta,\Lambda$ and boundary conditions) such that
\begin{equation}
\label{eq:h60_kp_contraction}
\delta\bigl(\gamma^{\langle B(\beta)\rangle}_{\beta,\Lambda}\bigr)\le q_\ast<1.
\end{equation}
In particular, the corresponding infinite-volume blocked state exists and exhibits exponential clustering.
\end{theorem}

\begin{proof}
\noindent\textbf{Step 1 (reference mixing at the one-shot scale; crossover input).}
By Proposition~\ref{prop:h51_kbeta_opt} and the definition of $B(\beta)$,
the majorant convolution power reaches the $\chi^2$-mixing level $\delta$ at time $B(\beta)^2/\beta\ge s_{\delta}$.
Because $s_{\delta}\ge S_\star$ by the standing choice of $\delta$, the same one-shot scale also meets the explicit KP activity budget of Proposition~\ref{prop:TBmix_to_f}:
\[
\frac{B(\beta)^2}{\beta}=\frac{|T_{B(\beta)}|}{\beta}\ge s_{\delta}\ge S_\star.
\]
Corollary~\ref{cor:TBmix_reference_budget_closed} therefore supplies the theorem-level QE1 reference budget at that same dyadic block size: Proposition~\ref{prop:TBmix_to_f} has already converted the crossover coefficient decay into KP-small $L^2$ activity by an explicit low-/high-representation split at scale $B(\beta)$:
\begin{equation}
\label{eq:h60_ref_small}
f_{B(\beta)}^{(\alpha)}\le f_{\star},
\end{equation}
with $f_{\star}$ chosen below the admissible threshold in the deterministic KP iteration.

\medskip
\noindent\textbf{Step 2 (exact-to-reference transfer at the one-shot scale).}
The finite-depth transfer control was strengthened in Proposition~\ref{prop:h55_planA_gate_closed} and Corollary~\ref{cor:h55_planA_gate_one_shot}.
In particular, under the pivot schedule determined by $\kappa_{\delta}$, the exact-to-reference comparison at the \emph{same} one-shot block size $L=B(\beta)$ satisfies the Plan-A conditional $\chi^2$ gate with defect $\varepsilon_\ast$, uniformly in finite volume, admissible boundary conditions, and all events $\Omega$ of reference probability at least $1/2$.
The gluing/transfer theorem (Theorem~\ref{thm:gluing_KP_entrance}) applies at this same block size and transfers KP-smallness of the reference specification to the exact blocked specification,
with constants independent of volume and boundary conditions.
Applying Theorem~\ref{thm:gluing_KP_entrance} at $L=B(\beta)$ (where \eqref{eq:h60_ref_small} holds) yields
\eqref{eq:h60_kp_contraction}.
By Proposition~\ref{prop:route1_block_finite_range}, the exact blocked specification is finite range on the entrance-block interaction graph, so Proposition~\ref{prop:dobrushin_standard} gives existence, uniqueness, and exponential clustering of the blocked infinite-volume state.
\end{proof}


\begin{corollary}[Bounded physical one-shot entrance scale along a fixed tuned dyadic trajectory]\label{cor:route1_bounded_physical_entrance}
Fix a target value $u_{\mathrm{ren}}\in(0,u_\ast]$ and let
\[
a_n:=\ell_0 2^{-n},\qquad \beta_n:=\beta(a_n;u_{\mathrm{ren}}),\qquad B_n:=B(\beta_n)
\]
be the associated tuned dyadic trajectory from Proposition~\ref{prop:h59_uv_tuning_exists}.
Then there exists a constant $C_{\mathrm{ent}}<\infty$, depending only on the group ledger and the one-shot threshold data, and an entry index $N_{\mathrm{ent}}(u_{\mathrm{ren}})\in\mathbb N$ such that
\[
 a_n B_n\le C_{\mathrm{ent}}\ell_0
 \qquad (n\ge N_{\mathrm{ent}}(u_{\mathrm{ren}})).
\]
In particular the one-shot entrance scale remains bounded in physical units along every tuned dyadic trajectory, with a uniform constant and a trajectory-dependent onset index.
\end{corollary}

\begin{proof}
By Proposition~\ref{prop:h51_kbeta_opt} one has
\[
B(\beta)\le 2\sqrt{\max\{1,\beta s_{\delta}\}}.
\]
Fix $u_{\mathrm{ren}}$ and apply Proposition~\ref{prop:h59_uv_tuning_exists} with $u_0=u_{\mathrm{ren}}$.
For every $n\ge N_0(u_{\mathrm{ren}})$, the explicit bound \eqref{eq:h59_beta_vs_n_bounds} gives
\[
\beta_n\le 3c_+ n.
\]
Hence for all $n\ge N_{\mathrm{ent}}(u_{\mathrm{ren}}):=N_0(u_{\mathrm{ren}})$,
\[
 a_n B_n
 \le 2\ell_0 2^{-n}\sqrt{\max\{1,3c_+s_{\delta}n\}}
 \le 2\ell_0\sqrt{\max\{1,3c_+s_{\delta}\}},
\]
using $\sqrt n\le 2^n$ for $n\ge 1$.
Thus one may take
\[
C_{\mathrm{ent}}:=2\sqrt{\max\{1,3c_+s_{\delta}\}},
\]
which is uniform in the chosen target value $u_{\mathrm{ren}}$.
\end{proof}

\subsubsection{Route~1 internal theorem chain: dense-core Euclidean clustering}
\label{sec:route1_dense_core_internal}

The load-bearing Route~1/Lane~B input is now recorded as an explicit theorem/proposition chain, rather than only as a front-facing theorem-package claim.
We work with the fixed countable positive-time local core from Lane~B3 and prove Euclidean exponential clustering directly on that core from the one-shot entrance-scale blocked theory.
The stronger fixed-$\beta$ microscopic descent statements that appear later remain auxiliary and are not used in the argument below.

\begin{definition}[Route~1 dense positive-time probe core]
\label{def:route1_dense_probe_core}
Set
\[
\mathcal G_{\mathrm{ent}}:=\mathcal A_{+,\mathrm{poly}}\subset\mathcal A_+,
\]
where $\mathcal A_{+,\mathrm{poly}}$ is the countable positive-time local $*$-algebra fixed in Lane~B3.
\end{definition}

\begin{lemma}[Route~1 probe core is OS-dense]
\label{lem:route1_probe_core_dense}
The span of OS-classes $\{[F]:F\in\mathcal G_{\mathrm{ent}}\}$ is dense in $\mathcal H$.
Consequently the centered classes $\{[\widetilde F]:F\in\mathcal G_{\mathrm{ent}}\}$ are dense in $\Omega^\perp$.
\end{lemma}

\begin{proof}
By Definition~\ref{def:route1_dense_probe_core}, $\mathcal G_{\mathrm{ent}}=\mathcal A_{+,\mathrm{poly}}$.
The first statement is exactly Lemma~\ref{lem:density_poly}; the second follows by subtracting the vacuum component, i.e. by orthogonal projection onto $\Omega^\perp$.
\end{proof}

\begin{lemma}[Entrance mesh shrinks along the tuned sequence]
\label{lem:route1_entrance_mesh_to_zero}
Fix a target value $u_{\mathrm{ren}}\in(0,u_\ast]$ and let $a_n:=\ell_0 2^{-n}$, $\beta_n:=\beta(a_n;u_{\mathrm{ren}})$, and $B_n:=B(\beta_n)$.
Then
\[
 a_n B_n\xrightarrow[n\to\infty]{}0.
\]
In particular $B_n\to\infty$ in lattice units, and every fixed bounded local cylinder observable is eventually contained in a uniformly bounded number of entrance blocks.
\end{lemma}

\begin{proof}
The same estimate as in the proof of Corollary~\ref{cor:route1_bounded_physical_entrance} gives
\[
 a_n B_n
 \le 2\ell_0 2^{-n}\sqrt{\max\{1,C_{\beta}s_{\delta}n\}},
\]
which tends to $0$ as $n\to\infty$.
Since each fixed local cylinder observable depends on finitely many fine links, the number of entrance blocks intersecting its support is therefore uniformly bounded for all sufficiently large $n$.
\end{proof}

\begin{remark}[Why a further fixed-physical-scale theorem is still needed]
Lemma~\ref{lem:route1_entrance_mesh_to_zero} shows that the physical size of one entrance block tends to zero along the tuned sequence.
Accordingly, any final fixed-physical-shift Euclidean clustering theorem must include an additional persistence or upgrade step beyond bare entrance-scale decay on the blocked interaction graph.
This is why the dense-core route recorded below is retained as a secondary bridge rather than as the primary front-facing closure path.
\end{remark}

\begin{proposition}[Entrance-scale block covariance decay in oscillation form]
\label{prop:route1_block_covariance}
Fix a target value $u_{\mathrm{ren}}\in(0,u_\ast]$, let $a_n:=\ell_0 2^{-n}$, $\beta_n:=\beta(a_n;u_{\mathrm{ren}})$, and $B_n:=B(\beta_n)$, and let $\omega_{a_n}^{\langle B_n\rangle}$ denote the unique infinite-volume exact blocked state at the one-shot entrance scale furnished by Theorem~\ref{thm:h60_one_shot_entrance}.
For a bounded blocked observable $A$, define the entrance-block oscillation seminorm
\[
\operatorname{osc}_{\mathcal Q}(A):=
\sup_{\eta,\eta':\ \eta_{\mathbb Z^4_{B_n}\setminus\{\mathcal Q\}}=\eta'_{\mathbb Z^4_{B_n}\setminus\{\mathcal Q\}}}
|A(\eta)-A(\eta')|,
\qquad
\supp_{\mathrm{osc}}(A):=\{\mathcal Q:\operatorname{osc}_{\mathcal Q}(A)\neq0\}.
\]
Then every pair of bounded blocked observables $A,B$ satisfies
\begin{equation}\label{eq:route1_block_covariance_densecore}
\begin{aligned}
\big|\mathrm{Cov}_{\omega_{a_n}^{\langle B_n\rangle}}(A,B)\big|
\le{}&\frac{e^{-m_{\mathrm{blk}}\,\mathrm{dist}_{\mathrm{blk}}(\supp_{\mathrm{osc}}(A),\supp_{\mathrm{osc}}(B))}}{1-q_\ast}
\Big(\sum_{\mathcal Q\in\supp_{\mathrm{osc}}(A)}\operatorname{osc}_{\mathcal Q}(A)\Big)\\
&\times\Big(\sum_{\mathcal R\in\supp_{\mathrm{osc}}(B)}\operatorname{osc}_{\mathcal R}(B)\Big),
\qquad m_{\mathrm{blk}}:=-\log q_\ast.
\end{aligned}
\end{equation}
where $\mathrm{dist}_{\mathrm{blk}}$ denotes graph distance on the interaction graph of Definition~\ref{def:route1_block_graph}.
In particular,
\begin{equation}\label{eq:route1_block_covariance_cardinality}
\big|\mathrm{Cov}_{\omega_{a_n}^{\langle B_n\rangle}}(A,B)\big|
\le
\frac{4\,|\supp_{\mathrm{osc}}(A)|\,|\supp_{\mathrm{osc}}(B)|}{1-q_\ast}\,
\|A\|_\infty\|B\|_\infty
 e^{-m_{\mathrm{blk}}\,\mathrm{dist}_{\mathrm{blk}}(\supp_{\mathrm{osc}}(A),\supp_{\mathrm{osc}}(B))}.
\end{equation}
\end{proposition}

\begin{proof}
By Theorem~\ref{thm:h60_one_shot_entrance}, the exact blocked specification at scale $B_n$ satisfies
\[
\delta\bigl(\gamma^{\langle B_n\rangle}_{\beta_n,\Lambda}\bigr)\le q_\ast<1
\]
uniformly in finite volume and admissible boundary conditions.
Proposition~\ref{prop:route1_block_finite_range} shows that the exact blocked specification is finite range on the interaction graph of Definition~\ref{def:route1_block_graph}, i.e. its influence matrix vanishes beyond one graph step.
Therefore Lemma~\ref{lem:dobrushin_to_block_decay} applies to the unique blocked infinite-volume state and yields the support-weighted estimate \eqref{eq:route1_block_covariance_densecore} with $m_{\mathrm{blk}}=-\log q_\ast$.
The cardinality bound \eqref{eq:route1_block_covariance_cardinality} follows from $\operatorname{osc}_{\mathcal Q}(A)\le 2\|A\|_\infty$ and $\operatorname{osc}_{\mathcal R}(B)\le 2\|B\|_\infty$.
\end{proof}



\begin{proposition}[Entrance-scale blocked representatives for arbitrary local probes]
\label{prop:route1_blocked_representatives}
Fix a target value $u_{\mathrm{ren}}\in(0,u_\ast]$.
Let $a_n:=\ell_0 2^{-n}$, $\beta_n:=\beta(a_n;u_{\mathrm{ren}})$, and $B_n:=B(\beta_n)$.
Then there exists $n_{\mathrm{repr}}(u_{\mathrm{ren}})<\infty$ such that for every $n\ge n_{\mathrm{repr}}(u_{\mathrm{ren}})$ and every bounded gauge-invariant local cylinder observable $H$ (not necessarily positive-time), there exists a bounded blocked observable $\Phi_{H,n}$ on the one-shot blocked state $\omega_{a_n}^{\langle B_n\rangle}$ such that:
\begin{enumerate}[leftmargin=*]
\item $\Phi_{H,n}$ depends only on the finite entrance-block hull $\mathsf B_n(H)$ of the support of $H$;
\item $\|\Phi_{H,n}\|_\infty\le \|H\|_\infty$;
\item if $H_1,H_2$ are bounded local cylinder observables with disjoint entrance-block hulls $\mathsf B_n(H_1)\cap \mathsf B_n(H_2)=\varnothing$, then
\begin{equation}\label{eq:route1_exact_core_cov_identity}
\omega_{a_n}(H_1H_2)
=
\omega_{a_n}^{\langle B_n\rangle}(\Phi_{H_1,n}\,\Phi_{H_2,n}),
\qquad
\omega_{a_n}(H_i)=\omega_{a_n}^{\langle B_n\rangle}(\Phi_{H_i,n})\quad(i=1,2),
\end{equation}
and therefore
\begin{equation}\label{eq:route1_exact_core_covariance}
\mathrm{Cov}_{\omega_{a_n}}(H_1,H_2)
=
\mathrm{Cov}_{\omega_{a_n}^{\langle B_n\rangle}}(\Phi_{H_1,n},\Phi_{H_2,n}).
\end{equation}
\end{enumerate}
No support-size restriction enters: at each fixed tuned lattice spacing, every bounded local cylinder observable admits such an exact entrance-scale representative.
\end{proposition}

\begin{proof}
Choose $n_{\mathrm{repr}}(u_{\mathrm{ren}})$ so that along the tuned dyadic trajectory the one-shot block depth
\[
k_n:=\log_2 B_n
\]
is realized by the exact dyadic Wilson recursion and Proposition~\ref{prop:route1_blocked_sigma_identification} applies for all $n\ge n_{\mathrm{repr}}(u_{\mathrm{ren}})$.
Fix such an $n$.

By Proposition~\ref{prop:route1_blocked_sigma_identification}, the sigma-algebra generated by the exact blocked variables $U^{\langle B_n\rangle}$ is exactly the terminal skeleton sigma-algebra $\mathcal F_{B_n}$ from Theorem~\ref{thm:wilson_exact_recursion} at depth $k_n$.
Define
\[
\Phi_{H,n}:=\mathbb E_{\omega_{a_n}}[H\mid U^{\langle B_n\rangle}]
=\mathbb E_{\omega_{a_n}}[H\mid \mathcal F_{B_n}].
\]
The $L^\infty$ bound in item~(2) is immediate from conditional expectation.

For item~(1), Theorem~\ref{thm:wilson_one_step_factorization} shows that at each dyadic step the conditional integration factors blockwise once the next skeleton is fixed, so links outside the parent block hull of the current support integrate out trivially.
Iterating through the $k_n$ dyadic steps therefore shows that $\Phi_{H,n}$ depends only on the entrance-block hull $\mathsf B_n(H)$.
Because $H$ is local, that hull is always finite, but its cardinality may now depend on both $H$ and $n$.

For item~(3), fix $H_1,H_2$ with disjoint entrance-block hulls.
By item~(1), each $H_i$ is measurable with respect to the fine links contained in the corresponding finite family of half-open entrance blocks $\mathsf B_n(H_i)$.
Set $j_0:=0$ and $m:=k_n=\log_2 B_n$ in Corollary~\ref{cor:wilson_terminal_block_product}; Proposition~\ref{prop:route1_blocked_sigma_identification} identifies its terminal skeleton sigma-algebra with $\sigma(U^{\langle B_n\rangle})$.
Because the block families $\mathsf B_n(H_1)$ and $\mathsf B_n(H_2)$ are disjoint, Corollary~\ref{cor:wilson_terminal_block_product} gives
\[
\mathbb E_{\omega_{a_n}}[H_1H_2\mid U^{\langle B_n\rangle}]
=
\mathbb E_{\omega_{a_n}}[H_1\mid U^{\langle B_n\rangle}]\,
\mathbb E_{\omega_{a_n}}[H_2\mid U^{\langle B_n\rangle}]
=
\Phi_{H_1,n}\,\Phi_{H_2,n}.
\]
Taking expectations gives \eqref{eq:route1_exact_core_cov_identity}, and subtracting the one-point identities yields \eqref{eq:route1_exact_core_covariance}.
\end{proof}



\begin{remark}[Status of the dense-core Route~1 bridge]\label{rem:route1_dense_core_status}
Proposition~\ref{prop:route1_block_covariance} and Proposition~\ref{prop:route1_blocked_representatives}
still isolate the microscopic lattice-side ingredients one would need for a \emph{direct}
fixed-physical-scale Euclidean-clustering proof from one-shot entrance.
However, Lemma~\ref{lem:route1_entrance_mesh_to_zero} shows that the entrance mesh shrinks in physical units, so such a direct route
would require an additional persistence/upgrade theorem beyond bare entrance-scale decay.
We therefore use the spectral-calculus bridge below:
once QE3 and Lane~A base-state compatibility identify the sharp-local OS Hamiltonian with the same gapped Hamiltonian, dense-core
Euclidean clustering follows immediately for that same state.
Thus Lane~B is a downstream reformulation of the front-facing Route~1 non-collapse proof,
while the block-covariance propositions remain as diagnostics toward an independent lattice-side derivation of the same input.
\end{remark}

\begin{theorem}[Secondary Route~1 dense-core bridge via QE3 and OS spectral calculus]
\label{thm:route1_E2_internal}
Fix a target value $u_{\mathrm{ren}}\in(0,u_\ast]$ and let $\omega^{(u_{\mathrm{ren}})}$ be the sharp-local continuum state
furnished by Corollary~\ref{cor:LaneA_full_scope} along the tuned dyadic sequence associated with $u_{\mathrm{ren}}$.
Let $(\mathcal H,\Omega,H)$ be its OS reconstruction.
Then the dense positive-time core
\[
\mathcal G_{\mathrm{ent}}=\mathcal A_{+,\mathrm{poly}}
\]
is OS-dense in $\mathcal H$ and satisfies Euclidean exponential clustering with rate
\begin{equation}\label{eq:route1_mstar_densecore}
 m=\mathfrak m_\ast:=\frac{m_{\mathrm{blk}}}{2C_{\mathrm{ent}}\ell_0}.
\end{equation}
More precisely, for every $F,G\in\mathcal G_{\mathrm{ent}}$ and every $s\ge 0$,
\begin{equation}\label{eq:route1_core_continuum_E2_cross}
 \big|\omega^{(u_{\mathrm{ren}})}\big(\Theta \widetilde F\,\tau_s\widetilde G\big)\big|
 \le e^{-\mathfrak m_\ast s}\,\|[\widetilde F]\|_{\mathcal H}\,\|[\widetilde G]\|_{\mathcal H}.
\end{equation}
In particular, for every $F\in\mathcal G_{\mathrm{ent}}$,
\begin{equation}\label{eq:route1_core_continuum_E2}
 \big|\omega^{(u_{\mathrm{ren}})}\big(\Theta \widetilde F\,\tau_s\widetilde F\big)\big|
 \le \|[\widetilde F]\|_{\mathcal H}^{2}\,e^{-\mathfrak m_\ast s}
 \qquad (s\ge 0).
\end{equation}
Hence the Lane~B Euclidean clustering input is verified for the same sharp-local state, now as a corollary of QE3 rather than as an independent replacement proof of QE3.
\end{theorem}

\begin{proof}
OS-density of the core is exactly Lemma~\ref{lem:route1_probe_core_dense}.
By Theorem~\ref{thm:gap_transfer}, the OS Hamiltonian reconstructed from $\omega^{(u_{\mathrm{ren}})}$ satisfies
\[
\sigma(H)\subset\{0\}\cup[\mathfrak m_\ast,\infty),
\qquad
\ker(H)=\mathbb C\Omega.
\]
Therefore, for every $\psi\in\Omega^\perp$ and every $s\ge 0$,
\[
\|e^{-sH}\psi\|\le e^{-\mathfrak m_\ast s}\,\|\psi\|.
\]
Now fix $F,G\in\mathcal G_{\mathrm{ent}}$.
Because $[\widetilde F],[\widetilde G]\in\Omega^\perp$, the OS semigroup identity gives
\[
\omega^{(u_{\mathrm{ren}})}(\Theta \widetilde F\,\tau_s\widetilde G)
=\langle [\widetilde F],e^{-sH}[\widetilde G]\rangle_{\mathcal H}.
\]
Applying Cauchy--Schwarz and the previous contraction bound yields
\[
\big|\omega^{(u_{\mathrm{ren}})}(\Theta \widetilde F\,\tau_s\widetilde G)\big|
\le \|[\widetilde F]\|_{\mathcal H}\,\|e^{-sH}[\widetilde G]\|_{\mathcal H}
\le e^{-\mathfrak m_\ast s}\,\|[\widetilde F]\|_{\mathcal H}\,\|[\widetilde G]\|_{\mathcal H},
\]
which is exactly \eqref{eq:route1_core_continuum_E2_cross}.
Taking $G=F$ gives \eqref{eq:route1_core_continuum_E2}.
\end{proof}


\begin{theorem}[Identification of the Route~1 bounded state and the Lane~A sharp-local state]\label{thm:route1_sharp_local_identification}
Fix a target value $u_{\mathrm{ren}}\in(0,u_\ast]$ and let
\[
a_n:=\ell_0 2^{-n},\qquad \beta_n:=\beta(a_n;u_{\mathrm{ren}})
\]
be the tuned dyadic sequence from Proposition~\ref{prop:h59_uv_tuning_exists}.
Let $\omega_{\mathrm{bd}}^{(u_{\mathrm{ren}})}$ and $\mathcal K_{u_{\mathrm{ren}}}(\cdot,\cdot;\cdot)$ be the bounded positive-time state and dyadic OS matrix-element limits furnished by Theorem~\ref{thm:continuum_tightness}, and let $\omega^{(u_{\mathrm{ren}})}$ be the full sharp-local continuum state from Corollary~\ref{cor:LaneA_full_scope} built over that same base state.
Then
\[
\omega_{\mathrm{bd}}^{(u_{\mathrm{ren}})} = \omega^{(u_{\mathrm{ren}})}\!\restriction_{\mathcal A_+},
\]
and, for every bounded $F,G\in\mathcal A_+$ and every dyadic shift $s\in\mathbb D_{\mathrm{dyad}}$,
\[
\omega^{(u_{\mathrm{ren}})}(\Theta F\,\tau_s G)=\mathcal K_{u_{\mathrm{ren}}}(F,G;s).
\]
Consequently any Route~1 estimate proved first for the bounded dyadic OS matrix elements transfers to the same sharp-local state that enters Theorem~\ref{thm:main}; by semigroup continuity the transfer then extends from dyadic $s$ to all $s\ge 0$.
\end{theorem}

\begin{proof}
The equality on $\mathcal A_+$ and the dyadic matrix-element identity are exactly the two compatibility clauses isolated in Corollary~\ref{cor:LaneA_bounded_state_restriction}.
In particular, that corollary is not the inductive-union theorem: the full sharp-local state exists by Corollary~\ref{cor:LaneA_full_scope}, while the capwise OS structure behind the later semigroup continuity comes from Theorem~\ref{thm:LaneA_closure}(b).
The final sentence is therefore just Theorem~\ref{thm:continuum_time_OS_upgrade}(iii): once OS reconstruction is available for $\omega^{(u_{\mathrm{ren}})}$, the matrix elements $\omega^{(u_{\mathrm{ren}})}(\Theta F\,\tau_s G)$ extend continuously from dyadic $s$ to all $s\ge 0$.
\end{proof}

\begin{remark}[Route~1 closure summary]\label{rem:route1_dense_core_summary}
Theorem~\ref{thm:h60_one_shot_entrance} and Corollary~\ref{cor:route1_bounded_physical_entrance} supply the entrance constants.
Lemma~\ref{lem:block_decay_to_micro}, Proposition~\ref{prop:lattice_OS_cyclicity_local_algebra}, Corollary~\ref{cor:uniform_phys_gap_from_bounded_entrance}, and Theorem~\ref{thm:physical_gap_from_gluing} close QE2$_{\mathrm{lat}}$ by exact entrance-scale covariance transfer plus the fixed-lattice OS gap theorem.
Theorem~\ref{thm:phys_gap_noncollapse} closes QE3 on the primary Route~1 spine, while Corollary~\ref{cor:route1_QE3_gap} records the resulting continuum vacuum spectral gap.
Proposition~\ref{prop:route1_block_covariance} and Proposition~\ref{prop:route1_blocked_representatives} are retained as quantitative diagnostics for the downstream dense-core reformulation.
Theorem~\ref{thm:route1_E2_internal} then converts the already established Route~1 spectral gap into the dense-core Euclidean bound used by Lane~B.
\end{remark}


\subsubsection{Consequence: physical mass-scale noncollapse from one-shot entrance}

The one-shot entrance theorem, together with the exact entrance-scale covariance transfer proved above, is sufficient to obtain a uniform physical mass-scale lower bound along the tuned UV sequence.
We record this as an explicit corollary, emphasizing that no infinite-step RG well-posedness is required once the one-shot entrance constants are available.

\begin{corollary}[Noncollapse of the physical mass scale from one-shot entrance]
\label{cor:h60_physical_gap_noncollapse}
Fix a target value $u_{\mathrm{ren}}\in(0,u_\ast]$ and let $a_n=\ell_0 2^{-n}$, $\beta_n:=\beta(a_n;u_{\mathrm{ren}})$, and $B_n:=B(\beta_n)$.
Then there exists $m_{\mathrm{phys},\ast}>0$ such that for all $n\ge N_{\mathrm{ent}}(u_{\mathrm{ren}})$,
the (infinite-volume) Yang--Mills theory at parameters $(a_n,\beta_n)$ has a physical mass gap
\begin{equation}
\label{eq:h60_mphys_bound}
m_{\mathrm{phys}}(a_n)\ \ge\ m_{\mathrm{phys},\ast}.
\end{equation}
\noindent\emph{Remark.} The finite dyadic overhead coming from the $\exp(C/(L\beta))$ factor in Proposition~\ref{prop:BM_dyadic} is absorbed into the universal constant $m_{\mathrm{phys},\ast}$; it does not affect uniformity or positivity.

\end{corollary}

\begin{proof}
Theorem~\ref{thm:physical_gap_from_gluing} yields a uniform lower bound on $\Delta_{\mathrm{lat}}(\beta_n)/a_n$.
By definition of the physical normalization used on the primary Route~1 spine, this is exactly the physical mass gap $m_{\mathrm{phys}}(a_n)$.
Take $m_{\mathrm{phys},\ast}:=\mathfrak m_\ast$.
\end{proof}

\begin{remark}[Interpretation]
Corollary~\ref{cor:h60_physical_gap_noncollapse} shows that, once the exact entrance-scale covariance transfer is taken into account, the only scale-dependent input entering the gap constant is the finite-depth
one-shot entrance to the KP basin. No global RG dynamical control or Banach-space stable manifold is required beyond the one-shot entrance itself.
Combined with Theorem~\ref{thm:continuum_time_OS_upgrade} and Theorem~\ref{thm:phys_gap_noncollapse}, this closes QE3 on the primary Route~1 spine.
\end{remark}




\begin{corFspecialmix}[Finite mixed bounded-family packaging on a finite truncation]\label{cor:finite_cap_mixed_family_packaging}
Fix $u_{\mathrm{ren}}\in(0,u_\ast]$ and $\Delta_{\max}\ge 0$.
Let
\[
\mathcal B_{\mathrm{cap}}^{(\le \Delta_{\max})}
:=
\ast\!\bigl(\mathcal A_{\mathrm{flow}}^{(\le \Delta_{\max})},\mathcal A_+\bigr)
\]
and let $\mathcal W\subset \mathcal B_{\mathrm{cap}}^{(\le \Delta_{\max})}$ be finite.
Write $\operatorname{Alg}(\mathcal W)$ for the unital $\ast$-subalgebra generated by $\mathcal W$.
Then:
\begin{enumerate}[label=(\roman*),leftmargin=2.3em]
\item For every finite family $\mathcal F\subset \operatorname{Alg}(\mathcal W)$, the tuned Wilson expectations of the members of $\mathcal F$ converge jointly along the dyadic trajectory associated with $u_{\mathrm{ren}}$.
Equivalently, the finite-cap evaluation functionals on $\mathcal F$ satisfy the same finite-observable-set bounds used in Theorem~\ref{thm:C3U}(a)--(b) and Corollary~\ref{cor:C3U_Aplus_tuned}.

\item Consequently there is a unique functional
\[
\omega_{\mathcal W,\mathrm{cap,bnd}}^{(u_{\mathrm{ren}})}
\]
on $\operatorname{Alg}(\mathcal W)$ satisfying
\[
\omega_{\mathcal W,\mathrm{cap,bnd}}^{(u_{\mathrm{ren}})}(1)=1,
\qquad
\omega_{\mathcal W,\mathrm{cap,bnd}}^{(u_{\mathrm{ren}})}(P^\ast P)\ge 0
\quad
(P\in \operatorname{Alg}(\mathcal W)),
\]
hence a state on $\operatorname{Alg}(\mathcal W)$.
Its restrictions to
$\operatorname{Alg}(\mathcal W)\cap \mathcal A_{\mathrm{flow}}^{(\le \Delta_{\max})}$ and
$\operatorname{Alg}(\mathcal W)\cap \mathcal A_+$ agree with
$\omega_{\mathrm{flow}}^{(u_{\mathrm{ren}})}$ and
$\omega_{\mathrm{bd}}^{(u_{\mathrm{ren}})}$, respectively.

\item If $\mathcal W_1\subseteq \mathcal W_2\subset \mathcal B_{\mathrm{cap}}^{(\le \Delta_{\max})}$ are finite, then
\[
\omega_{\mathcal W_2,\mathrm{cap,bnd}}^{(u_{\mathrm{ren}})}\big|_{\operatorname{Alg}(\mathcal W_1)}
=
\omega_{\mathcal W_1,\mathrm{cap,bnd}}^{(u_{\mathrm{ren}})}.
\]
In particular, whenever
\[
X\in \operatorname{Alg}(\mathcal W_1)\cap \operatorname{Alg}(\mathcal W_2),
\]
one has
\[
\omega_{\mathcal W_1,\mathrm{cap,bnd}}^{(u_{\mathrm{ren}})}(X)
=
\omega_{\mathcal W_2,\mathrm{cap,bnd}}^{(u_{\mathrm{ren}})}(X).
\]

\item For every Euclidean time translation $\tau_T$ such that the translated family
\[
\tau_T\mathcal W:=\{\tau_TY:Y\in \mathcal W\}
\]
again lies in $\mathcal B_{\mathrm{cap}}^{(\le \Delta_{\max})}$, the packaged state attached to $\tau_T\mathcal W$ satisfies
\[
\omega_{\tau_T\mathcal W,\mathrm{cap,bnd}}^{(u_{\mathrm{ren}})}(\tau_TX)
=
\omega_{\mathcal W,\mathrm{cap,bnd}}^{(u_{\mathrm{ren}})}(X)
\qquad
\bigl(X\in \operatorname{Alg}(\mathcal W)\bigr).
\]
Only this translation-covariance clause is used later in Theorem~\ref{thm:finite_cap_flowed_sector_bridge}(i).
\end{enumerate}
\end{corFspecialmix}

\begin{proof}
\emph{Finite-family algebra step.}
By definition, $\operatorname{Alg}(\mathcal W)$ consists of finite linear combinations of noncommutative monomials in $\mathcal W\cup \mathcal W^\ast$.
Fix a finite family $\mathcal F\subset \operatorname{Alg}(\mathcal W)$.
After expanding each element of $\mathcal F$ into such monomials, there is a finite bounded observable set
\[
\mathcal G_{\mathcal F}\subset
\mathcal A_{\mathrm{flow}}^{(\le \Delta_{\max})}\cup \mathcal A_+
\]
containing every bounded flowed factor and every bounded positive-time factor appearing in those monomials.
The observable-evaluation Lipschitz step and the Wilson-to-chart-truncated transport theorem used in Theorem~\ref{thm:C3U}(a)--(b) and Corollary~\ref{cor:C3U_Aplus_tuned} are formulated for arbitrary finite bounded observable sets.
Hence the tuned Wilson expectations satisfy on $\mathcal G_{\mathcal F}$ the same finite-observable-set bounds used in that packaging mechanism.
Because products, adjoints, and finite linear combinations of bounded observables built from a fixed finite list of generators are again bounded observables depending only on that same list, the same joint convergence holds for the members of $\mathcal F$ themselves.
This proves~(i).

\emph{Packaging, positivity, and normalization.}
Varying $\mathcal F$ and using compatibility on overlaps packages these dyadic limits into a unique unital functional
$\omega_{\mathcal W,\mathrm{cap,bnd}}^{(u_{\mathrm{ren}})}$ on $\operatorname{Alg}(\mathcal W)$.
To record positivity explicitly, let $P\in \operatorname{Alg}(\mathcal W)$ and choose a finite family $\mathcal F_P\subset \operatorname{Alg}(\mathcal W)$ containing $1$, $P$, and $P^\ast P$.
For every finite-cap Wilson functional along the tuned dyadic trajectory one has
\[
\omega_{a_n}(1)=1,
\qquad
\omega_{a_n}(P^\ast P)\ge 0.
\]
Passing to the packaged dyadic limit on $\mathcal F_P$ gives
\[
\omega_{\mathcal W,\mathrm{cap,bnd}}^{(u_{\mathrm{ren}})}(1)=1,
\qquad
\omega_{\mathcal W,\mathrm{cap,bnd}}^{(u_{\mathrm{ren}})}(P^\ast P)
=
\lim_{n\to\infty}\omega_{a_n}(P^\ast P)
\ge 0.
\]
Hence $\omega_{\mathcal W,\mathrm{cap,bnd}}^{(u_{\mathrm{ren}})}$ is a state on $\operatorname{Alg}(\mathcal W)$.
The restriction statement follows because, when an element of $\operatorname{Alg}(\mathcal W)$ already lies in $\mathcal A_{\mathrm{flow}}^{(\le \Delta_{\max})}$ or in $\mathcal A_+$, the packaged limit is exactly the one already fixed by Corollary~\ref{cor:flow_state} or Corollary~\ref{cor:C3U_Aplus_tuned}.

\emph{Overlap compatibility.}
If $\mathcal W_1\subseteq \mathcal W_2$, then every $X\in \operatorname{Alg}(\mathcal W_1)$ is evaluated from the same finite-family tuned Wilson limits whether one packages inside $\operatorname{Alg}(\mathcal W_1)$ or inside $\operatorname{Alg}(\mathcal W_2)$.
Therefore
\[
\omega_{\mathcal W_2,\mathrm{cap,bnd}}^{(u_{\mathrm{ren}})}\big|_{\operatorname{Alg}(\mathcal W_1)}
=
\omega_{\mathcal W_1,\mathrm{cap,bnd}}^{(u_{\mathrm{ren}})}.
\]
This is exactly the overlap-compatibility statement in~(iii).

\emph{Symmetry inheritance.}
Finally, let $\tau_T$ be as in~(iv).
For every finite family $\mathcal F\subset \operatorname{Alg}(\mathcal W)$, the translated family $\tau_T\mathcal F$ lies in $\operatorname{Alg}(\tau_T\mathcal W)$, and for every $Y\in \mathcal F$ the finite-cap Wilson expectations satisfy
\[
\omega_{a_n}(\tau_TY)=\omega_{a_n}(Y)
\qquad
(Y\in \mathcal F)
\]
by Euclidean covariance of the finite-$a_n$ Wilson expectations together with translation covariance of the flowed representatives.
Passing to the packaged dyadic limits yields
\[
\omega_{\tau_T\mathcal W,\mathrm{cap,bnd}}^{(u_{\mathrm{ren}})}(\tau_TX)
=
\omega_{\mathcal W,\mathrm{cap,bnd}}^{(u_{\mathrm{ren}})}(X)
\qquad
\bigl(X\in \operatorname{Alg}(\mathcal W)\bigr),
\]
which is the translation-covariance clause used later in Theorem~\ref{thm:finite_cap_flowed_sector_bridge}(i).
\end{proof}

\begin{theorem}[Finite-cap mixed-correlator extension from finite-family packaged bounded capwise states]\label{thm:finite_cap_flowed_sector_bridge}
Fix $u_{\mathrm{ren}}\in(0,u_\ast]$ and $\Delta_{\max}\ge 0$.
Let
\[
\mathcal B_{\mathrm{cap}}^{(\le \Delta_{\max})}
:=
\ast\!\bigl(\mathcal A_{\mathrm{flow}}^{(\le \Delta_{\max})},\mathcal A_+\bigr)
\]
denote the unital bounded capwise algebra generated by the capwise flowed subalgebra and the bounded positive-time algebra, and write
\[
\mathcal B_{+,\mathrm{cap}}^{(\le \Delta_{\max})}
:=
\ast\!\bigl(\mathcal A_{+,\mathrm{flow}}^{(\le \Delta_{\max})},\mathcal A_+\bigr)
\subset \mathcal B_{\mathrm{cap}}^{(\le \Delta_{\max})}
\]
for its positive-time subalgebra.
For every finite family $\mathcal W\subset \mathcal B_{\mathrm{cap}}^{(\le \Delta_{\max})}$, Corollary~\ref{cor:finite_cap_mixed_family_packaging} furnishes a packaged state
\[
\omega_{\mathcal W,\mathrm{cap,bnd}}^{(u_{\mathrm{ren}})}
\]
on the generated algebra $\operatorname{Alg}(\mathcal W)$.
In every argument below, first fix a finite family $\mathcal W\subset \mathcal B_{\mathrm{cap}}^{(\le \Delta_{\max})}$ containing all bounded words appearing in the formula under consideration and work inside $\operatorname{Alg}(\mathcal W)$.
Once this ambient finite family has been fixed, we abbreviate
\[
\omega_{\mathcal W,\mathrm{cap,bnd}}^{(u_{\mathrm{ren}})}
\]
by
\[
\omega_{\mathrm{cap,bnd}}^{(u_{\mathrm{ren}})}
\]
for that calculation.
The compatibility statement in Corollary~\ref{cor:finite_cap_mixed_family_packaging} guarantees that the resulting value is independent of the chosen admissible finite family.

Let $\widehat{\mathfrak A}_{\mathrm{loc}}^{(\le \Delta_{\max})}$ be the unital $\ast$-algebra generated by
$\mathcal B_{\mathrm{cap}}^{(\le \Delta_{\max})}$ and the formal centered capwise local symbols
\[
\widehat O_j^{(\le \Delta_{\max})}(f;\mu)
\qquad
(1\le j\le N,\ f\in C_c^\infty(\mathbb R^4)).
\]
For an abstract word
\[
W
=
B_0\,\widehat O_{j_1}^{(\le \Delta_{\max})}(f_1;\mu)\,B_1\cdots
\widehat O_{j_m}^{(\le \Delta_{\max})}(f_m;\mu)\,B_m
\qquad
\bigl(B_r\in\mathcal B_{\mathrm{cap}}^{(\le \Delta_{\max})}\bigr),
\]
define its flowed representative at scales $\mathbf t=(t_1,\dots,t_m)$ by
\[
W(\mathbf t)
:=
B_0\,\mathcal O_{j_1}^{\mathrm{app},(\le \Delta_{\max})}(t_1;f_1;\mu)\,B_1\cdots
\mathcal O_{j_m}^{\mathrm{app},(\le \Delta_{\max})}(t_m;f_m;\mu)\,B_m
\in \mathcal B_{\mathrm{cap}}^{(\le \Delta_{\max})}.
\]
Write also
\[
\|W\|_{\mathrm{bnd}}:=\prod_{r=0}^{m}\|B_r\|_{\infty},
\]
and let $K_W\subset\mathbb R^4$ be a fixed compact neighborhood containing the supports of the test functions $f_r$ together with the supports of every bounded cylinder / flowed factor appearing in the chosen representatives of $B_0,\dots,B_m$.
For theorem-facing bookkeeping, let $C_W^{\mathrm{ins}}$ denote any fixed constant dominating the single-slot bounded-word insertion constants attached to $W$.
Concretely, in each individual application of Corollary~\ref{cor:ins_to_OS_bounded_word} / Proposition~\ref{prop:finite_cap_bounded_factor_insertion}, one may take
\[
C_W^{\mathrm{ins}}
=
C_{\mathrm{word}}(K_{W,f};W^{\sharp})\,\|W\|_{\mathrm{bnd}},
\]
for the relevant slotwise positive-time neighborhood $K_{W,f}$ and normalized bounded word $W^{\sharp}$, with the trivial convention $C_W^{\mathrm{ins}}=0$ if some $\|B_r\|_{\infty}=0$; when finitely many such channels arise in Lemma~\ref{lem:finite_cap_mixed_channel_reduction}, enlarge $C_W^{\mathrm{ins}}$ once and for all to dominate them all.
In the mixed-channel proof of part~(i), the constants attached to the fixed word $W$ are recorded through $K_W$, $C_W^{\mathrm{ins}}$, and the finite truncation data $(m,\Delta_{\max},\mu)$.
Then:
\begin{enumerate}[label=(\roman*),leftmargin=2.3em]
\item For every abstract word $W\in\widehat{\mathfrak A}_{\mathrm{loc}}^{(\le \Delta_{\max})}$ the limit
\[
\widetilde\omega_{\mathrm{cap}}(W)
:=
\lim_{\mathbf t\downarrow 0}
\omega_{\mathrm{cap,bnd}}^{(u_{\mathrm{ren}})}\bigl(W(\mathbf t)\bigr)
\]
exists and is independent of the order or relative rates with which the coordinates of $\mathbf t$ tend to $0$.
\item The rule $W\mapsto\widetilde\omega_{\mathrm{cap}}(W)$ defines a unique unital linear functional on
$\widehat{\mathfrak A}_{\mathrm{loc}}^{(\le \Delta_{\max})}$ extending the compatible finite-family package limits
$\omega_{\mathcal W,\mathrm{cap,bnd}}^{(u_{\mathrm{ren}})}$ furnished by Corollary~\ref{cor:finite_cap_mixed_family_packaging}.
In particular,
\[
\widetilde\omega_{\mathrm{cap}}\!\restriction_{\mathcal A_{\mathrm{flow}}^{(\le \Delta_{\max})}}
=
\omega_{\mathrm{flow}}^{(u_{\mathrm{ren}})}\!\restriction_{\mathcal A_{\mathrm{flow}}^{(\le \Delta_{\max})}}
\qquad\text{and}\qquad
\widetilde\omega_{\mathrm{cap}}\!\restriction_{\mathcal A_+}
=
\omega_{\mathrm{bd}}^{(u_{\mathrm{ren}})}.
\]
\item For every $P\in\widehat{\mathfrak A}_{\mathrm{loc}}^{(\le \Delta_{\max})}$ one has
\[
\widetilde\omega_{\mathrm{cap}}(P^\ast P)\ge 0,
\]
so $\widetilde\omega_{\mathrm{cap}}$ is a state on the capwise mixed algebra.
\item If $H,G\in\mathcal A_{+,\mathrm{flow}}^{(\le \Delta_{\max})}$ and all test functions $f_r$ have support in $\{x_0>0\}$, then, writing
\[
\widehat{\mathbf O}
:=
\widehat O_{j_1}^{(\le \Delta_{\max})}(f_1;\mu)\cdots
\widehat O_{j_m}^{(\le \Delta_{\max})}(f_m;\mu),
\]
one has
\[
\widetilde\omega_{\mathrm{cap}}\bigl(\Theta H\,\widehat{\mathbf O}\,G\bigr)
=
\big\langle [H],\widehat{\mathbf O}[G]\big\rangle_{\mathrm{cap}}.
\]
Here the local fields on the right are the graph-norm limits furnished by
Theorem~\ref{thm:field_algebra_truncation}, identified with the Lane~A inverse-flow representatives by
Corollary~\ref{cor:finite_truncation_SFTE}, on the common core of
Lemma~\ref{lem:finite_cap_common_core}.
Hence the capwise mixed-correlator extension agrees with the finite-cap positive-time OS operator realization on its natural domain.
\end{enumerate}
\end{theorem}

\begin{proof}
Corollary~\ref{cor:finite_cap_mixed_family_packaging} is the explicit packaging node for the bounded capwise mixed words used here.
Whenever a formula below involves finitely many bounded representatives, fix a finite family $\mathcal W\subset \mathcal B_{\mathrm{cap}}^{(\le \Delta_{\max})}$ containing them and evaluate with the packaged state $\omega_{\mathcal W,\mathrm{cap,bnd}}^{(u_{\mathrm{ren}})}$ furnished by that corollary.
When convenient we then abbreviate this packaged state by $\omega_{\mathrm{cap,bnd}}^{(u_{\mathrm{ren}})}$ only for that local calculation; overlap compatibility ensures that this shorthand does not change any value.

For part~(i), fix an abstract word $W$ as above and two scale vectors $\mathbf t,\mathbf s$, and retain the notation $K_W$, $\|W\|_{\mathrm{bnd}}$, and $C_W^{\mathrm{ins}}$ from the theorem statement.

\paragraph{Primary mixed-channel proof chain for part~(i).}
The mixed-channel Cauchy step is proved here once and for all through the theorem chain
\[
\text{Theorem~\ref{thm:ins_to_OS}}
\Longrightarrow
\text{Corollary~\ref{cor:ins_to_OS_bounded_word}}
\Longrightarrow
\text{Proposition~\ref{prop:finite_cap_bounded_factor_insertion}}
\]
together with
\[
\text{Theorem~\ref{thm:finite_truncation_inverse_control}}
\Longrightarrow
\text{Lemma~\ref{lem:finite_cap_mixed_channel_reduction}}
\Longrightarrow
\text{Theorem~\ref{thm:finite_cap_flowed_sector_bridge}(i)}.
\]
No later Lane~A statement reproves this step; all downstream uses cite the present theorem.

Because the word is finite and every bounded factor / local test function entering it has compact support, one may choose a common Euclidean time translation $\tau_T$ with $T\gg 1$ such that, after translating every bounded factor and every test function in $W$ by $\tau_T$, all translated bounded factors lie in
$\mathcal B_{+,\mathrm{cap}}^{(\le \Delta_{\max})}$ and all translated test functions are supported in $\{x_0>0\}$.
Write $(\tau_T W)(\mathbf t)$ for the corresponding translated bounded representative.
Write
\[
(\tau_T W)(\mathbf t)-(\tau_T W)(\mathbf s)=\sum_{r=1}^m T_r^T(\mathbf t,\mathbf s)
\]
for the multilinear telescoping decomposition of the translated word.
Lemma~\ref{lem:finite_cap_mixed_channel_reduction}, applied to $\tau_T W$, reduces each literal telescoping summand $T_r^T(\mathbf t,\mathbf s)$ to finitely many uniformly bounded positive-time word families multiplying one single-slot difference channel, and its displayed estimate \eqref{eq:finite_cap_literal_summand_bound} gives the corresponding vacuum-channel envelope.
There the coefficient growth is recorded explicitly in \eqref{eq:finite_cap_literal_summand_coeff_env} via Theorem~\ref{thm:finite_truncation_inverse_control}, while Proposition~\ref{prop:finite_cap_bounded_factor_insertion}, specialized to $\Psi=\Phi=\Omega_{\mathrm{cap}}$, supplies the factor $t_r^{\eta_{\Delta_{\max}}}+s_r^{\eta_{\Delta_{\max}}}$ after subtracting the centered finite-cap SFTE identities in the $r$th slot.
The single-slot bounded-word contribution from Proposition~\ref{prop:finite_cap_bounded_factor_insertion} is first packaged into a fixed insertion constant for the translated word $\tau_TW$; enlarging $C_W^{\mathrm{ins}}$ once more if necessary, we record that translated-word contribution again under the shorthand $C_W^{\mathrm{ins}}$, while Theorem~\ref{thm:finite_truncation_inverse_control} controls the remaining coefficient growth.
Equation~\eqref{eq:finite_cap_literal_summand_o1} records the corresponding slotwise final $o(1)$ handoff.
Thus, for each $r$, there exists a nonnegative function $\Xi_{W,r}(\mathbf t,\mathbf s)$, depending on the fixed word through $K_W$, $C_W^{\mathrm{ins}}$, and the finite truncation data, such that
\[
\big|
\big\langle \Omega_{\mathrm{cap}},T_r^T(\mathbf t,\mathbf s)\,\Omega_{\mathrm{cap}}\big\rangle_{\mathrm{cap}}
\big|
\le
\Xi_{W,r}(\mathbf t,\mathbf s),
\qquad
\Xi_{W,r}(\mathbf t,\mathbf s)\xrightarrow[\mathbf t,\mathbf s\downarrow 0]{}0.
\]
Summing over the finitely many slots gives
\[
\big|
\big\langle \Omega_{\mathrm{cap}},\bigl((\tau_T W)(\mathbf t)-(\tau_T W)(\mathbf s)\bigr)\Omega_{\mathrm{cap}}\big\rangle_{\mathrm{cap}}
\big|
\le
\sum_{r=1}^m \Xi_{W,r}(\mathbf t,\mathbf s),
\qquad
\sum_{r=1}^m \Xi_{W,r}(\mathbf t,\mathbf s)
\xrightarrow[\mathbf t,\mathbf s\downarrow 0]{}0.
\]
To expose the capwise packaging step on the page, fix once and for all a finite set
\[
\mathcal F\subset \mathcal A_{+,\mathrm{flow}}^{(\le \Delta_{\max})},
\qquad
1\in\mathcal F,
\]
and write
\[
p_{\mathcal F}(X):=\|X\|_\infty
\qquad
(X\in \mathcal F).
\]
For each $X\in\mathcal F$, define the finite-cap difference functional
\[
\Lambda^{W,\mathcal F}_{\mathbf t,\mathbf s}(X)
:=
\big\langle
\Omega_{\mathrm{cap}},
\bigl((\tau_T W)(\mathbf t)-(\tau_T W)(\mathbf s)\bigr)X\,\Omega_{\mathrm{cap}}
\big\rangle_{\mathrm{cap}}.
\]
Because $\mathcal F$ is finite and each $X\in\mathcal F$ belongs to the common core, Proposition~\ref{prop:finite_cap_bounded_factor_insertion}, applied with $\Psi=\Omega_{\mathrm{cap}}$ and $\Phi=[X]$ to each telescoping summand together with Lemma~\ref{lem:finite_cap_mixed_channel_reduction}, gives a constant $C_{W,\mathcal F}<\infty$ such that
\[
\bigl|\Lambda^{W,\mathcal F}_{\mathbf t,\mathbf s}(X)\bigr|
\le
C_{W,\mathcal F}\,p_{\mathcal F}(X)\,
\sum_{r=1}^m \Xi_{W,r}(\mathbf t,\mathbf s)
\qquad
(X\in\mathcal F),
\]
where the fixed-word dependence is recorded through $K_W$, $C_W^{\mathrm{ins}}$, and the finite truncation data.
Equivalently,
\begin{equation}\label{eq:finite_cap_flowed_sector_bridge_lambda_o1}
\bigl|\Lambda^{W,\mathcal F}_{\mathbf t,\mathbf s}(X)\bigr|
\le
C_{W,\mathcal F}\,p_{\mathcal F}(X)\,
\sum_{r=1}^m \Xi_{W,r}(\mathbf t,\mathbf s),
\qquad
\sum_{r=1}^m \Xi_{W,r}(\mathbf t,\mathbf s)
\xrightarrow[\mathbf t,\mathbf s\downarrow 0]{}0,
\end{equation}
so the finite-cap functional difference is displayed on the page as a finite sum of literal slotwise $o(1)$ handoffs coming from \eqref{eq:finite_cap_literal_summand_o1}.
Define the finite bounded family
\[
\mathcal W_{W,\mathcal F,\mathbf t,\mathbf s}
:=
\Bigl\{
\bigl((\tau_T W)(\mathbf t)-(\tau_T W)(\mathbf s)\bigr)X
:
X\in\mathcal F
\Bigr\}
\cup
\{(\tau_T W)(\mathbf t),(\tau_T W)(\mathbf s)\}
\subset
\mathcal B_{+,\mathrm{cap}}^{(\le \Delta_{\max})}.
\]
Corollary~\ref{cor:finite_cap_mixed_family_packaging} therefore furnishes a packaged state
\[
\omega_{W,\mathcal F,\mathbf t,\mathbf s}^{\mathrm{pkg}}
:=
\omega_{\mathcal W_{W,\mathcal F,\mathbf t,\mathbf s},\mathrm{cap,bnd}}^{(u_{\mathrm{ren}})}
\]
on $\operatorname{Alg}(\mathcal W_{W,\mathcal F,\mathbf t,\mathbf s})$.
In particular, since $1\in\mathcal F$, this packaged state satisfies
\[
\bigl|
\omega_{W,\mathcal F,\mathbf t,\mathbf s}^{\mathrm{pkg}}\bigl((\tau_T W)(\mathbf t)-(\tau_T W)(\mathbf s)\bigr)
\bigr|
\le
C_{W,\mathcal F}\sum_{r=1}^m \Xi_{W,r}(\mathbf t,\mathbf s)
\xrightarrow[\mathbf t,\mathbf s\downarrow 0]{}0.
\]
To compare with the unshifted word, set
\[
\mathcal W_{W,\mathbf t,\mathbf s}^{\mathrm{tr}}
:=
\{(\tau_TW)(\mathbf t),(\tau_TW)(\mathbf s)\},
\qquad
\mathcal W_{W,\mathbf t,\mathbf s}
:=
\{W(\mathbf t),W(\mathbf s)\}.
\]
By the compatibility clause of Corollary~\ref{cor:finite_cap_mixed_family_packaging}, the restriction of $\omega_{W,\mathcal F,\mathbf t,\mathbf s}^{\mathrm{pkg}}$ to $\operatorname{Alg}(\mathcal W_{W,\mathbf t,\mathbf s}^{\mathrm{tr}})$ is
\[
\omega_{\mathcal W_{W,\mathbf t,\mathbf s}^{\mathrm{tr}},\mathrm{cap,bnd}}^{(u_{\mathrm{ren}})}.
\]
Hence
\[
\bigl|
\omega_{\mathcal W_{W,\mathbf t,\mathbf s}^{\mathrm{tr}},\mathrm{cap,bnd}}^{(u_{\mathrm{ren}})}\bigl((\tau_T W)(\mathbf t)-(\tau_T W)(\mathbf s)\bigr)
\bigr|
\le
C_{W,\mathcal F}\sum_{r=1}^m \Xi_{W,r}(\mathbf t,\mathbf s)
\xrightarrow[\mathbf t,\mathbf s\downarrow 0]{}0.
\]
Because $\tau_T\mathcal W_{W,\mathbf t,\mathbf s}=\mathcal W_{W,\mathbf t,\mathbf s}^{\mathrm{tr}}$, the translation-covariance clause of Corollary~\ref{cor:finite_cap_mixed_family_packaging}, together with translation covariance of the inverse-flow representatives, gives
\[
\omega_{\mathcal W_{W,\mathbf t,\mathbf s}^{\mathrm{tr}},\mathrm{cap,bnd}}^{(u_{\mathrm{ren}})}\bigl((\tau_T W)(\mathbf t)-(\tau_T W)(\mathbf s)\bigr)
=
\omega_{\mathcal W_{W,\mathbf t,\mathbf s},\mathrm{cap,bnd}}^{(u_{\mathrm{ren}})}\bigl(W(\mathbf t)-W(\mathbf s)\bigr).
\]
Therefore
\[
\bigl|
\omega_{\mathcal W_{W,\mathbf t,\mathbf s},\mathrm{cap,bnd}}^{(u_{\mathrm{ren}})}\bigl(W(\mathbf t)-W(\mathbf s)\bigr)
\bigr|
\le
C_{W,\mathcal F}\sum_{r=1}^m \Xi_{W,r}(\mathbf t,\mathbf s)
\xrightarrow[\mathbf t,\mathbf s\downarrow 0]{}0.
\]
By the local shorthand convention fixed in the theorem statement, this is exactly the Cauchy estimate for $\omega_{\mathrm{cap,bnd}}^{(u_{\mathrm{ren}})}\bigl(W(\mathbf t)\bigr)$.
Hence the net $\omega_{\mathrm{cap,bnd}}^{(u_{\mathrm{ren}})}\bigl(W(\mathbf t)\bigr)$ is Cauchy, proving existence of the limit.
Since the displayed Cauchy estimate was obtained for an arbitrary pair of scale vectors $\mathbf t,\mathbf s\downarrow 0$, the limit is independent of the order or relative rates with which the components of $\mathbf t$ approach $0$.

\paragraph*{Local certification for the mixed-channel seam.}
\phantomsection\label{par:finite_cap_mixed_channel_certification}
For theorem-facing use, the mixed-channel seam is certified locally here as follows.
The single-slot bounded-word channel is the chain
\[
\text{Theorem~\ref{thm:ins_to_OS}}
\Longrightarrow
\text{Corollary~\ref{cor:ins_to_OS_bounded_word}}
\Longrightarrow
\text{Proposition~\ref{prop:finite_cap_bounded_factor_insertion}},
\]
and the literal telescoping reduction is the chain
\[
\text{Theorem~\ref{thm:finite_truncation_inverse_control}}
\Longrightarrow
\text{Lemma~\ref{lem:finite_cap_mixed_channel_reduction}}
\Longrightarrow
\text{Theorem~\ref{thm:finite_cap_flowed_sector_bridge}(i)}.
\]
The mixed bounded-family packaging node used here is Corollary~\ref{cor:finite_cap_mixed_family_packaging}: for each finite family of mixed bounded words it packages the corresponding capwise limits on the generated finite-family algebra, and the notation
\(
\omega_{\mathrm{cap,bnd}}^{(u_{\mathrm{ren}})}
\)
in Theorem~\ref{thm:finite_cap_flowed_sector_bridge} is only the local shorthand for
\(
\omega_{\mathcal W,\mathrm{cap,bnd}}^{(u_{\mathrm{ren}})}
\)
after the relevant finite family $\mathcal W$ has been fixed.
This packaging is invoked only after the state-free estimate of Proposition~\ref{prop:finite_cap_bounded_factor_insertion} and the literal summand reduction of Lemma~\ref{lem:finite_cap_mixed_channel_reduction} are in place.
The dependence on the fixed bounded word is recorded through the support ledger $K_W$, the packaged insertion constant $C_W^{\mathrm{ins}}$, and the finite truncation data.
Concretely, in any single-slot channel one may take
\[
C_W^{\mathrm{ins}}
=
C_{\mathrm{word}}(K_{W,f};W^{\sharp})\,\|W\|_{\mathrm{bnd}},
\qquad
\|W\|_{\mathrm{bnd}}:=\prod_r \|B_r\|_\infty,
\]
for the relevant slotwise neighborhood $K_{W,f}$ and normalized bounded word $W^{\sharp}$ from Corollary~\ref{cor:ins_to_OS_bounded_word}; when finitely many such channels appear, $C_W^{\mathrm{ins}}$ is enlarged once and for all to dominate them all.
The displayed slotwise-to-functional $o(1)$ handoff is given by
\eqref{eq:finite_cap_literal_summand_o1} together with
\eqref{eq:finite_cap_flowed_sector_bridge_lambda_o1}: the former converts the literal summand envelope into slotwise $o(1)$, and the latter records the resulting finite-cap functional difference estimate on the page.
Appendix~\ref{app:referee_audit}, \S\ref{subsec:packet6_mixed_channel_note}, records the matching theorem-location guide for this same seam.

Part~(ii) is immediate from the linearity of the approximant words and the uniqueness of the limit just proved.
If no formal local symbol is present, then $W(\mathbf t)=W\in\mathcal B_{\mathrm{cap}}^{(\le \Delta_{\max})}$ for every $\mathbf t$, so the restriction statement follows directly.

For part~(iii), let $P\in\widehat{\mathfrak A}_{\mathrm{loc}}^{(\le \Delta_{\max})}$ and choose one scale parameter for every formal local symbol appearing in $P$.
Denote by $P(\mathbf t)\in\mathcal B_{\mathrm{cap}}^{(\le \Delta_{\max})}$ the corresponding bounded representative, and set
\[
\mathcal W_{P,\mathbf t}:=\{P(\mathbf t)\}\subset \mathcal B_{\mathrm{cap}}^{(\le \Delta_{\max})}.
\]
Corollary~\ref{cor:finite_cap_mixed_family_packaging} furnishes a state
\[
\omega_{\mathcal W_{P,\mathbf t},\mathrm{cap,bnd}}^{(u_{\mathrm{ren}})}
\]
on $\operatorname{Alg}(\mathcal W_{P,\mathbf t})$, hence
\[
\omega_{\mathcal W_{P,\mathbf t},\mathrm{cap,bnd}}^{(u_{\mathrm{ren}})}\bigl(P(\mathbf t)^\ast P(\mathbf t)\bigr)\ge 0
\qquad\text{for all }\mathbf t.
\]
The approximant for $P^\ast P$ is precisely $P(\mathbf t)^\ast P(\mathbf t)$ when the same scale parameters are used on the two factors, so part~(i) allows one to pass to the limit and obtain
\[
\widetilde\omega_{\mathrm{cap}}(P^\ast P)
=
\lim_{\mathbf t\downarrow 0}
\omega_{\mathcal W_{P,\mathbf t},\mathrm{cap,bnd}}^{(u_{\mathrm{ren}})}\bigl(P(\mathbf t)^\ast P(\mathbf t)\bigr)
\ge 0.
\]
Thus $\widetilde\omega_{\mathrm{cap}}$ is positive and normalized, hence a state.

For part~(iv), fix $H,G\in\mathcal A_{+,\mathrm{flow}}^{(\le \Delta_{\max})}$ and choose any sequence of dyadic scales
$t_r^{(n)}\downarrow 0$ for each slot $r$.
Set
\[
\widehat{\mathbf O}
:=
\widehat O_{j_1}^{(\le \Delta_{\max})}(f_1;\mu)\cdots
\widehat O_{j_m}^{(\le \Delta_{\max})}(f_m;\mu)
\]
and
\[
\mathbf O_n^{\mathrm{app}}
:=
\mathcal O_{j_1}^{\mathrm{app}}(t_1^{(n)};f_1;\mu)\cdots
\mathcal O_{j_m}^{\mathrm{app}}(t_m^{(n)};f_m;\mu).
\]
Theorem~\ref{thm:field_algebra_truncation} and Lemma~\ref{lem:finite_cap_common_core} realize the local fields as graph-norm limits on the common core, and
Corollary~\ref{cor:finite_truncation_SFTE} identifies those graph-norm limits with the Lane~A inverse-flow representatives.
Therefore
\[
\big\langle [H],\widehat{\mathbf O}[G]\big\rangle_{\mathrm{cap}}
=
\lim_{n\to\infty}\big\langle [H],\mathbf O_n^{\mathrm{app}}[G]\big\rangle_{\mathrm{cap}}.
\]
For each $n$, set
\[
\mathcal W_n^{\mathrm{OS}}:=\{\Theta H\,\mathbf O_n^{\mathrm{app}}\,G\}\subset \mathcal B_{\mathrm{cap}}^{(\le \Delta_{\max})}.
\]
Corollary~\ref{cor:finite_cap_mixed_family_packaging} furnishes a packaged state
\[
\omega_{\mathcal W_n^{\mathrm{OS}},\mathrm{cap,bnd}}^{(u_{\mathrm{ren}})}
\]
on $\operatorname{Alg}(\mathcal W_n^{\mathrm{OS}})$.
Since $\mathbf O_n^{\mathrm{app}}G$ is a bounded positive-time observable and $H\in\mathcal A_{+,\mathrm{flow}}^{(\le \Delta_{\max})}$, the definition of the finite-cap OS inner product on bounded flowed positive-time classes gives
\[
\big\langle [H],\mathbf O_n^{\mathrm{app}}[G]\big\rangle_{\mathrm{cap}}
=
\omega_{\mathcal W_n^{\mathrm{OS}},\mathrm{cap,bnd}}^{(u_{\mathrm{ren}})}\bigl(\Theta H\,\mathbf O_n^{\mathrm{app}}\,G\bigr).
\]
Passing to the limit and invoking part~(i), together with the local shorthand convention fixed in the theorem statement, gives the claimed identity.
\end{proof}

\begin{center}
\fbox{\begin{minipage}{0.94\textwidth}
\textbf{End of Appendix F (primary mass-gap spine).} Theorems and propositions in this appendix close QE1, QE2$_{\mathrm{lat}}$, and QE3 on the primary Route~1 spine. Lane~B remains a downstream reformulation.
\end{minipage}}
\end{center}

